
\documentclass[11pt,a4paper]{article}
\usepackage[utf8]{inputenc}
\usepackage[T1]{fontenc}
\usepackage[spanish]{babel}
\usepackage{amsmath,amssymb}
\usepackage{geometry}
\usepackage{enumitem}
\usepackage{xcolor}
\usepackage{fancyhdr}
\usepackage{hyperref}
\geometry{margin=2cm}
\hypersetup{colorlinks=true, linkcolor=black, urlcolor=blue}
\setlength{\parindent}{0pt}
\setlength{\parskip}{0.45em}
\setlength{\headheight}{14pt}
\sloppy
\pagestyle{fancy}
\fancyhf{}
\lhead{Practico de factorizacion}
\rhead{Algebra}
\cfoot{\thepage}
\newcommand{\nivel}[2]{\subsection*{Nivel #1: #2}\addcontentsline{toc}{subsection}{Nivel #1: #2}}
\newcommand{\metodo}[1]{\section{#1}}
\newcommand{\ejercicio}[3]{\item \textbf{[#1]} \textit{#2}. $ #3 $}
\begin{document}
\begin{titlepage}
\centering
{\Huge\bfseries Practico Extenso de Factorizacion\par}
\vspace{0.7cm}
{\Large De nivel facil a nivel imposible\par}
\vspace{1cm}
{\large Factor comun, diferencia de cuadrados, aspas, suma y diferencia de potencias y Ruffini\par}
\vfill
\begin{tabular}{rl}
\textbf{Materia:} & Algebra \\
\textbf{Tema:} & Metodos de factorizacion \\
\textbf{Formato:} & Ejercicios + solucionario paso a paso \\
\textbf{Cantidad:} & 500 ejercicios \\
\end{tabular}
\vfill
{\large Elaborado en \LaTeX\par}
\end{titlepage}
\tableofcontents
\newpage
\section*{Como usar este practico}
\addcontentsline{toc}{section}{Como usar este practico}
Este material esta ordenado de menor a mayor dificultad. La recomendacion es intentar primero los ejercicios sin mirar el solucionario. Si te trabas, revisa solamente el primer paso de la solucion y vuelve a intentar.

Cada metodo tiene cinco niveles:
\begin{enumerate}
\item Nivel 1: reconocimiento directo.
\item Nivel 2: signos, coeficientes y pequenos cambios.
\item Nivel 3: expresiones con mas estructura.
\item Nivel 4: combinaciones algebraicas mas exigentes.
\item Nivel 5: nivel imposible, pensado para dominar el metodo con seguridad.
\end{enumerate}

\newpage

\part*{Ejercicios propuestos}\addcontentsline{toc}{part}{Ejercicios propuestos}
\metodo{Factor comun normal y por agrupacion}
\nivel{1}{basico}
\begin{enumerate}
\ejercicio{1}{Factor comun normal}{4 x^{3} y + 6 x^{2} y}
\ejercicio{2}{Factor comun normal}{9 x^{4} + 12 x^{3}}
\ejercicio{3}{Factor comun normal}{16 x^{2} y + 20 x y}
\ejercicio{4}{Factor comun normal}{25 x^{3} + 30 x^{2}}
\ejercicio{5}{Factor comun normal}{36 x^{4} y + 42 x^{3} y}
\ejercicio{6}{Factor comun normal}{49 x^{2} + 56 x}
\ejercicio{7}{Factor comun normal}{64 x^{3} y + 72 x^{2} y}
\ejercicio{8}{Factor comun normal}{81 x^{4} + 90 x^{3}}
\ejercicio{9}{Factor comun normal}{100 x^{2} y + 110 x y}
\ejercicio{10}{Factor comun normal}{121 x^{3} + 132 x^{2}}
\ejercicio{11}{Factor comun por agrupacion}{3 x y + 3 x + 6 y + 6}
\ejercicio{12}{Factor comun por agrupacion}{4 x y + 5 x + 12 y + 15}
\ejercicio{13}{Factor comun por agrupacion}{5 x y + 7 x + 20 y + 28}
\ejercicio{14}{Factor comun por agrupacion}{6 x y + 9 x + 30 y + 45}
\ejercicio{15}{Factor comun por agrupacion}{7 x y + 11 x + 42 y + 66}
\ejercicio{16}{Factor comun por agrupacion}{8 x y + 13 x + 56 y + 91}
\ejercicio{17}{Factor comun por agrupacion}{9 x y + 15 x + 72 y + 120}
\ejercicio{18}{Factor comun por agrupacion}{10 x y + 17 x + 90 y + 153}
\ejercicio{19}{Factor comun por agrupacion}{11 x y + 19 x + 110 y + 190}
\ejercicio{20}{Factor comun por agrupacion}{12 x y + 21 x + 132 y + 231}
\end{enumerate}
\nivel{2}{intermedio inicial}
\begin{enumerate}
\ejercicio{21}{Factor comun normal}{9 x^{4} y + 6 x^{3} y^{2} - 12 x^{3} y}
\ejercicio{22}{Factor comun normal}{16 x^{5} + 12 x^{4} y - 20 x^{4}}
\ejercicio{23}{Factor comun normal}{25 x^{3} y + 20 x^{2} y^{2} - 30 x^{2} y}
\ejercicio{24}{Factor comun normal}{36 x^{4} + 30 x^{3} y - 42 x^{3}}
\ejercicio{25}{Factor comun normal}{49 x^{5} y + 42 x^{4} y^{2} - 56 x^{4} y}
\ejercicio{26}{Factor comun normal}{64 x^{3} + 56 x^{2} y - 72 x^{2}}
\ejercicio{27}{Factor comun normal}{81 x^{4} y + 72 x^{3} y^{2} - 90 x^{3} y}
\ejercicio{28}{Factor comun normal}{100 x^{5} + 90 x^{4} y - 110 x^{4}}
\ejercicio{29}{Factor comun normal}{121 x^{3} y + 110 x^{2} y^{2} - 132 x^{2} y}
\ejercicio{30}{Factor comun normal}{144 x^{4} + 132 x^{3} y - 156 x^{3}}
\ejercicio{31}{Factor comun por agrupacion}{9 x y + 12 x - 6 y - 8}
\ejercicio{32}{Factor comun por agrupacion}{16 x y + 24 x - 12 y - 18}
\ejercicio{33}{Factor comun por agrupacion}{25 x y + 40 x - 20 y - 32}
\ejercicio{34}{Factor comun por agrupacion}{36 x y + 60 x - 30 y - 50}
\ejercicio{35}{Factor comun por agrupacion}{49 x y + 84 x - 42 y - 72}
\ejercicio{36}{Factor comun por agrupacion}{64 x y + 112 x - 56 y - 98}
\ejercicio{37}{Factor comun por agrupacion}{81 x y + 144 x - 72 y - 128}
\ejercicio{38}{Factor comun por agrupacion}{100 x y + 180 x - 90 y - 162}
\ejercicio{39}{Factor comun por agrupacion}{121 x y + 220 x - 110 y - 200}
\ejercicio{40}{Factor comun por agrupacion}{144 x y + 264 x - 132 y - 242}
\end{enumerate}
\nivel{3}{intermedio}
\begin{enumerate}
\ejercicio{41}{Factor comun normal}{8 x^{6} y^{2} - 16 x^{4} y^{3} + 12 x^{4} y^{2}}
\ejercicio{42}{Factor comun normal}{15 x^{7} y - 25 x^{5} y^{2} + 25 x^{5} y}
\ejercicio{43}{Factor comun normal}{24 x^{5} y^{2} - 36 x^{3} y^{3} + 42 x^{3} y^{2}}
\ejercicio{44}{Factor comun normal}{35 x^{6} y - 49 x^{4} y^{2} + 63 x^{4} y}
\ejercicio{45}{Factor comun normal}{48 x^{7} y^{2} - 64 x^{5} y^{3} + 88 x^{5} y^{2}}
\ejercicio{46}{Factor comun normal}{63 x^{5} y - 81 x^{3} y^{2} + 117 x^{3} y}
\ejercicio{47}{Factor comun normal}{80 x^{6} y^{2} - 100 x^{4} y^{3} + 150 x^{4} y^{2}}
\ejercicio{48}{Factor comun normal}{99 x^{7} y - 121 x^{5} y^{2} + 187 x^{5} y}
\ejercicio{49}{Factor comun normal}{120 x^{5} y^{2} - 144 x^{3} y^{3} + 228 x^{3} y^{2}}
\ejercicio{50}{Factor comun normal}{143 x^{6} y - 169 x^{4} y^{2} + 273 x^{4} y}
\ejercicio{51}{Factor comun por agrupacion}{12 x^{2} + 6 x y - 20 x - 10 y}
\ejercicio{52}{Factor comun por agrupacion}{20 x^{2} + 12 x y - 35 x - 21 y}
\ejercicio{53}{Factor comun por agrupacion}{30 x^{2} + 20 x y - 54 x - 36 y}
\ejercicio{54}{Factor comun por agrupacion}{42 x^{2} + 30 x y - 77 x - 55 y}
\ejercicio{55}{Factor comun por agrupacion}{56 x^{2} + 42 x y - 104 x - 78 y}
\ejercicio{56}{Factor comun por agrupacion}{72 x^{2} + 56 x y - 135 x - 105 y}
\ejercicio{57}{Factor comun por agrupacion}{90 x^{2} + 72 x y - 170 x - 136 y}
\ejercicio{58}{Factor comun por agrupacion}{110 x^{2} + 90 x y - 209 x - 171 y}
\ejercicio{59}{Factor comun por agrupacion}{132 x^{2} + 110 x y - 252 x - 210 y}
\ejercicio{60}{Factor comun por agrupacion}{156 x^{2} + 132 x y - 299 x - 253 y}
\end{enumerate}
\nivel{4}{avanzado}
\begin{enumerate}
\ejercicio{61}{Factor comun normal}{15 x^{7} y^{3} - 10 x^{6} y^{2} + 25 x^{5} y^{4} - 15 x^{5} y^{2}}
\ejercicio{62}{Factor comun normal}{24 x^{8} y^{2} - 18 x^{7} y + 36 x^{6} y^{3} - 24 x^{6} y}
\ejercicio{63}{Factor comun normal}{35 x^{6} y^{3} - 28 x^{5} y^{2} + 49 x^{4} y^{4} - 35 x^{4} y^{2}}
\ejercicio{64}{Factor comun normal}{48 x^{7} y^{2} - 40 x^{6} y + 64 x^{5} y^{3} - 48 x^{5} y}
\ejercicio{65}{Factor comun normal}{63 x^{8} y^{3} - 54 x^{7} y^{2} + 81 x^{6} y^{4} - 63 x^{6} y^{2}}
\ejercicio{66}{Factor comun normal}{80 x^{6} y^{2} - 70 x^{5} y + 100 x^{4} y^{3} - 80 x^{4} y}
\ejercicio{67}{Factor comun normal}{99 x^{7} y^{3} - 88 x^{6} y^{2} + 121 x^{5} y^{4} - 99 x^{5} y^{2}}
\ejercicio{68}{Factor comun normal}{120 x^{8} y^{2} - 108 x^{7} y + 144 x^{6} y^{3} - 120 x^{6} y}
\ejercicio{69}{Factor comun normal}{143 x^{6} y^{3} - 130 x^{5} y^{2} + 169 x^{4} y^{4} - 143 x^{4} y^{2}}
\ejercicio{70}{Factor comun normal}{168 x^{7} y^{2} - 154 x^{6} y + 196 x^{5} y^{3} - 168 x^{5} y}
\ejercicio{71}{Factor comun por agrupacion}{15 x^{3} + 30 x^{2} y - 6 x y - 12 y^{2}}
\ejercicio{72}{Factor comun por agrupacion}{24 x^{3} + 48 x^{2} y - 12 x y - 24 y^{2}}
\ejercicio{73}{Factor comun por agrupacion}{35 x^{3} + 70 x^{2} y - 20 x y - 40 y^{2}}
\ejercicio{74}{Factor comun por agrupacion}{48 x^{3} + 96 x^{2} y - 30 x y - 60 y^{2}}
\ejercicio{75}{Factor comun por agrupacion}{63 x^{3} + 126 x^{2} y - 42 x y - 84 y^{2}}
\ejercicio{76}{Factor comun por agrupacion}{80 x^{3} + 160 x^{2} y - 56 x y - 112 y^{2}}
\ejercicio{77}{Factor comun por agrupacion}{99 x^{3} + 198 x^{2} y - 72 x y - 144 y^{2}}
\ejercicio{78}{Factor comun por agrupacion}{120 x^{3} + 240 x^{2} y - 90 x y - 180 y^{2}}
\ejercicio{79}{Factor comun por agrupacion}{143 x^{3} + 286 x^{2} y - 110 x y - 220 y^{2}}
\ejercicio{80}{Factor comun por agrupacion}{168 x^{3} + 336 x^{2} y - 132 x y - 264 y^{2}}
\end{enumerate}
\nivel{5}{imposible}
\begin{enumerate}
\ejercicio{81}{Factor comun normal}{12 x^{9} y^{4} + 24 x^{8} y^{3} - 18 x^{7} y^{5} - 30 x^{6} y^{4} + 36 x^{6} y^{3}}
\ejercicio{82}{Factor comun normal}{21 x^{10} y^{3} + 35 x^{9} y^{2} - 28 x^{8} y^{4} - 42 x^{7} y^{3} + 49 x^{7} y^{2}}
\ejercicio{83}{Factor comun normal}{32 x^{8} y^{4} + 48 x^{7} y^{3} - 40 x^{6} y^{5} - 56 x^{5} y^{4} + 64 x^{5} y^{3}}
\ejercicio{84}{Factor comun normal}{45 x^{9} y^{3} + 63 x^{8} y^{2} - 54 x^{7} y^{4} - 72 x^{6} y^{3} + 81 x^{6} y^{2}}
\ejercicio{85}{Factor comun normal}{60 x^{10} y^{4} + 80 x^{9} y^{3} - 70 x^{8} y^{5} - 90 x^{7} y^{4} + 100 x^{7} y^{3}}
\ejercicio{86}{Factor comun normal}{77 x^{8} y^{3} + 99 x^{7} y^{2} - 88 x^{6} y^{4} - 110 x^{5} y^{3} + 121 x^{5} y^{2}}
\ejercicio{87}{Factor comun normal}{96 x^{9} y^{4} + 120 x^{8} y^{3} - 108 x^{7} y^{5} - 132 x^{6} y^{4} + 144 x^{6} y^{3}}
\ejercicio{88}{Factor comun normal}{117 x^{10} y^{3} + 143 x^{9} y^{2} - 130 x^{8} y^{4} - 156 x^{7} y^{3} + 169 x^{7} y^{2}}
\ejercicio{89}{Factor comun normal}{140 x^{8} y^{4} + 168 x^{7} y^{3} - 154 x^{6} y^{5} - 182 x^{5} y^{4} + 196 x^{5} y^{3}}
\ejercicio{90}{Factor comun normal}{165 x^{9} y^{3} + 195 x^{8} y^{2} - 180 x^{7} y^{4} - 210 x^{6} y^{3} + 225 x^{6} y^{2}}
\ejercicio{91}{Factor comun por agrupacion}{18 x^{3} y^{2} + 42 x^{2} y^{3} - 30 x^{2} y - 11 x y^{2} + 8 x + 7 y^{3} - 4 y}
\ejercicio{92}{Factor comun por agrupacion}{28 x^{3} y^{2} + 63 x^{2} y^{3} - 47 x^{2} y - 23 x y^{2} + 15 x + 9 y^{3} - 5 y}
\ejercicio{93}{Factor comun por agrupacion}{40 x^{3} y^{2} + 88 x^{2} y^{3} - 68 x^{2} y - 39 x y^{2} + 24 x + 11 y^{3} - 6 y}
\ejercicio{94}{Factor comun por agrupacion}{54 x^{3} y^{2} + 117 x^{2} y^{3} - 93 x^{2} y - 59 x y^{2} + 35 x + 13 y^{3} - 7 y}
\ejercicio{95}{Factor comun por agrupacion}{70 x^{3} y^{2} + 150 x^{2} y^{3} - 122 x^{2} y - 83 x y^{2} + 48 x + 15 y^{3} - 8 y}
\ejercicio{96}{Factor comun por agrupacion}{88 x^{3} y^{2} + 187 x^{2} y^{3} - 155 x^{2} y - 111 x y^{2} + 63 x + 17 y^{3} - 9 y}
\ejercicio{97}{Factor comun por agrupacion}{108 x^{3} y^{2} + 228 x^{2} y^{3} - 192 x^{2} y - 143 x y^{2} + 80 x + 19 y^{3} - 10 y}
\ejercicio{98}{Factor comun por agrupacion}{130 x^{3} y^{2} + 273 x^{2} y^{3} - 233 x^{2} y - 179 x y^{2} + 99 x + 21 y^{3} - 11 y}
\ejercicio{99}{Factor comun por agrupacion}{154 x^{3} y^{2} + 322 x^{2} y^{3} - 278 x^{2} y - 219 x y^{2} + 120 x + 23 y^{3} - 12 y}
\ejercicio{100}{Factor comun por agrupacion}{180 x^{3} y^{2} + 375 x^{2} y^{3} - 327 x^{2} y - 263 x y^{2} + 143 x + 25 y^{3} - 13 y}
\end{enumerate}
\metodo{Diferencia de cuadrados}
\nivel{1}{basico}
\begin{enumerate}
\ejercicio{101}{Diferencia de cuadrados}{4 x^{2} - 4}
\ejercicio{102}{Diferencia de cuadrados}{9 x^{2} - 9}
\ejercicio{103}{Diferencia de cuadrados}{16 x^{2} - 16}
\ejercicio{104}{Diferencia de cuadrados}{25 x^{2} - 25}
\ejercicio{105}{Diferencia de cuadrados}{x^{2} - 36}
\ejercicio{106}{Diferencia de cuadrados}{4 x^{2} - 49}
\ejercicio{107}{Diferencia de cuadrados}{9 x^{2} - 64}
\ejercicio{108}{Diferencia de cuadrados}{16 x^{2} - 81}
\ejercicio{109}{Diferencia de cuadrados}{25 x^{2} - 100}
\ejercicio{110}{Diferencia de cuadrados}{x^{2} - 121}
\ejercicio{111}{Diferencia de cuadrados}{4 x^{2} - 144}
\ejercicio{112}{Diferencia de cuadrados}{9 x^{2} - 169}
\ejercicio{113}{Diferencia de cuadrados}{16 x^{2} - 196}
\ejercicio{114}{Diferencia de cuadrados}{25 x^{2} - 225}
\ejercicio{115}{Diferencia de cuadrados}{x^{2} - 256}
\ejercicio{116}{Diferencia de cuadrados}{4 x^{2} - 289}
\ejercicio{117}{Diferencia de cuadrados}{9 x^{2} - 324}
\ejercicio{118}{Diferencia de cuadrados}{16 x^{2} - 361}
\ejercicio{119}{Diferencia de cuadrados}{25 x^{2} - 400}
\ejercicio{120}{Diferencia de cuadrados}{x^{2} - 441}
\end{enumerate}
\nivel{2}{intermedio inicial}
\begin{enumerate}
\ejercicio{121}{Diferencia de cuadrados}{9 x^{4} - 4 y^{2}}
\ejercicio{122}{Diferencia de cuadrados}{16 x^{2} - 9 y^{2}}
\ejercicio{123}{Diferencia de cuadrados}{25 x^{4} - 16 y^{2}}
\ejercicio{124}{Diferencia de cuadrados}{4 x^{2} - 25 y^{2}}
\ejercicio{125}{Diferencia de cuadrados}{9 x^{4} - y^{2}}
\ejercicio{126}{Diferencia de cuadrados}{16 x^{2} - 4 y^{2}}
\ejercicio{127}{Diferencia de cuadrados}{25 x^{4} - 9 y^{2}}
\ejercicio{128}{Diferencia de cuadrados}{4 x^{2} - 16 y^{2}}
\ejercicio{129}{Diferencia de cuadrados}{9 x^{4} - 25 y^{2}}
\ejercicio{130}{Diferencia de cuadrados}{16 x^{2} - y^{2}}
\ejercicio{131}{Diferencia de cuadrados}{25 x^{4} - 4 y^{2}}
\ejercicio{132}{Diferencia de cuadrados}{4 x^{2} - 9 y^{2}}
\ejercicio{133}{Diferencia de cuadrados}{9 x^{4} - 16 y^{2}}
\ejercicio{134}{Diferencia de cuadrados}{16 x^{2} - 25 y^{2}}
\ejercicio{135}{Diferencia de cuadrados}{25 x^{4} - y^{2}}
\ejercicio{136}{Diferencia de cuadrados}{4 x^{2} - 4 y^{2}}
\ejercicio{137}{Diferencia de cuadrados}{9 x^{4} - 9 y^{2}}
\ejercicio{138}{Diferencia de cuadrados}{16 x^{2} - 16 y^{2}}
\ejercicio{139}{Diferencia de cuadrados}{25 x^{4} - 25 y^{2}}
\ejercicio{140}{Diferencia de cuadrados}{4 x^{2} - y^{2}}
\end{enumerate}
\nivel{3}{intermedio}
\begin{enumerate}
\ejercicio{141}{Diferencia de cuadrados}{4 x^{2} + 8 x - 9 y^{2} + 4}
\ejercicio{142}{Diferencia de cuadrados}{9 x^{2} + 18 x - 16 y^{2} + 9}
\ejercicio{143}{Diferencia de cuadrados}{x^{2} + 8 x - 25 y^{2} + 16}
\ejercicio{144}{Diferencia de cuadrados}{4 x^{2} + 4 x - 36 y^{2} + 1}
\ejercicio{145}{Diferencia de cuadrados}{9 x^{2} + 12 x - 4 y^{2} + 4}
\ejercicio{146}{Diferencia de cuadrados}{x^{2} + 6 x - 9 y^{2} + 9}
\ejercicio{147}{Diferencia de cuadrados}{4 x^{2} + 16 x - 16 y^{2} + 16}
\ejercicio{148}{Diferencia de cuadrados}{9 x^{2} + 6 x - 25 y^{2} + 1}
\ejercicio{149}{Diferencia de cuadrados}{x^{2} + 4 x - 36 y^{2} + 4}
\ejercicio{150}{Diferencia de cuadrados}{4 x^{2} + 12 x - 4 y^{2} + 9}
\ejercicio{151}{Diferencia de cuadrados}{9 x^{2} + 24 x - 9 y^{2} + 16}
\ejercicio{152}{Diferencia de cuadrados}{x^{2} + 2 x - 16 y^{2} + 1}
\ejercicio{153}{Diferencia de cuadrados}{4 x^{2} + 8 x - 25 y^{2} + 4}
\ejercicio{154}{Diferencia de cuadrados}{9 x^{2} + 18 x - 36 y^{2} + 9}
\ejercicio{155}{Diferencia de cuadrados}{x^{2} + 8 x - 4 y^{2} + 16}
\ejercicio{156}{Diferencia de cuadrados}{4 x^{2} + 4 x - 9 y^{2} + 1}
\ejercicio{157}{Diferencia de cuadrados}{9 x^{2} + 12 x - 16 y^{2} + 4}
\ejercicio{158}{Diferencia de cuadrados}{x^{2} + 6 x - 25 y^{2} + 9}
\ejercicio{159}{Diferencia de cuadrados}{4 x^{2} + 16 x - 36 y^{2} + 16}
\ejercicio{160}{Diferencia de cuadrados}{9 x^{2} + 6 x - 4 y^{2} + 1}
\end{enumerate}
\nivel{4}{avanzado}
\begin{enumerate}
\ejercicio{161}{Diferencia de cuadrados}{5 x^{2} - 12 x y - 12 x + 4 y^{2} - 9}
\ejercicio{162}{Diferencia de cuadrados}{7 x^{2} - 24 x y - 12 x + 9 y^{2} - 4}
\ejercicio{163}{Diferencia de cuadrados}{9 x^{2} - 10 x y - 24 x + y^{2} - 9}
\ejercicio{164}{Diferencia de cuadrados}{- 21 x^{2} - 8 x y - 20 x + 4 y^{2} - 4}
\ejercicio{165}{Diferencia de cuadrados}{8 x^{2} - 18 x y - 6 x + 9 y^{2} - 9}
\ejercicio{166}{Diferencia de cuadrados}{12 x^{2} - 8 x y - 8 x + y^{2} - 4}
\ejercicio{167}{Diferencia de cuadrados}{16 x^{2} - 20 x y - 18 x + 4 y^{2} - 9}
\ejercicio{168}{Diferencia de cuadrados}{- 12 x^{2} - 12 x y - 16 x + 9 y^{2} - 4}
\ejercicio{169}{Diferencia de cuadrados}{- 16 x^{2} - 6 x y - 30 x + y^{2} - 9}
\ejercicio{170}{Diferencia de cuadrados}{15 x^{2} - 16 x y - 4 x + 4 y^{2} - 4}
\ejercicio{171}{Diferencia de cuadrados}{21 x^{2} - 30 x y - 12 x + 9 y^{2} - 9}
\ejercicio{172}{Diferencia de cuadrados}{- 5 x^{2} - 4 x y - 12 x + y^{2} - 4}
\ejercicio{173}{Diferencia de cuadrados}{- 7 x^{2} - 12 x y - 24 x + 4 y^{2} - 9}
\ejercicio{174}{Diferencia de cuadrados}{- 9 x^{2} - 24 x y - 20 x + 9 y^{2} - 4}
\ejercicio{175}{Diferencia de cuadrados}{24 x^{2} - 10 x y - 6 x + y^{2} - 9}
\ejercicio{176}{Diferencia de cuadrados}{- 8 x y - 8 x + 4 y^{2} - 4}
\ejercicio{177}{Diferencia de cuadrados}{- 18 x y - 18 x + 9 y^{2} - 9}
\ejercicio{178}{Diferencia de cuadrados}{- 8 x y - 16 x + y^{2} - 4}
\ejercicio{179}{Diferencia de cuadrados}{- 20 x y - 30 x + 4 y^{2} - 9}
\ejercicio{180}{Diferencia de cuadrados}{3 x^{2} - 12 x y - 4 x + 9 y^{2} - 4}
\end{enumerate}
\nivel{5}{imposible}
\begin{enumerate}
\ejercicio{181}{Diferencia de cuadrados}{4 x^{4} - 8 x^{3} y + 4 x^{2} y^{2} + 4 x^{2} - 12 x y^{2} - 8 x y + 8 x - 9 y^{4} + 12 y^{2}}
\ejercicio{182}{Diferencia de cuadrados}{9 x^{4} - 18 x^{3} y + 9 x^{2} y^{2} + 9 x^{2} - 24 x y^{2} - 18 x y + 6 x - 16 y^{4} + 8 y^{2} + 8}
\ejercicio{183}{Diferencia de cuadrados}{x^{4} - 8 x^{3} y + 16 x^{2} y^{2} + 7 x^{2} - 10 x y^{2} - 32 x y + 4 x - 25 y^{4} + 20 y^{2} + 12}
\ejercicio{184}{Diferencia de cuadrados}{4 x^{4} - 4 x^{3} y + x^{2} y^{2} + 16 x^{2} - 8 x y^{2} - 10 x y + 4 x - 4 y^{4} + 4 y^{2} + 24}
\ejercicio{185}{Diferencia de cuadrados}{9 x^{4} - 12 x^{3} y + 4 x^{2} y^{2} - 3 x^{2} - 18 x y^{2} - 4 x y + 12 x - 9 y^{4} + 12 y^{2} - 3}
\ejercicio{186}{Diferencia de cuadrados}{x^{4} - 6 x^{3} y + 9 x^{2} y^{2} + 3 x^{2} - 8 x y^{2} - 12 x y + 2 x - 16 y^{4} + 8 y^{2} + 3}
\ejercicio{187}{Diferencia de cuadrados}{4 x^{4} - 16 x^{3} y + 16 x^{2} y^{2} + 8 x^{2} - 20 x y^{2} - 24 x y + 8 x - 25 y^{4} + 20 y^{2} + 5}
\ejercicio{188}{Diferencia de cuadrados}{9 x^{4} - 6 x^{3} y + x^{2} y^{2} + 15 x^{2} - 12 x y^{2} - 8 x y + 6 x - 4 y^{4} + 4 y^{2} + 15}
\ejercicio{189}{Diferencia de cuadrados}{x^{4} - 4 x^{3} y + 4 x^{2} y^{2} + 9 x^{2} - 6 x y^{2} - 20 x y + 4 x - 9 y^{4} + 12 y^{2} + 21}
\ejercicio{190}{Diferencia de cuadrados}{4 x^{4} - 12 x^{3} y + 9 x^{2} y^{2} - 16 x y^{2} - 6 x y + 4 x - 16 y^{4} + 8 y^{2}}
\ejercicio{191}{Diferencia de cuadrados}{9 x^{4} - 24 x^{3} y + 16 x^{2} y^{2} + 3 x^{2} - 30 x y^{2} - 16 x y + 12 x - 25 y^{4} + 20 y^{2}}
\ejercicio{192}{Diferencia de cuadrados}{x^{4} - 2 x^{3} y + x^{2} y^{2} + 5 x^{2} - 4 x y^{2} - 6 x y + 2 x - 4 y^{4} + 4 y^{2} + 8}
\ejercicio{193}{Diferencia de cuadrados}{4 x^{4} - 8 x^{3} y + 4 x^{2} y^{2} + 12 x^{2} - 12 x y^{2} - 16 x y + 8 x - 9 y^{4} + 12 y^{2} + 12}
\ejercicio{194}{Diferencia de cuadrados}{9 x^{4} - 18 x^{3} y + 9 x^{2} y^{2} + 21 x^{2} - 24 x y^{2} - 30 x y + 6 x - 16 y^{4} + 8 y^{2} + 24}
\ejercicio{195}{Diferencia de cuadrados}{x^{4} - 8 x^{3} y + 16 x^{2} y^{2} + x^{2} - 10 x y^{2} - 8 x y + 4 x - 25 y^{4} + 20 y^{2} - 3}
\ejercicio{196}{Diferencia de cuadrados}{4 x^{4} - 4 x^{3} y + x^{2} y^{2} + 4 x^{2} - 8 x y^{2} - 4 x y + 4 x - 4 y^{4} + 4 y^{2} + 3}
\ejercicio{197}{Diferencia de cuadrados}{9 x^{4} - 12 x^{3} y + 4 x^{2} y^{2} + 9 x^{2} - 18 x y^{2} - 12 x y + 12 x - 9 y^{4} + 12 y^{2} + 5}
\ejercicio{198}{Diferencia de cuadrados}{x^{4} - 6 x^{3} y + 9 x^{2} y^{2} + 7 x^{2} - 8 x y^{2} - 24 x y + 2 x - 16 y^{4} + 8 y^{2} + 15}
\ejercicio{199}{Diferencia de cuadrados}{4 x^{4} - 16 x^{3} y + 16 x^{2} y^{2} + 16 x^{2} - 20 x y^{2} - 40 x y + 8 x - 25 y^{4} + 20 y^{2} + 21}
\ejercicio{200}{Diferencia de cuadrados}{9 x^{4} - 6 x^{3} y + x^{2} y^{2} - 3 x^{2} - 12 x y^{2} - 2 x y + 6 x - 4 y^{4} + 4 y^{2}}
\end{enumerate}
\metodo{Metodo de aspas}
\nivel{1}{basico}
\begin{enumerate}
\ejercicio{201}{Aspas}{x^{2} + 5 x + 6}
\ejercicio{202}{Aspas}{x^{2} + 7 x + 12}
\ejercicio{203}{Aspas}{x^{2} + 9 x + 20}
\ejercicio{204}{Aspas}{x^{2} + 11 x + 30}
\ejercicio{205}{Aspas}{x^{2} + 13 x + 42}
\ejercicio{206}{Aspas}{x^{2} + 9 x + 8}
\ejercicio{207}{Aspas}{x^{2} + 4 x + 4}
\ejercicio{208}{Aspas}{x^{2} + 6 x + 9}
\ejercicio{209}{Aspas}{x^{2} + 8 x + 16}
\ejercicio{210}{Aspas}{x^{2} + 10 x + 25}
\ejercicio{211}{Aspas}{x^{2} + 12 x + 36}
\ejercicio{212}{Aspas}{x^{2} + 8 x + 7}
\ejercicio{213}{Aspas}{x^{2} + 10 x + 16}
\ejercicio{214}{Aspas}{x^{2} + 5 x + 6}
\ejercicio{215}{Aspas}{x^{2} + 7 x + 12}
\ejercicio{216}{Aspas}{x^{2} + 9 x + 20}
\ejercicio{217}{Aspas}{x^{2} + 11 x + 30}
\ejercicio{218}{Aspas}{x^{2} + 7 x + 6}
\ejercicio{219}{Aspas}{x^{2} + 9 x + 14}
\ejercicio{220}{Aspas}{x^{2} + 11 x + 24}
\end{enumerate}
\nivel{2}{intermedio inicial}
\begin{enumerate}
\ejercicio{221}{Aspas}{x^{2} + x - 6}
\ejercicio{222}{Aspas}{x^{2} - x - 12}
\ejercicio{223}{Aspas}{x^{2} + x - 20}
\ejercicio{224}{Aspas}{x^{2} - x - 30}
\ejercicio{225}{Aspas}{x^{2} + x - 42}
\ejercicio{226}{Aspas}{x^{2} - 7 x - 8}
\ejercicio{227}{Aspas}{x^{2} - 4}
\ejercicio{228}{Aspas}{x^{2} - 9}
\ejercicio{229}{Aspas}{x^{2} - 16}
\ejercicio{230}{Aspas}{x^{2} - 25}
\ejercicio{231}{Aspas}{x^{2} - 36}
\ejercicio{232}{Aspas}{x^{2} - 6 x - 7}
\ejercicio{233}{Aspas}{x^{2} + 6 x - 16}
\ejercicio{234}{Aspas}{x^{2} + x - 6}
\ejercicio{235}{Aspas}{x^{2} - x - 12}
\ejercicio{236}{Aspas}{x^{2} + x - 20}
\ejercicio{237}{Aspas}{x^{2} - x - 30}
\ejercicio{238}{Aspas}{x^{2} - 5 x - 6}
\ejercicio{239}{Aspas}{x^{2} + 5 x - 14}
\ejercicio{240}{Aspas}{x^{2} - 5 x - 24}
\end{enumerate}
\nivel{3}{intermedio}
\begin{enumerate}
\ejercicio{241}{Aspas}{6 x^{2} + 5 x - 6}
\ejercicio{242}{Aspas}{12 x^{2} - 7 x - 12}
\ejercicio{243}{Aspas}{5 x^{2} + 21 x - 20}
\ejercicio{244}{Aspas}{4 x^{2} - 2 x - 30}
\ejercicio{245}{Aspas}{9 x^{2} + 18 x - 7}
\ejercicio{246}{Aspas}{4 x^{2} - 6 x - 4}
\ejercicio{247}{Aspas}{10 x^{2} + 9 x - 9}
\ejercicio{248}{Aspas}{6 x^{2} + 4 x - 16}
\ejercicio{249}{Aspas}{3 x^{2} + 10 x - 25}
\ejercicio{250}{Aspas}{8 x^{2} - 22 x - 6}
\ejercicio{251}{Aspas}{15 x^{2} + 29 x - 14}
\ejercicio{252}{Aspas}{2 x^{2} - x - 6}
\ejercicio{253}{Aspas}{6 x^{2} + x - 12}
\ejercicio{254}{Aspas}{12 x^{2} - x - 20}
\ejercicio{255}{Aspas}{5 x^{2} + 24 x - 5}
\ejercicio{256}{Aspas}{4 x^{2} - 8 x - 12}
\ejercicio{257}{Aspas}{9 x^{2} + 12 x - 21}
\ejercicio{258}{Aspas}{4 x^{2} - 4 x - 8}
\ejercicio{259}{Aspas}{10 x^{2} + 5 x - 15}
\ejercicio{260}{Aspas}{6 x^{2} - 5 x - 4}
\end{enumerate}
\nivel{4}{avanzado}
\begin{enumerate}
\ejercicio{261}{Aspas}{6 x^{2} + 5 x y - 6 y^{2}}
\ejercicio{262}{Aspas}{12 x^{2} - 7 x y - 12 y^{2}}
\ejercicio{263}{Aspas}{5 x^{2} + 21 x y - 20 y^{2}}
\ejercicio{264}{Aspas}{4 x^{2} - 2 x y - 30 y^{2}}
\ejercicio{265}{Aspas}{9 x^{2} + 18 x y - 7 y^{2}}
\ejercicio{266}{Aspas}{4 x^{2} - 6 x y - 4 y^{2}}
\ejercicio{267}{Aspas}{10 x^{2} + 9 x y - 9 y^{2}}
\ejercicio{268}{Aspas}{6 x^{2} + 4 x y - 16 y^{2}}
\ejercicio{269}{Aspas}{3 x^{2} + 10 x y - 25 y^{2}}
\ejercicio{270}{Aspas}{8 x^{2} - 22 x y - 6 y^{2}}
\ejercicio{271}{Aspas}{15 x^{2} + 29 x y - 14 y^{2}}
\ejercicio{272}{Aspas}{2 x^{2} - x y - 6 y^{2}}
\ejercicio{273}{Aspas}{6 x^{2} + x y - 12 y^{2}}
\ejercicio{274}{Aspas}{12 x^{2} - x y - 20 y^{2}}
\ejercicio{275}{Aspas}{5 x^{2} + 24 x y - 5 y^{2}}
\ejercicio{276}{Aspas}{4 x^{2} - 8 x y - 12 y^{2}}
\ejercicio{277}{Aspas}{9 x^{2} + 12 x y - 21 y^{2}}
\ejercicio{278}{Aspas}{4 x^{2} - 4 x y - 8 y^{2}}
\ejercicio{279}{Aspas}{10 x^{2} + 5 x y - 15 y^{2}}
\ejercicio{280}{Aspas}{6 x^{2} - 5 x y - 4 y^{2}}
\end{enumerate}
\nivel{5}{imposible}
\begin{enumerate}
\ejercicio{281}{Aspas}{9 x^{4} + 6 x^{2} y - 8 y^{2}}
\ejercicio{282}{Aspas}{16 x^{4} - 8 x^{2} y - 15 y^{2}}
\ejercicio{283}{Aspas}{25 x^{4} + 10 x^{2} y - 24 y^{2}}
\ejercicio{284}{Aspas}{12 x^{4} - 32 x^{2} y - 35 y^{2}}
\ejercicio{285}{Aspas}{6 x^{4} - 2 x^{2} y - 48 y^{2}}
\ejercicio{286}{Aspas}{12 x^{4} - 23 x^{2} y - 9 y^{2}}
\ejercicio{287}{Aspas}{20 x^{4} + 2 x^{2} y - 6 y^{2}}
\ejercicio{288}{Aspas}{10 x^{4} - 14 x^{2} y - 12 y^{2}}
\ejercicio{289}{Aspas}{18 x^{4} + 18 x^{2} y - 20 y^{2}}
\ejercicio{290}{Aspas}{8 x^{4} + 8 x^{2} y - 30 y^{2}}
\ejercicio{291}{Aspas}{15 x^{4} - 9 x^{2} y - 42 y^{2}}
\ejercicio{292}{Aspas}{8 x^{4} - 30 x^{2} y - 8 y^{2}}
\ejercicio{293}{Aspas}{15 x^{4} + 39 x^{2} y - 18 y^{2}}
\ejercicio{294}{Aspas}{24 x^{4} - 6 x^{2} y - 9 y^{2}}
\ejercicio{295}{Aspas}{10 x^{4} - 12 x^{2} y - 16 y^{2}}
\ejercicio{296}{Aspas}{6 x^{4} - 5 x^{2} y - 25 y^{2}}
\ejercicio{297}{Aspas}{12 x^{4} + 6 x^{2} y - 36 y^{2}}
\ejercicio{298}{Aspas}{20 x^{4} - 31 x^{2} y - 7 y^{2}}
\ejercicio{299}{Aspas}{30 x^{4} + 38 x^{2} y - 16 y^{2}}
\ejercicio{300}{Aspas}{4 x^{4} - 12 x^{2} y - 27 y^{2}}
\end{enumerate}
\metodo{Suma y diferencia de potencias}
\nivel{1}{basico}
\begin{enumerate}
\ejercicio{301}{Diferencia de potencias}{\left(x\right)^{3}-\left(2\right)^{3}}
\ejercicio{302}{Suma de potencias}{\left(x\right)^{3}+\left(3\right)^{3}}
\ejercicio{303}{Diferencia de potencias}{\left(x\right)^{3}-\left(4\right)^{3}}
\ejercicio{304}{Suma de potencias}{\left(x\right)^{3}+\left(5\right)^{3}}
\ejercicio{305}{Diferencia de potencias}{\left(x\right)^{3}-\left(1\right)^{3}}
\ejercicio{306}{Suma de potencias}{\left(x\right)^{3}+\left(2\right)^{3}}
\ejercicio{307}{Diferencia de potencias}{\left(x\right)^{3}-\left(3\right)^{3}}
\ejercicio{308}{Suma de potencias}{\left(x\right)^{3}+\left(4\right)^{3}}
\ejercicio{309}{Diferencia de potencias}{\left(x\right)^{3}-\left(5\right)^{3}}
\ejercicio{310}{Suma de potencias}{\left(x\right)^{3}+\left(1\right)^{3}}
\ejercicio{311}{Diferencia de potencias}{\left(x\right)^{3}-\left(2\right)^{3}}
\ejercicio{312}{Suma de potencias}{\left(x\right)^{3}+\left(3\right)^{3}}
\ejercicio{313}{Diferencia de potencias}{\left(x\right)^{3}-\left(4\right)^{3}}
\ejercicio{314}{Suma de potencias}{\left(x\right)^{3}+\left(5\right)^{3}}
\ejercicio{315}{Diferencia de potencias}{\left(x\right)^{3}-\left(1\right)^{3}}
\ejercicio{316}{Suma de potencias}{\left(x\right)^{3}+\left(2\right)^{3}}
\ejercicio{317}{Diferencia de potencias}{\left(x\right)^{3}-\left(3\right)^{3}}
\ejercicio{318}{Suma de potencias}{\left(x\right)^{3}+\left(4\right)^{3}}
\ejercicio{319}{Diferencia de potencias}{\left(x\right)^{3}-\left(5\right)^{3}}
\ejercicio{320}{Suma de potencias}{\left(x\right)^{3}+\left(1\right)^{3}}
\end{enumerate}
\nivel{2}{intermedio inicial}
\begin{enumerate}
\ejercicio{321}{Diferencia de potencias}{\left(2 x\right)^{3}-\left(2 y\right)^{3}}
\ejercicio{322}{Suma de potencias}{\left(3 x\right)^{5}+\left(3 y\right)^{5}}
\ejercicio{323}{Diferencia de potencias}{\left(x\right)^{4}-\left(4 y\right)^{4}}
\ejercicio{324}{Suma de potencias}{\left(2 x\right)^{5}+\left(y\right)^{5}}
\ejercicio{325}{Diferencia de potencias}{\left(3 x\right)^{3}-\left(2 y\right)^{3}}
\ejercicio{326}{Suma de potencias}{\left(x\right)^{5}+\left(3 y\right)^{5}}
\ejercicio{327}{Diferencia de potencias}{\left(2 x\right)^{3}-\left(4 y\right)^{3}}
\ejercicio{328}{Suma de potencias}{\left(3 x\right)^{5}+\left(y\right)^{5}}
\ejercicio{329}{Diferencia de potencias}{\left(x\right)^{4}-\left(2 y\right)^{4}}
\ejercicio{330}{Suma de potencias}{\left(2 x\right)^{5}+\left(3 y\right)^{5}}
\ejercicio{331}{Diferencia de potencias}{\left(3 x\right)^{3}-\left(4 y\right)^{3}}
\ejercicio{332}{Suma de potencias}{\left(x\right)^{5}+\left(y\right)^{5}}
\ejercicio{333}{Diferencia de potencias}{\left(2 x\right)^{3}-\left(2 y\right)^{3}}
\ejercicio{334}{Suma de potencias}{\left(3 x\right)^{5}+\left(3 y\right)^{5}}
\ejercicio{335}{Diferencia de potencias}{\left(x\right)^{4}-\left(4 y\right)^{4}}
\ejercicio{336}{Suma de potencias}{\left(2 x\right)^{5}+\left(y\right)^{5}}
\ejercicio{337}{Diferencia de potencias}{\left(3 x\right)^{3}-\left(2 y\right)^{3}}
\ejercicio{338}{Suma de potencias}{\left(x\right)^{5}+\left(3 y\right)^{5}}
\ejercicio{339}{Diferencia de potencias}{\left(2 x\right)^{3}-\left(4 y\right)^{3}}
\ejercicio{340}{Suma de potencias}{\left(3 x\right)^{5}+\left(y\right)^{5}}
\end{enumerate}
\nivel{3}{intermedio}
\begin{enumerate}
\ejercicio{341}{Diferencia de potencias}{\left(x + 2\right)^{4}-\left(2 y\right)^{4}}
\ejercicio{342}{Suma de potencias}{\left(x + 3\right)^{5}+\left(3 y\right)^{5}}
\ejercicio{343}{Diferencia de potencias}{\left(x + 4\right)^{4}-\left(y\right)^{4}}
\ejercicio{344}{Suma de potencias}{\left(x + 1\right)^{5}+\left(2 y\right)^{5}}
\ejercicio{345}{Diferencia de potencias}{\left(x + 2\right)^{4}-\left(3 y\right)^{4}}
\ejercicio{346}{Suma de potencias}{\left(x + 3\right)^{5}+\left(y\right)^{5}}
\ejercicio{347}{Diferencia de potencias}{\left(x + 4\right)^{4}-\left(2 y\right)^{4}}
\ejercicio{348}{Suma de potencias}{\left(x + 1\right)^{5}+\left(3 y\right)^{5}}
\ejercicio{349}{Diferencia de potencias}{\left(x + 2\right)^{4}-\left(y\right)^{4}}
\ejercicio{350}{Suma de potencias}{\left(x + 3\right)^{5}+\left(2 y\right)^{5}}
\ejercicio{351}{Diferencia de potencias}{\left(x + 4\right)^{4}-\left(3 y\right)^{4}}
\ejercicio{352}{Suma de potencias}{\left(x + 1\right)^{5}+\left(y\right)^{5}}
\ejercicio{353}{Diferencia de potencias}{\left(x + 2\right)^{4}-\left(2 y\right)^{4}}
\ejercicio{354}{Suma de potencias}{\left(x + 3\right)^{5}+\left(3 y\right)^{5}}
\ejercicio{355}{Diferencia de potencias}{\left(x + 4\right)^{4}-\left(y\right)^{4}}
\ejercicio{356}{Suma de potencias}{\left(x + 1\right)^{5}+\left(2 y\right)^{5}}
\ejercicio{357}{Diferencia de potencias}{\left(x + 2\right)^{4}-\left(3 y\right)^{4}}
\ejercicio{358}{Suma de potencias}{\left(x + 3\right)^{5}+\left(y\right)^{5}}
\ejercicio{359}{Diferencia de potencias}{\left(x + 4\right)^{4}-\left(2 y\right)^{4}}
\ejercicio{360}{Suma de potencias}{\left(x + 1\right)^{5}+\left(3 y\right)^{5}}
\end{enumerate}
\nivel{4}{avanzado}
\begin{enumerate}
\ejercicio{361}{Diferencia de potencias}{\left(2 x - 2 y\right)^{6}-\left(2 x + 2\right)^{6}}
\ejercicio{362}{Suma de potencias}{\left(3 x - 3 y\right)^{5}+\left(3 x + 1\right)^{5}}
\ejercicio{363}{Diferencia de potencias}{\left(x - 4 y\right)^{6}-\left(4 x + 2\right)^{6}}
\ejercicio{364}{Suma de potencias}{\left(2 x - y\right)^{5}+\left(5 x + 1\right)^{5}}
\ejercicio{365}{Diferencia de potencias}{\left(3 x - 2 y\right)^{8}-\left(x + 2\right)^{8}}
\ejercicio{366}{Suma de potencias}{\left(x - 3 y\right)^{5}+\left(2 x + 1\right)^{5}}
\ejercicio{367}{Diferencia de potencias}{\left(2 x - 4 y\right)^{6}-\left(3 x + 2\right)^{6}}
\ejercicio{368}{Suma de potencias}{\left(3 x - y\right)^{5}+\left(4 x + 1\right)^{5}}
\ejercicio{369}{Diferencia de potencias}{\left(x - 2 y\right)^{6}-\left(5 x + 2\right)^{6}}
\ejercicio{370}{Suma de potencias}{\left(2 x - 3 y\right)^{5}+\left(x + 1\right)^{5}}
\ejercicio{371}{Diferencia de potencias}{\left(3 x - 4 y\right)^{6}-\left(2 x + 2\right)^{6}}
\ejercicio{372}{Suma de potencias}{\left(x - y\right)^{5}+\left(3 x + 1\right)^{5}}
\ejercicio{373}{Diferencia de potencias}{\left(2 x - 2 y\right)^{6}-\left(4 x + 2\right)^{6}}
\ejercicio{374}{Suma de potencias}{\left(3 x - 3 y\right)^{5}+\left(5 x + 1\right)^{5}}
\ejercicio{375}{Diferencia de potencias}{\left(x - 4 y\right)^{8}-\left(x + 2\right)^{8}}
\ejercicio{376}{Suma de potencias}{\left(2 x - y\right)^{5}+\left(2 x + 1\right)^{5}}
\ejercicio{377}{Diferencia de potencias}{\left(3 x - 2 y\right)^{6}-\left(3 x + 2\right)^{6}}
\ejercicio{378}{Suma de potencias}{\left(x - 3 y\right)^{5}+\left(4 x + 1\right)^{5}}
\ejercicio{379}{Diferencia de potencias}{\left(2 x - 4 y\right)^{6}-\left(5 x + 2\right)^{6}}
\ejercicio{380}{Suma de potencias}{\left(3 x - y\right)^{5}+\left(x + 1\right)^{5}}
\end{enumerate}
\nivel{5}{imposible}
\begin{enumerate}
\ejercicio{381}{Diferencia de potencias}{\left(2 x^{2} - 2 y\right)^{8}-\left(3 x y + 2\right)^{8}}
\ejercicio{382}{Suma de potencias}{\left(3 x^{2} - 3 y\right)^{7}+\left(4 x y + 3\right)^{7}}
\ejercicio{383}{Diferencia de potencias}{\left(x^{2} - 4 y\right)^{9}-\left(5 x y + 4\right)^{9}}
\ejercicio{384}{Suma de potencias}{\left(2 x^{2} - y\right)^{7}+\left(2 x y + 5\right)^{7}}
\ejercicio{385}{Diferencia de potencias}{\left(3 x^{2} - 2 y\right)^{8}-\left(3 x y + 1\right)^{8}}
\ejercicio{386}{Suma de potencias}{\left(x^{2} - 3 y\right)^{7}+\left(4 x y + 2\right)^{7}}
\ejercicio{387}{Diferencia de potencias}{\left(2 x^{2} - 4 y\right)^{8}-\left(5 x y + 3\right)^{8}}
\ejercicio{388}{Suma de potencias}{\left(3 x^{2} - y\right)^{7}+\left(2 x y + 4\right)^{7}}
\ejercicio{389}{Diferencia de potencias}{\left(x^{2} - 2 y\right)^{9}-\left(3 x y + 5\right)^{9}}
\ejercicio{390}{Suma de potencias}{\left(2 x^{2} - 3 y\right)^{7}+\left(4 x y + 1\right)^{7}}
\ejercicio{391}{Diferencia de potencias}{\left(3 x^{2} - 4 y\right)^{8}-\left(5 x y + 2\right)^{8}}
\ejercicio{392}{Suma de potencias}{\left(x^{2} - y\right)^{7}+\left(2 x y + 3\right)^{7}}
\ejercicio{393}{Diferencia de potencias}{\left(2 x^{2} - 2 y\right)^{8}-\left(3 x y + 4\right)^{8}}
\ejercicio{394}{Suma de potencias}{\left(3 x^{2} - 3 y\right)^{7}+\left(4 x y + 5\right)^{7}}
\ejercicio{395}{Diferencia de potencias}{\left(x^{2} - 4 y\right)^{9}-\left(5 x y + 1\right)^{9}}
\ejercicio{396}{Suma de potencias}{\left(2 x^{2} - y\right)^{7}+\left(2 x y + 2\right)^{7}}
\ejercicio{397}{Diferencia de potencias}{\left(3 x^{2} - 2 y\right)^{8}-\left(3 x y + 3\right)^{8}}
\ejercicio{398}{Suma de potencias}{\left(x^{2} - 3 y\right)^{7}+\left(4 x y + 4\right)^{7}}
\ejercicio{399}{Diferencia de potencias}{\left(2 x^{2} - 4 y\right)^{8}-\left(5 x y + 5\right)^{8}}
\ejercicio{400}{Suma de potencias}{\left(3 x^{2} - y\right)^{7}+\left(2 x y + 1\right)^{7}}
\end{enumerate}
\metodo{Ruffini}
\nivel{1}{basico}
\begin{enumerate}
\ejercicio{401}{Ruffini}{x^{3} + 3 x^{2} - 4 x - 12}
\ejercicio{402}{Ruffini}{x^{3} + 4 x^{2} - 9 x - 36}
\ejercicio{403}{Ruffini}{x^{3} + 2 x^{2} - 19 x - 20}
\ejercicio{404}{Ruffini}{x^{3} - x^{2} - 16 x - 20}
\ejercicio{405}{Ruffini}{x^{3} + 5 x^{2} + 3 x - 9}
\ejercicio{406}{Ruffini}{x^{3} + 3 x^{2} - 6 x - 8}
\ejercicio{407}{Ruffini}{x^{3} + 4 x^{2} - 11 x - 30}
\ejercicio{408}{Ruffini}{x^{3} + x^{2} - 14 x - 24}
\ejercicio{409}{Ruffini}{x^{3} - x^{2} - 17 x - 15}
\ejercicio{410}{Ruffini}{x^{3} + 5 x^{2} + 2 x - 8}
\ejercicio{411}{Ruffini}{x^{3} + 6 x^{2} - x - 30}
\ejercicio{412}{Ruffini}{x^{3} - 7 x - 6}
\ejercicio{413}{Ruffini}{x^{3} + x^{2} - 14 x - 24}
\ejercicio{414}{Ruffini}{x^{3} + 2 x^{2} - 23 x - 60}
\ejercicio{415}{Ruffini}{x^{3} + 5 x^{2} - x - 5}
\ejercicio{416}{Ruffini}{x^{3} + 2 x^{2} - 4 x - 8}
\ejercicio{417}{Ruffini}{x^{3} + 3 x^{2} - 9 x - 27}
\ejercicio{418}{Ruffini}{x^{3} + x^{2} - 16 x - 16}
\ejercicio{419}{Ruffini}{x^{3} + 2 x^{2} - 25 x - 50}
\ejercicio{420}{Ruffini}{x^{3} + 4 x^{2} + x - 6}
\end{enumerate}
\nivel{2}{intermedio inicial}
\begin{enumerate}
\ejercicio{421}{Ruffini}{x^{3} + 3 x^{2} - 4 x - 12}
\ejercicio{422}{Ruffini}{x^{3} + 4 x^{2} - 9 x - 36}
\ejercicio{423}{Ruffini}{x^{3} + 2 x^{2} - 16 x - 32}
\ejercicio{424}{Ruffini}{x^{3} + 7 x^{2} + 7 x - 15}
\ejercicio{425}{Ruffini}{x^{3} + 3 x^{2} - 6 x - 8}
\ejercicio{426}{Ruffini}{x^{3} + x^{2} - 8 x - 12}
\ejercicio{427}{Ruffini}{x^{3} + 2 x^{2} - 15 x - 36}
\ejercicio{428}{Ruffini}{x^{3} + 7 x^{2} + 8 x - 16}
\ejercicio{429}{Ruffini}{x^{3} + 5 x^{2} - 4 x - 20}
\ejercicio{430}{Ruffini}{x^{3} + x^{2} - 9 x - 9}
\ejercicio{431}{Ruffini}{x^{3} + 2 x^{2} - 16 x - 32}
\ejercicio{432}{Ruffini}{x^{3} + 4 x^{2} + x - 6}
\ejercicio{433}{Ruffini}{x^{3} + 5 x^{2} - 2 x - 24}
\ejercicio{434}{Ruffini}{x^{3} + 6 x^{2} - 7 x - 60}
\ejercicio{435}{Ruffini}{x^{3} - x^{2} - 10 x - 8}
\ejercicio{436}{Ruffini}{x^{3} + 4 x^{2} + x - 6}
\ejercicio{437}{Ruffini}{x^{3} + 5 x^{2} - 2 x - 24}
\ejercicio{438}{Ruffini}{x^{3} + 3 x^{2} - 10 x - 24}
\ejercicio{439}{Ruffini}{x^{3} + 4 x^{2} - 17 x - 60}
\ejercicio{440}{Ruffini}{x^{3} + 4 x^{2} - x - 4}
\end{enumerate}
\nivel{3}{intermedio}
\begin{enumerate}
\ejercicio{441}{Ruffini}{3 x^{3} - 9 x^{2} - 12 x + 36}
\ejercicio{442}{Ruffini}{4 x^{3} - 16 x^{2} - 36 x + 144}
\ejercicio{443}{Ruffini}{2 x^{3} - 4 x^{2} - 32 x + 64}
\ejercicio{444}{Ruffini}{3 x^{3} + 3 x^{2} - 51 x + 45}
\ejercicio{445}{Ruffini}{4 x^{3} - 20 x^{2} + 8 x + 32}
\ejercicio{446}{Ruffini}{2 x^{3} - 6 x^{2} - 8 x + 24}
\ejercicio{447}{Ruffini}{3 x^{3} - 12 x^{2} - 27 x + 108}
\ejercicio{448}{Ruffini}{4 x^{3} - 4 x^{2} - 64 x + 64}
\ejercicio{449}{Ruffini}{2 x^{3} + 2 x^{2} - 32 x + 40}
\ejercicio{450}{Ruffini}{3 x^{3} - 15 x^{2} + 9 x + 27}
\ejercicio{451}{Ruffini}{4 x^{3} - 24 x^{2} + 128}
\ejercicio{452}{Ruffini}{2 x^{3} - 14 x + 12}
\ejercicio{453}{Ruffini}{3 x^{3} - 3 x^{2} - 42 x + 72}
\ejercicio{454}{Ruffini}{4 x^{3} - 8 x^{2} - 92 x + 240}
\ejercicio{455}{Ruffini}{2 x^{3} - 10 x^{2} + 4 x + 16}
\ejercicio{456}{Ruffini}{3 x^{3} - 6 x^{2} - 15 x + 18}
\ejercicio{457}{Ruffini}{4 x^{3} - 12 x^{2} - 40 x + 96}
\ejercicio{458}{Ruffini}{2 x^{3} - 2 x^{2} - 28 x + 48}
\ejercicio{459}{Ruffini}{3 x^{3} - 6 x^{2} - 69 x + 180}
\ejercicio{460}{Ruffini}{4 x^{3} - 16 x^{2} - 4 x + 16}
\end{enumerate}
\nivel{4}{avanzado}
\begin{enumerate}
\ejercicio{461}{Ruffini}{3 x^{4} - 39 x^{2} + 108}
\ejercicio{462}{Ruffini}{4 x^{4} - 8 x^{3} - 68 x^{2} + 72 x + 288}
\ejercicio{463}{Ruffini}{5 x^{4} + 5 x^{3} - 110 x^{2} - 80 x + 480}
\ejercicio{464}{Ruffini}{2 x^{4} - 10 x^{3} - 14 x^{2} + 58 x + 60}
\ejercicio{465}{Ruffini}{3 x^{4} - 45 x^{2} - 30 x + 72}
\ejercicio{466}{Ruffini}{4 x^{4} + 4 x^{3} - 40 x^{2} - 16 x + 96}
\ejercicio{467}{Ruffini}{5 x^{4} + 5 x^{3} - 105 x^{2} - 45 x + 540}
\ejercicio{468}{Ruffini}{2 x^{4} - 10 x^{3} - 12 x^{2} + 64 x + 64}
\ejercicio{469}{Ruffini}{3 x^{4} - 6 x^{3} - 57 x^{2} + 24 x + 180}
\ejercicio{470}{Ruffini}{4 x^{4} + 4 x^{3} - 44 x^{2} - 36 x + 72}
\ejercicio{471}{Ruffini}{5 x^{4} + 5 x^{3} - 110 x^{2} - 80 x + 480}
\ejercicio{472}{Ruffini}{2 x^{4} - 4 x^{3} - 14 x^{2} + 16 x + 24}
\ejercicio{473}{Ruffini}{3 x^{4} - 6 x^{3} - 51 x^{2} + 54 x + 216}
\ejercicio{474}{Ruffini}{4 x^{4} - 16 x^{3} - 76 x^{2} + 184 x + 480}
\ejercicio{475}{Ruffini}{5 x^{4} + 20 x^{3} - 35 x^{2} - 110 x + 120}
\ejercicio{476}{Ruffini}{2 x^{4} - 4 x^{3} - 14 x^{2} + 16 x + 24}
\ejercicio{477}{Ruffini}{3 x^{4} - 6 x^{3} - 51 x^{2} + 54 x + 216}
\ejercicio{478}{Ruffini}{4 x^{4} - 4 x^{3} - 64 x^{2} + 16 x + 192}
\ejercicio{479}{Ruffini}{5 x^{4} - 5 x^{3} - 145 x^{2} + 45 x + 900}
\ejercicio{480}{Ruffini}{2 x^{4} - 4 x^{3} - 18 x^{2} + 4 x + 16}
\end{enumerate}
\nivel{5}{imposible}
\begin{enumerate}
\ejercicio{481}{Ruffini}{3 x^{5} - 9 x^{4} - 39 x^{3} + 117 x^{2} + 108 x - 324}
\ejercicio{482}{Ruffini}{4 x^{5} - 8 x^{4} - 100 x^{3} + 200 x^{2} + 576 x - 1152}
\ejercicio{483}{Ruffini}{5 x^{5} - 30 x^{4} - 85 x^{3} + 630 x^{2} + 80 x - 2400}
\ejercicio{484}{Ruffini}{6 x^{5} - 6 x^{4} - 198 x^{3} + 222 x^{2} + 1200 x - 1800}
\ejercicio{485}{Ruffini}{2 x^{5} - 14 x^{4} - 22 x^{3} + 246 x^{2} - 180 x - 432}
\ejercicio{486}{Ruffini}{3 x^{5} - 9 x^{4} - 30 x^{3} + 60 x^{2} + 72 x - 96}
\ejercicio{487}{Ruffini}{4 x^{5} - 16 x^{4} - 80 x^{3} + 264 x^{2} + 396 x - 1080}
\ejercicio{488}{Ruffini}{5 x^{5} + 5 x^{4} - 120 x^{3} + 20 x^{2} + 800 x - 960}
\ejercicio{489}{Ruffini}{6 x^{5} - 18 x^{4} - 162 x^{3} + 570 x^{2} + 468 x - 2160}
\ejercicio{490}{Ruffini}{2 x^{5} - 14 x^{4} - 6 x^{3} + 158 x^{2} - 92 x - 240}
\ejercicio{491}{Ruffini}{3 x^{5} - 24 x^{4} - 39 x^{3} + 528 x^{2} - 108 x - 2160}
\ejercicio{492}{Ruffini}{4 x^{5} - 44 x^{3} + 24 x^{2} + 112 x - 96}
\ejercicio{493}{Ruffini}{5 x^{5} - 5 x^{4} - 115 x^{3} + 165 x^{2} + 630 x - 1080}
\ejercicio{494}{Ruffini}{6 x^{5} - 210 x^{3} + 180 x^{2} + 1824 x - 2880}
\ejercicio{495}{Ruffini}{2 x^{5} - 18 x^{4} + 26 x^{3} + 114 x^{2} - 172 x - 240}
\ejercicio{496}{Ruffini}{3 x^{5} - 12 x^{4} - 45 x^{3} + 138 x^{2} + 132 x - 360}
\ejercicio{497}{Ruffini}{4 x^{5} - 20 x^{4} - 108 x^{3} + 468 x^{2} + 648 x - 2592}
\ejercicio{498}{Ruffini}{5 x^{5} - 5 x^{4} - 100 x^{3} + 100 x^{2} + 320 x - 320}
\ejercicio{499}{Ruffini}{6 x^{5} - 12 x^{4} - 204 x^{3} + 408 x^{2} + 1350 x - 2700}
\ejercicio{500}{Ruffini}{2 x^{5} - 4 x^{4} - 30 x^{3} + 80 x^{2} + 8 x - 96}
\end{enumerate}
\newpage
\part*{Solucionario desarrollado paso a paso}\addcontentsline{toc}{part}{Solucionario desarrollado paso a paso}
\section*{Nota para estudiar las soluciones}\addcontentsline{toc}{section}{Nota para estudiar las soluciones}
Las soluciones estan explicadas con pasos pequenos. La idea no es memorizar el resultado, sino aprender a reconocer la forma del ejercicio, elegir el metodo y justificar cada transformacion.
\section{Soluciones: Factor comun normal y por agrupacion}
\nivel{1}{basico}
\subsubsection*{Ejercicio 1 - Factor comun normal}
\textbf{Enunciado:} $ 4 x^{3} y + 6 x^{2} y $\par
\begin{enumerate}[leftmargin=2em]
\item Miramos todos los terminos y buscamos que parte se repite en cada uno.
\item El factor comun es \(2 x^{2} y\). Eso significa que todos los terminos se pueden dividir entre \(2 x^{2} y\).
\item Al dividir el polinomio entre el factor comun queda \(2 x + 3\).
\item Sacamos el factor comun fuera del parentesis. Resultado: \[\boxed{2 x^{2} y\left(2 x + 3\right)}\]
\end{enumerate}

\vspace{0.2cm}
\subsubsection*{Ejercicio 2 - Factor comun normal}
\textbf{Enunciado:} $ 9 x^{4} + 12 x^{3} $\par
\begin{enumerate}[leftmargin=2em]
\item Miramos todos los terminos y buscamos que parte se repite en cada uno.
\item El factor comun es \(3 x^{3}\). Eso significa que todos los terminos se pueden dividir entre \(3 x^{3}\).
\item Al dividir el polinomio entre el factor comun queda \(3 x + 4\).
\item Sacamos el factor comun fuera del parentesis. Resultado: \[\boxed{3 x^{3}\left(3 x + 4\right)}\]
\end{enumerate}

\vspace{0.2cm}
\subsubsection*{Ejercicio 3 - Factor comun normal}
\textbf{Enunciado:} $ 16 x^{2} y + 20 x y $\par
\begin{enumerate}[leftmargin=2em]
\item Miramos todos los terminos y buscamos que parte se repite en cada uno.
\item El factor comun es \(4 x y\). Eso significa que todos los terminos se pueden dividir entre \(4 x y\).
\item Al dividir el polinomio entre el factor comun queda \(4 x + 5\).
\item Sacamos el factor comun fuera del parentesis. Resultado: \[\boxed{4 x y\left(4 x + 5\right)}\]
\end{enumerate}

\vspace{0.2cm}
\subsubsection*{Ejercicio 4 - Factor comun normal}
\textbf{Enunciado:} $ 25 x^{3} + 30 x^{2} $\par
\begin{enumerate}[leftmargin=2em]
\item Miramos todos los terminos y buscamos que parte se repite en cada uno.
\item El factor comun es \(5 x^{2}\). Eso significa que todos los terminos se pueden dividir entre \(5 x^{2}\).
\item Al dividir el polinomio entre el factor comun queda \(5 x + 6\).
\item Sacamos el factor comun fuera del parentesis. Resultado: \[\boxed{5 x^{2}\left(5 x + 6\right)}\]
\end{enumerate}

\vspace{0.2cm}
\subsubsection*{Ejercicio 5 - Factor comun normal}
\textbf{Enunciado:} $ 36 x^{4} y + 42 x^{3} y $\par
\begin{enumerate}[leftmargin=2em]
\item Miramos todos los terminos y buscamos que parte se repite en cada uno.
\item El factor comun es \(6 x^{3} y\). Eso significa que todos los terminos se pueden dividir entre \(6 x^{3} y\).
\item Al dividir el polinomio entre el factor comun queda \(6 x + 7\).
\item Sacamos el factor comun fuera del parentesis. Resultado: \[\boxed{6 x^{3} y\left(6 x + 7\right)}\]
\end{enumerate}

\vspace{0.2cm}
\subsubsection*{Ejercicio 6 - Factor comun normal}
\textbf{Enunciado:} $ 49 x^{2} + 56 x $\par
\begin{enumerate}[leftmargin=2em]
\item Miramos todos los terminos y buscamos que parte se repite en cada uno.
\item El factor comun es \(7 x\). Eso significa que todos los terminos se pueden dividir entre \(7 x\).
\item Al dividir el polinomio entre el factor comun queda \(7 x + 8\).
\item Sacamos el factor comun fuera del parentesis. Resultado: \[\boxed{7 x\left(7 x + 8\right)}\]
\end{enumerate}

\vspace{0.2cm}
\subsubsection*{Ejercicio 7 - Factor comun normal}
\textbf{Enunciado:} $ 64 x^{3} y + 72 x^{2} y $\par
\begin{enumerate}[leftmargin=2em]
\item Miramos todos los terminos y buscamos que parte se repite en cada uno.
\item El factor comun es \(8 x^{2} y\). Eso significa que todos los terminos se pueden dividir entre \(8 x^{2} y\).
\item Al dividir el polinomio entre el factor comun queda \(8 x + 9\).
\item Sacamos el factor comun fuera del parentesis. Resultado: \[\boxed{8 x^{2} y\left(8 x + 9\right)}\]
\end{enumerate}

\vspace{0.2cm}
\subsubsection*{Ejercicio 8 - Factor comun normal}
\textbf{Enunciado:} $ 81 x^{4} + 90 x^{3} $\par
\begin{enumerate}[leftmargin=2em]
\item Miramos todos los terminos y buscamos que parte se repite en cada uno.
\item El factor comun es \(9 x^{3}\). Eso significa que todos los terminos se pueden dividir entre \(9 x^{3}\).
\item Al dividir el polinomio entre el factor comun queda \(9 x + 10\).
\item Sacamos el factor comun fuera del parentesis. Resultado: \[\boxed{9 x^{3}\left(9 x + 10\right)}\]
\end{enumerate}

\vspace{0.2cm}
\subsubsection*{Ejercicio 9 - Factor comun normal}
\textbf{Enunciado:} $ 100 x^{2} y + 110 x y $\par
\begin{enumerate}[leftmargin=2em]
\item Miramos todos los terminos y buscamos que parte se repite en cada uno.
\item El factor comun es \(10 x y\). Eso significa que todos los terminos se pueden dividir entre \(10 x y\).
\item Al dividir el polinomio entre el factor comun queda \(10 x + 11\).
\item Sacamos el factor comun fuera del parentesis. Resultado: \[\boxed{10 x y\left(10 x + 11\right)}\]
\end{enumerate}

\vspace{0.2cm}
\subsubsection*{Ejercicio 10 - Factor comun normal}
\textbf{Enunciado:} $ 121 x^{3} + 132 x^{2} $\par
\begin{enumerate}[leftmargin=2em]
\item Miramos todos los terminos y buscamos que parte se repite en cada uno.
\item El factor comun es \(11 x^{2}\). Eso significa que todos los terminos se pueden dividir entre \(11 x^{2}\).
\item Al dividir el polinomio entre el factor comun queda \(11 x + 12\).
\item Sacamos el factor comun fuera del parentesis. Resultado: \[\boxed{11 x^{2}\left(11 x + 12\right)}\]
\end{enumerate}

\vspace{0.2cm}
\subsubsection*{Ejercicio 11 - Factor comun por agrupacion}
\textbf{Enunciado:} $ 3 x y + 3 x + 6 y + 6 $\par
\begin{enumerate}[leftmargin=2em]
\item En agrupacion no siempre hay un unico factor visible al inicio. Por eso intentamos formar grupos.
\item Los grupos se pueden ordenar para que aparezca el mismo parentesis \((3 y + 3)\).
\item Cuando ese parentesis se repite, lo sacamos como factor comun. Lo que queda afuera forma el otro factor \((x + 2)\).
\item Resultado: \[\boxed{\left(x + 2\right)\left(3 y + 3\right)}\]
\end{enumerate}

\vspace{0.2cm}
\subsubsection*{Ejercicio 12 - Factor comun por agrupacion}
\textbf{Enunciado:} $ 4 x y + 5 x + 12 y + 15 $\par
\begin{enumerate}[leftmargin=2em]
\item En agrupacion no siempre hay un unico factor visible al inicio. Por eso intentamos formar grupos.
\item Los grupos se pueden ordenar para que aparezca el mismo parentesis \((4 y + 5)\).
\item Cuando ese parentesis se repite, lo sacamos como factor comun. Lo que queda afuera forma el otro factor \((x + 3)\).
\item Resultado: \[\boxed{\left(x + 3\right)\left(4 y + 5\right)}\]
\end{enumerate}

\vspace{0.2cm}
\subsubsection*{Ejercicio 13 - Factor comun por agrupacion}
\textbf{Enunciado:} $ 5 x y + 7 x + 20 y + 28 $\par
\begin{enumerate}[leftmargin=2em]
\item En agrupacion no siempre hay un unico factor visible al inicio. Por eso intentamos formar grupos.
\item Los grupos se pueden ordenar para que aparezca el mismo parentesis \((5 y + 7)\).
\item Cuando ese parentesis se repite, lo sacamos como factor comun. Lo que queda afuera forma el otro factor \((x + 4)\).
\item Resultado: \[\boxed{\left(x + 4\right)\left(5 y + 7\right)}\]
\end{enumerate}

\vspace{0.2cm}
\subsubsection*{Ejercicio 14 - Factor comun por agrupacion}
\textbf{Enunciado:} $ 6 x y + 9 x + 30 y + 45 $\par
\begin{enumerate}[leftmargin=2em]
\item En agrupacion no siempre hay un unico factor visible al inicio. Por eso intentamos formar grupos.
\item Los grupos se pueden ordenar para que aparezca el mismo parentesis \((6 y + 9)\).
\item Cuando ese parentesis se repite, lo sacamos como factor comun. Lo que queda afuera forma el otro factor \((x + 5)\).
\item Resultado: \[\boxed{\left(x + 5\right)\left(6 y + 9\right)}\]
\end{enumerate}

\vspace{0.2cm}
\subsubsection*{Ejercicio 15 - Factor comun por agrupacion}
\textbf{Enunciado:} $ 7 x y + 11 x + 42 y + 66 $\par
\begin{enumerate}[leftmargin=2em]
\item En agrupacion no siempre hay un unico factor visible al inicio. Por eso intentamos formar grupos.
\item Los grupos se pueden ordenar para que aparezca el mismo parentesis \((7 y + 11)\).
\item Cuando ese parentesis se repite, lo sacamos como factor comun. Lo que queda afuera forma el otro factor \((x + 6)\).
\item Resultado: \[\boxed{\left(x + 6\right)\left(7 y + 11\right)}\]
\end{enumerate}

\vspace{0.2cm}
\subsubsection*{Ejercicio 16 - Factor comun por agrupacion}
\textbf{Enunciado:} $ 8 x y + 13 x + 56 y + 91 $\par
\begin{enumerate}[leftmargin=2em]
\item En agrupacion no siempre hay un unico factor visible al inicio. Por eso intentamos formar grupos.
\item Los grupos se pueden ordenar para que aparezca el mismo parentesis \((8 y + 13)\).
\item Cuando ese parentesis se repite, lo sacamos como factor comun. Lo que queda afuera forma el otro factor \((x + 7)\).
\item Resultado: \[\boxed{\left(x + 7\right)\left(8 y + 13\right)}\]
\end{enumerate}

\vspace{0.2cm}
\subsubsection*{Ejercicio 17 - Factor comun por agrupacion}
\textbf{Enunciado:} $ 9 x y + 15 x + 72 y + 120 $\par
\begin{enumerate}[leftmargin=2em]
\item En agrupacion no siempre hay un unico factor visible al inicio. Por eso intentamos formar grupos.
\item Los grupos se pueden ordenar para que aparezca el mismo parentesis \((9 y + 15)\).
\item Cuando ese parentesis se repite, lo sacamos como factor comun. Lo que queda afuera forma el otro factor \((x + 8)\).
\item Resultado: \[\boxed{\left(x + 8\right)\left(9 y + 15\right)}\]
\end{enumerate}

\vspace{0.2cm}
\subsubsection*{Ejercicio 18 - Factor comun por agrupacion}
\textbf{Enunciado:} $ 10 x y + 17 x + 90 y + 153 $\par
\begin{enumerate}[leftmargin=2em]
\item En agrupacion no siempre hay un unico factor visible al inicio. Por eso intentamos formar grupos.
\item Los grupos se pueden ordenar para que aparezca el mismo parentesis \((10 y + 17)\).
\item Cuando ese parentesis se repite, lo sacamos como factor comun. Lo que queda afuera forma el otro factor \((x + 9)\).
\item Resultado: \[\boxed{\left(x + 9\right)\left(10 y + 17\right)}\]
\end{enumerate}

\vspace{0.2cm}
\subsubsection*{Ejercicio 19 - Factor comun por agrupacion}
\textbf{Enunciado:} $ 11 x y + 19 x + 110 y + 190 $\par
\begin{enumerate}[leftmargin=2em]
\item En agrupacion no siempre hay un unico factor visible al inicio. Por eso intentamos formar grupos.
\item Los grupos se pueden ordenar para que aparezca el mismo parentesis \((11 y + 19)\).
\item Cuando ese parentesis se repite, lo sacamos como factor comun. Lo que queda afuera forma el otro factor \((x + 10)\).
\item Resultado: \[\boxed{\left(x + 10\right)\left(11 y + 19\right)}\]
\end{enumerate}

\vspace{0.2cm}
\subsubsection*{Ejercicio 20 - Factor comun por agrupacion}
\textbf{Enunciado:} $ 12 x y + 21 x + 132 y + 231 $\par
\begin{enumerate}[leftmargin=2em]
\item En agrupacion no siempre hay un unico factor visible al inicio. Por eso intentamos formar grupos.
\item Los grupos se pueden ordenar para que aparezca el mismo parentesis \((12 y + 21)\).
\item Cuando ese parentesis se repite, lo sacamos como factor comun. Lo que queda afuera forma el otro factor \((x + 11)\).
\item Resultado: \[\boxed{\left(x + 11\right)\left(12 y + 21\right)}\]
\end{enumerate}

\vspace{0.2cm}
\nivel{2}{intermedio inicial}
\subsubsection*{Ejercicio 21 - Factor comun normal}
\textbf{Enunciado:} $ 9 x^{4} y + 6 x^{3} y^{2} - 12 x^{3} y $\par
\begin{enumerate}[leftmargin=2em]
\item Miramos todos los terminos y buscamos que parte se repite en cada uno.
\item El factor comun es \(3 x^{3} y\). Eso significa que todos los terminos se pueden dividir entre \(3 x^{3} y\).
\item Al dividir el polinomio entre el factor comun queda \(3 x + 2 y - 4\).
\item Sacamos el factor comun fuera del parentesis. Resultado: \[\boxed{3 x^{3} y\left(3 x + 2 y - 4\right)}\]
\end{enumerate}

\vspace{0.2cm}
\subsubsection*{Ejercicio 22 - Factor comun normal}
\textbf{Enunciado:} $ 16 x^{5} + 12 x^{4} y - 20 x^{4} $\par
\begin{enumerate}[leftmargin=2em]
\item Miramos todos los terminos y buscamos que parte se repite en cada uno.
\item El factor comun es \(4 x^{4}\). Eso significa que todos los terminos se pueden dividir entre \(4 x^{4}\).
\item Al dividir el polinomio entre el factor comun queda \(4 x + 3 y - 5\).
\item Sacamos el factor comun fuera del parentesis. Resultado: \[\boxed{4 x^{4}\left(4 x + 3 y - 5\right)}\]
\end{enumerate}

\vspace{0.2cm}
\subsubsection*{Ejercicio 23 - Factor comun normal}
\textbf{Enunciado:} $ 25 x^{3} y + 20 x^{2} y^{2} - 30 x^{2} y $\par
\begin{enumerate}[leftmargin=2em]
\item Miramos todos los terminos y buscamos que parte se repite en cada uno.
\item El factor comun es \(5 x^{2} y\). Eso significa que todos los terminos se pueden dividir entre \(5 x^{2} y\).
\item Al dividir el polinomio entre el factor comun queda \(5 x + 4 y - 6\).
\item Sacamos el factor comun fuera del parentesis. Resultado: \[\boxed{5 x^{2} y\left(5 x + 4 y - 6\right)}\]
\end{enumerate}

\vspace{0.2cm}
\subsubsection*{Ejercicio 24 - Factor comun normal}
\textbf{Enunciado:} $ 36 x^{4} + 30 x^{3} y - 42 x^{3} $\par
\begin{enumerate}[leftmargin=2em]
\item Miramos todos los terminos y buscamos que parte se repite en cada uno.
\item El factor comun es \(6 x^{3}\). Eso significa que todos los terminos se pueden dividir entre \(6 x^{3}\).
\item Al dividir el polinomio entre el factor comun queda \(6 x + 5 y - 7\).
\item Sacamos el factor comun fuera del parentesis. Resultado: \[\boxed{6 x^{3}\left(6 x + 5 y - 7\right)}\]
\end{enumerate}

\vspace{0.2cm}
\subsubsection*{Ejercicio 25 - Factor comun normal}
\textbf{Enunciado:} $ 49 x^{5} y + 42 x^{4} y^{2} - 56 x^{4} y $\par
\begin{enumerate}[leftmargin=2em]
\item Miramos todos los terminos y buscamos que parte se repite en cada uno.
\item El factor comun es \(7 x^{4} y\). Eso significa que todos los terminos se pueden dividir entre \(7 x^{4} y\).
\item Al dividir el polinomio entre el factor comun queda \(7 x + 6 y - 8\).
\item Sacamos el factor comun fuera del parentesis. Resultado: \[\boxed{7 x^{4} y\left(7 x + 6 y - 8\right)}\]
\end{enumerate}

\vspace{0.2cm}
\subsubsection*{Ejercicio 26 - Factor comun normal}
\textbf{Enunciado:} $ 64 x^{3} + 56 x^{2} y - 72 x^{2} $\par
\begin{enumerate}[leftmargin=2em]
\item Miramos todos los terminos y buscamos que parte se repite en cada uno.
\item El factor comun es \(8 x^{2}\). Eso significa que todos los terminos se pueden dividir entre \(8 x^{2}\).
\item Al dividir el polinomio entre el factor comun queda \(8 x + 7 y - 9\).
\item Sacamos el factor comun fuera del parentesis. Resultado: \[\boxed{8 x^{2}\left(8 x + 7 y - 9\right)}\]
\end{enumerate}

\vspace{0.2cm}
\subsubsection*{Ejercicio 27 - Factor comun normal}
\textbf{Enunciado:} $ 81 x^{4} y + 72 x^{3} y^{2} - 90 x^{3} y $\par
\begin{enumerate}[leftmargin=2em]
\item Miramos todos los terminos y buscamos que parte se repite en cada uno.
\item El factor comun es \(9 x^{3} y\). Eso significa que todos los terminos se pueden dividir entre \(9 x^{3} y\).
\item Al dividir el polinomio entre el factor comun queda \(9 x + 8 y - 10\).
\item Sacamos el factor comun fuera del parentesis. Resultado: \[\boxed{9 x^{3} y\left(9 x + 8 y - 10\right)}\]
\end{enumerate}

\vspace{0.2cm}
\subsubsection*{Ejercicio 28 - Factor comun normal}
\textbf{Enunciado:} $ 100 x^{5} + 90 x^{4} y - 110 x^{4} $\par
\begin{enumerate}[leftmargin=2em]
\item Miramos todos los terminos y buscamos que parte se repite en cada uno.
\item El factor comun es \(10 x^{4}\). Eso significa que todos los terminos se pueden dividir entre \(10 x^{4}\).
\item Al dividir el polinomio entre el factor comun queda \(10 x + 9 y - 11\).
\item Sacamos el factor comun fuera del parentesis. Resultado: \[\boxed{10 x^{4}\left(10 x + 9 y - 11\right)}\]
\end{enumerate}

\vspace{0.2cm}
\subsubsection*{Ejercicio 29 - Factor comun normal}
\textbf{Enunciado:} $ 121 x^{3} y + 110 x^{2} y^{2} - 132 x^{2} y $\par
\begin{enumerate}[leftmargin=2em]
\item Miramos todos los terminos y buscamos que parte se repite en cada uno.
\item El factor comun es \(11 x^{2} y\). Eso significa que todos los terminos se pueden dividir entre \(11 x^{2} y\).
\item Al dividir el polinomio entre el factor comun queda \(11 x + 10 y - 12\).
\item Sacamos el factor comun fuera del parentesis. Resultado: \[\boxed{11 x^{2} y\left(11 x + 10 y - 12\right)}\]
\end{enumerate}

\vspace{0.2cm}
\subsubsection*{Ejercicio 30 - Factor comun normal}
\textbf{Enunciado:} $ 144 x^{4} + 132 x^{3} y - 156 x^{3} $\par
\begin{enumerate}[leftmargin=2em]
\item Miramos todos los terminos y buscamos que parte se repite en cada uno.
\item El factor comun es \(12 x^{3}\). Eso significa que todos los terminos se pueden dividir entre \(12 x^{3}\).
\item Al dividir el polinomio entre el factor comun queda \(12 x + 11 y - 13\).
\item Sacamos el factor comun fuera del parentesis. Resultado: \[\boxed{12 x^{3}\left(12 x + 11 y - 13\right)}\]
\end{enumerate}

\vspace{0.2cm}
\subsubsection*{Ejercicio 31 - Factor comun por agrupacion}
\textbf{Enunciado:} $ 9 x y + 12 x - 6 y - 8 $\par
\begin{enumerate}[leftmargin=2em]
\item En agrupacion no siempre hay un unico factor visible al inicio. Por eso intentamos formar grupos.
\item Los grupos se pueden ordenar para que aparezca el mismo parentesis \((3 y + 4)\).
\item Cuando ese parentesis se repite, lo sacamos como factor comun. Lo que queda afuera forma el otro factor \((3 x - 2)\).
\item Resultado: \[\boxed{\left(3 x - 2\right)\left(3 y + 4\right)}\]
\end{enumerate}

\vspace{0.2cm}
\subsubsection*{Ejercicio 32 - Factor comun por agrupacion}
\textbf{Enunciado:} $ 16 x y + 24 x - 12 y - 18 $\par
\begin{enumerate}[leftmargin=2em]
\item En agrupacion no siempre hay un unico factor visible al inicio. Por eso intentamos formar grupos.
\item Los grupos se pueden ordenar para que aparezca el mismo parentesis \((4 y + 6)\).
\item Cuando ese parentesis se repite, lo sacamos como factor comun. Lo que queda afuera forma el otro factor \((4 x - 3)\).
\item Resultado: \[\boxed{\left(4 x - 3\right)\left(4 y + 6\right)}\]
\end{enumerate}

\vspace{0.2cm}
\subsubsection*{Ejercicio 33 - Factor comun por agrupacion}
\textbf{Enunciado:} $ 25 x y + 40 x - 20 y - 32 $\par
\begin{enumerate}[leftmargin=2em]
\item En agrupacion no siempre hay un unico factor visible al inicio. Por eso intentamos formar grupos.
\item Los grupos se pueden ordenar para que aparezca el mismo parentesis \((5 y + 8)\).
\item Cuando ese parentesis se repite, lo sacamos como factor comun. Lo que queda afuera forma el otro factor \((5 x - 4)\).
\item Resultado: \[\boxed{\left(5 x - 4\right)\left(5 y + 8\right)}\]
\end{enumerate}

\vspace{0.2cm}
\subsubsection*{Ejercicio 34 - Factor comun por agrupacion}
\textbf{Enunciado:} $ 36 x y + 60 x - 30 y - 50 $\par
\begin{enumerate}[leftmargin=2em]
\item En agrupacion no siempre hay un unico factor visible al inicio. Por eso intentamos formar grupos.
\item Los grupos se pueden ordenar para que aparezca el mismo parentesis \((6 y + 10)\).
\item Cuando ese parentesis se repite, lo sacamos como factor comun. Lo que queda afuera forma el otro factor \((6 x - 5)\).
\item Resultado: \[\boxed{\left(6 x - 5\right)\left(6 y + 10\right)}\]
\end{enumerate}

\vspace{0.2cm}
\subsubsection*{Ejercicio 35 - Factor comun por agrupacion}
\textbf{Enunciado:} $ 49 x y + 84 x - 42 y - 72 $\par
\begin{enumerate}[leftmargin=2em]
\item En agrupacion no siempre hay un unico factor visible al inicio. Por eso intentamos formar grupos.
\item Los grupos se pueden ordenar para que aparezca el mismo parentesis \((7 y + 12)\).
\item Cuando ese parentesis se repite, lo sacamos como factor comun. Lo que queda afuera forma el otro factor \((7 x - 6)\).
\item Resultado: \[\boxed{\left(7 x - 6\right)\left(7 y + 12\right)}\]
\end{enumerate}

\vspace{0.2cm}
\subsubsection*{Ejercicio 36 - Factor comun por agrupacion}
\textbf{Enunciado:} $ 64 x y + 112 x - 56 y - 98 $\par
\begin{enumerate}[leftmargin=2em]
\item En agrupacion no siempre hay un unico factor visible al inicio. Por eso intentamos formar grupos.
\item Los grupos se pueden ordenar para que aparezca el mismo parentesis \((8 y + 14)\).
\item Cuando ese parentesis se repite, lo sacamos como factor comun. Lo que queda afuera forma el otro factor \((8 x - 7)\).
\item Resultado: \[\boxed{\left(8 x - 7\right)\left(8 y + 14\right)}\]
\end{enumerate}

\vspace{0.2cm}
\subsubsection*{Ejercicio 37 - Factor comun por agrupacion}
\textbf{Enunciado:} $ 81 x y + 144 x - 72 y - 128 $\par
\begin{enumerate}[leftmargin=2em]
\item En agrupacion no siempre hay un unico factor visible al inicio. Por eso intentamos formar grupos.
\item Los grupos se pueden ordenar para que aparezca el mismo parentesis \((9 y + 16)\).
\item Cuando ese parentesis se repite, lo sacamos como factor comun. Lo que queda afuera forma el otro factor \((9 x - 8)\).
\item Resultado: \[\boxed{\left(9 x - 8\right)\left(9 y + 16\right)}\]
\end{enumerate}

\vspace{0.2cm}
\subsubsection*{Ejercicio 38 - Factor comun por agrupacion}
\textbf{Enunciado:} $ 100 x y + 180 x - 90 y - 162 $\par
\begin{enumerate}[leftmargin=2em]
\item En agrupacion no siempre hay un unico factor visible al inicio. Por eso intentamos formar grupos.
\item Los grupos se pueden ordenar para que aparezca el mismo parentesis \((10 y + 18)\).
\item Cuando ese parentesis se repite, lo sacamos como factor comun. Lo que queda afuera forma el otro factor \((10 x - 9)\).
\item Resultado: \[\boxed{\left(10 x - 9\right)\left(10 y + 18\right)}\]
\end{enumerate}

\vspace{0.2cm}
\subsubsection*{Ejercicio 39 - Factor comun por agrupacion}
\textbf{Enunciado:} $ 121 x y + 220 x - 110 y - 200 $\par
\begin{enumerate}[leftmargin=2em]
\item En agrupacion no siempre hay un unico factor visible al inicio. Por eso intentamos formar grupos.
\item Los grupos se pueden ordenar para que aparezca el mismo parentesis \((11 y + 20)\).
\item Cuando ese parentesis se repite, lo sacamos como factor comun. Lo que queda afuera forma el otro factor \((11 x - 10)\).
\item Resultado: \[\boxed{\left(11 x - 10\right)\left(11 y + 20\right)}\]
\end{enumerate}

\vspace{0.2cm}
\subsubsection*{Ejercicio 40 - Factor comun por agrupacion}
\textbf{Enunciado:} $ 144 x y + 264 x - 132 y - 242 $\par
\begin{enumerate}[leftmargin=2em]
\item En agrupacion no siempre hay un unico factor visible al inicio. Por eso intentamos formar grupos.
\item Los grupos se pueden ordenar para que aparezca el mismo parentesis \((12 y + 22)\).
\item Cuando ese parentesis se repite, lo sacamos como factor comun. Lo que queda afuera forma el otro factor \((12 x - 11)\).
\item Resultado: \[\boxed{\left(12 x - 11\right)\left(12 y + 22\right)}\]
\end{enumerate}

\vspace{0.2cm}
\nivel{3}{intermedio}
\subsubsection*{Ejercicio 41 - Factor comun normal}
\textbf{Enunciado:} $ 8 x^{6} y^{2} - 16 x^{4} y^{3} + 12 x^{4} y^{2} $\par
\begin{enumerate}[leftmargin=2em]
\item Miramos todos los terminos y buscamos que parte se repite en cada uno.
\item El factor comun es \(4 x^{4} y^{2}\). Eso significa que todos los terminos se pueden dividir entre \(4 x^{4} y^{2}\).
\item Al dividir el polinomio entre el factor comun queda \(2 x^{2} - 4 y + 3\).
\item Sacamos el factor comun fuera del parentesis. Resultado: \[\boxed{4 x^{4} y^{2}\left(2 x^{2} - 4 y + 3\right)}\]
\end{enumerate}

\vspace{0.2cm}
\subsubsection*{Ejercicio 42 - Factor comun normal}
\textbf{Enunciado:} $ 15 x^{7} y - 25 x^{5} y^{2} + 25 x^{5} y $\par
\begin{enumerate}[leftmargin=2em]
\item Miramos todos los terminos y buscamos que parte se repite en cada uno.
\item El factor comun es \(5 x^{5} y\). Eso significa que todos los terminos se pueden dividir entre \(5 x^{5} y\).
\item Al dividir el polinomio entre el factor comun queda \(3 x^{2} - 5 y + 5\).
\item Sacamos el factor comun fuera del parentesis. Resultado: \[\boxed{5 x^{5} y\left(3 x^{2} - 5 y + 5\right)}\]
\end{enumerate}

\vspace{0.2cm}
\subsubsection*{Ejercicio 43 - Factor comun normal}
\textbf{Enunciado:} $ 24 x^{5} y^{2} - 36 x^{3} y^{3} + 42 x^{3} y^{2} $\par
\begin{enumerate}[leftmargin=2em]
\item Miramos todos los terminos y buscamos que parte se repite en cada uno.
\item El factor comun es \(6 x^{3} y^{2}\). Eso significa que todos los terminos se pueden dividir entre \(6 x^{3} y^{2}\).
\item Al dividir el polinomio entre el factor comun queda \(4 x^{2} - 6 y + 7\).
\item Sacamos el factor comun fuera del parentesis. Resultado: \[\boxed{6 x^{3} y^{2}\left(4 x^{2} - 6 y + 7\right)}\]
\end{enumerate}

\vspace{0.2cm}
\subsubsection*{Ejercicio 44 - Factor comun normal}
\textbf{Enunciado:} $ 35 x^{6} y - 49 x^{4} y^{2} + 63 x^{4} y $\par
\begin{enumerate}[leftmargin=2em]
\item Miramos todos los terminos y buscamos que parte se repite en cada uno.
\item El factor comun es \(7 x^{4} y\). Eso significa que todos los terminos se pueden dividir entre \(7 x^{4} y\).
\item Al dividir el polinomio entre el factor comun queda \(5 x^{2} - 7 y + 9\).
\item Sacamos el factor comun fuera del parentesis. Resultado: \[\boxed{7 x^{4} y\left(5 x^{2} - 7 y + 9\right)}\]
\end{enumerate}

\vspace{0.2cm}
\subsubsection*{Ejercicio 45 - Factor comun normal}
\textbf{Enunciado:} $ 48 x^{7} y^{2} - 64 x^{5} y^{3} + 88 x^{5} y^{2} $\par
\begin{enumerate}[leftmargin=2em]
\item Miramos todos los terminos y buscamos que parte se repite en cada uno.
\item El factor comun es \(8 x^{5} y^{2}\). Eso significa que todos los terminos se pueden dividir entre \(8 x^{5} y^{2}\).
\item Al dividir el polinomio entre el factor comun queda \(6 x^{2} - 8 y + 11\).
\item Sacamos el factor comun fuera del parentesis. Resultado: \[\boxed{8 x^{5} y^{2}\left(6 x^{2} - 8 y + 11\right)}\]
\end{enumerate}

\vspace{0.2cm}
\subsubsection*{Ejercicio 46 - Factor comun normal}
\textbf{Enunciado:} $ 63 x^{5} y - 81 x^{3} y^{2} + 117 x^{3} y $\par
\begin{enumerate}[leftmargin=2em]
\item Miramos todos los terminos y buscamos que parte se repite en cada uno.
\item El factor comun es \(9 x^{3} y\). Eso significa que todos los terminos se pueden dividir entre \(9 x^{3} y\).
\item Al dividir el polinomio entre el factor comun queda \(7 x^{2} - 9 y + 13\).
\item Sacamos el factor comun fuera del parentesis. Resultado: \[\boxed{9 x^{3} y\left(7 x^{2} - 9 y + 13\right)}\]
\end{enumerate}

\vspace{0.2cm}
\subsubsection*{Ejercicio 47 - Factor comun normal}
\textbf{Enunciado:} $ 80 x^{6} y^{2} - 100 x^{4} y^{3} + 150 x^{4} y^{2} $\par
\begin{enumerate}[leftmargin=2em]
\item Miramos todos los terminos y buscamos que parte se repite en cada uno.
\item El factor comun es \(10 x^{4} y^{2}\). Eso significa que todos los terminos se pueden dividir entre \(10 x^{4} y^{2}\).
\item Al dividir el polinomio entre el factor comun queda \(8 x^{2} - 10 y + 15\).
\item Sacamos el factor comun fuera del parentesis. Resultado: \[\boxed{10 x^{4} y^{2}\left(8 x^{2} - 10 y + 15\right)}\]
\end{enumerate}

\vspace{0.2cm}
\subsubsection*{Ejercicio 48 - Factor comun normal}
\textbf{Enunciado:} $ 99 x^{7} y - 121 x^{5} y^{2} + 187 x^{5} y $\par
\begin{enumerate}[leftmargin=2em]
\item Miramos todos los terminos y buscamos que parte se repite en cada uno.
\item El factor comun es \(11 x^{5} y\). Eso significa que todos los terminos se pueden dividir entre \(11 x^{5} y\).
\item Al dividir el polinomio entre el factor comun queda \(9 x^{2} - 11 y + 17\).
\item Sacamos el factor comun fuera del parentesis. Resultado: \[\boxed{11 x^{5} y\left(9 x^{2} - 11 y + 17\right)}\]
\end{enumerate}

\vspace{0.2cm}
\subsubsection*{Ejercicio 49 - Factor comun normal}
\textbf{Enunciado:} $ 120 x^{5} y^{2} - 144 x^{3} y^{3} + 228 x^{3} y^{2} $\par
\begin{enumerate}[leftmargin=2em]
\item Miramos todos los terminos y buscamos que parte se repite en cada uno.
\item El factor comun es \(12 x^{3} y^{2}\). Eso significa que todos los terminos se pueden dividir entre \(12 x^{3} y^{2}\).
\item Al dividir el polinomio entre el factor comun queda \(10 x^{2} - 12 y + 19\).
\item Sacamos el factor comun fuera del parentesis. Resultado: \[\boxed{12 x^{3} y^{2}\left(10 x^{2} - 12 y + 19\right)}\]
\end{enumerate}

\vspace{0.2cm}
\subsubsection*{Ejercicio 50 - Factor comun normal}
\textbf{Enunciado:} $ 143 x^{6} y - 169 x^{4} y^{2} + 273 x^{4} y $\par
\begin{enumerate}[leftmargin=2em]
\item Miramos todos los terminos y buscamos que parte se repite en cada uno.
\item El factor comun es \(13 x^{4} y\). Eso significa que todos los terminos se pueden dividir entre \(13 x^{4} y\).
\item Al dividir el polinomio entre el factor comun queda \(11 x^{2} - 13 y + 21\).
\item Sacamos el factor comun fuera del parentesis. Resultado: \[\boxed{13 x^{4} y\left(11 x^{2} - 13 y + 21\right)}\]
\end{enumerate}

\vspace{0.2cm}
\subsubsection*{Ejercicio 51 - Factor comun por agrupacion}
\textbf{Enunciado:} $ 12 x^{2} + 6 x y - 20 x - 10 y $\par
\begin{enumerate}[leftmargin=2em]
\item En agrupacion no siempre hay un unico factor visible al inicio. Por eso intentamos formar grupos.
\item Los grupos se pueden ordenar para que aparezca el mismo parentesis \((3 x - 5)\).
\item Cuando ese parentesis se repite, lo sacamos como factor comun. Lo que queda afuera forma el otro factor \((4 x + 2 y)\).
\item Resultado: \[\boxed{\left(4 x + 2 y\right)\left(3 x - 5\right)}\]
\end{enumerate}

\vspace{0.2cm}
\subsubsection*{Ejercicio 52 - Factor comun por agrupacion}
\textbf{Enunciado:} $ 20 x^{2} + 12 x y - 35 x - 21 y $\par
\begin{enumerate}[leftmargin=2em]
\item En agrupacion no siempre hay un unico factor visible al inicio. Por eso intentamos formar grupos.
\item Los grupos se pueden ordenar para que aparezca el mismo parentesis \((4 x - 7)\).
\item Cuando ese parentesis se repite, lo sacamos como factor comun. Lo que queda afuera forma el otro factor \((5 x + 3 y)\).
\item Resultado: \[\boxed{\left(5 x + 3 y\right)\left(4 x - 7\right)}\]
\end{enumerate}

\vspace{0.2cm}
\subsubsection*{Ejercicio 53 - Factor comun por agrupacion}
\textbf{Enunciado:} $ 30 x^{2} + 20 x y - 54 x - 36 y $\par
\begin{enumerate}[leftmargin=2em]
\item En agrupacion no siempre hay un unico factor visible al inicio. Por eso intentamos formar grupos.
\item Los grupos se pueden ordenar para que aparezca el mismo parentesis \((5 x - 9)\).
\item Cuando ese parentesis se repite, lo sacamos como factor comun. Lo que queda afuera forma el otro factor \((6 x + 4 y)\).
\item Resultado: \[\boxed{\left(6 x + 4 y\right)\left(5 x - 9\right)}\]
\end{enumerate}

\vspace{0.2cm}
\subsubsection*{Ejercicio 54 - Factor comun por agrupacion}
\textbf{Enunciado:} $ 42 x^{2} + 30 x y - 77 x - 55 y $\par
\begin{enumerate}[leftmargin=2em]
\item En agrupacion no siempre hay un unico factor visible al inicio. Por eso intentamos formar grupos.
\item Los grupos se pueden ordenar para que aparezca el mismo parentesis \((6 x - 11)\).
\item Cuando ese parentesis se repite, lo sacamos como factor comun. Lo que queda afuera forma el otro factor \((7 x + 5 y)\).
\item Resultado: \[\boxed{\left(7 x + 5 y\right)\left(6 x - 11\right)}\]
\end{enumerate}

\vspace{0.2cm}
\subsubsection*{Ejercicio 55 - Factor comun por agrupacion}
\textbf{Enunciado:} $ 56 x^{2} + 42 x y - 104 x - 78 y $\par
\begin{enumerate}[leftmargin=2em]
\item En agrupacion no siempre hay un unico factor visible al inicio. Por eso intentamos formar grupos.
\item Los grupos se pueden ordenar para que aparezca el mismo parentesis \((7 x - 13)\).
\item Cuando ese parentesis se repite, lo sacamos como factor comun. Lo que queda afuera forma el otro factor \((8 x + 6 y)\).
\item Resultado: \[\boxed{\left(8 x + 6 y\right)\left(7 x - 13\right)}\]
\end{enumerate}

\vspace{0.2cm}
\subsubsection*{Ejercicio 56 - Factor comun por agrupacion}
\textbf{Enunciado:} $ 72 x^{2} + 56 x y - 135 x - 105 y $\par
\begin{enumerate}[leftmargin=2em]
\item En agrupacion no siempre hay un unico factor visible al inicio. Por eso intentamos formar grupos.
\item Los grupos se pueden ordenar para que aparezca el mismo parentesis \((8 x - 15)\).
\item Cuando ese parentesis se repite, lo sacamos como factor comun. Lo que queda afuera forma el otro factor \((9 x + 7 y)\).
\item Resultado: \[\boxed{\left(9 x + 7 y\right)\left(8 x - 15\right)}\]
\end{enumerate}

\vspace{0.2cm}
\subsubsection*{Ejercicio 57 - Factor comun por agrupacion}
\textbf{Enunciado:} $ 90 x^{2} + 72 x y - 170 x - 136 y $\par
\begin{enumerate}[leftmargin=2em]
\item En agrupacion no siempre hay un unico factor visible al inicio. Por eso intentamos formar grupos.
\item Los grupos se pueden ordenar para que aparezca el mismo parentesis \((9 x - 17)\).
\item Cuando ese parentesis se repite, lo sacamos como factor comun. Lo que queda afuera forma el otro factor \((10 x + 8 y)\).
\item Resultado: \[\boxed{\left(10 x + 8 y\right)\left(9 x - 17\right)}\]
\end{enumerate}

\vspace{0.2cm}
\subsubsection*{Ejercicio 58 - Factor comun por agrupacion}
\textbf{Enunciado:} $ 110 x^{2} + 90 x y - 209 x - 171 y $\par
\begin{enumerate}[leftmargin=2em]
\item En agrupacion no siempre hay un unico factor visible al inicio. Por eso intentamos formar grupos.
\item Los grupos se pueden ordenar para que aparezca el mismo parentesis \((10 x - 19)\).
\item Cuando ese parentesis se repite, lo sacamos como factor comun. Lo que queda afuera forma el otro factor \((11 x + 9 y)\).
\item Resultado: \[\boxed{\left(11 x + 9 y\right)\left(10 x - 19\right)}\]
\end{enumerate}

\vspace{0.2cm}
\subsubsection*{Ejercicio 59 - Factor comun por agrupacion}
\textbf{Enunciado:} $ 132 x^{2} + 110 x y - 252 x - 210 y $\par
\begin{enumerate}[leftmargin=2em]
\item En agrupacion no siempre hay un unico factor visible al inicio. Por eso intentamos formar grupos.
\item Los grupos se pueden ordenar para que aparezca el mismo parentesis \((11 x - 21)\).
\item Cuando ese parentesis se repite, lo sacamos como factor comun. Lo que queda afuera forma el otro factor \((12 x + 10 y)\).
\item Resultado: \[\boxed{\left(12 x + 10 y\right)\left(11 x - 21\right)}\]
\end{enumerate}

\vspace{0.2cm}
\subsubsection*{Ejercicio 60 - Factor comun por agrupacion}
\textbf{Enunciado:} $ 156 x^{2} + 132 x y - 299 x - 253 y $\par
\begin{enumerate}[leftmargin=2em]
\item En agrupacion no siempre hay un unico factor visible al inicio. Por eso intentamos formar grupos.
\item Los grupos se pueden ordenar para que aparezca el mismo parentesis \((12 x - 23)\).
\item Cuando ese parentesis se repite, lo sacamos como factor comun. Lo que queda afuera forma el otro factor \((13 x + 11 y)\).
\item Resultado: \[\boxed{\left(13 x + 11 y\right)\left(12 x - 23\right)}\]
\end{enumerate}

\vspace{0.2cm}
\nivel{4}{avanzado}
\subsubsection*{Ejercicio 61 - Factor comun normal}
\textbf{Enunciado:} $ 15 x^{7} y^{3} - 10 x^{6} y^{2} + 25 x^{5} y^{4} - 15 x^{5} y^{2} $\par
\begin{enumerate}[leftmargin=2em]
\item Miramos todos los terminos y buscamos que parte se repite en cada uno.
\item El factor comun es \(5 x^{5} y^{2}\). Eso significa que todos los terminos se pueden dividir entre \(5 x^{5} y^{2}\).
\item Al dividir el polinomio entre el factor comun queda \(3 x^{2} y - 2 x + 5 y^{2} - 3\).
\item Sacamos el factor comun fuera del parentesis. Resultado: \[\boxed{5 x^{5} y^{2}\left(3 x^{2} y - 2 x + 5 y^{2} - 3\right)}\]
\end{enumerate}

\vspace{0.2cm}
\subsubsection*{Ejercicio 62 - Factor comun normal}
\textbf{Enunciado:} $ 24 x^{8} y^{2} - 18 x^{7} y + 36 x^{6} y^{3} - 24 x^{6} y $\par
\begin{enumerate}[leftmargin=2em]
\item Miramos todos los terminos y buscamos que parte se repite en cada uno.
\item El factor comun es \(6 x^{6} y\). Eso significa que todos los terminos se pueden dividir entre \(6 x^{6} y\).
\item Al dividir el polinomio entre el factor comun queda \(4 x^{2} y - 3 x + 6 y^{2} - 4\).
\item Sacamos el factor comun fuera del parentesis. Resultado: \[\boxed{6 x^{6} y\left(4 x^{2} y - 3 x + 6 y^{2} - 4\right)}\]
\end{enumerate}

\vspace{0.2cm}
\subsubsection*{Ejercicio 63 - Factor comun normal}
\textbf{Enunciado:} $ 35 x^{6} y^{3} - 28 x^{5} y^{2} + 49 x^{4} y^{4} - 35 x^{4} y^{2} $\par
\begin{enumerate}[leftmargin=2em]
\item Miramos todos los terminos y buscamos que parte se repite en cada uno.
\item El factor comun es \(7 x^{4} y^{2}\). Eso significa que todos los terminos se pueden dividir entre \(7 x^{4} y^{2}\).
\item Al dividir el polinomio entre el factor comun queda \(5 x^{2} y - 4 x + 7 y^{2} - 5\).
\item Sacamos el factor comun fuera del parentesis. Resultado: \[\boxed{7 x^{4} y^{2}\left(5 x^{2} y - 4 x + 7 y^{2} - 5\right)}\]
\end{enumerate}

\vspace{0.2cm}
\subsubsection*{Ejercicio 64 - Factor comun normal}
\textbf{Enunciado:} $ 48 x^{7} y^{2} - 40 x^{6} y + 64 x^{5} y^{3} - 48 x^{5} y $\par
\begin{enumerate}[leftmargin=2em]
\item Miramos todos los terminos y buscamos que parte se repite en cada uno.
\item El factor comun es \(8 x^{5} y\). Eso significa que todos los terminos se pueden dividir entre \(8 x^{5} y\).
\item Al dividir el polinomio entre el factor comun queda \(6 x^{2} y - 5 x + 8 y^{2} - 6\).
\item Sacamos el factor comun fuera del parentesis. Resultado: \[\boxed{8 x^{5} y\left(6 x^{2} y - 5 x + 8 y^{2} - 6\right)}\]
\end{enumerate}

\vspace{0.2cm}
\subsubsection*{Ejercicio 65 - Factor comun normal}
\textbf{Enunciado:} $ 63 x^{8} y^{3} - 54 x^{7} y^{2} + 81 x^{6} y^{4} - 63 x^{6} y^{2} $\par
\begin{enumerate}[leftmargin=2em]
\item Miramos todos los terminos y buscamos que parte se repite en cada uno.
\item El factor comun es \(9 x^{6} y^{2}\). Eso significa que todos los terminos se pueden dividir entre \(9 x^{6} y^{2}\).
\item Al dividir el polinomio entre el factor comun queda \(7 x^{2} y - 6 x + 9 y^{2} - 7\).
\item Sacamos el factor comun fuera del parentesis. Resultado: \[\boxed{9 x^{6} y^{2}\left(7 x^{2} y - 6 x + 9 y^{2} - 7\right)}\]
\end{enumerate}

\vspace{0.2cm}
\subsubsection*{Ejercicio 66 - Factor comun normal}
\textbf{Enunciado:} $ 80 x^{6} y^{2} - 70 x^{5} y + 100 x^{4} y^{3} - 80 x^{4} y $\par
\begin{enumerate}[leftmargin=2em]
\item Miramos todos los terminos y buscamos que parte se repite en cada uno.
\item El factor comun es \(10 x^{4} y\). Eso significa que todos los terminos se pueden dividir entre \(10 x^{4} y\).
\item Al dividir el polinomio entre el factor comun queda \(8 x^{2} y - 7 x + 10 y^{2} - 8\).
\item Sacamos el factor comun fuera del parentesis. Resultado: \[\boxed{10 x^{4} y\left(8 x^{2} y - 7 x + 10 y^{2} - 8\right)}\]
\end{enumerate}

\vspace{0.2cm}
\subsubsection*{Ejercicio 67 - Factor comun normal}
\textbf{Enunciado:} $ 99 x^{7} y^{3} - 88 x^{6} y^{2} + 121 x^{5} y^{4} - 99 x^{5} y^{2} $\par
\begin{enumerate}[leftmargin=2em]
\item Miramos todos los terminos y buscamos que parte se repite en cada uno.
\item El factor comun es \(11 x^{5} y^{2}\). Eso significa que todos los terminos se pueden dividir entre \(11 x^{5} y^{2}\).
\item Al dividir el polinomio entre el factor comun queda \(9 x^{2} y - 8 x + 11 y^{2} - 9\).
\item Sacamos el factor comun fuera del parentesis. Resultado: \[\boxed{11 x^{5} y^{2}\left(9 x^{2} y - 8 x + 11 y^{2} - 9\right)}\]
\end{enumerate}

\vspace{0.2cm}
\subsubsection*{Ejercicio 68 - Factor comun normal}
\textbf{Enunciado:} $ 120 x^{8} y^{2} - 108 x^{7} y + 144 x^{6} y^{3} - 120 x^{6} y $\par
\begin{enumerate}[leftmargin=2em]
\item Miramos todos los terminos y buscamos que parte se repite en cada uno.
\item El factor comun es \(12 x^{6} y\). Eso significa que todos los terminos se pueden dividir entre \(12 x^{6} y\).
\item Al dividir el polinomio entre el factor comun queda \(10 x^{2} y - 9 x + 12 y^{2} - 10\).
\item Sacamos el factor comun fuera del parentesis. Resultado: \[\boxed{12 x^{6} y\left(10 x^{2} y - 9 x + 12 y^{2} - 10\right)}\]
\end{enumerate}

\vspace{0.2cm}
\subsubsection*{Ejercicio 69 - Factor comun normal}
\textbf{Enunciado:} $ 143 x^{6} y^{3} - 130 x^{5} y^{2} + 169 x^{4} y^{4} - 143 x^{4} y^{2} $\par
\begin{enumerate}[leftmargin=2em]
\item Miramos todos los terminos y buscamos que parte se repite en cada uno.
\item El factor comun es \(13 x^{4} y^{2}\). Eso significa que todos los terminos se pueden dividir entre \(13 x^{4} y^{2}\).
\item Al dividir el polinomio entre el factor comun queda \(11 x^{2} y - 10 x + 13 y^{2} - 11\).
\item Sacamos el factor comun fuera del parentesis. Resultado: \[\boxed{13 x^{4} y^{2}\left(11 x^{2} y - 10 x + 13 y^{2} - 11\right)}\]
\end{enumerate}

\vspace{0.2cm}
\subsubsection*{Ejercicio 70 - Factor comun normal}
\textbf{Enunciado:} $ 168 x^{7} y^{2} - 154 x^{6} y + 196 x^{5} y^{3} - 168 x^{5} y $\par
\begin{enumerate}[leftmargin=2em]
\item Miramos todos los terminos y buscamos que parte se repite en cada uno.
\item El factor comun es \(14 x^{5} y\). Eso significa que todos los terminos se pueden dividir entre \(14 x^{5} y\).
\item Al dividir el polinomio entre el factor comun queda \(12 x^{2} y - 11 x + 14 y^{2} - 12\).
\item Sacamos el factor comun fuera del parentesis. Resultado: \[\boxed{14 x^{5} y\left(12 x^{2} y - 11 x + 14 y^{2} - 12\right)}\]
\end{enumerate}

\vspace{0.2cm}
\subsubsection*{Ejercicio 71 - Factor comun por agrupacion}
\textbf{Enunciado:} $ 15 x^{3} + 30 x^{2} y - 6 x y - 12 y^{2} $\par
\begin{enumerate}[leftmargin=2em]
\item En agrupacion no siempre hay un unico factor visible al inicio. Por eso intentamos formar grupos.
\item Los grupos se pueden ordenar para que aparezca el mismo parentesis \((3 x + 6 y)\).
\item Cuando ese parentesis se repite, lo sacamos como factor comun. Lo que queda afuera forma el otro factor \((5 x^{2} - 2 y)\).
\item Resultado: \[\boxed{\left(5 x^{2} - 2 y\right)\left(3 x + 6 y\right)}\]
\end{enumerate}

\vspace{0.2cm}
\subsubsection*{Ejercicio 72 - Factor comun por agrupacion}
\textbf{Enunciado:} $ 24 x^{3} + 48 x^{2} y - 12 x y - 24 y^{2} $\par
\begin{enumerate}[leftmargin=2em]
\item En agrupacion no siempre hay un unico factor visible al inicio. Por eso intentamos formar grupos.
\item Los grupos se pueden ordenar para que aparezca el mismo parentesis \((4 x + 8 y)\).
\item Cuando ese parentesis se repite, lo sacamos como factor comun. Lo que queda afuera forma el otro factor \((6 x^{2} - 3 y)\).
\item Resultado: \[\boxed{\left(6 x^{2} - 3 y\right)\left(4 x + 8 y\right)}\]
\end{enumerate}

\vspace{0.2cm}
\subsubsection*{Ejercicio 73 - Factor comun por agrupacion}
\textbf{Enunciado:} $ 35 x^{3} + 70 x^{2} y - 20 x y - 40 y^{2} $\par
\begin{enumerate}[leftmargin=2em]
\item En agrupacion no siempre hay un unico factor visible al inicio. Por eso intentamos formar grupos.
\item Los grupos se pueden ordenar para que aparezca el mismo parentesis \((5 x + 10 y)\).
\item Cuando ese parentesis se repite, lo sacamos como factor comun. Lo que queda afuera forma el otro factor \((7 x^{2} - 4 y)\).
\item Resultado: \[\boxed{\left(7 x^{2} - 4 y\right)\left(5 x + 10 y\right)}\]
\end{enumerate}

\vspace{0.2cm}
\subsubsection*{Ejercicio 74 - Factor comun por agrupacion}
\textbf{Enunciado:} $ 48 x^{3} + 96 x^{2} y - 30 x y - 60 y^{2} $\par
\begin{enumerate}[leftmargin=2em]
\item En agrupacion no siempre hay un unico factor visible al inicio. Por eso intentamos formar grupos.
\item Los grupos se pueden ordenar para que aparezca el mismo parentesis \((6 x + 12 y)\).
\item Cuando ese parentesis se repite, lo sacamos como factor comun. Lo que queda afuera forma el otro factor \((8 x^{2} - 5 y)\).
\item Resultado: \[\boxed{\left(8 x^{2} - 5 y\right)\left(6 x + 12 y\right)}\]
\end{enumerate}

\vspace{0.2cm}
\subsubsection*{Ejercicio 75 - Factor comun por agrupacion}
\textbf{Enunciado:} $ 63 x^{3} + 126 x^{2} y - 42 x y - 84 y^{2} $\par
\begin{enumerate}[leftmargin=2em]
\item En agrupacion no siempre hay un unico factor visible al inicio. Por eso intentamos formar grupos.
\item Los grupos se pueden ordenar para que aparezca el mismo parentesis \((7 x + 14 y)\).
\item Cuando ese parentesis se repite, lo sacamos como factor comun. Lo que queda afuera forma el otro factor \((9 x^{2} - 6 y)\).
\item Resultado: \[\boxed{\left(9 x^{2} - 6 y\right)\left(7 x + 14 y\right)}\]
\end{enumerate}

\vspace{0.2cm}
\subsubsection*{Ejercicio 76 - Factor comun por agrupacion}
\textbf{Enunciado:} $ 80 x^{3} + 160 x^{2} y - 56 x y - 112 y^{2} $\par
\begin{enumerate}[leftmargin=2em]
\item En agrupacion no siempre hay un unico factor visible al inicio. Por eso intentamos formar grupos.
\item Los grupos se pueden ordenar para que aparezca el mismo parentesis \((8 x + 16 y)\).
\item Cuando ese parentesis se repite, lo sacamos como factor comun. Lo que queda afuera forma el otro factor \((10 x^{2} - 7 y)\).
\item Resultado: \[\boxed{\left(10 x^{2} - 7 y\right)\left(8 x + 16 y\right)}\]
\end{enumerate}

\vspace{0.2cm}
\subsubsection*{Ejercicio 77 - Factor comun por agrupacion}
\textbf{Enunciado:} $ 99 x^{3} + 198 x^{2} y - 72 x y - 144 y^{2} $\par
\begin{enumerate}[leftmargin=2em]
\item En agrupacion no siempre hay un unico factor visible al inicio. Por eso intentamos formar grupos.
\item Los grupos se pueden ordenar para que aparezca el mismo parentesis \((9 x + 18 y)\).
\item Cuando ese parentesis se repite, lo sacamos como factor comun. Lo que queda afuera forma el otro factor \((11 x^{2} - 8 y)\).
\item Resultado: \[\boxed{\left(11 x^{2} - 8 y\right)\left(9 x + 18 y\right)}\]
\end{enumerate}

\vspace{0.2cm}
\subsubsection*{Ejercicio 78 - Factor comun por agrupacion}
\textbf{Enunciado:} $ 120 x^{3} + 240 x^{2} y - 90 x y - 180 y^{2} $\par
\begin{enumerate}[leftmargin=2em]
\item En agrupacion no siempre hay un unico factor visible al inicio. Por eso intentamos formar grupos.
\item Los grupos se pueden ordenar para que aparezca el mismo parentesis \((10 x + 20 y)\).
\item Cuando ese parentesis se repite, lo sacamos como factor comun. Lo que queda afuera forma el otro factor \((12 x^{2} - 9 y)\).
\item Resultado: \[\boxed{\left(12 x^{2} - 9 y\right)\left(10 x + 20 y\right)}\]
\end{enumerate}

\vspace{0.2cm}
\subsubsection*{Ejercicio 79 - Factor comun por agrupacion}
\textbf{Enunciado:} $ 143 x^{3} + 286 x^{2} y - 110 x y - 220 y^{2} $\par
\begin{enumerate}[leftmargin=2em]
\item En agrupacion no siempre hay un unico factor visible al inicio. Por eso intentamos formar grupos.
\item Los grupos se pueden ordenar para que aparezca el mismo parentesis \((11 x + 22 y)\).
\item Cuando ese parentesis se repite, lo sacamos como factor comun. Lo que queda afuera forma el otro factor \((13 x^{2} - 10 y)\).
\item Resultado: \[\boxed{\left(13 x^{2} - 10 y\right)\left(11 x + 22 y\right)}\]
\end{enumerate}

\vspace{0.2cm}
\subsubsection*{Ejercicio 80 - Factor comun por agrupacion}
\textbf{Enunciado:} $ 168 x^{3} + 336 x^{2} y - 132 x y - 264 y^{2} $\par
\begin{enumerate}[leftmargin=2em]
\item En agrupacion no siempre hay un unico factor visible al inicio. Por eso intentamos formar grupos.
\item Los grupos se pueden ordenar para que aparezca el mismo parentesis \((12 x + 24 y)\).
\item Cuando ese parentesis se repite, lo sacamos como factor comun. Lo que queda afuera forma el otro factor \((14 x^{2} - 11 y)\).
\item Resultado: \[\boxed{\left(14 x^{2} - 11 y\right)\left(12 x + 24 y\right)}\]
\end{enumerate}

\vspace{0.2cm}
\nivel{5}{imposible}
\subsubsection*{Ejercicio 81 - Factor comun normal}
\textbf{Enunciado:} $ 12 x^{9} y^{4} + 24 x^{8} y^{3} - 18 x^{7} y^{5} - 30 x^{6} y^{4} + 36 x^{6} y^{3} $\par
\begin{enumerate}[leftmargin=2em]
\item Miramos todos los terminos y buscamos que parte se repite en cada uno.
\item El factor comun es \(6 x^{6} y^{3}\). Eso significa que todos los terminos se pueden dividir entre \(6 x^{6} y^{3}\).
\item Al dividir el polinomio entre el factor comun queda \(2 x^{3} y + 4 x^{2} - 3 x y^{2} - 5 y + 6\).
\item Sacamos el factor comun fuera del parentesis. Resultado: \[\boxed{6 x^{6} y^{3}\left(2 x^{3} y + 4 x^{2} - 3 x y^{2} - 5 y + 6\right)}\]
\end{enumerate}

\vspace{0.2cm}
\subsubsection*{Ejercicio 82 - Factor comun normal}
\textbf{Enunciado:} $ 21 x^{10} y^{3} + 35 x^{9} y^{2} - 28 x^{8} y^{4} - 42 x^{7} y^{3} + 49 x^{7} y^{2} $\par
\begin{enumerate}[leftmargin=2em]
\item Miramos todos los terminos y buscamos que parte se repite en cada uno.
\item El factor comun es \(7 x^{7} y^{2}\). Eso significa que todos los terminos se pueden dividir entre \(7 x^{7} y^{2}\).
\item Al dividir el polinomio entre el factor comun queda \(3 x^{3} y + 5 x^{2} - 4 x y^{2} - 6 y + 7\).
\item Sacamos el factor comun fuera del parentesis. Resultado: \[\boxed{7 x^{7} y^{2}\left(3 x^{3} y + 5 x^{2} - 4 x y^{2} - 6 y + 7\right)}\]
\end{enumerate}

\vspace{0.2cm}
\subsubsection*{Ejercicio 83 - Factor comun normal}
\textbf{Enunciado:} $ 32 x^{8} y^{4} + 48 x^{7} y^{3} - 40 x^{6} y^{5} - 56 x^{5} y^{4} + 64 x^{5} y^{3} $\par
\begin{enumerate}[leftmargin=2em]
\item Miramos todos los terminos y buscamos que parte se repite en cada uno.
\item El factor comun es \(8 x^{5} y^{3}\). Eso significa que todos los terminos se pueden dividir entre \(8 x^{5} y^{3}\).
\item Al dividir el polinomio entre el factor comun queda \(4 x^{3} y + 6 x^{2} - 5 x y^{2} - 7 y + 8\).
\item Sacamos el factor comun fuera del parentesis. Resultado: \[\boxed{8 x^{5} y^{3}\left(4 x^{3} y + 6 x^{2} - 5 x y^{2} - 7 y + 8\right)}\]
\end{enumerate}

\vspace{0.2cm}
\subsubsection*{Ejercicio 84 - Factor comun normal}
\textbf{Enunciado:} $ 45 x^{9} y^{3} + 63 x^{8} y^{2} - 54 x^{7} y^{4} - 72 x^{6} y^{3} + 81 x^{6} y^{2} $\par
\begin{enumerate}[leftmargin=2em]
\item Miramos todos los terminos y buscamos que parte se repite en cada uno.
\item El factor comun es \(9 x^{6} y^{2}\). Eso significa que todos los terminos se pueden dividir entre \(9 x^{6} y^{2}\).
\item Al dividir el polinomio entre el factor comun queda \(5 x^{3} y + 7 x^{2} - 6 x y^{2} - 8 y + 9\).
\item Sacamos el factor comun fuera del parentesis. Resultado: \[\boxed{9 x^{6} y^{2}\left(5 x^{3} y + 7 x^{2} - 6 x y^{2} - 8 y + 9\right)}\]
\end{enumerate}

\vspace{0.2cm}
\subsubsection*{Ejercicio 85 - Factor comun normal}
\textbf{Enunciado:} $ 60 x^{10} y^{4} + 80 x^{9} y^{3} - 70 x^{8} y^{5} - 90 x^{7} y^{4} + 100 x^{7} y^{3} $\par
\begin{enumerate}[leftmargin=2em]
\item Miramos todos los terminos y buscamos que parte se repite en cada uno.
\item El factor comun es \(10 x^{7} y^{3}\). Eso significa que todos los terminos se pueden dividir entre \(10 x^{7} y^{3}\).
\item Al dividir el polinomio entre el factor comun queda \(6 x^{3} y + 8 x^{2} - 7 x y^{2} - 9 y + 10\).
\item Sacamos el factor comun fuera del parentesis. Resultado: \[\boxed{10 x^{7} y^{3}\left(6 x^{3} y + 8 x^{2} - 7 x y^{2} - 9 y + 10\right)}\]
\end{enumerate}

\vspace{0.2cm}
\subsubsection*{Ejercicio 86 - Factor comun normal}
\textbf{Enunciado:} $ 77 x^{8} y^{3} + 99 x^{7} y^{2} - 88 x^{6} y^{4} - 110 x^{5} y^{3} + 121 x^{5} y^{2} $\par
\begin{enumerate}[leftmargin=2em]
\item Miramos todos los terminos y buscamos que parte se repite en cada uno.
\item El factor comun es \(11 x^{5} y^{2}\). Eso significa que todos los terminos se pueden dividir entre \(11 x^{5} y^{2}\).
\item Al dividir el polinomio entre el factor comun queda \(7 x^{3} y + 9 x^{2} - 8 x y^{2} - 10 y + 11\).
\item Sacamos el factor comun fuera del parentesis. Resultado: \[\boxed{11 x^{5} y^{2}\left(7 x^{3} y + 9 x^{2} - 8 x y^{2} - 10 y + 11\right)}\]
\end{enumerate}

\vspace{0.2cm}
\subsubsection*{Ejercicio 87 - Factor comun normal}
\textbf{Enunciado:} $ 96 x^{9} y^{4} + 120 x^{8} y^{3} - 108 x^{7} y^{5} - 132 x^{6} y^{4} + 144 x^{6} y^{3} $\par
\begin{enumerate}[leftmargin=2em]
\item Miramos todos los terminos y buscamos que parte se repite en cada uno.
\item El factor comun es \(12 x^{6} y^{3}\). Eso significa que todos los terminos se pueden dividir entre \(12 x^{6} y^{3}\).
\item Al dividir el polinomio entre el factor comun queda \(8 x^{3} y + 10 x^{2} - 9 x y^{2} - 11 y + 12\).
\item Sacamos el factor comun fuera del parentesis. Resultado: \[\boxed{12 x^{6} y^{3}\left(8 x^{3} y + 10 x^{2} - 9 x y^{2} - 11 y + 12\right)}\]
\end{enumerate}

\vspace{0.2cm}
\subsubsection*{Ejercicio 88 - Factor comun normal}
\textbf{Enunciado:} $ 117 x^{10} y^{3} + 143 x^{9} y^{2} - 130 x^{8} y^{4} - 156 x^{7} y^{3} + 169 x^{7} y^{2} $\par
\begin{enumerate}[leftmargin=2em]
\item Miramos todos los terminos y buscamos que parte se repite en cada uno.
\item El factor comun es \(13 x^{7} y^{2}\). Eso significa que todos los terminos se pueden dividir entre \(13 x^{7} y^{2}\).
\item Al dividir el polinomio entre el factor comun queda \(9 x^{3} y + 11 x^{2} - 10 x y^{2} - 12 y + 13\).
\item Sacamos el factor comun fuera del parentesis. Resultado: \[\boxed{13 x^{7} y^{2}\left(9 x^{3} y + 11 x^{2} - 10 x y^{2} - 12 y + 13\right)}\]
\end{enumerate}

\vspace{0.2cm}
\subsubsection*{Ejercicio 89 - Factor comun normal}
\textbf{Enunciado:} $ 140 x^{8} y^{4} + 168 x^{7} y^{3} - 154 x^{6} y^{5} - 182 x^{5} y^{4} + 196 x^{5} y^{3} $\par
\begin{enumerate}[leftmargin=2em]
\item Miramos todos los terminos y buscamos que parte se repite en cada uno.
\item El factor comun es \(14 x^{5} y^{3}\). Eso significa que todos los terminos se pueden dividir entre \(14 x^{5} y^{3}\).
\item Al dividir el polinomio entre el factor comun queda \(10 x^{3} y + 12 x^{2} - 11 x y^{2} - 13 y + 14\).
\item Sacamos el factor comun fuera del parentesis. Resultado: \[\boxed{14 x^{5} y^{3}\left(10 x^{3} y + 12 x^{2} - 11 x y^{2} - 13 y + 14\right)}\]
\end{enumerate}

\vspace{0.2cm}
\subsubsection*{Ejercicio 90 - Factor comun normal}
\textbf{Enunciado:} $ 165 x^{9} y^{3} + 195 x^{8} y^{2} - 180 x^{7} y^{4} - 210 x^{6} y^{3} + 225 x^{6} y^{2} $\par
\begin{enumerate}[leftmargin=2em]
\item Miramos todos los terminos y buscamos que parte se repite en cada uno.
\item El factor comun es \(15 x^{6} y^{2}\). Eso significa que todos los terminos se pueden dividir entre \(15 x^{6} y^{2}\).
\item Al dividir el polinomio entre el factor comun queda \(11 x^{3} y + 13 x^{2} - 12 x y^{2} - 14 y + 15\).
\item Sacamos el factor comun fuera del parentesis. Resultado: \[\boxed{15 x^{6} y^{2}\left(11 x^{3} y + 13 x^{2} - 12 x y^{2} - 14 y + 15\right)}\]
\end{enumerate}

\vspace{0.2cm}
\subsubsection*{Ejercicio 91 - Factor comun por agrupacion}
\textbf{Enunciado:} $ 18 x^{3} y^{2} + 42 x^{2} y^{3} - 30 x^{2} y - 11 x y^{2} + 8 x + 7 y^{3} - 4 y $\par
\begin{enumerate}[leftmargin=2em]
\item En agrupacion no siempre hay un unico factor visible al inicio. Por eso intentamos formar grupos.
\item Los grupos se pueden ordenar para que aparezca el mismo parentesis \((3 x y + 7 y^{2} - 4)\).
\item Cuando ese parentesis se repite, lo sacamos como factor comun. Lo que queda afuera forma el otro factor \((6 x^{2} y - 2 x + y)\).
\item Resultado: \[\boxed{\left(6 x^{2} y - 2 x + y\right)\left(3 x y + 7 y^{2} - 4\right)}\]
\end{enumerate}

\vspace{0.2cm}
\subsubsection*{Ejercicio 92 - Factor comun por agrupacion}
\textbf{Enunciado:} $ 28 x^{3} y^{2} + 63 x^{2} y^{3} - 47 x^{2} y - 23 x y^{2} + 15 x + 9 y^{3} - 5 y $\par
\begin{enumerate}[leftmargin=2em]
\item En agrupacion no siempre hay un unico factor visible al inicio. Por eso intentamos formar grupos.
\item Los grupos se pueden ordenar para que aparezca el mismo parentesis \((4 x y + 9 y^{2} - 5)\).
\item Cuando ese parentesis se repite, lo sacamos como factor comun. Lo que queda afuera forma el otro factor \((7 x^{2} y - 3 x + y)\).
\item Resultado: \[\boxed{\left(7 x^{2} y - 3 x + y\right)\left(4 x y + 9 y^{2} - 5\right)}\]
\end{enumerate}

\vspace{0.2cm}
\subsubsection*{Ejercicio 93 - Factor comun por agrupacion}
\textbf{Enunciado:} $ 40 x^{3} y^{2} + 88 x^{2} y^{3} - 68 x^{2} y - 39 x y^{2} + 24 x + 11 y^{3} - 6 y $\par
\begin{enumerate}[leftmargin=2em]
\item En agrupacion no siempre hay un unico factor visible al inicio. Por eso intentamos formar grupos.
\item Los grupos se pueden ordenar para que aparezca el mismo parentesis \((5 x y + 11 y^{2} - 6)\).
\item Cuando ese parentesis se repite, lo sacamos como factor comun. Lo que queda afuera forma el otro factor \((8 x^{2} y - 4 x + y)\).
\item Resultado: \[\boxed{\left(8 x^{2} y - 4 x + y\right)\left(5 x y + 11 y^{2} - 6\right)}\]
\end{enumerate}

\vspace{0.2cm}
\subsubsection*{Ejercicio 94 - Factor comun por agrupacion}
\textbf{Enunciado:} $ 54 x^{3} y^{2} + 117 x^{2} y^{3} - 93 x^{2} y - 59 x y^{2} + 35 x + 13 y^{3} - 7 y $\par
\begin{enumerate}[leftmargin=2em]
\item En agrupacion no siempre hay un unico factor visible al inicio. Por eso intentamos formar grupos.
\item Los grupos se pueden ordenar para que aparezca el mismo parentesis \((6 x y + 13 y^{2} - 7)\).
\item Cuando ese parentesis se repite, lo sacamos como factor comun. Lo que queda afuera forma el otro factor \((9 x^{2} y - 5 x + y)\).
\item Resultado: \[\boxed{\left(9 x^{2} y - 5 x + y\right)\left(6 x y + 13 y^{2} - 7\right)}\]
\end{enumerate}

\vspace{0.2cm}
\subsubsection*{Ejercicio 95 - Factor comun por agrupacion}
\textbf{Enunciado:} $ 70 x^{3} y^{2} + 150 x^{2} y^{3} - 122 x^{2} y - 83 x y^{2} + 48 x + 15 y^{3} - 8 y $\par
\begin{enumerate}[leftmargin=2em]
\item En agrupacion no siempre hay un unico factor visible al inicio. Por eso intentamos formar grupos.
\item Los grupos se pueden ordenar para que aparezca el mismo parentesis \((7 x y + 15 y^{2} - 8)\).
\item Cuando ese parentesis se repite, lo sacamos como factor comun. Lo que queda afuera forma el otro factor \((10 x^{2} y - 6 x + y)\).
\item Resultado: \[\boxed{\left(10 x^{2} y - 6 x + y\right)\left(7 x y + 15 y^{2} - 8\right)}\]
\end{enumerate}

\vspace{0.2cm}
\subsubsection*{Ejercicio 96 - Factor comun por agrupacion}
\textbf{Enunciado:} $ 88 x^{3} y^{2} + 187 x^{2} y^{3} - 155 x^{2} y - 111 x y^{2} + 63 x + 17 y^{3} - 9 y $\par
\begin{enumerate}[leftmargin=2em]
\item En agrupacion no siempre hay un unico factor visible al inicio. Por eso intentamos formar grupos.
\item Los grupos se pueden ordenar para que aparezca el mismo parentesis \((8 x y + 17 y^{2} - 9)\).
\item Cuando ese parentesis se repite, lo sacamos como factor comun. Lo que queda afuera forma el otro factor \((11 x^{2} y - 7 x + y)\).
\item Resultado: \[\boxed{\left(11 x^{2} y - 7 x + y\right)\left(8 x y + 17 y^{2} - 9\right)}\]
\end{enumerate}

\vspace{0.2cm}
\subsubsection*{Ejercicio 97 - Factor comun por agrupacion}
\textbf{Enunciado:} $ 108 x^{3} y^{2} + 228 x^{2} y^{3} - 192 x^{2} y - 143 x y^{2} + 80 x + 19 y^{3} - 10 y $\par
\begin{enumerate}[leftmargin=2em]
\item En agrupacion no siempre hay un unico factor visible al inicio. Por eso intentamos formar grupos.
\item Los grupos se pueden ordenar para que aparezca el mismo parentesis \((9 x y + 19 y^{2} - 10)\).
\item Cuando ese parentesis se repite, lo sacamos como factor comun. Lo que queda afuera forma el otro factor \((12 x^{2} y - 8 x + y)\).
\item Resultado: \[\boxed{\left(12 x^{2} y - 8 x + y\right)\left(9 x y + 19 y^{2} - 10\right)}\]
\end{enumerate}

\vspace{0.2cm}
\subsubsection*{Ejercicio 98 - Factor comun por agrupacion}
\textbf{Enunciado:} $ 130 x^{3} y^{2} + 273 x^{2} y^{3} - 233 x^{2} y - 179 x y^{2} + 99 x + 21 y^{3} - 11 y $\par
\begin{enumerate}[leftmargin=2em]
\item En agrupacion no siempre hay un unico factor visible al inicio. Por eso intentamos formar grupos.
\item Los grupos se pueden ordenar para que aparezca el mismo parentesis \((10 x y + 21 y^{2} - 11)\).
\item Cuando ese parentesis se repite, lo sacamos como factor comun. Lo que queda afuera forma el otro factor \((13 x^{2} y - 9 x + y)\).
\item Resultado: \[\boxed{\left(13 x^{2} y - 9 x + y\right)\left(10 x y + 21 y^{2} - 11\right)}\]
\end{enumerate}

\vspace{0.2cm}
\subsubsection*{Ejercicio 99 - Factor comun por agrupacion}
\textbf{Enunciado:} $ 154 x^{3} y^{2} + 322 x^{2} y^{3} - 278 x^{2} y - 219 x y^{2} + 120 x + 23 y^{3} - 12 y $\par
\begin{enumerate}[leftmargin=2em]
\item En agrupacion no siempre hay un unico factor visible al inicio. Por eso intentamos formar grupos.
\item Los grupos se pueden ordenar para que aparezca el mismo parentesis \((11 x y + 23 y^{2} - 12)\).
\item Cuando ese parentesis se repite, lo sacamos como factor comun. Lo que queda afuera forma el otro factor \((14 x^{2} y - 10 x + y)\).
\item Resultado: \[\boxed{\left(14 x^{2} y - 10 x + y\right)\left(11 x y + 23 y^{2} - 12\right)}\]
\end{enumerate}

\vspace{0.2cm}
\subsubsection*{Ejercicio 100 - Factor comun por agrupacion}
\textbf{Enunciado:} $ 180 x^{3} y^{2} + 375 x^{2} y^{3} - 327 x^{2} y - 263 x y^{2} + 143 x + 25 y^{3} - 13 y $\par
\begin{enumerate}[leftmargin=2em]
\item En agrupacion no siempre hay un unico factor visible al inicio. Por eso intentamos formar grupos.
\item Los grupos se pueden ordenar para que aparezca el mismo parentesis \((12 x y + 25 y^{2} - 13)\).
\item Cuando ese parentesis se repite, lo sacamos como factor comun. Lo que queda afuera forma el otro factor \((15 x^{2} y - 11 x + y)\).
\item Resultado: \[\boxed{\left(15 x^{2} y - 11 x + y\right)\left(12 x y + 25 y^{2} - 13\right)}\]
\end{enumerate}

\vspace{0.2cm}
\section{Soluciones: Diferencia de cuadrados}
\nivel{1}{basico}
\subsubsection*{Ejercicio 101 - Diferencia de cuadrados}
\textbf{Enunciado:} $ 4 x^{2} - 4 $\par
\begin{enumerate}[leftmargin=2em]
\item La diferencia de cuadrados siempre tiene esta forma: \(A^2-B^2\).
\item Aqui identificamos \(A=2 x\) y \(B=2\).
\item Aplicamos la formula: \(A^2-B^2=(A-B)(A+B)\).
\item Resultado: \[\boxed{\left(\left(2 x\right)-\left(2\right)\right)\left(\left(2 x\right)+\left(2\right)\right)}\]
\end{enumerate}

\vspace{0.2cm}
\subsubsection*{Ejercicio 102 - Diferencia de cuadrados}
\textbf{Enunciado:} $ 9 x^{2} - 9 $\par
\begin{enumerate}[leftmargin=2em]
\item La diferencia de cuadrados siempre tiene esta forma: \(A^2-B^2\).
\item Aqui identificamos \(A=3 x\) y \(B=3\).
\item Aplicamos la formula: \(A^2-B^2=(A-B)(A+B)\).
\item Resultado: \[\boxed{\left(\left(3 x\right)-\left(3\right)\right)\left(\left(3 x\right)+\left(3\right)\right)}\]
\end{enumerate}

\vspace{0.2cm}
\subsubsection*{Ejercicio 103 - Diferencia de cuadrados}
\textbf{Enunciado:} $ 16 x^{2} - 16 $\par
\begin{enumerate}[leftmargin=2em]
\item La diferencia de cuadrados siempre tiene esta forma: \(A^2-B^2\).
\item Aqui identificamos \(A=4 x\) y \(B=4\).
\item Aplicamos la formula: \(A^2-B^2=(A-B)(A+B)\).
\item Resultado: \[\boxed{\left(\left(4 x\right)-\left(4\right)\right)\left(\left(4 x\right)+\left(4\right)\right)}\]
\end{enumerate}

\vspace{0.2cm}
\subsubsection*{Ejercicio 104 - Diferencia de cuadrados}
\textbf{Enunciado:} $ 25 x^{2} - 25 $\par
\begin{enumerate}[leftmargin=2em]
\item La diferencia de cuadrados siempre tiene esta forma: \(A^2-B^2\).
\item Aqui identificamos \(A=5 x\) y \(B=5\).
\item Aplicamos la formula: \(A^2-B^2=(A-B)(A+B)\).
\item Resultado: \[\boxed{\left(\left(5 x\right)-\left(5\right)\right)\left(\left(5 x\right)+\left(5\right)\right)}\]
\end{enumerate}

\vspace{0.2cm}
\subsubsection*{Ejercicio 105 - Diferencia de cuadrados}
\textbf{Enunciado:} $ x^{2} - 36 $\par
\begin{enumerate}[leftmargin=2em]
\item La diferencia de cuadrados siempre tiene esta forma: \(A^2-B^2\).
\item Aqui identificamos \(A=x\) y \(B=6\).
\item Aplicamos la formula: \(A^2-B^2=(A-B)(A+B)\).
\item Resultado: \[\boxed{\left(\left(x\right)-\left(6\right)\right)\left(\left(x\right)+\left(6\right)\right)}\]
\end{enumerate}

\vspace{0.2cm}
\subsubsection*{Ejercicio 106 - Diferencia de cuadrados}
\textbf{Enunciado:} $ 4 x^{2} - 49 $\par
\begin{enumerate}[leftmargin=2em]
\item La diferencia de cuadrados siempre tiene esta forma: \(A^2-B^2\).
\item Aqui identificamos \(A=2 x\) y \(B=7\).
\item Aplicamos la formula: \(A^2-B^2=(A-B)(A+B)\).
\item Resultado: \[\boxed{\left(\left(2 x\right)-\left(7\right)\right)\left(\left(2 x\right)+\left(7\right)\right)}\]
\end{enumerate}

\vspace{0.2cm}
\subsubsection*{Ejercicio 107 - Diferencia de cuadrados}
\textbf{Enunciado:} $ 9 x^{2} - 64 $\par
\begin{enumerate}[leftmargin=2em]
\item La diferencia de cuadrados siempre tiene esta forma: \(A^2-B^2\).
\item Aqui identificamos \(A=3 x\) y \(B=8\).
\item Aplicamos la formula: \(A^2-B^2=(A-B)(A+B)\).
\item Resultado: \[\boxed{\left(\left(3 x\right)-\left(8\right)\right)\left(\left(3 x\right)+\left(8\right)\right)}\]
\end{enumerate}

\vspace{0.2cm}
\subsubsection*{Ejercicio 108 - Diferencia de cuadrados}
\textbf{Enunciado:} $ 16 x^{2} - 81 $\par
\begin{enumerate}[leftmargin=2em]
\item La diferencia de cuadrados siempre tiene esta forma: \(A^2-B^2\).
\item Aqui identificamos \(A=4 x\) y \(B=9\).
\item Aplicamos la formula: \(A^2-B^2=(A-B)(A+B)\).
\item Resultado: \[\boxed{\left(\left(4 x\right)-\left(9\right)\right)\left(\left(4 x\right)+\left(9\right)\right)}\]
\end{enumerate}

\vspace{0.2cm}
\subsubsection*{Ejercicio 109 - Diferencia de cuadrados}
\textbf{Enunciado:} $ 25 x^{2} - 100 $\par
\begin{enumerate}[leftmargin=2em]
\item La diferencia de cuadrados siempre tiene esta forma: \(A^2-B^2\).
\item Aqui identificamos \(A=5 x\) y \(B=10\).
\item Aplicamos la formula: \(A^2-B^2=(A-B)(A+B)\).
\item Resultado: \[\boxed{\left(\left(5 x\right)-\left(10\right)\right)\left(\left(5 x\right)+\left(10\right)\right)}\]
\end{enumerate}

\vspace{0.2cm}
\subsubsection*{Ejercicio 110 - Diferencia de cuadrados}
\textbf{Enunciado:} $ x^{2} - 121 $\par
\begin{enumerate}[leftmargin=2em]
\item La diferencia de cuadrados siempre tiene esta forma: \(A^2-B^2\).
\item Aqui identificamos \(A=x\) y \(B=11\).
\item Aplicamos la formula: \(A^2-B^2=(A-B)(A+B)\).
\item Resultado: \[\boxed{\left(\left(x\right)-\left(11\right)\right)\left(\left(x\right)+\left(11\right)\right)}\]
\end{enumerate}

\vspace{0.2cm}
\subsubsection*{Ejercicio 111 - Diferencia de cuadrados}
\textbf{Enunciado:} $ 4 x^{2} - 144 $\par
\begin{enumerate}[leftmargin=2em]
\item La diferencia de cuadrados siempre tiene esta forma: \(A^2-B^2\).
\item Aqui identificamos \(A=2 x\) y \(B=12\).
\item Aplicamos la formula: \(A^2-B^2=(A-B)(A+B)\).
\item Resultado: \[\boxed{\left(\left(2 x\right)-\left(12\right)\right)\left(\left(2 x\right)+\left(12\right)\right)}\]
\end{enumerate}

\vspace{0.2cm}
\subsubsection*{Ejercicio 112 - Diferencia de cuadrados}
\textbf{Enunciado:} $ 9 x^{2} - 169 $\par
\begin{enumerate}[leftmargin=2em]
\item La diferencia de cuadrados siempre tiene esta forma: \(A^2-B^2\).
\item Aqui identificamos \(A=3 x\) y \(B=13\).
\item Aplicamos la formula: \(A^2-B^2=(A-B)(A+B)\).
\item Resultado: \[\boxed{\left(\left(3 x\right)-\left(13\right)\right)\left(\left(3 x\right)+\left(13\right)\right)}\]
\end{enumerate}

\vspace{0.2cm}
\subsubsection*{Ejercicio 113 - Diferencia de cuadrados}
\textbf{Enunciado:} $ 16 x^{2} - 196 $\par
\begin{enumerate}[leftmargin=2em]
\item La diferencia de cuadrados siempre tiene esta forma: \(A^2-B^2\).
\item Aqui identificamos \(A=4 x\) y \(B=14\).
\item Aplicamos la formula: \(A^2-B^2=(A-B)(A+B)\).
\item Resultado: \[\boxed{\left(\left(4 x\right)-\left(14\right)\right)\left(\left(4 x\right)+\left(14\right)\right)}\]
\end{enumerate}

\vspace{0.2cm}
\subsubsection*{Ejercicio 114 - Diferencia de cuadrados}
\textbf{Enunciado:} $ 25 x^{2} - 225 $\par
\begin{enumerate}[leftmargin=2em]
\item La diferencia de cuadrados siempre tiene esta forma: \(A^2-B^2\).
\item Aqui identificamos \(A=5 x\) y \(B=15\).
\item Aplicamos la formula: \(A^2-B^2=(A-B)(A+B)\).
\item Resultado: \[\boxed{\left(\left(5 x\right)-\left(15\right)\right)\left(\left(5 x\right)+\left(15\right)\right)}\]
\end{enumerate}

\vspace{0.2cm}
\subsubsection*{Ejercicio 115 - Diferencia de cuadrados}
\textbf{Enunciado:} $ x^{2} - 256 $\par
\begin{enumerate}[leftmargin=2em]
\item La diferencia de cuadrados siempre tiene esta forma: \(A^2-B^2\).
\item Aqui identificamos \(A=x\) y \(B=16\).
\item Aplicamos la formula: \(A^2-B^2=(A-B)(A+B)\).
\item Resultado: \[\boxed{\left(\left(x\right)-\left(16\right)\right)\left(\left(x\right)+\left(16\right)\right)}\]
\end{enumerate}

\vspace{0.2cm}
\subsubsection*{Ejercicio 116 - Diferencia de cuadrados}
\textbf{Enunciado:} $ 4 x^{2} - 289 $\par
\begin{enumerate}[leftmargin=2em]
\item La diferencia de cuadrados siempre tiene esta forma: \(A^2-B^2\).
\item Aqui identificamos \(A=2 x\) y \(B=17\).
\item Aplicamos la formula: \(A^2-B^2=(A-B)(A+B)\).
\item Resultado: \[\boxed{\left(\left(2 x\right)-\left(17\right)\right)\left(\left(2 x\right)+\left(17\right)\right)}\]
\end{enumerate}

\vspace{0.2cm}
\subsubsection*{Ejercicio 117 - Diferencia de cuadrados}
\textbf{Enunciado:} $ 9 x^{2} - 324 $\par
\begin{enumerate}[leftmargin=2em]
\item La diferencia de cuadrados siempre tiene esta forma: \(A^2-B^2\).
\item Aqui identificamos \(A=3 x\) y \(B=18\).
\item Aplicamos la formula: \(A^2-B^2=(A-B)(A+B)\).
\item Resultado: \[\boxed{\left(\left(3 x\right)-\left(18\right)\right)\left(\left(3 x\right)+\left(18\right)\right)}\]
\end{enumerate}

\vspace{0.2cm}
\subsubsection*{Ejercicio 118 - Diferencia de cuadrados}
\textbf{Enunciado:} $ 16 x^{2} - 361 $\par
\begin{enumerate}[leftmargin=2em]
\item La diferencia de cuadrados siempre tiene esta forma: \(A^2-B^2\).
\item Aqui identificamos \(A=4 x\) y \(B=19\).
\item Aplicamos la formula: \(A^2-B^2=(A-B)(A+B)\).
\item Resultado: \[\boxed{\left(\left(4 x\right)-\left(19\right)\right)\left(\left(4 x\right)+\left(19\right)\right)}\]
\end{enumerate}

\vspace{0.2cm}
\subsubsection*{Ejercicio 119 - Diferencia de cuadrados}
\textbf{Enunciado:} $ 25 x^{2} - 400 $\par
\begin{enumerate}[leftmargin=2em]
\item La diferencia de cuadrados siempre tiene esta forma: \(A^2-B^2\).
\item Aqui identificamos \(A=5 x\) y \(B=20\).
\item Aplicamos la formula: \(A^2-B^2=(A-B)(A+B)\).
\item Resultado: \[\boxed{\left(\left(5 x\right)-\left(20\right)\right)\left(\left(5 x\right)+\left(20\right)\right)}\]
\end{enumerate}

\vspace{0.2cm}
\subsubsection*{Ejercicio 120 - Diferencia de cuadrados}
\textbf{Enunciado:} $ x^{2} - 441 $\par
\begin{enumerate}[leftmargin=2em]
\item La diferencia de cuadrados siempre tiene esta forma: \(A^2-B^2\).
\item Aqui identificamos \(A=x\) y \(B=21\).
\item Aplicamos la formula: \(A^2-B^2=(A-B)(A+B)\).
\item Resultado: \[\boxed{\left(\left(x\right)-\left(21\right)\right)\left(\left(x\right)+\left(21\right)\right)}\]
\end{enumerate}

\vspace{0.2cm}
\nivel{2}{intermedio inicial}
\subsubsection*{Ejercicio 121 - Diferencia de cuadrados}
\textbf{Enunciado:} $ 9 x^{4} - 4 y^{2} $\par
\begin{enumerate}[leftmargin=2em]
\item La diferencia de cuadrados siempre tiene esta forma: \(A^2-B^2\).
\item Aqui identificamos \(A=3 x^{2}\) y \(B=2 y\).
\item Aplicamos la formula: \(A^2-B^2=(A-B)(A+B)\).
\item Resultado: \[\boxed{\left(\left(3 x^{2}\right)-\left(2 y\right)\right)\left(\left(3 x^{2}\right)+\left(2 y\right)\right)}\]
\end{enumerate}

\vspace{0.2cm}
\subsubsection*{Ejercicio 122 - Diferencia de cuadrados}
\textbf{Enunciado:} $ 16 x^{2} - 9 y^{2} $\par
\begin{enumerate}[leftmargin=2em]
\item La diferencia de cuadrados siempre tiene esta forma: \(A^2-B^2\).
\item Aqui identificamos \(A=4 x\) y \(B=3 y\).
\item Aplicamos la formula: \(A^2-B^2=(A-B)(A+B)\).
\item Resultado: \[\boxed{\left(\left(4 x\right)-\left(3 y\right)\right)\left(\left(4 x\right)+\left(3 y\right)\right)}\]
\end{enumerate}

\vspace{0.2cm}
\subsubsection*{Ejercicio 123 - Diferencia de cuadrados}
\textbf{Enunciado:} $ 25 x^{4} - 16 y^{2} $\par
\begin{enumerate}[leftmargin=2em]
\item La diferencia de cuadrados siempre tiene esta forma: \(A^2-B^2\).
\item Aqui identificamos \(A=5 x^{2}\) y \(B=4 y\).
\item Aplicamos la formula: \(A^2-B^2=(A-B)(A+B)\).
\item Resultado: \[\boxed{\left(\left(5 x^{2}\right)-\left(4 y\right)\right)\left(\left(5 x^{2}\right)+\left(4 y\right)\right)}\]
\end{enumerate}

\vspace{0.2cm}
\subsubsection*{Ejercicio 124 - Diferencia de cuadrados}
\textbf{Enunciado:} $ 4 x^{2} - 25 y^{2} $\par
\begin{enumerate}[leftmargin=2em]
\item La diferencia de cuadrados siempre tiene esta forma: \(A^2-B^2\).
\item Aqui identificamos \(A=2 x\) y \(B=5 y\).
\item Aplicamos la formula: \(A^2-B^2=(A-B)(A+B)\).
\item Resultado: \[\boxed{\left(\left(2 x\right)-\left(5 y\right)\right)\left(\left(2 x\right)+\left(5 y\right)\right)}\]
\end{enumerate}

\vspace{0.2cm}
\subsubsection*{Ejercicio 125 - Diferencia de cuadrados}
\textbf{Enunciado:} $ 9 x^{4} - y^{2} $\par
\begin{enumerate}[leftmargin=2em]
\item La diferencia de cuadrados siempre tiene esta forma: \(A^2-B^2\).
\item Aqui identificamos \(A=3 x^{2}\) y \(B=y\).
\item Aplicamos la formula: \(A^2-B^2=(A-B)(A+B)\).
\item Resultado: \[\boxed{\left(\left(3 x^{2}\right)-\left(y\right)\right)\left(\left(3 x^{2}\right)+\left(y\right)\right)}\]
\end{enumerate}

\vspace{0.2cm}
\subsubsection*{Ejercicio 126 - Diferencia de cuadrados}
\textbf{Enunciado:} $ 16 x^{2} - 4 y^{2} $\par
\begin{enumerate}[leftmargin=2em]
\item La diferencia de cuadrados siempre tiene esta forma: \(A^2-B^2\).
\item Aqui identificamos \(A=4 x\) y \(B=2 y\).
\item Aplicamos la formula: \(A^2-B^2=(A-B)(A+B)\).
\item Resultado: \[\boxed{\left(\left(4 x\right)-\left(2 y\right)\right)\left(\left(4 x\right)+\left(2 y\right)\right)}\]
\end{enumerate}

\vspace{0.2cm}
\subsubsection*{Ejercicio 127 - Diferencia de cuadrados}
\textbf{Enunciado:} $ 25 x^{4} - 9 y^{2} $\par
\begin{enumerate}[leftmargin=2em]
\item La diferencia de cuadrados siempre tiene esta forma: \(A^2-B^2\).
\item Aqui identificamos \(A=5 x^{2}\) y \(B=3 y\).
\item Aplicamos la formula: \(A^2-B^2=(A-B)(A+B)\).
\item Resultado: \[\boxed{\left(\left(5 x^{2}\right)-\left(3 y\right)\right)\left(\left(5 x^{2}\right)+\left(3 y\right)\right)}\]
\end{enumerate}

\vspace{0.2cm}
\subsubsection*{Ejercicio 128 - Diferencia de cuadrados}
\textbf{Enunciado:} $ 4 x^{2} - 16 y^{2} $\par
\begin{enumerate}[leftmargin=2em]
\item La diferencia de cuadrados siempre tiene esta forma: \(A^2-B^2\).
\item Aqui identificamos \(A=2 x\) y \(B=4 y\).
\item Aplicamos la formula: \(A^2-B^2=(A-B)(A+B)\).
\item Resultado: \[\boxed{\left(\left(2 x\right)-\left(4 y\right)\right)\left(\left(2 x\right)+\left(4 y\right)\right)}\]
\end{enumerate}

\vspace{0.2cm}
\subsubsection*{Ejercicio 129 - Diferencia de cuadrados}
\textbf{Enunciado:} $ 9 x^{4} - 25 y^{2} $\par
\begin{enumerate}[leftmargin=2em]
\item La diferencia de cuadrados siempre tiene esta forma: \(A^2-B^2\).
\item Aqui identificamos \(A=3 x^{2}\) y \(B=5 y\).
\item Aplicamos la formula: \(A^2-B^2=(A-B)(A+B)\).
\item Resultado: \[\boxed{\left(\left(3 x^{2}\right)-\left(5 y\right)\right)\left(\left(3 x^{2}\right)+\left(5 y\right)\right)}\]
\end{enumerate}

\vspace{0.2cm}
\subsubsection*{Ejercicio 130 - Diferencia de cuadrados}
\textbf{Enunciado:} $ 16 x^{2} - y^{2} $\par
\begin{enumerate}[leftmargin=2em]
\item La diferencia de cuadrados siempre tiene esta forma: \(A^2-B^2\).
\item Aqui identificamos \(A=4 x\) y \(B=y\).
\item Aplicamos la formula: \(A^2-B^2=(A-B)(A+B)\).
\item Resultado: \[\boxed{\left(\left(4 x\right)-\left(y\right)\right)\left(\left(4 x\right)+\left(y\right)\right)}\]
\end{enumerate}

\vspace{0.2cm}
\subsubsection*{Ejercicio 131 - Diferencia de cuadrados}
\textbf{Enunciado:} $ 25 x^{4} - 4 y^{2} $\par
\begin{enumerate}[leftmargin=2em]
\item La diferencia de cuadrados siempre tiene esta forma: \(A^2-B^2\).
\item Aqui identificamos \(A=5 x^{2}\) y \(B=2 y\).
\item Aplicamos la formula: \(A^2-B^2=(A-B)(A+B)\).
\item Resultado: \[\boxed{\left(\left(5 x^{2}\right)-\left(2 y\right)\right)\left(\left(5 x^{2}\right)+\left(2 y\right)\right)}\]
\end{enumerate}

\vspace{0.2cm}
\subsubsection*{Ejercicio 132 - Diferencia de cuadrados}
\textbf{Enunciado:} $ 4 x^{2} - 9 y^{2} $\par
\begin{enumerate}[leftmargin=2em]
\item La diferencia de cuadrados siempre tiene esta forma: \(A^2-B^2\).
\item Aqui identificamos \(A=2 x\) y \(B=3 y\).
\item Aplicamos la formula: \(A^2-B^2=(A-B)(A+B)\).
\item Resultado: \[\boxed{\left(\left(2 x\right)-\left(3 y\right)\right)\left(\left(2 x\right)+\left(3 y\right)\right)}\]
\end{enumerate}

\vspace{0.2cm}
\subsubsection*{Ejercicio 133 - Diferencia de cuadrados}
\textbf{Enunciado:} $ 9 x^{4} - 16 y^{2} $\par
\begin{enumerate}[leftmargin=2em]
\item La diferencia de cuadrados siempre tiene esta forma: \(A^2-B^2\).
\item Aqui identificamos \(A=3 x^{2}\) y \(B=4 y\).
\item Aplicamos la formula: \(A^2-B^2=(A-B)(A+B)\).
\item Resultado: \[\boxed{\left(\left(3 x^{2}\right)-\left(4 y\right)\right)\left(\left(3 x^{2}\right)+\left(4 y\right)\right)}\]
\end{enumerate}

\vspace{0.2cm}
\subsubsection*{Ejercicio 134 - Diferencia de cuadrados}
\textbf{Enunciado:} $ 16 x^{2} - 25 y^{2} $\par
\begin{enumerate}[leftmargin=2em]
\item La diferencia de cuadrados siempre tiene esta forma: \(A^2-B^2\).
\item Aqui identificamos \(A=4 x\) y \(B=5 y\).
\item Aplicamos la formula: \(A^2-B^2=(A-B)(A+B)\).
\item Resultado: \[\boxed{\left(\left(4 x\right)-\left(5 y\right)\right)\left(\left(4 x\right)+\left(5 y\right)\right)}\]
\end{enumerate}

\vspace{0.2cm}
\subsubsection*{Ejercicio 135 - Diferencia de cuadrados}
\textbf{Enunciado:} $ 25 x^{4} - y^{2} $\par
\begin{enumerate}[leftmargin=2em]
\item La diferencia de cuadrados siempre tiene esta forma: \(A^2-B^2\).
\item Aqui identificamos \(A=5 x^{2}\) y \(B=y\).
\item Aplicamos la formula: \(A^2-B^2=(A-B)(A+B)\).
\item Resultado: \[\boxed{\left(\left(5 x^{2}\right)-\left(y\right)\right)\left(\left(5 x^{2}\right)+\left(y\right)\right)}\]
\end{enumerate}

\vspace{0.2cm}
\subsubsection*{Ejercicio 136 - Diferencia de cuadrados}
\textbf{Enunciado:} $ 4 x^{2} - 4 y^{2} $\par
\begin{enumerate}[leftmargin=2em]
\item La diferencia de cuadrados siempre tiene esta forma: \(A^2-B^2\).
\item Aqui identificamos \(A=2 x\) y \(B=2 y\).
\item Aplicamos la formula: \(A^2-B^2=(A-B)(A+B)\).
\item Resultado: \[\boxed{\left(\left(2 x\right)-\left(2 y\right)\right)\left(\left(2 x\right)+\left(2 y\right)\right)}\]
\end{enumerate}

\vspace{0.2cm}
\subsubsection*{Ejercicio 137 - Diferencia de cuadrados}
\textbf{Enunciado:} $ 9 x^{4} - 9 y^{2} $\par
\begin{enumerate}[leftmargin=2em]
\item La diferencia de cuadrados siempre tiene esta forma: \(A^2-B^2\).
\item Aqui identificamos \(A=3 x^{2}\) y \(B=3 y\).
\item Aplicamos la formula: \(A^2-B^2=(A-B)(A+B)\).
\item Resultado: \[\boxed{\left(\left(3 x^{2}\right)-\left(3 y\right)\right)\left(\left(3 x^{2}\right)+\left(3 y\right)\right)}\]
\end{enumerate}

\vspace{0.2cm}
\subsubsection*{Ejercicio 138 - Diferencia de cuadrados}
\textbf{Enunciado:} $ 16 x^{2} - 16 y^{2} $\par
\begin{enumerate}[leftmargin=2em]
\item La diferencia de cuadrados siempre tiene esta forma: \(A^2-B^2\).
\item Aqui identificamos \(A=4 x\) y \(B=4 y\).
\item Aplicamos la formula: \(A^2-B^2=(A-B)(A+B)\).
\item Resultado: \[\boxed{\left(\left(4 x\right)-\left(4 y\right)\right)\left(\left(4 x\right)+\left(4 y\right)\right)}\]
\end{enumerate}

\vspace{0.2cm}
\subsubsection*{Ejercicio 139 - Diferencia de cuadrados}
\textbf{Enunciado:} $ 25 x^{4} - 25 y^{2} $\par
\begin{enumerate}[leftmargin=2em]
\item La diferencia de cuadrados siempre tiene esta forma: \(A^2-B^2\).
\item Aqui identificamos \(A=5 x^{2}\) y \(B=5 y\).
\item Aplicamos la formula: \(A^2-B^2=(A-B)(A+B)\).
\item Resultado: \[\boxed{\left(\left(5 x^{2}\right)-\left(5 y\right)\right)\left(\left(5 x^{2}\right)+\left(5 y\right)\right)}\]
\end{enumerate}

\vspace{0.2cm}
\subsubsection*{Ejercicio 140 - Diferencia de cuadrados}
\textbf{Enunciado:} $ 4 x^{2} - y^{2} $\par
\begin{enumerate}[leftmargin=2em]
\item La diferencia de cuadrados siempre tiene esta forma: \(A^2-B^2\).
\item Aqui identificamos \(A=2 x\) y \(B=y\).
\item Aplicamos la formula: \(A^2-B^2=(A-B)(A+B)\).
\item Resultado: \[\boxed{\left(\left(2 x\right)-\left(y\right)\right)\left(\left(2 x\right)+\left(y\right)\right)}\]
\end{enumerate}

\vspace{0.2cm}
\nivel{3}{intermedio}
\subsubsection*{Ejercicio 141 - Diferencia de cuadrados}
\textbf{Enunciado:} $ 4 x^{2} + 8 x - 9 y^{2} + 4 $\par
\begin{enumerate}[leftmargin=2em]
\item La diferencia de cuadrados siempre tiene esta forma: \(A^2-B^2\).
\item Aqui identificamos \(A=2 x + 2\) y \(B=3 y\).
\item Aplicamos la formula: \(A^2-B^2=(A-B)(A+B)\).
\item Resultado: \[\boxed{\left(\left(2 x + 2\right)-\left(3 y\right)\right)\left(\left(2 x + 2\right)+\left(3 y\right)\right)}\]
\end{enumerate}

\vspace{0.2cm}
\subsubsection*{Ejercicio 142 - Diferencia de cuadrados}
\textbf{Enunciado:} $ 9 x^{2} + 18 x - 16 y^{2} + 9 $\par
\begin{enumerate}[leftmargin=2em]
\item La diferencia de cuadrados siempre tiene esta forma: \(A^2-B^2\).
\item Aqui identificamos \(A=3 x + 3\) y \(B=4 y\).
\item Aplicamos la formula: \(A^2-B^2=(A-B)(A+B)\).
\item Resultado: \[\boxed{\left(\left(3 x + 3\right)-\left(4 y\right)\right)\left(\left(3 x + 3\right)+\left(4 y\right)\right)}\]
\end{enumerate}

\vspace{0.2cm}
\subsubsection*{Ejercicio 143 - Diferencia de cuadrados}
\textbf{Enunciado:} $ x^{2} + 8 x - 25 y^{2} + 16 $\par
\begin{enumerate}[leftmargin=2em]
\item La diferencia de cuadrados siempre tiene esta forma: \(A^2-B^2\).
\item Aqui identificamos \(A=x + 4\) y \(B=5 y\).
\item Aplicamos la formula: \(A^2-B^2=(A-B)(A+B)\).
\item Resultado: \[\boxed{\left(\left(x + 4\right)-\left(5 y\right)\right)\left(\left(x + 4\right)+\left(5 y\right)\right)}\]
\end{enumerate}

\vspace{0.2cm}
\subsubsection*{Ejercicio 144 - Diferencia de cuadrados}
\textbf{Enunciado:} $ 4 x^{2} + 4 x - 36 y^{2} + 1 $\par
\begin{enumerate}[leftmargin=2em]
\item La diferencia de cuadrados siempre tiene esta forma: \(A^2-B^2\).
\item Aqui identificamos \(A=2 x + 1\) y \(B=6 y\).
\item Aplicamos la formula: \(A^2-B^2=(A-B)(A+B)\).
\item Resultado: \[\boxed{\left(\left(2 x + 1\right)-\left(6 y\right)\right)\left(\left(2 x + 1\right)+\left(6 y\right)\right)}\]
\end{enumerate}

\vspace{0.2cm}
\subsubsection*{Ejercicio 145 - Diferencia de cuadrados}
\textbf{Enunciado:} $ 9 x^{2} + 12 x - 4 y^{2} + 4 $\par
\begin{enumerate}[leftmargin=2em]
\item La diferencia de cuadrados siempre tiene esta forma: \(A^2-B^2\).
\item Aqui identificamos \(A=3 x + 2\) y \(B=2 y\).
\item Aplicamos la formula: \(A^2-B^2=(A-B)(A+B)\).
\item Resultado: \[\boxed{\left(\left(3 x + 2\right)-\left(2 y\right)\right)\left(\left(3 x + 2\right)+\left(2 y\right)\right)}\]
\end{enumerate}

\vspace{0.2cm}
\subsubsection*{Ejercicio 146 - Diferencia de cuadrados}
\textbf{Enunciado:} $ x^{2} + 6 x - 9 y^{2} + 9 $\par
\begin{enumerate}[leftmargin=2em]
\item La diferencia de cuadrados siempre tiene esta forma: \(A^2-B^2\).
\item Aqui identificamos \(A=x + 3\) y \(B=3 y\).
\item Aplicamos la formula: \(A^2-B^2=(A-B)(A+B)\).
\item Resultado: \[\boxed{\left(\left(x + 3\right)-\left(3 y\right)\right)\left(\left(x + 3\right)+\left(3 y\right)\right)}\]
\end{enumerate}

\vspace{0.2cm}
\subsubsection*{Ejercicio 147 - Diferencia de cuadrados}
\textbf{Enunciado:} $ 4 x^{2} + 16 x - 16 y^{2} + 16 $\par
\begin{enumerate}[leftmargin=2em]
\item La diferencia de cuadrados siempre tiene esta forma: \(A^2-B^2\).
\item Aqui identificamos \(A=2 x + 4\) y \(B=4 y\).
\item Aplicamos la formula: \(A^2-B^2=(A-B)(A+B)\).
\item Resultado: \[\boxed{\left(\left(2 x + 4\right)-\left(4 y\right)\right)\left(\left(2 x + 4\right)+\left(4 y\right)\right)}\]
\end{enumerate}

\vspace{0.2cm}
\subsubsection*{Ejercicio 148 - Diferencia de cuadrados}
\textbf{Enunciado:} $ 9 x^{2} + 6 x - 25 y^{2} + 1 $\par
\begin{enumerate}[leftmargin=2em]
\item La diferencia de cuadrados siempre tiene esta forma: \(A^2-B^2\).
\item Aqui identificamos \(A=3 x + 1\) y \(B=5 y\).
\item Aplicamos la formula: \(A^2-B^2=(A-B)(A+B)\).
\item Resultado: \[\boxed{\left(\left(3 x + 1\right)-\left(5 y\right)\right)\left(\left(3 x + 1\right)+\left(5 y\right)\right)}\]
\end{enumerate}

\vspace{0.2cm}
\subsubsection*{Ejercicio 149 - Diferencia de cuadrados}
\textbf{Enunciado:} $ x^{2} + 4 x - 36 y^{2} + 4 $\par
\begin{enumerate}[leftmargin=2em]
\item La diferencia de cuadrados siempre tiene esta forma: \(A^2-B^2\).
\item Aqui identificamos \(A=x + 2\) y \(B=6 y\).
\item Aplicamos la formula: \(A^2-B^2=(A-B)(A+B)\).
\item Resultado: \[\boxed{\left(\left(x + 2\right)-\left(6 y\right)\right)\left(\left(x + 2\right)+\left(6 y\right)\right)}\]
\end{enumerate}

\vspace{0.2cm}
\subsubsection*{Ejercicio 150 - Diferencia de cuadrados}
\textbf{Enunciado:} $ 4 x^{2} + 12 x - 4 y^{2} + 9 $\par
\begin{enumerate}[leftmargin=2em]
\item La diferencia de cuadrados siempre tiene esta forma: \(A^2-B^2\).
\item Aqui identificamos \(A=2 x + 3\) y \(B=2 y\).
\item Aplicamos la formula: \(A^2-B^2=(A-B)(A+B)\).
\item Resultado: \[\boxed{\left(\left(2 x + 3\right)-\left(2 y\right)\right)\left(\left(2 x + 3\right)+\left(2 y\right)\right)}\]
\end{enumerate}

\vspace{0.2cm}
\subsubsection*{Ejercicio 151 - Diferencia de cuadrados}
\textbf{Enunciado:} $ 9 x^{2} + 24 x - 9 y^{2} + 16 $\par
\begin{enumerate}[leftmargin=2em]
\item La diferencia de cuadrados siempre tiene esta forma: \(A^2-B^2\).
\item Aqui identificamos \(A=3 x + 4\) y \(B=3 y\).
\item Aplicamos la formula: \(A^2-B^2=(A-B)(A+B)\).
\item Resultado: \[\boxed{\left(\left(3 x + 4\right)-\left(3 y\right)\right)\left(\left(3 x + 4\right)+\left(3 y\right)\right)}\]
\end{enumerate}

\vspace{0.2cm}
\subsubsection*{Ejercicio 152 - Diferencia de cuadrados}
\textbf{Enunciado:} $ x^{2} + 2 x - 16 y^{2} + 1 $\par
\begin{enumerate}[leftmargin=2em]
\item La diferencia de cuadrados siempre tiene esta forma: \(A^2-B^2\).
\item Aqui identificamos \(A=x + 1\) y \(B=4 y\).
\item Aplicamos la formula: \(A^2-B^2=(A-B)(A+B)\).
\item Resultado: \[\boxed{\left(\left(x + 1\right)-\left(4 y\right)\right)\left(\left(x + 1\right)+\left(4 y\right)\right)}\]
\end{enumerate}

\vspace{0.2cm}
\subsubsection*{Ejercicio 153 - Diferencia de cuadrados}
\textbf{Enunciado:} $ 4 x^{2} + 8 x - 25 y^{2} + 4 $\par
\begin{enumerate}[leftmargin=2em]
\item La diferencia de cuadrados siempre tiene esta forma: \(A^2-B^2\).
\item Aqui identificamos \(A=2 x + 2\) y \(B=5 y\).
\item Aplicamos la formula: \(A^2-B^2=(A-B)(A+B)\).
\item Resultado: \[\boxed{\left(\left(2 x + 2\right)-\left(5 y\right)\right)\left(\left(2 x + 2\right)+\left(5 y\right)\right)}\]
\end{enumerate}

\vspace{0.2cm}
\subsubsection*{Ejercicio 154 - Diferencia de cuadrados}
\textbf{Enunciado:} $ 9 x^{2} + 18 x - 36 y^{2} + 9 $\par
\begin{enumerate}[leftmargin=2em]
\item La diferencia de cuadrados siempre tiene esta forma: \(A^2-B^2\).
\item Aqui identificamos \(A=3 x + 3\) y \(B=6 y\).
\item Aplicamos la formula: \(A^2-B^2=(A-B)(A+B)\).
\item Resultado: \[\boxed{\left(\left(3 x + 3\right)-\left(6 y\right)\right)\left(\left(3 x + 3\right)+\left(6 y\right)\right)}\]
\end{enumerate}

\vspace{0.2cm}
\subsubsection*{Ejercicio 155 - Diferencia de cuadrados}
\textbf{Enunciado:} $ x^{2} + 8 x - 4 y^{2} + 16 $\par
\begin{enumerate}[leftmargin=2em]
\item La diferencia de cuadrados siempre tiene esta forma: \(A^2-B^2\).
\item Aqui identificamos \(A=x + 4\) y \(B=2 y\).
\item Aplicamos la formula: \(A^2-B^2=(A-B)(A+B)\).
\item Resultado: \[\boxed{\left(\left(x + 4\right)-\left(2 y\right)\right)\left(\left(x + 4\right)+\left(2 y\right)\right)}\]
\end{enumerate}

\vspace{0.2cm}
\subsubsection*{Ejercicio 156 - Diferencia de cuadrados}
\textbf{Enunciado:} $ 4 x^{2} + 4 x - 9 y^{2} + 1 $\par
\begin{enumerate}[leftmargin=2em]
\item La diferencia de cuadrados siempre tiene esta forma: \(A^2-B^2\).
\item Aqui identificamos \(A=2 x + 1\) y \(B=3 y\).
\item Aplicamos la formula: \(A^2-B^2=(A-B)(A+B)\).
\item Resultado: \[\boxed{\left(\left(2 x + 1\right)-\left(3 y\right)\right)\left(\left(2 x + 1\right)+\left(3 y\right)\right)}\]
\end{enumerate}

\vspace{0.2cm}
\subsubsection*{Ejercicio 157 - Diferencia de cuadrados}
\textbf{Enunciado:} $ 9 x^{2} + 12 x - 16 y^{2} + 4 $\par
\begin{enumerate}[leftmargin=2em]
\item La diferencia de cuadrados siempre tiene esta forma: \(A^2-B^2\).
\item Aqui identificamos \(A=3 x + 2\) y \(B=4 y\).
\item Aplicamos la formula: \(A^2-B^2=(A-B)(A+B)\).
\item Resultado: \[\boxed{\left(\left(3 x + 2\right)-\left(4 y\right)\right)\left(\left(3 x + 2\right)+\left(4 y\right)\right)}\]
\end{enumerate}

\vspace{0.2cm}
\subsubsection*{Ejercicio 158 - Diferencia de cuadrados}
\textbf{Enunciado:} $ x^{2} + 6 x - 25 y^{2} + 9 $\par
\begin{enumerate}[leftmargin=2em]
\item La diferencia de cuadrados siempre tiene esta forma: \(A^2-B^2\).
\item Aqui identificamos \(A=x + 3\) y \(B=5 y\).
\item Aplicamos la formula: \(A^2-B^2=(A-B)(A+B)\).
\item Resultado: \[\boxed{\left(\left(x + 3\right)-\left(5 y\right)\right)\left(\left(x + 3\right)+\left(5 y\right)\right)}\]
\end{enumerate}

\vspace{0.2cm}
\subsubsection*{Ejercicio 159 - Diferencia de cuadrados}
\textbf{Enunciado:} $ 4 x^{2} + 16 x - 36 y^{2} + 16 $\par
\begin{enumerate}[leftmargin=2em]
\item La diferencia de cuadrados siempre tiene esta forma: \(A^2-B^2\).
\item Aqui identificamos \(A=2 x + 4\) y \(B=6 y\).
\item Aplicamos la formula: \(A^2-B^2=(A-B)(A+B)\).
\item Resultado: \[\boxed{\left(\left(2 x + 4\right)-\left(6 y\right)\right)\left(\left(2 x + 4\right)+\left(6 y\right)\right)}\]
\end{enumerate}

\vspace{0.2cm}
\subsubsection*{Ejercicio 160 - Diferencia de cuadrados}
\textbf{Enunciado:} $ 9 x^{2} + 6 x - 4 y^{2} + 1 $\par
\begin{enumerate}[leftmargin=2em]
\item La diferencia de cuadrados siempre tiene esta forma: \(A^2-B^2\).
\item Aqui identificamos \(A=3 x + 1\) y \(B=2 y\).
\item Aplicamos la formula: \(A^2-B^2=(A-B)(A+B)\).
\item Resultado: \[\boxed{\left(\left(3 x + 1\right)-\left(2 y\right)\right)\left(\left(3 x + 1\right)+\left(2 y\right)\right)}\]
\end{enumerate}

\vspace{0.2cm}
\nivel{4}{avanzado}
\subsubsection*{Ejercicio 161 - Diferencia de cuadrados}
\textbf{Enunciado:} $ 5 x^{2} - 12 x y - 12 x + 4 y^{2} - 9 $\par
\begin{enumerate}[leftmargin=2em]
\item La diferencia de cuadrados siempre tiene esta forma: \(A^2-B^2\).
\item Aqui identificamos \(A=3 x - 2 y\) y \(B=2 x + 3\).
\item Aplicamos la formula: \(A^2-B^2=(A-B)(A+B)\).
\item Resultado: \[\boxed{\left(\left(3 x - 2 y\right)-\left(2 x + 3\right)\right)\left(\left(3 x - 2 y\right)+\left(2 x + 3\right)\right)}\]
\end{enumerate}

\vspace{0.2cm}
\subsubsection*{Ejercicio 162 - Diferencia de cuadrados}
\textbf{Enunciado:} $ 7 x^{2} - 24 x y - 12 x + 9 y^{2} - 4 $\par
\begin{enumerate}[leftmargin=2em]
\item La diferencia de cuadrados siempre tiene esta forma: \(A^2-B^2\).
\item Aqui identificamos \(A=4 x - 3 y\) y \(B=3 x + 2\).
\item Aplicamos la formula: \(A^2-B^2=(A-B)(A+B)\).
\item Resultado: \[\boxed{\left(\left(4 x - 3 y\right)-\left(3 x + 2\right)\right)\left(\left(4 x - 3 y\right)+\left(3 x + 2\right)\right)}\]
\end{enumerate}

\vspace{0.2cm}
\subsubsection*{Ejercicio 163 - Diferencia de cuadrados}
\textbf{Enunciado:} $ 9 x^{2} - 10 x y - 24 x + y^{2} - 9 $\par
\begin{enumerate}[leftmargin=2em]
\item La diferencia de cuadrados siempre tiene esta forma: \(A^2-B^2\).
\item Aqui identificamos \(A=5 x - y\) y \(B=4 x + 3\).
\item Aplicamos la formula: \(A^2-B^2=(A-B)(A+B)\).
\item Resultado: \[\boxed{\left(\left(5 x - y\right)-\left(4 x + 3\right)\right)\left(\left(5 x - y\right)+\left(4 x + 3\right)\right)}\]
\end{enumerate}

\vspace{0.2cm}
\subsubsection*{Ejercicio 164 - Diferencia de cuadrados}
\textbf{Enunciado:} $ - 21 x^{2} - 8 x y - 20 x + 4 y^{2} - 4 $\par
\begin{enumerate}[leftmargin=2em]
\item La diferencia de cuadrados siempre tiene esta forma: \(A^2-B^2\).
\item Aqui identificamos \(A=2 x - 2 y\) y \(B=5 x + 2\).
\item Aplicamos la formula: \(A^2-B^2=(A-B)(A+B)\).
\item Resultado: \[\boxed{\left(\left(2 x - 2 y\right)-\left(5 x + 2\right)\right)\left(\left(2 x - 2 y\right)+\left(5 x + 2\right)\right)}\]
\end{enumerate}

\vspace{0.2cm}
\subsubsection*{Ejercicio 165 - Diferencia de cuadrados}
\textbf{Enunciado:} $ 8 x^{2} - 18 x y - 6 x + 9 y^{2} - 9 $\par
\begin{enumerate}[leftmargin=2em]
\item La diferencia de cuadrados siempre tiene esta forma: \(A^2-B^2\).
\item Aqui identificamos \(A=3 x - 3 y\) y \(B=x + 3\).
\item Aplicamos la formula: \(A^2-B^2=(A-B)(A+B)\).
\item Resultado: \[\boxed{\left(\left(3 x - 3 y\right)-\left(x + 3\right)\right)\left(\left(3 x - 3 y\right)+\left(x + 3\right)\right)}\]
\end{enumerate}

\vspace{0.2cm}
\subsubsection*{Ejercicio 166 - Diferencia de cuadrados}
\textbf{Enunciado:} $ 12 x^{2} - 8 x y - 8 x + y^{2} - 4 $\par
\begin{enumerate}[leftmargin=2em]
\item La diferencia de cuadrados siempre tiene esta forma: \(A^2-B^2\).
\item Aqui identificamos \(A=4 x - y\) y \(B=2 x + 2\).
\item Aplicamos la formula: \(A^2-B^2=(A-B)(A+B)\).
\item Resultado: \[\boxed{\left(\left(4 x - y\right)-\left(2 x + 2\right)\right)\left(\left(4 x - y\right)+\left(2 x + 2\right)\right)}\]
\end{enumerate}

\vspace{0.2cm}
\subsubsection*{Ejercicio 167 - Diferencia de cuadrados}
\textbf{Enunciado:} $ 16 x^{2} - 20 x y - 18 x + 4 y^{2} - 9 $\par
\begin{enumerate}[leftmargin=2em]
\item La diferencia de cuadrados siempre tiene esta forma: \(A^2-B^2\).
\item Aqui identificamos \(A=5 x - 2 y\) y \(B=3 x + 3\).
\item Aplicamos la formula: \(A^2-B^2=(A-B)(A+B)\).
\item Resultado: \[\boxed{\left(\left(5 x - 2 y\right)-\left(3 x + 3\right)\right)\left(\left(5 x - 2 y\right)+\left(3 x + 3\right)\right)}\]
\end{enumerate}

\vspace{0.2cm}
\subsubsection*{Ejercicio 168 - Diferencia de cuadrados}
\textbf{Enunciado:} $ - 12 x^{2} - 12 x y - 16 x + 9 y^{2} - 4 $\par
\begin{enumerate}[leftmargin=2em]
\item La diferencia de cuadrados siempre tiene esta forma: \(A^2-B^2\).
\item Aqui identificamos \(A=2 x - 3 y\) y \(B=4 x + 2\).
\item Aplicamos la formula: \(A^2-B^2=(A-B)(A+B)\).
\item Resultado: \[\boxed{\left(\left(2 x - 3 y\right)-\left(4 x + 2\right)\right)\left(\left(2 x - 3 y\right)+\left(4 x + 2\right)\right)}\]
\end{enumerate}

\vspace{0.2cm}
\subsubsection*{Ejercicio 169 - Diferencia de cuadrados}
\textbf{Enunciado:} $ - 16 x^{2} - 6 x y - 30 x + y^{2} - 9 $\par
\begin{enumerate}[leftmargin=2em]
\item La diferencia de cuadrados siempre tiene esta forma: \(A^2-B^2\).
\item Aqui identificamos \(A=3 x - y\) y \(B=5 x + 3\).
\item Aplicamos la formula: \(A^2-B^2=(A-B)(A+B)\).
\item Resultado: \[\boxed{\left(\left(3 x - y\right)-\left(5 x + 3\right)\right)\left(\left(3 x - y\right)+\left(5 x + 3\right)\right)}\]
\end{enumerate}

\vspace{0.2cm}
\subsubsection*{Ejercicio 170 - Diferencia de cuadrados}
\textbf{Enunciado:} $ 15 x^{2} - 16 x y - 4 x + 4 y^{2} - 4 $\par
\begin{enumerate}[leftmargin=2em]
\item La diferencia de cuadrados siempre tiene esta forma: \(A^2-B^2\).
\item Aqui identificamos \(A=4 x - 2 y\) y \(B=x + 2\).
\item Aplicamos la formula: \(A^2-B^2=(A-B)(A+B)\).
\item Resultado: \[\boxed{\left(\left(4 x - 2 y\right)-\left(x + 2\right)\right)\left(\left(4 x - 2 y\right)+\left(x + 2\right)\right)}\]
\end{enumerate}

\vspace{0.2cm}
\subsubsection*{Ejercicio 171 - Diferencia de cuadrados}
\textbf{Enunciado:} $ 21 x^{2} - 30 x y - 12 x + 9 y^{2} - 9 $\par
\begin{enumerate}[leftmargin=2em]
\item La diferencia de cuadrados siempre tiene esta forma: \(A^2-B^2\).
\item Aqui identificamos \(A=5 x - 3 y\) y \(B=2 x + 3\).
\item Aplicamos la formula: \(A^2-B^2=(A-B)(A+B)\).
\item Resultado: \[\boxed{\left(\left(5 x - 3 y\right)-\left(2 x + 3\right)\right)\left(\left(5 x - 3 y\right)+\left(2 x + 3\right)\right)}\]
\end{enumerate}

\vspace{0.2cm}
\subsubsection*{Ejercicio 172 - Diferencia de cuadrados}
\textbf{Enunciado:} $ - 5 x^{2} - 4 x y - 12 x + y^{2} - 4 $\par
\begin{enumerate}[leftmargin=2em]
\item La diferencia de cuadrados siempre tiene esta forma: \(A^2-B^2\).
\item Aqui identificamos \(A=2 x - y\) y \(B=3 x + 2\).
\item Aplicamos la formula: \(A^2-B^2=(A-B)(A+B)\).
\item Resultado: \[\boxed{\left(\left(2 x - y\right)-\left(3 x + 2\right)\right)\left(\left(2 x - y\right)+\left(3 x + 2\right)\right)}\]
\end{enumerate}

\vspace{0.2cm}
\subsubsection*{Ejercicio 173 - Diferencia de cuadrados}
\textbf{Enunciado:} $ - 7 x^{2} - 12 x y - 24 x + 4 y^{2} - 9 $\par
\begin{enumerate}[leftmargin=2em]
\item La diferencia de cuadrados siempre tiene esta forma: \(A^2-B^2\).
\item Aqui identificamos \(A=3 x - 2 y\) y \(B=4 x + 3\).
\item Aplicamos la formula: \(A^2-B^2=(A-B)(A+B)\).
\item Resultado: \[\boxed{\left(\left(3 x - 2 y\right)-\left(4 x + 3\right)\right)\left(\left(3 x - 2 y\right)+\left(4 x + 3\right)\right)}\]
\end{enumerate}

\vspace{0.2cm}
\subsubsection*{Ejercicio 174 - Diferencia de cuadrados}
\textbf{Enunciado:} $ - 9 x^{2} - 24 x y - 20 x + 9 y^{2} - 4 $\par
\begin{enumerate}[leftmargin=2em]
\item La diferencia de cuadrados siempre tiene esta forma: \(A^2-B^2\).
\item Aqui identificamos \(A=4 x - 3 y\) y \(B=5 x + 2\).
\item Aplicamos la formula: \(A^2-B^2=(A-B)(A+B)\).
\item Resultado: \[\boxed{\left(\left(4 x - 3 y\right)-\left(5 x + 2\right)\right)\left(\left(4 x - 3 y\right)+\left(5 x + 2\right)\right)}\]
\end{enumerate}

\vspace{0.2cm}
\subsubsection*{Ejercicio 175 - Diferencia de cuadrados}
\textbf{Enunciado:} $ 24 x^{2} - 10 x y - 6 x + y^{2} - 9 $\par
\begin{enumerate}[leftmargin=2em]
\item La diferencia de cuadrados siempre tiene esta forma: \(A^2-B^2\).
\item Aqui identificamos \(A=5 x - y\) y \(B=x + 3\).
\item Aplicamos la formula: \(A^2-B^2=(A-B)(A+B)\).
\item Resultado: \[\boxed{\left(\left(5 x - y\right)-\left(x + 3\right)\right)\left(\left(5 x - y\right)+\left(x + 3\right)\right)}\]
\end{enumerate}

\vspace{0.2cm}
\subsubsection*{Ejercicio 176 - Diferencia de cuadrados}
\textbf{Enunciado:} $ - 8 x y - 8 x + 4 y^{2} - 4 $\par
\begin{enumerate}[leftmargin=2em]
\item La diferencia de cuadrados siempre tiene esta forma: \(A^2-B^2\).
\item Aqui identificamos \(A=2 x - 2 y\) y \(B=2 x + 2\).
\item Aplicamos la formula: \(A^2-B^2=(A-B)(A+B)\).
\item Resultado: \[\boxed{\left(\left(2 x - 2 y\right)-\left(2 x + 2\right)\right)\left(\left(2 x - 2 y\right)+\left(2 x + 2\right)\right)}\]
\end{enumerate}

\vspace{0.2cm}
\subsubsection*{Ejercicio 177 - Diferencia de cuadrados}
\textbf{Enunciado:} $ - 18 x y - 18 x + 9 y^{2} - 9 $\par
\begin{enumerate}[leftmargin=2em]
\item La diferencia de cuadrados siempre tiene esta forma: \(A^2-B^2\).
\item Aqui identificamos \(A=3 x - 3 y\) y \(B=3 x + 3\).
\item Aplicamos la formula: \(A^2-B^2=(A-B)(A+B)\).
\item Resultado: \[\boxed{\left(\left(3 x - 3 y\right)-\left(3 x + 3\right)\right)\left(\left(3 x - 3 y\right)+\left(3 x + 3\right)\right)}\]
\end{enumerate}

\vspace{0.2cm}
\subsubsection*{Ejercicio 178 - Diferencia de cuadrados}
\textbf{Enunciado:} $ - 8 x y - 16 x + y^{2} - 4 $\par
\begin{enumerate}[leftmargin=2em]
\item La diferencia de cuadrados siempre tiene esta forma: \(A^2-B^2\).
\item Aqui identificamos \(A=4 x - y\) y \(B=4 x + 2\).
\item Aplicamos la formula: \(A^2-B^2=(A-B)(A+B)\).
\item Resultado: \[\boxed{\left(\left(4 x - y\right)-\left(4 x + 2\right)\right)\left(\left(4 x - y\right)+\left(4 x + 2\right)\right)}\]
\end{enumerate}

\vspace{0.2cm}
\subsubsection*{Ejercicio 179 - Diferencia de cuadrados}
\textbf{Enunciado:} $ - 20 x y - 30 x + 4 y^{2} - 9 $\par
\begin{enumerate}[leftmargin=2em]
\item La diferencia de cuadrados siempre tiene esta forma: \(A^2-B^2\).
\item Aqui identificamos \(A=5 x - 2 y\) y \(B=5 x + 3\).
\item Aplicamos la formula: \(A^2-B^2=(A-B)(A+B)\).
\item Resultado: \[\boxed{\left(\left(5 x - 2 y\right)-\left(5 x + 3\right)\right)\left(\left(5 x - 2 y\right)+\left(5 x + 3\right)\right)}\]
\end{enumerate}

\vspace{0.2cm}
\subsubsection*{Ejercicio 180 - Diferencia de cuadrados}
\textbf{Enunciado:} $ 3 x^{2} - 12 x y - 4 x + 9 y^{2} - 4 $\par
\begin{enumerate}[leftmargin=2em]
\item La diferencia de cuadrados siempre tiene esta forma: \(A^2-B^2\).
\item Aqui identificamos \(A=2 x - 3 y\) y \(B=x + 2\).
\item Aplicamos la formula: \(A^2-B^2=(A-B)(A+B)\).
\item Resultado: \[\boxed{\left(\left(2 x - 3 y\right)-\left(x + 2\right)\right)\left(\left(2 x - 3 y\right)+\left(x + 2\right)\right)}\]
\end{enumerate}

\vspace{0.2cm}
\nivel{5}{imposible}
\subsubsection*{Ejercicio 181 - Diferencia de cuadrados}
\textbf{Enunciado:} $ 4 x^{4} - 8 x^{3} y + 4 x^{2} y^{2} + 4 x^{2} - 12 x y^{2} - 8 x y + 8 x - 9 y^{4} + 12 y^{2} $\par
\begin{enumerate}[leftmargin=2em]
\item La diferencia de cuadrados siempre tiene esta forma: \(A^2-B^2\).
\item Aqui identificamos \(A=2 x^{2} - 2 x y + 2\) y \(B=2 x + 3 y^{2} - 2\).
\item Aplicamos la formula: \(A^2-B^2=(A-B)(A+B)\).
\item Resultado: \[\boxed{\left(\left(2 x^{2} - 2 x y + 2\right)-\left(2 x + 3 y^{2} - 2\right)\right)\left(\left(2 x^{2} - 2 x y + 2\right)+\left(2 x + 3 y^{2} - 2\right)\right)}\]
\end{enumerate}

\vspace{0.2cm}
\subsubsection*{Ejercicio 182 - Diferencia de cuadrados}
\textbf{Enunciado:} $ 9 x^{4} - 18 x^{3} y + 9 x^{2} y^{2} + 9 x^{2} - 24 x y^{2} - 18 x y + 6 x - 16 y^{4} + 8 y^{2} + 8 $\par
\begin{enumerate}[leftmargin=2em]
\item La diferencia de cuadrados siempre tiene esta forma: \(A^2-B^2\).
\item Aqui identificamos \(A=3 x^{2} - 3 x y + 3\) y \(B=3 x + 4 y^{2} - 1\).
\item Aplicamos la formula: \(A^2-B^2=(A-B)(A+B)\).
\item Resultado: \[\boxed{\left(\left(3 x^{2} - 3 x y + 3\right)-\left(3 x + 4 y^{2} - 1\right)\right)\left(\left(3 x^{2} - 3 x y + 3\right)+\left(3 x + 4 y^{2} - 1\right)\right)}\]
\end{enumerate}

\vspace{0.2cm}
\subsubsection*{Ejercicio 183 - Diferencia de cuadrados}
\textbf{Enunciado:} $ x^{4} - 8 x^{3} y + 16 x^{2} y^{2} + 7 x^{2} - 10 x y^{2} - 32 x y + 4 x - 25 y^{4} + 20 y^{2} + 12 $\par
\begin{enumerate}[leftmargin=2em]
\item La diferencia de cuadrados siempre tiene esta forma: \(A^2-B^2\).
\item Aqui identificamos \(A=x^{2} - 4 x y + 4\) y \(B=x + 5 y^{2} - 2\).
\item Aplicamos la formula: \(A^2-B^2=(A-B)(A+B)\).
\item Resultado: \[\boxed{\left(\left(x^{2} - 4 x y + 4\right)-\left(x + 5 y^{2} - 2\right)\right)\left(\left(x^{2} - 4 x y + 4\right)+\left(x + 5 y^{2} - 2\right)\right)}\]
\end{enumerate}

\vspace{0.2cm}
\subsubsection*{Ejercicio 184 - Diferencia de cuadrados}
\textbf{Enunciado:} $ 4 x^{4} - 4 x^{3} y + x^{2} y^{2} + 16 x^{2} - 8 x y^{2} - 10 x y + 4 x - 4 y^{4} + 4 y^{2} + 24 $\par
\begin{enumerate}[leftmargin=2em]
\item La diferencia de cuadrados siempre tiene esta forma: \(A^2-B^2\).
\item Aqui identificamos \(A=2 x^{2} - x y + 5\) y \(B=2 x + 2 y^{2} - 1\).
\item Aplicamos la formula: \(A^2-B^2=(A-B)(A+B)\).
\item Resultado: \[\boxed{\left(\left(2 x^{2} - x y + 5\right)-\left(2 x + 2 y^{2} - 1\right)\right)\left(\left(2 x^{2} - x y + 5\right)+\left(2 x + 2 y^{2} - 1\right)\right)}\]
\end{enumerate}

\vspace{0.2cm}
\subsubsection*{Ejercicio 185 - Diferencia de cuadrados}
\textbf{Enunciado:} $ 9 x^{4} - 12 x^{3} y + 4 x^{2} y^{2} - 3 x^{2} - 18 x y^{2} - 4 x y + 12 x - 9 y^{4} + 12 y^{2} - 3 $\par
\begin{enumerate}[leftmargin=2em]
\item La diferencia de cuadrados siempre tiene esta forma: \(A^2-B^2\).
\item Aqui identificamos \(A=3 x^{2} - 2 x y + 1\) y \(B=3 x + 3 y^{2} - 2\).
\item Aplicamos la formula: \(A^2-B^2=(A-B)(A+B)\).
\item Resultado: \[\boxed{\left(\left(3 x^{2} - 2 x y + 1\right)-\left(3 x + 3 y^{2} - 2\right)\right)\left(\left(3 x^{2} - 2 x y + 1\right)+\left(3 x + 3 y^{2} - 2\right)\right)}\]
\end{enumerate}

\vspace{0.2cm}
\subsubsection*{Ejercicio 186 - Diferencia de cuadrados}
\textbf{Enunciado:} $ x^{4} - 6 x^{3} y + 9 x^{2} y^{2} + 3 x^{2} - 8 x y^{2} - 12 x y + 2 x - 16 y^{4} + 8 y^{2} + 3 $\par
\begin{enumerate}[leftmargin=2em]
\item La diferencia de cuadrados siempre tiene esta forma: \(A^2-B^2\).
\item Aqui identificamos \(A=x^{2} - 3 x y + 2\) y \(B=x + 4 y^{2} - 1\).
\item Aplicamos la formula: \(A^2-B^2=(A-B)(A+B)\).
\item Resultado: \[\boxed{\left(\left(x^{2} - 3 x y + 2\right)-\left(x + 4 y^{2} - 1\right)\right)\left(\left(x^{2} - 3 x y + 2\right)+\left(x + 4 y^{2} - 1\right)\right)}\]
\end{enumerate}

\vspace{0.2cm}
\subsubsection*{Ejercicio 187 - Diferencia de cuadrados}
\textbf{Enunciado:} $ 4 x^{4} - 16 x^{3} y + 16 x^{2} y^{2} + 8 x^{2} - 20 x y^{2} - 24 x y + 8 x - 25 y^{4} + 20 y^{2} + 5 $\par
\begin{enumerate}[leftmargin=2em]
\item La diferencia de cuadrados siempre tiene esta forma: \(A^2-B^2\).
\item Aqui identificamos \(A=2 x^{2} - 4 x y + 3\) y \(B=2 x + 5 y^{2} - 2\).
\item Aplicamos la formula: \(A^2-B^2=(A-B)(A+B)\).
\item Resultado: \[\boxed{\left(\left(2 x^{2} - 4 x y + 3\right)-\left(2 x + 5 y^{2} - 2\right)\right)\left(\left(2 x^{2} - 4 x y + 3\right)+\left(2 x + 5 y^{2} - 2\right)\right)}\]
\end{enumerate}

\vspace{0.2cm}
\subsubsection*{Ejercicio 188 - Diferencia de cuadrados}
\textbf{Enunciado:} $ 9 x^{4} - 6 x^{3} y + x^{2} y^{2} + 15 x^{2} - 12 x y^{2} - 8 x y + 6 x - 4 y^{4} + 4 y^{2} + 15 $\par
\begin{enumerate}[leftmargin=2em]
\item La diferencia de cuadrados siempre tiene esta forma: \(A^2-B^2\).
\item Aqui identificamos \(A=3 x^{2} - x y + 4\) y \(B=3 x + 2 y^{2} - 1\).
\item Aplicamos la formula: \(A^2-B^2=(A-B)(A+B)\).
\item Resultado: \[\boxed{\left(\left(3 x^{2} - x y + 4\right)-\left(3 x + 2 y^{2} - 1\right)\right)\left(\left(3 x^{2} - x y + 4\right)+\left(3 x + 2 y^{2} - 1\right)\right)}\]
\end{enumerate}

\vspace{0.2cm}
\subsubsection*{Ejercicio 189 - Diferencia de cuadrados}
\textbf{Enunciado:} $ x^{4} - 4 x^{3} y + 4 x^{2} y^{2} + 9 x^{2} - 6 x y^{2} - 20 x y + 4 x - 9 y^{4} + 12 y^{2} + 21 $\par
\begin{enumerate}[leftmargin=2em]
\item La diferencia de cuadrados siempre tiene esta forma: \(A^2-B^2\).
\item Aqui identificamos \(A=x^{2} - 2 x y + 5\) y \(B=x + 3 y^{2} - 2\).
\item Aplicamos la formula: \(A^2-B^2=(A-B)(A+B)\).
\item Resultado: \[\boxed{\left(\left(x^{2} - 2 x y + 5\right)-\left(x + 3 y^{2} - 2\right)\right)\left(\left(x^{2} - 2 x y + 5\right)+\left(x + 3 y^{2} - 2\right)\right)}\]
\end{enumerate}

\vspace{0.2cm}
\subsubsection*{Ejercicio 190 - Diferencia de cuadrados}
\textbf{Enunciado:} $ 4 x^{4} - 12 x^{3} y + 9 x^{2} y^{2} - 16 x y^{2} - 6 x y + 4 x - 16 y^{4} + 8 y^{2} $\par
\begin{enumerate}[leftmargin=2em]
\item La diferencia de cuadrados siempre tiene esta forma: \(A^2-B^2\).
\item Aqui identificamos \(A=2 x^{2} - 3 x y + 1\) y \(B=2 x + 4 y^{2} - 1\).
\item Aplicamos la formula: \(A^2-B^2=(A-B)(A+B)\).
\item Resultado: \[\boxed{\left(\left(2 x^{2} - 3 x y + 1\right)-\left(2 x + 4 y^{2} - 1\right)\right)\left(\left(2 x^{2} - 3 x y + 1\right)+\left(2 x + 4 y^{2} - 1\right)\right)}\]
\end{enumerate}

\vspace{0.2cm}
\subsubsection*{Ejercicio 191 - Diferencia de cuadrados}
\textbf{Enunciado:} $ 9 x^{4} - 24 x^{3} y + 16 x^{2} y^{2} + 3 x^{2} - 30 x y^{2} - 16 x y + 12 x - 25 y^{4} + 20 y^{2} $\par
\begin{enumerate}[leftmargin=2em]
\item La diferencia de cuadrados siempre tiene esta forma: \(A^2-B^2\).
\item Aqui identificamos \(A=3 x^{2} - 4 x y + 2\) y \(B=3 x + 5 y^{2} - 2\).
\item Aplicamos la formula: \(A^2-B^2=(A-B)(A+B)\).
\item Resultado: \[\boxed{\left(\left(3 x^{2} - 4 x y + 2\right)-\left(3 x + 5 y^{2} - 2\right)\right)\left(\left(3 x^{2} - 4 x y + 2\right)+\left(3 x + 5 y^{2} - 2\right)\right)}\]
\end{enumerate}

\vspace{0.2cm}
\subsubsection*{Ejercicio 192 - Diferencia de cuadrados}
\textbf{Enunciado:} $ x^{4} - 2 x^{3} y + x^{2} y^{2} + 5 x^{2} - 4 x y^{2} - 6 x y + 2 x - 4 y^{4} + 4 y^{2} + 8 $\par
\begin{enumerate}[leftmargin=2em]
\item La diferencia de cuadrados siempre tiene esta forma: \(A^2-B^2\).
\item Aqui identificamos \(A=x^{2} - x y + 3\) y \(B=x + 2 y^{2} - 1\).
\item Aplicamos la formula: \(A^2-B^2=(A-B)(A+B)\).
\item Resultado: \[\boxed{\left(\left(x^{2} - x y + 3\right)-\left(x + 2 y^{2} - 1\right)\right)\left(\left(x^{2} - x y + 3\right)+\left(x + 2 y^{2} - 1\right)\right)}\]
\end{enumerate}

\vspace{0.2cm}
\subsubsection*{Ejercicio 193 - Diferencia de cuadrados}
\textbf{Enunciado:} $ 4 x^{4} - 8 x^{3} y + 4 x^{2} y^{2} + 12 x^{2} - 12 x y^{2} - 16 x y + 8 x - 9 y^{4} + 12 y^{2} + 12 $\par
\begin{enumerate}[leftmargin=2em]
\item La diferencia de cuadrados siempre tiene esta forma: \(A^2-B^2\).
\item Aqui identificamos \(A=2 x^{2} - 2 x y + 4\) y \(B=2 x + 3 y^{2} - 2\).
\item Aplicamos la formula: \(A^2-B^2=(A-B)(A+B)\).
\item Resultado: \[\boxed{\left(\left(2 x^{2} - 2 x y + 4\right)-\left(2 x + 3 y^{2} - 2\right)\right)\left(\left(2 x^{2} - 2 x y + 4\right)+\left(2 x + 3 y^{2} - 2\right)\right)}\]
\end{enumerate}

\vspace{0.2cm}
\subsubsection*{Ejercicio 194 - Diferencia de cuadrados}
\textbf{Enunciado:} $ 9 x^{4} - 18 x^{3} y + 9 x^{2} y^{2} + 21 x^{2} - 24 x y^{2} - 30 x y + 6 x - 16 y^{4} + 8 y^{2} + 24 $\par
\begin{enumerate}[leftmargin=2em]
\item La diferencia de cuadrados siempre tiene esta forma: \(A^2-B^2\).
\item Aqui identificamos \(A=3 x^{2} - 3 x y + 5\) y \(B=3 x + 4 y^{2} - 1\).
\item Aplicamos la formula: \(A^2-B^2=(A-B)(A+B)\).
\item Resultado: \[\boxed{\left(\left(3 x^{2} - 3 x y + 5\right)-\left(3 x + 4 y^{2} - 1\right)\right)\left(\left(3 x^{2} - 3 x y + 5\right)+\left(3 x + 4 y^{2} - 1\right)\right)}\]
\end{enumerate}

\vspace{0.2cm}
\subsubsection*{Ejercicio 195 - Diferencia de cuadrados}
\textbf{Enunciado:} $ x^{4} - 8 x^{3} y + 16 x^{2} y^{2} + x^{2} - 10 x y^{2} - 8 x y + 4 x - 25 y^{4} + 20 y^{2} - 3 $\par
\begin{enumerate}[leftmargin=2em]
\item La diferencia de cuadrados siempre tiene esta forma: \(A^2-B^2\).
\item Aqui identificamos \(A=x^{2} - 4 x y + 1\) y \(B=x + 5 y^{2} - 2\).
\item Aplicamos la formula: \(A^2-B^2=(A-B)(A+B)\).
\item Resultado: \[\boxed{\left(\left(x^{2} - 4 x y + 1\right)-\left(x + 5 y^{2} - 2\right)\right)\left(\left(x^{2} - 4 x y + 1\right)+\left(x + 5 y^{2} - 2\right)\right)}\]
\end{enumerate}

\vspace{0.2cm}
\subsubsection*{Ejercicio 196 - Diferencia de cuadrados}
\textbf{Enunciado:} $ 4 x^{4} - 4 x^{3} y + x^{2} y^{2} + 4 x^{2} - 8 x y^{2} - 4 x y + 4 x - 4 y^{4} + 4 y^{2} + 3 $\par
\begin{enumerate}[leftmargin=2em]
\item La diferencia de cuadrados siempre tiene esta forma: \(A^2-B^2\).
\item Aqui identificamos \(A=2 x^{2} - x y + 2\) y \(B=2 x + 2 y^{2} - 1\).
\item Aplicamos la formula: \(A^2-B^2=(A-B)(A+B)\).
\item Resultado: \[\boxed{\left(\left(2 x^{2} - x y + 2\right)-\left(2 x + 2 y^{2} - 1\right)\right)\left(\left(2 x^{2} - x y + 2\right)+\left(2 x + 2 y^{2} - 1\right)\right)}\]
\end{enumerate}

\vspace{0.2cm}
\subsubsection*{Ejercicio 197 - Diferencia de cuadrados}
\textbf{Enunciado:} $ 9 x^{4} - 12 x^{3} y + 4 x^{2} y^{2} + 9 x^{2} - 18 x y^{2} - 12 x y + 12 x - 9 y^{4} + 12 y^{2} + 5 $\par
\begin{enumerate}[leftmargin=2em]
\item La diferencia de cuadrados siempre tiene esta forma: \(A^2-B^2\).
\item Aqui identificamos \(A=3 x^{2} - 2 x y + 3\) y \(B=3 x + 3 y^{2} - 2\).
\item Aplicamos la formula: \(A^2-B^2=(A-B)(A+B)\).
\item Resultado: \[\boxed{\left(\left(3 x^{2} - 2 x y + 3\right)-\left(3 x + 3 y^{2} - 2\right)\right)\left(\left(3 x^{2} - 2 x y + 3\right)+\left(3 x + 3 y^{2} - 2\right)\right)}\]
\end{enumerate}

\vspace{0.2cm}
\subsubsection*{Ejercicio 198 - Diferencia de cuadrados}
\textbf{Enunciado:} $ x^{4} - 6 x^{3} y + 9 x^{2} y^{2} + 7 x^{2} - 8 x y^{2} - 24 x y + 2 x - 16 y^{4} + 8 y^{2} + 15 $\par
\begin{enumerate}[leftmargin=2em]
\item La diferencia de cuadrados siempre tiene esta forma: \(A^2-B^2\).
\item Aqui identificamos \(A=x^{2} - 3 x y + 4\) y \(B=x + 4 y^{2} - 1\).
\item Aplicamos la formula: \(A^2-B^2=(A-B)(A+B)\).
\item Resultado: \[\boxed{\left(\left(x^{2} - 3 x y + 4\right)-\left(x + 4 y^{2} - 1\right)\right)\left(\left(x^{2} - 3 x y + 4\right)+\left(x + 4 y^{2} - 1\right)\right)}\]
\end{enumerate}

\vspace{0.2cm}
\subsubsection*{Ejercicio 199 - Diferencia de cuadrados}
\textbf{Enunciado:} $ 4 x^{4} - 16 x^{3} y + 16 x^{2} y^{2} + 16 x^{2} - 20 x y^{2} - 40 x y + 8 x - 25 y^{4} + 20 y^{2} + 21 $\par
\begin{enumerate}[leftmargin=2em]
\item La diferencia de cuadrados siempre tiene esta forma: \(A^2-B^2\).
\item Aqui identificamos \(A=2 x^{2} - 4 x y + 5\) y \(B=2 x + 5 y^{2} - 2\).
\item Aplicamos la formula: \(A^2-B^2=(A-B)(A+B)\).
\item Resultado: \[\boxed{\left(\left(2 x^{2} - 4 x y + 5\right)-\left(2 x + 5 y^{2} - 2\right)\right)\left(\left(2 x^{2} - 4 x y + 5\right)+\left(2 x + 5 y^{2} - 2\right)\right)}\]
\end{enumerate}

\vspace{0.2cm}
\subsubsection*{Ejercicio 200 - Diferencia de cuadrados}
\textbf{Enunciado:} $ 9 x^{4} - 6 x^{3} y + x^{2} y^{2} - 3 x^{2} - 12 x y^{2} - 2 x y + 6 x - 4 y^{4} + 4 y^{2} $\par
\begin{enumerate}[leftmargin=2em]
\item La diferencia de cuadrados siempre tiene esta forma: \(A^2-B^2\).
\item Aqui identificamos \(A=3 x^{2} - x y + 1\) y \(B=3 x + 2 y^{2} - 1\).
\item Aplicamos la formula: \(A^2-B^2=(A-B)(A+B)\).
\item Resultado: \[\boxed{\left(\left(3 x^{2} - x y + 1\right)-\left(3 x + 2 y^{2} - 1\right)\right)\left(\left(3 x^{2} - x y + 1\right)+\left(3 x + 2 y^{2} - 1\right)\right)}\]
\end{enumerate}

\vspace{0.2cm}
\section{Soluciones: Metodo de aspas}
\nivel{1}{basico}
\subsubsection*{Ejercicio 201 - Aspas}
\textbf{Enunciado:} $ x^{2} + 5 x + 6 $\par
\begin{enumerate}[leftmargin=2em]
\item El metodo de aspas busca dos binomios. Los extremos deben producir el primer y el ultimo termino.
\item Probamos con \((x + 2)\) y \((x + 3)\).
\item Verificamos el termino del medio con los productos cruzados: \(3 x+2 x=5 x\).
\item Como los productos cruzados coinciden con el termino central, el resultado es: \[\boxed{\left(x + 2\right)\left(x + 3\right)}\]
\end{enumerate}

\vspace{0.2cm}
\subsubsection*{Ejercicio 202 - Aspas}
\textbf{Enunciado:} $ x^{2} + 7 x + 12 $\par
\begin{enumerate}[leftmargin=2em]
\item El metodo de aspas busca dos binomios. Los extremos deben producir el primer y el ultimo termino.
\item Probamos con \((x + 3)\) y \((x + 4)\).
\item Verificamos el termino del medio con los productos cruzados: \(4 x+3 x=7 x\).
\item Como los productos cruzados coinciden con el termino central, el resultado es: \[\boxed{\left(x + 3\right)\left(x + 4\right)}\]
\end{enumerate}

\vspace{0.2cm}
\subsubsection*{Ejercicio 203 - Aspas}
\textbf{Enunciado:} $ x^{2} + 9 x + 20 $\par
\begin{enumerate}[leftmargin=2em]
\item El metodo de aspas busca dos binomios. Los extremos deben producir el primer y el ultimo termino.
\item Probamos con \((x + 4)\) y \((x + 5)\).
\item Verificamos el termino del medio con los productos cruzados: \(5 x+4 x=9 x\).
\item Como los productos cruzados coinciden con el termino central, el resultado es: \[\boxed{\left(x + 4\right)\left(x + 5\right)}\]
\end{enumerate}

\vspace{0.2cm}
\subsubsection*{Ejercicio 204 - Aspas}
\textbf{Enunciado:} $ x^{2} + 11 x + 30 $\par
\begin{enumerate}[leftmargin=2em]
\item El metodo de aspas busca dos binomios. Los extremos deben producir el primer y el ultimo termino.
\item Probamos con \((x + 5)\) y \((x + 6)\).
\item Verificamos el termino del medio con los productos cruzados: \(6 x+5 x=11 x\).
\item Como los productos cruzados coinciden con el termino central, el resultado es: \[\boxed{\left(x + 5\right)\left(x + 6\right)}\]
\end{enumerate}

\vspace{0.2cm}
\subsubsection*{Ejercicio 205 - Aspas}
\textbf{Enunciado:} $ x^{2} + 13 x + 42 $\par
\begin{enumerate}[leftmargin=2em]
\item El metodo de aspas busca dos binomios. Los extremos deben producir el primer y el ultimo termino.
\item Probamos con \((x + 6)\) y \((x + 7)\).
\item Verificamos el termino del medio con los productos cruzados: \(7 x+6 x=13 x\).
\item Como los productos cruzados coinciden con el termino central, el resultado es: \[\boxed{\left(x + 6\right)\left(x + 7\right)}\]
\end{enumerate}

\vspace{0.2cm}
\subsubsection*{Ejercicio 206 - Aspas}
\textbf{Enunciado:} $ x^{2} + 9 x + 8 $\par
\begin{enumerate}[leftmargin=2em]
\item El metodo de aspas busca dos binomios. Los extremos deben producir el primer y el ultimo termino.
\item Probamos con \((x + 1)\) y \((x + 8)\).
\item Verificamos el termino del medio con los productos cruzados: \(8 x+x=9 x\).
\item Como los productos cruzados coinciden con el termino central, el resultado es: \[\boxed{\left(x + 1\right)\left(x + 8\right)}\]
\end{enumerate}

\vspace{0.2cm}
\subsubsection*{Ejercicio 207 - Aspas}
\textbf{Enunciado:} $ x^{2} + 4 x + 4 $\par
\begin{enumerate}[leftmargin=2em]
\item El metodo de aspas busca dos binomios. Los extremos deben producir el primer y el ultimo termino.
\item Probamos con \((x + 2)\) y \((x + 2)\).
\item Verificamos el termino del medio con los productos cruzados: \(2 x+2 x=4 x\).
\item Como los productos cruzados coinciden con el termino central, el resultado es: \[\boxed{\left(x + 2\right)\left(x + 2\right)}\]
\end{enumerate}

\vspace{0.2cm}
\subsubsection*{Ejercicio 208 - Aspas}
\textbf{Enunciado:} $ x^{2} + 6 x + 9 $\par
\begin{enumerate}[leftmargin=2em]
\item El metodo de aspas busca dos binomios. Los extremos deben producir el primer y el ultimo termino.
\item Probamos con \((x + 3)\) y \((x + 3)\).
\item Verificamos el termino del medio con los productos cruzados: \(3 x+3 x=6 x\).
\item Como los productos cruzados coinciden con el termino central, el resultado es: \[\boxed{\left(x + 3\right)\left(x + 3\right)}\]
\end{enumerate}

\vspace{0.2cm}
\subsubsection*{Ejercicio 209 - Aspas}
\textbf{Enunciado:} $ x^{2} + 8 x + 16 $\par
\begin{enumerate}[leftmargin=2em]
\item El metodo de aspas busca dos binomios. Los extremos deben producir el primer y el ultimo termino.
\item Probamos con \((x + 4)\) y \((x + 4)\).
\item Verificamos el termino del medio con los productos cruzados: \(4 x+4 x=8 x\).
\item Como los productos cruzados coinciden con el termino central, el resultado es: \[\boxed{\left(x + 4\right)\left(x + 4\right)}\]
\end{enumerate}

\vspace{0.2cm}
\subsubsection*{Ejercicio 210 - Aspas}
\textbf{Enunciado:} $ x^{2} + 10 x + 25 $\par
\begin{enumerate}[leftmargin=2em]
\item El metodo de aspas busca dos binomios. Los extremos deben producir el primer y el ultimo termino.
\item Probamos con \((x + 5)\) y \((x + 5)\).
\item Verificamos el termino del medio con los productos cruzados: \(5 x+5 x=10 x\).
\item Como los productos cruzados coinciden con el termino central, el resultado es: \[\boxed{\left(x + 5\right)\left(x + 5\right)}\]
\end{enumerate}

\vspace{0.2cm}
\subsubsection*{Ejercicio 211 - Aspas}
\textbf{Enunciado:} $ x^{2} + 12 x + 36 $\par
\begin{enumerate}[leftmargin=2em]
\item El metodo de aspas busca dos binomios. Los extremos deben producir el primer y el ultimo termino.
\item Probamos con \((x + 6)\) y \((x + 6)\).
\item Verificamos el termino del medio con los productos cruzados: \(6 x+6 x=12 x\).
\item Como los productos cruzados coinciden con el termino central, el resultado es: \[\boxed{\left(x + 6\right)\left(x + 6\right)}\]
\end{enumerate}

\vspace{0.2cm}
\subsubsection*{Ejercicio 212 - Aspas}
\textbf{Enunciado:} $ x^{2} + 8 x + 7 $\par
\begin{enumerate}[leftmargin=2em]
\item El metodo de aspas busca dos binomios. Los extremos deben producir el primer y el ultimo termino.
\item Probamos con \((x + 1)\) y \((x + 7)\).
\item Verificamos el termino del medio con los productos cruzados: \(7 x+x=8 x\).
\item Como los productos cruzados coinciden con el termino central, el resultado es: \[\boxed{\left(x + 1\right)\left(x + 7\right)}\]
\end{enumerate}

\vspace{0.2cm}
\subsubsection*{Ejercicio 213 - Aspas}
\textbf{Enunciado:} $ x^{2} + 10 x + 16 $\par
\begin{enumerate}[leftmargin=2em]
\item El metodo de aspas busca dos binomios. Los extremos deben producir el primer y el ultimo termino.
\item Probamos con \((x + 2)\) y \((x + 8)\).
\item Verificamos el termino del medio con los productos cruzados: \(8 x+2 x=10 x\).
\item Como los productos cruzados coinciden con el termino central, el resultado es: \[\boxed{\left(x + 2\right)\left(x + 8\right)}\]
\end{enumerate}

\vspace{0.2cm}
\subsubsection*{Ejercicio 214 - Aspas}
\textbf{Enunciado:} $ x^{2} + 5 x + 6 $\par
\begin{enumerate}[leftmargin=2em]
\item El metodo de aspas busca dos binomios. Los extremos deben producir el primer y el ultimo termino.
\item Probamos con \((x + 3)\) y \((x + 2)\).
\item Verificamos el termino del medio con los productos cruzados: \(2 x+3 x=5 x\).
\item Como los productos cruzados coinciden con el termino central, el resultado es: \[\boxed{\left(x + 3\right)\left(x + 2\right)}\]
\end{enumerate}

\vspace{0.2cm}
\subsubsection*{Ejercicio 215 - Aspas}
\textbf{Enunciado:} $ x^{2} + 7 x + 12 $\par
\begin{enumerate}[leftmargin=2em]
\item El metodo de aspas busca dos binomios. Los extremos deben producir el primer y el ultimo termino.
\item Probamos con \((x + 4)\) y \((x + 3)\).
\item Verificamos el termino del medio con los productos cruzados: \(3 x+4 x=7 x\).
\item Como los productos cruzados coinciden con el termino central, el resultado es: \[\boxed{\left(x + 4\right)\left(x + 3\right)}\]
\end{enumerate}

\vspace{0.2cm}
\subsubsection*{Ejercicio 216 - Aspas}
\textbf{Enunciado:} $ x^{2} + 9 x + 20 $\par
\begin{enumerate}[leftmargin=2em]
\item El metodo de aspas busca dos binomios. Los extremos deben producir el primer y el ultimo termino.
\item Probamos con \((x + 5)\) y \((x + 4)\).
\item Verificamos el termino del medio con los productos cruzados: \(4 x+5 x=9 x\).
\item Como los productos cruzados coinciden con el termino central, el resultado es: \[\boxed{\left(x + 5\right)\left(x + 4\right)}\]
\end{enumerate}

\vspace{0.2cm}
\subsubsection*{Ejercicio 217 - Aspas}
\textbf{Enunciado:} $ x^{2} + 11 x + 30 $\par
\begin{enumerate}[leftmargin=2em]
\item El metodo de aspas busca dos binomios. Los extremos deben producir el primer y el ultimo termino.
\item Probamos con \((x + 6)\) y \((x + 5)\).
\item Verificamos el termino del medio con los productos cruzados: \(5 x+6 x=11 x\).
\item Como los productos cruzados coinciden con el termino central, el resultado es: \[\boxed{\left(x + 6\right)\left(x + 5\right)}\]
\end{enumerate}

\vspace{0.2cm}
\subsubsection*{Ejercicio 218 - Aspas}
\textbf{Enunciado:} $ x^{2} + 7 x + 6 $\par
\begin{enumerate}[leftmargin=2em]
\item El metodo de aspas busca dos binomios. Los extremos deben producir el primer y el ultimo termino.
\item Probamos con \((x + 1)\) y \((x + 6)\).
\item Verificamos el termino del medio con los productos cruzados: \(6 x+x=7 x\).
\item Como los productos cruzados coinciden con el termino central, el resultado es: \[\boxed{\left(x + 1\right)\left(x + 6\right)}\]
\end{enumerate}

\vspace{0.2cm}
\subsubsection*{Ejercicio 219 - Aspas}
\textbf{Enunciado:} $ x^{2} + 9 x + 14 $\par
\begin{enumerate}[leftmargin=2em]
\item El metodo de aspas busca dos binomios. Los extremos deben producir el primer y el ultimo termino.
\item Probamos con \((x + 2)\) y \((x + 7)\).
\item Verificamos el termino del medio con los productos cruzados: \(7 x+2 x=9 x\).
\item Como los productos cruzados coinciden con el termino central, el resultado es: \[\boxed{\left(x + 2\right)\left(x + 7\right)}\]
\end{enumerate}

\vspace{0.2cm}
\subsubsection*{Ejercicio 220 - Aspas}
\textbf{Enunciado:} $ x^{2} + 11 x + 24 $\par
\begin{enumerate}[leftmargin=2em]
\item El metodo de aspas busca dos binomios. Los extremos deben producir el primer y el ultimo termino.
\item Probamos con \((x + 3)\) y \((x + 8)\).
\item Verificamos el termino del medio con los productos cruzados: \(8 x+3 x=11 x\).
\item Como los productos cruzados coinciden con el termino central, el resultado es: \[\boxed{\left(x + 3\right)\left(x + 8\right)}\]
\end{enumerate}

\vspace{0.2cm}
\nivel{2}{intermedio inicial}
\subsubsection*{Ejercicio 221 - Aspas}
\textbf{Enunciado:} $ x^{2} + x - 6 $\par
\begin{enumerate}[leftmargin=2em]
\item El metodo de aspas busca dos binomios. Los extremos deben producir el primer y el ultimo termino.
\item Probamos con \((x - 2)\) y \((x + 3)\).
\item Verificamos el termino del medio con los productos cruzados: \(3 x- 2 x=x\).
\item Como los productos cruzados coinciden con el termino central, el resultado es: \[\boxed{\left(x - 2\right)\left(x + 3\right)}\]
\end{enumerate}

\vspace{0.2cm}
\subsubsection*{Ejercicio 222 - Aspas}
\textbf{Enunciado:} $ x^{2} - x - 12 $\par
\begin{enumerate}[leftmargin=2em]
\item El metodo de aspas busca dos binomios. Los extremos deben producir el primer y el ultimo termino.
\item Probamos con \((x + 3)\) y \((x - 4)\).
\item Verificamos el termino del medio con los productos cruzados: \(- 4 x+3 x=- x\).
\item Como los productos cruzados coinciden con el termino central, el resultado es: \[\boxed{\left(x + 3\right)\left(x - 4\right)}\]
\end{enumerate}

\vspace{0.2cm}
\subsubsection*{Ejercicio 223 - Aspas}
\textbf{Enunciado:} $ x^{2} + x - 20 $\par
\begin{enumerate}[leftmargin=2em]
\item El metodo de aspas busca dos binomios. Los extremos deben producir el primer y el ultimo termino.
\item Probamos con \((x - 4)\) y \((x + 5)\).
\item Verificamos el termino del medio con los productos cruzados: \(5 x- 4 x=x\).
\item Como los productos cruzados coinciden con el termino central, el resultado es: \[\boxed{\left(x - 4\right)\left(x + 5\right)}\]
\end{enumerate}

\vspace{0.2cm}
\subsubsection*{Ejercicio 224 - Aspas}
\textbf{Enunciado:} $ x^{2} - x - 30 $\par
\begin{enumerate}[leftmargin=2em]
\item El metodo de aspas busca dos binomios. Los extremos deben producir el primer y el ultimo termino.
\item Probamos con \((x + 5)\) y \((x - 6)\).
\item Verificamos el termino del medio con los productos cruzados: \(- 6 x+5 x=- x\).
\item Como los productos cruzados coinciden con el termino central, el resultado es: \[\boxed{\left(x + 5\right)\left(x - 6\right)}\]
\end{enumerate}

\vspace{0.2cm}
\subsubsection*{Ejercicio 225 - Aspas}
\textbf{Enunciado:} $ x^{2} + x - 42 $\par
\begin{enumerate}[leftmargin=2em]
\item El metodo de aspas busca dos binomios. Los extremos deben producir el primer y el ultimo termino.
\item Probamos con \((x - 6)\) y \((x + 7)\).
\item Verificamos el termino del medio con los productos cruzados: \(7 x- 6 x=x\).
\item Como los productos cruzados coinciden con el termino central, el resultado es: \[\boxed{\left(x - 6\right)\left(x + 7\right)}\]
\end{enumerate}

\vspace{0.2cm}
\subsubsection*{Ejercicio 226 - Aspas}
\textbf{Enunciado:} $ x^{2} - 7 x - 8 $\par
\begin{enumerate}[leftmargin=2em]
\item El metodo de aspas busca dos binomios. Los extremos deben producir el primer y el ultimo termino.
\item Probamos con \((x + 1)\) y \((x - 8)\).
\item Verificamos el termino del medio con los productos cruzados: \(- 8 x+x=- 7 x\).
\item Como los productos cruzados coinciden con el termino central, el resultado es: \[\boxed{\left(x + 1\right)\left(x - 8\right)}\]
\end{enumerate}

\vspace{0.2cm}
\subsubsection*{Ejercicio 227 - Aspas}
\textbf{Enunciado:} $ x^{2} - 4 $\par
\begin{enumerate}[leftmargin=2em]
\item El metodo de aspas busca dos binomios. Los extremos deben producir el primer y el ultimo termino.
\item Probamos con \((x - 2)\) y \((x + 2)\).
\item Verificamos el termino del medio con los productos cruzados: \(2 x- 2 x=0\).
\item Como los productos cruzados coinciden con el termino central, el resultado es: \[\boxed{\left(x - 2\right)\left(x + 2\right)}\]
\end{enumerate}

\vspace{0.2cm}
\subsubsection*{Ejercicio 228 - Aspas}
\textbf{Enunciado:} $ x^{2} - 9 $\par
\begin{enumerate}[leftmargin=2em]
\item El metodo de aspas busca dos binomios. Los extremos deben producir el primer y el ultimo termino.
\item Probamos con \((x + 3)\) y \((x - 3)\).
\item Verificamos el termino del medio con los productos cruzados: \(- 3 x+3 x=0\).
\item Como los productos cruzados coinciden con el termino central, el resultado es: \[\boxed{\left(x + 3\right)\left(x - 3\right)}\]
\end{enumerate}

\vspace{0.2cm}
\subsubsection*{Ejercicio 229 - Aspas}
\textbf{Enunciado:} $ x^{2} - 16 $\par
\begin{enumerate}[leftmargin=2em]
\item El metodo de aspas busca dos binomios. Los extremos deben producir el primer y el ultimo termino.
\item Probamos con \((x - 4)\) y \((x + 4)\).
\item Verificamos el termino del medio con los productos cruzados: \(4 x- 4 x=0\).
\item Como los productos cruzados coinciden con el termino central, el resultado es: \[\boxed{\left(x - 4\right)\left(x + 4\right)}\]
\end{enumerate}

\vspace{0.2cm}
\subsubsection*{Ejercicio 230 - Aspas}
\textbf{Enunciado:} $ x^{2} - 25 $\par
\begin{enumerate}[leftmargin=2em]
\item El metodo de aspas busca dos binomios. Los extremos deben producir el primer y el ultimo termino.
\item Probamos con \((x + 5)\) y \((x - 5)\).
\item Verificamos el termino del medio con los productos cruzados: \(- 5 x+5 x=0\).
\item Como los productos cruzados coinciden con el termino central, el resultado es: \[\boxed{\left(x + 5\right)\left(x - 5\right)}\]
\end{enumerate}

\vspace{0.2cm}
\subsubsection*{Ejercicio 231 - Aspas}
\textbf{Enunciado:} $ x^{2} - 36 $\par
\begin{enumerate}[leftmargin=2em]
\item El metodo de aspas busca dos binomios. Los extremos deben producir el primer y el ultimo termino.
\item Probamos con \((x - 6)\) y \((x + 6)\).
\item Verificamos el termino del medio con los productos cruzados: \(6 x- 6 x=0\).
\item Como los productos cruzados coinciden con el termino central, el resultado es: \[\boxed{\left(x - 6\right)\left(x + 6\right)}\]
\end{enumerate}

\vspace{0.2cm}
\subsubsection*{Ejercicio 232 - Aspas}
\textbf{Enunciado:} $ x^{2} - 6 x - 7 $\par
\begin{enumerate}[leftmargin=2em]
\item El metodo de aspas busca dos binomios. Los extremos deben producir el primer y el ultimo termino.
\item Probamos con \((x + 1)\) y \((x - 7)\).
\item Verificamos el termino del medio con los productos cruzados: \(- 7 x+x=- 6 x\).
\item Como los productos cruzados coinciden con el termino central, el resultado es: \[\boxed{\left(x + 1\right)\left(x - 7\right)}\]
\end{enumerate}

\vspace{0.2cm}
\subsubsection*{Ejercicio 233 - Aspas}
\textbf{Enunciado:} $ x^{2} + 6 x - 16 $\par
\begin{enumerate}[leftmargin=2em]
\item El metodo de aspas busca dos binomios. Los extremos deben producir el primer y el ultimo termino.
\item Probamos con \((x - 2)\) y \((x + 8)\).
\item Verificamos el termino del medio con los productos cruzados: \(8 x- 2 x=6 x\).
\item Como los productos cruzados coinciden con el termino central, el resultado es: \[\boxed{\left(x - 2\right)\left(x + 8\right)}\]
\end{enumerate}

\vspace{0.2cm}
\subsubsection*{Ejercicio 234 - Aspas}
\textbf{Enunciado:} $ x^{2} + x - 6 $\par
\begin{enumerate}[leftmargin=2em]
\item El metodo de aspas busca dos binomios. Los extremos deben producir el primer y el ultimo termino.
\item Probamos con \((x + 3)\) y \((x - 2)\).
\item Verificamos el termino del medio con los productos cruzados: \(- 2 x+3 x=x\).
\item Como los productos cruzados coinciden con el termino central, el resultado es: \[\boxed{\left(x + 3\right)\left(x - 2\right)}\]
\end{enumerate}

\vspace{0.2cm}
\subsubsection*{Ejercicio 235 - Aspas}
\textbf{Enunciado:} $ x^{2} - x - 12 $\par
\begin{enumerate}[leftmargin=2em]
\item El metodo de aspas busca dos binomios. Los extremos deben producir el primer y el ultimo termino.
\item Probamos con \((x - 4)\) y \((x + 3)\).
\item Verificamos el termino del medio con los productos cruzados: \(3 x- 4 x=- x\).
\item Como los productos cruzados coinciden con el termino central, el resultado es: \[\boxed{\left(x - 4\right)\left(x + 3\right)}\]
\end{enumerate}

\vspace{0.2cm}
\subsubsection*{Ejercicio 236 - Aspas}
\textbf{Enunciado:} $ x^{2} + x - 20 $\par
\begin{enumerate}[leftmargin=2em]
\item El metodo de aspas busca dos binomios. Los extremos deben producir el primer y el ultimo termino.
\item Probamos con \((x + 5)\) y \((x - 4)\).
\item Verificamos el termino del medio con los productos cruzados: \(- 4 x+5 x=x\).
\item Como los productos cruzados coinciden con el termino central, el resultado es: \[\boxed{\left(x + 5\right)\left(x - 4\right)}\]
\end{enumerate}

\vspace{0.2cm}
\subsubsection*{Ejercicio 237 - Aspas}
\textbf{Enunciado:} $ x^{2} - x - 30 $\par
\begin{enumerate}[leftmargin=2em]
\item El metodo de aspas busca dos binomios. Los extremos deben producir el primer y el ultimo termino.
\item Probamos con \((x - 6)\) y \((x + 5)\).
\item Verificamos el termino del medio con los productos cruzados: \(5 x- 6 x=- x\).
\item Como los productos cruzados coinciden con el termino central, el resultado es: \[\boxed{\left(x - 6\right)\left(x + 5\right)}\]
\end{enumerate}

\vspace{0.2cm}
\subsubsection*{Ejercicio 238 - Aspas}
\textbf{Enunciado:} $ x^{2} - 5 x - 6 $\par
\begin{enumerate}[leftmargin=2em]
\item El metodo de aspas busca dos binomios. Los extremos deben producir el primer y el ultimo termino.
\item Probamos con \((x + 1)\) y \((x - 6)\).
\item Verificamos el termino del medio con los productos cruzados: \(- 6 x+x=- 5 x\).
\item Como los productos cruzados coinciden con el termino central, el resultado es: \[\boxed{\left(x + 1\right)\left(x - 6\right)}\]
\end{enumerate}

\vspace{0.2cm}
\subsubsection*{Ejercicio 239 - Aspas}
\textbf{Enunciado:} $ x^{2} + 5 x - 14 $\par
\begin{enumerate}[leftmargin=2em]
\item El metodo de aspas busca dos binomios. Los extremos deben producir el primer y el ultimo termino.
\item Probamos con \((x - 2)\) y \((x + 7)\).
\item Verificamos el termino del medio con los productos cruzados: \(7 x- 2 x=5 x\).
\item Como los productos cruzados coinciden con el termino central, el resultado es: \[\boxed{\left(x - 2\right)\left(x + 7\right)}\]
\end{enumerate}

\vspace{0.2cm}
\subsubsection*{Ejercicio 240 - Aspas}
\textbf{Enunciado:} $ x^{2} - 5 x - 24 $\par
\begin{enumerate}[leftmargin=2em]
\item El metodo de aspas busca dos binomios. Los extremos deben producir el primer y el ultimo termino.
\item Probamos con \((x + 3)\) y \((x - 8)\).
\item Verificamos el termino del medio con los productos cruzados: \(- 8 x+3 x=- 5 x\).
\item Como los productos cruzados coinciden con el termino central, el resultado es: \[\boxed{\left(x + 3\right)\left(x - 8\right)}\]
\end{enumerate}

\vspace{0.2cm}
\nivel{3}{intermedio}
\subsubsection*{Ejercicio 241 - Aspas}
\textbf{Enunciado:} $ 6 x^{2} + 5 x - 6 $\par
\begin{enumerate}[leftmargin=2em]
\item El metodo de aspas busca dos binomios. Los extremos deben producir el primer y el ultimo termino.
\item Probamos con \((3 x - 2)\) y \((2 x + 3)\).
\item Verificamos el termino del medio con los productos cruzados: \(9 x- 4 x=5 x\).
\item Como los productos cruzados coinciden con el termino central, el resultado es: \[\boxed{\left(3 x - 2\right)\left(2 x + 3\right)}\]
\end{enumerate}

\vspace{0.2cm}
\subsubsection*{Ejercicio 242 - Aspas}
\textbf{Enunciado:} $ 12 x^{2} - 7 x - 12 $\par
\begin{enumerate}[leftmargin=2em]
\item El metodo de aspas busca dos binomios. Los extremos deben producir el primer y el ultimo termino.
\item Probamos con \((4 x + 3)\) y \((3 x - 4)\).
\item Verificamos el termino del medio con los productos cruzados: \(- 16 x+9 x=- 7 x\).
\item Como los productos cruzados coinciden con el termino central, el resultado es: \[\boxed{\left(4 x + 3\right)\left(3 x - 4\right)}\]
\end{enumerate}

\vspace{0.2cm}
\subsubsection*{Ejercicio 243 - Aspas}
\textbf{Enunciado:} $ 5 x^{2} + 21 x - 20 $\par
\begin{enumerate}[leftmargin=2em]
\item El metodo de aspas busca dos binomios. Los extremos deben producir el primer y el ultimo termino.
\item Probamos con \((5 x - 4)\) y \((x + 5)\).
\item Verificamos el termino del medio con los productos cruzados: \(25 x- 4 x=21 x\).
\item Como los productos cruzados coinciden con el termino central, el resultado es: \[\boxed{\left(5 x - 4\right)\left(x + 5\right)}\]
\end{enumerate}

\vspace{0.2cm}
\subsubsection*{Ejercicio 244 - Aspas}
\textbf{Enunciado:} $ 4 x^{2} - 2 x - 30 $\par
\begin{enumerate}[leftmargin=2em]
\item El metodo de aspas busca dos binomios. Los extremos deben producir el primer y el ultimo termino.
\item Probamos con \((2 x + 5)\) y \((2 x - 6)\).
\item Verificamos el termino del medio con los productos cruzados: \(- 12 x+10 x=- 2 x\).
\item Como los productos cruzados coinciden con el termino central, el resultado es: \[\boxed{\left(2 x + 5\right)\left(2 x - 6\right)}\]
\end{enumerate}

\vspace{0.2cm}
\subsubsection*{Ejercicio 245 - Aspas}
\textbf{Enunciado:} $ 9 x^{2} + 18 x - 7 $\par
\begin{enumerate}[leftmargin=2em]
\item El metodo de aspas busca dos binomios. Los extremos deben producir el primer y el ultimo termino.
\item Probamos con \((3 x - 1)\) y \((3 x + 7)\).
\item Verificamos el termino del medio con los productos cruzados: \(21 x- 3 x=18 x\).
\item Como los productos cruzados coinciden con el termino central, el resultado es: \[\boxed{\left(3 x - 1\right)\left(3 x + 7\right)}\]
\end{enumerate}

\vspace{0.2cm}
\subsubsection*{Ejercicio 246 - Aspas}
\textbf{Enunciado:} $ 4 x^{2} - 6 x - 4 $\par
\begin{enumerate}[leftmargin=2em]
\item El metodo de aspas busca dos binomios. Los extremos deben producir el primer y el ultimo termino.
\item Probamos con \((4 x + 2)\) y \((x - 2)\).
\item Verificamos el termino del medio con los productos cruzados: \(- 8 x+2 x=- 6 x\).
\item Como los productos cruzados coinciden con el termino central, el resultado es: \[\boxed{\left(4 x + 2\right)\left(x - 2\right)}\]
\end{enumerate}

\vspace{0.2cm}
\subsubsection*{Ejercicio 247 - Aspas}
\textbf{Enunciado:} $ 10 x^{2} + 9 x - 9 $\par
\begin{enumerate}[leftmargin=2em]
\item El metodo de aspas busca dos binomios. Los extremos deben producir el primer y el ultimo termino.
\item Probamos con \((5 x - 3)\) y \((2 x + 3)\).
\item Verificamos el termino del medio con los productos cruzados: \(15 x- 6 x=9 x\).
\item Como los productos cruzados coinciden con el termino central, el resultado es: \[\boxed{\left(5 x - 3\right)\left(2 x + 3\right)}\]
\end{enumerate}

\vspace{0.2cm}
\subsubsection*{Ejercicio 248 - Aspas}
\textbf{Enunciado:} $ 6 x^{2} + 4 x - 16 $\par
\begin{enumerate}[leftmargin=2em]
\item El metodo de aspas busca dos binomios. Los extremos deben producir el primer y el ultimo termino.
\item Probamos con \((2 x + 4)\) y \((3 x - 4)\).
\item Verificamos el termino del medio con los productos cruzados: \(- 8 x+12 x=4 x\).
\item Como los productos cruzados coinciden con el termino central, el resultado es: \[\boxed{\left(2 x + 4\right)\left(3 x - 4\right)}\]
\end{enumerate}

\vspace{0.2cm}
\subsubsection*{Ejercicio 249 - Aspas}
\textbf{Enunciado:} $ 3 x^{2} + 10 x - 25 $\par
\begin{enumerate}[leftmargin=2em]
\item El metodo de aspas busca dos binomios. Los extremos deben producir el primer y el ultimo termino.
\item Probamos con \((3 x - 5)\) y \((x + 5)\).
\item Verificamos el termino del medio con los productos cruzados: \(15 x- 5 x=10 x\).
\item Como los productos cruzados coinciden con el termino central, el resultado es: \[\boxed{\left(3 x - 5\right)\left(x + 5\right)}\]
\end{enumerate}

\vspace{0.2cm}
\subsubsection*{Ejercicio 250 - Aspas}
\textbf{Enunciado:} $ 8 x^{2} - 22 x - 6 $\par
\begin{enumerate}[leftmargin=2em]
\item El metodo de aspas busca dos binomios. Los extremos deben producir el primer y el ultimo termino.
\item Probamos con \((4 x + 1)\) y \((2 x - 6)\).
\item Verificamos el termino del medio con los productos cruzados: \(- 24 x+2 x=- 22 x\).
\item Como los productos cruzados coinciden con el termino central, el resultado es: \[\boxed{\left(4 x + 1\right)\left(2 x - 6\right)}\]
\end{enumerate}

\vspace{0.2cm}
\subsubsection*{Ejercicio 251 - Aspas}
\textbf{Enunciado:} $ 15 x^{2} + 29 x - 14 $\par
\begin{enumerate}[leftmargin=2em]
\item El metodo de aspas busca dos binomios. Los extremos deben producir el primer y el ultimo termino.
\item Probamos con \((5 x - 2)\) y \((3 x + 7)\).
\item Verificamos el termino del medio con los productos cruzados: \(35 x- 6 x=29 x\).
\item Como los productos cruzados coinciden con el termino central, el resultado es: \[\boxed{\left(5 x - 2\right)\left(3 x + 7\right)}\]
\end{enumerate}

\vspace{0.2cm}
\subsubsection*{Ejercicio 252 - Aspas}
\textbf{Enunciado:} $ 2 x^{2} - x - 6 $\par
\begin{enumerate}[leftmargin=2em]
\item El metodo de aspas busca dos binomios. Los extremos deben producir el primer y el ultimo termino.
\item Probamos con \((2 x + 3)\) y \((x - 2)\).
\item Verificamos el termino del medio con los productos cruzados: \(- 4 x+3 x=- x\).
\item Como los productos cruzados coinciden con el termino central, el resultado es: \[\boxed{\left(2 x + 3\right)\left(x - 2\right)}\]
\end{enumerate}

\vspace{0.2cm}
\subsubsection*{Ejercicio 253 - Aspas}
\textbf{Enunciado:} $ 6 x^{2} + x - 12 $\par
\begin{enumerate}[leftmargin=2em]
\item El metodo de aspas busca dos binomios. Los extremos deben producir el primer y el ultimo termino.
\item Probamos con \((3 x - 4)\) y \((2 x + 3)\).
\item Verificamos el termino del medio con los productos cruzados: \(9 x- 8 x=x\).
\item Como los productos cruzados coinciden con el termino central, el resultado es: \[\boxed{\left(3 x - 4\right)\left(2 x + 3\right)}\]
\end{enumerate}

\vspace{0.2cm}
\subsubsection*{Ejercicio 254 - Aspas}
\textbf{Enunciado:} $ 12 x^{2} - x - 20 $\par
\begin{enumerate}[leftmargin=2em]
\item El metodo de aspas busca dos binomios. Los extremos deben producir el primer y el ultimo termino.
\item Probamos con \((4 x + 5)\) y \((3 x - 4)\).
\item Verificamos el termino del medio con los productos cruzados: \(- 16 x+15 x=- x\).
\item Como los productos cruzados coinciden con el termino central, el resultado es: \[\boxed{\left(4 x + 5\right)\left(3 x - 4\right)}\]
\end{enumerate}

\vspace{0.2cm}
\subsubsection*{Ejercicio 255 - Aspas}
\textbf{Enunciado:} $ 5 x^{2} + 24 x - 5 $\par
\begin{enumerate}[leftmargin=2em]
\item El metodo de aspas busca dos binomios. Los extremos deben producir el primer y el ultimo termino.
\item Probamos con \((5 x - 1)\) y \((x + 5)\).
\item Verificamos el termino del medio con los productos cruzados: \(25 x- x=24 x\).
\item Como los productos cruzados coinciden con el termino central, el resultado es: \[\boxed{\left(5 x - 1\right)\left(x + 5\right)}\]
\end{enumerate}

\vspace{0.2cm}
\subsubsection*{Ejercicio 256 - Aspas}
\textbf{Enunciado:} $ 4 x^{2} - 8 x - 12 $\par
\begin{enumerate}[leftmargin=2em]
\item El metodo de aspas busca dos binomios. Los extremos deben producir el primer y el ultimo termino.
\item Probamos con \((2 x + 2)\) y \((2 x - 6)\).
\item Verificamos el termino del medio con los productos cruzados: \(- 12 x+4 x=- 8 x\).
\item Como los productos cruzados coinciden con el termino central, el resultado es: \[\boxed{\left(2 x + 2\right)\left(2 x - 6\right)}\]
\end{enumerate}

\vspace{0.2cm}
\subsubsection*{Ejercicio 257 - Aspas}
\textbf{Enunciado:} $ 9 x^{2} + 12 x - 21 $\par
\begin{enumerate}[leftmargin=2em]
\item El metodo de aspas busca dos binomios. Los extremos deben producir el primer y el ultimo termino.
\item Probamos con \((3 x - 3)\) y \((3 x + 7)\).
\item Verificamos el termino del medio con los productos cruzados: \(21 x- 9 x=12 x\).
\item Como los productos cruzados coinciden con el termino central, el resultado es: \[\boxed{\left(3 x - 3\right)\left(3 x + 7\right)}\]
\end{enumerate}

\vspace{0.2cm}
\subsubsection*{Ejercicio 258 - Aspas}
\textbf{Enunciado:} $ 4 x^{2} - 4 x - 8 $\par
\begin{enumerate}[leftmargin=2em]
\item El metodo de aspas busca dos binomios. Los extremos deben producir el primer y el ultimo termino.
\item Probamos con \((4 x + 4)\) y \((x - 2)\).
\item Verificamos el termino del medio con los productos cruzados: \(- 8 x+4 x=- 4 x\).
\item Como los productos cruzados coinciden con el termino central, el resultado es: \[\boxed{\left(4 x + 4\right)\left(x - 2\right)}\]
\end{enumerate}

\vspace{0.2cm}
\subsubsection*{Ejercicio 259 - Aspas}
\textbf{Enunciado:} $ 10 x^{2} + 5 x - 15 $\par
\begin{enumerate}[leftmargin=2em]
\item El metodo de aspas busca dos binomios. Los extremos deben producir el primer y el ultimo termino.
\item Probamos con \((5 x - 5)\) y \((2 x + 3)\).
\item Verificamos el termino del medio con los productos cruzados: \(15 x- 10 x=5 x\).
\item Como los productos cruzados coinciden con el termino central, el resultado es: \[\boxed{\left(5 x - 5\right)\left(2 x + 3\right)}\]
\end{enumerate}

\vspace{0.2cm}
\subsubsection*{Ejercicio 260 - Aspas}
\textbf{Enunciado:} $ 6 x^{2} - 5 x - 4 $\par
\begin{enumerate}[leftmargin=2em]
\item El metodo de aspas busca dos binomios. Los extremos deben producir el primer y el ultimo termino.
\item Probamos con \((2 x + 1)\) y \((3 x - 4)\).
\item Verificamos el termino del medio con los productos cruzados: \(- 8 x+3 x=- 5 x\).
\item Como los productos cruzados coinciden con el termino central, el resultado es: \[\boxed{\left(2 x + 1\right)\left(3 x - 4\right)}\]
\end{enumerate}

\vspace{0.2cm}
\nivel{4}{avanzado}
\subsubsection*{Ejercicio 261 - Aspas}
\textbf{Enunciado:} $ 6 x^{2} + 5 x y - 6 y^{2} $\par
\begin{enumerate}[leftmargin=2em]
\item El metodo de aspas busca dos binomios. Los extremos deben producir el primer y el ultimo termino.
\item Probamos con \((3 x - 2 y)\) y \((2 x + 3 y)\).
\item Verificamos el termino del medio con los productos cruzados: \(- 6 y^{2}+6 x^{2}=6 x^{2} - 6 y^{2}\).
\item Como los productos cruzados coinciden con el termino central, el resultado es: \[\boxed{\left(3 x - 2 y\right)\left(2 x + 3 y\right)}\]
\end{enumerate}

\vspace{0.2cm}
\subsubsection*{Ejercicio 262 - Aspas}
\textbf{Enunciado:} $ 12 x^{2} - 7 x y - 12 y^{2} $\par
\begin{enumerate}[leftmargin=2em]
\item El metodo de aspas busca dos binomios. Los extremos deben producir el primer y el ultimo termino.
\item Probamos con \((4 x + 3 y)\) y \((3 x - 4 y)\).
\item Verificamos el termino del medio con los productos cruzados: \(9 x y- 16 x y=- 7 x y\).
\item Como los productos cruzados coinciden con el termino central, el resultado es: \[\boxed{\left(4 x + 3 y\right)\left(3 x - 4 y\right)}\]
\end{enumerate}

\vspace{0.2cm}
\subsubsection*{Ejercicio 263 - Aspas}
\textbf{Enunciado:} $ 5 x^{2} + 21 x y - 20 y^{2} $\par
\begin{enumerate}[leftmargin=2em]
\item El metodo de aspas busca dos binomios. Los extremos deben producir el primer y el ultimo termino.
\item Probamos con \((5 x - 4 y)\) y \((x + 5 y)\).
\item Verificamos el termino del medio con los productos cruzados: \(- 20 y^{2}+5 x^{2}=5 x^{2} - 20 y^{2}\).
\item Como los productos cruzados coinciden con el termino central, el resultado es: \[\boxed{\left(5 x - 4 y\right)\left(x + 5 y\right)}\]
\end{enumerate}

\vspace{0.2cm}
\subsubsection*{Ejercicio 264 - Aspas}
\textbf{Enunciado:} $ 4 x^{2} - 2 x y - 30 y^{2} $\par
\begin{enumerate}[leftmargin=2em]
\item El metodo de aspas busca dos binomios. Los extremos deben producir el primer y el ultimo termino.
\item Probamos con \((2 x + 5 y)\) y \((2 x - 6 y)\).
\item Verificamos el termino del medio con los productos cruzados: \(4 x^{2}- 30 y^{2}=4 x^{2} - 30 y^{2}\).
\item Como los productos cruzados coinciden con el termino central, el resultado es: \[\boxed{\left(2 x + 5 y\right)\left(2 x - 6 y\right)}\]
\end{enumerate}

\vspace{0.2cm}
\subsubsection*{Ejercicio 265 - Aspas}
\textbf{Enunciado:} $ 9 x^{2} + 18 x y - 7 y^{2} $\par
\begin{enumerate}[leftmargin=2em]
\item El metodo de aspas busca dos binomios. Los extremos deben producir el primer y el ultimo termino.
\item Probamos con \((3 x - y)\) y \((3 x + 7 y)\).
\item Verificamos el termino del medio con los productos cruzados: \(- 7 y^{2}+9 x^{2}=9 x^{2} - 7 y^{2}\).
\item Como los productos cruzados coinciden con el termino central, el resultado es: \[\boxed{\left(3 x - y\right)\left(3 x + 7 y\right)}\]
\end{enumerate}

\vspace{0.2cm}
\subsubsection*{Ejercicio 266 - Aspas}
\textbf{Enunciado:} $ 4 x^{2} - 6 x y - 4 y^{2} $\par
\begin{enumerate}[leftmargin=2em]
\item El metodo de aspas busca dos binomios. Los extremos deben producir el primer y el ultimo termino.
\item Probamos con \((4 x + 2 y)\) y \((x - 2 y)\).
\item Verificamos el termino del medio con los productos cruzados: \(- 4 y^{2}+4 x^{2}=4 x^{2} - 4 y^{2}\).
\item Como los productos cruzados coinciden con el termino central, el resultado es: \[\boxed{\left(4 x + 2 y\right)\left(x - 2 y\right)}\]
\end{enumerate}

\vspace{0.2cm}
\subsubsection*{Ejercicio 267 - Aspas}
\textbf{Enunciado:} $ 10 x^{2} + 9 x y - 9 y^{2} $\par
\begin{enumerate}[leftmargin=2em]
\item El metodo de aspas busca dos binomios. Los extremos deben producir el primer y el ultimo termino.
\item Probamos con \((5 x - 3 y)\) y \((2 x + 3 y)\).
\item Verificamos el termino del medio con los productos cruzados: \(- 9 y^{2}+10 x^{2}=10 x^{2} - 9 y^{2}\).
\item Como los productos cruzados coinciden con el termino central, el resultado es: \[\boxed{\left(5 x - 3 y\right)\left(2 x + 3 y\right)}\]
\end{enumerate}

\vspace{0.2cm}
\subsubsection*{Ejercicio 268 - Aspas}
\textbf{Enunciado:} $ 6 x^{2} + 4 x y - 16 y^{2} $\par
\begin{enumerate}[leftmargin=2em]
\item El metodo de aspas busca dos binomios. Los extremos deben producir el primer y el ultimo termino.
\item Probamos con \((2 x + 4 y)\) y \((3 x - 4 y)\).
\item Verificamos el termino del medio con los productos cruzados: \(6 x^{2}- 16 y^{2}=6 x^{2} - 16 y^{2}\).
\item Como los productos cruzados coinciden con el termino central, el resultado es: \[\boxed{\left(2 x + 4 y\right)\left(3 x - 4 y\right)}\]
\end{enumerate}

\vspace{0.2cm}
\subsubsection*{Ejercicio 269 - Aspas}
\textbf{Enunciado:} $ 3 x^{2} + 10 x y - 25 y^{2} $\par
\begin{enumerate}[leftmargin=2em]
\item El metodo de aspas busca dos binomios. Los extremos deben producir el primer y el ultimo termino.
\item Probamos con \((3 x - 5 y)\) y \((x + 5 y)\).
\item Verificamos el termino del medio con los productos cruzados: \(- 25 y^{2}+3 x^{2}=3 x^{2} - 25 y^{2}\).
\item Como los productos cruzados coinciden con el termino central, el resultado es: \[\boxed{\left(3 x - 5 y\right)\left(x + 5 y\right)}\]
\end{enumerate}

\vspace{0.2cm}
\subsubsection*{Ejercicio 270 - Aspas}
\textbf{Enunciado:} $ 8 x^{2} - 22 x y - 6 y^{2} $\par
\begin{enumerate}[leftmargin=2em]
\item El metodo de aspas busca dos binomios. Los extremos deben producir el primer y el ultimo termino.
\item Probamos con \((4 x + y)\) y \((2 x - 6 y)\).
\item Verificamos el termino del medio con los productos cruzados: \(2 x y- 24 x y=- 22 x y\).
\item Como los productos cruzados coinciden con el termino central, el resultado es: \[\boxed{\left(4 x + y\right)\left(2 x - 6 y\right)}\]
\end{enumerate}

\vspace{0.2cm}
\subsubsection*{Ejercicio 271 - Aspas}
\textbf{Enunciado:} $ 15 x^{2} + 29 x y - 14 y^{2} $\par
\begin{enumerate}[leftmargin=2em]
\item El metodo de aspas busca dos binomios. Los extremos deben producir el primer y el ultimo termino.
\item Probamos con \((5 x - 2 y)\) y \((3 x + 7 y)\).
\item Verificamos el termino del medio con los productos cruzados: \(- 14 y^{2}+15 x^{2}=15 x^{2} - 14 y^{2}\).
\item Como los productos cruzados coinciden con el termino central, el resultado es: \[\boxed{\left(5 x - 2 y\right)\left(3 x + 7 y\right)}\]
\end{enumerate}

\vspace{0.2cm}
\subsubsection*{Ejercicio 272 - Aspas}
\textbf{Enunciado:} $ 2 x^{2} - x y - 6 y^{2} $\par
\begin{enumerate}[leftmargin=2em]
\item El metodo de aspas busca dos binomios. Los extremos deben producir el primer y el ultimo termino.
\item Probamos con \((2 x + 3 y)\) y \((x - 2 y)\).
\item Verificamos el termino del medio con los productos cruzados: \(- 4 x y+3 x y=- x y\).
\item Como los productos cruzados coinciden con el termino central, el resultado es: \[\boxed{\left(2 x + 3 y\right)\left(x - 2 y\right)}\]
\end{enumerate}

\vspace{0.2cm}
\subsubsection*{Ejercicio 273 - Aspas}
\textbf{Enunciado:} $ 6 x^{2} + x y - 12 y^{2} $\par
\begin{enumerate}[leftmargin=2em]
\item El metodo de aspas busca dos binomios. Los extremos deben producir el primer y el ultimo termino.
\item Probamos con \((3 x - 4 y)\) y \((2 x + 3 y)\).
\item Verificamos el termino del medio con los productos cruzados: \(- 12 y^{2}+6 x^{2}=6 x^{2} - 12 y^{2}\).
\item Como los productos cruzados coinciden con el termino central, el resultado es: \[\boxed{\left(3 x - 4 y\right)\left(2 x + 3 y\right)}\]
\end{enumerate}

\vspace{0.2cm}
\subsubsection*{Ejercicio 274 - Aspas}
\textbf{Enunciado:} $ 12 x^{2} - x y - 20 y^{2} $\par
\begin{enumerate}[leftmargin=2em]
\item El metodo de aspas busca dos binomios. Los extremos deben producir el primer y el ultimo termino.
\item Probamos con \((4 x + 5 y)\) y \((3 x - 4 y)\).
\item Verificamos el termino del medio con los productos cruzados: \(12 x^{2}- 20 y^{2}=12 x^{2} - 20 y^{2}\).
\item Como los productos cruzados coinciden con el termino central, el resultado es: \[\boxed{\left(4 x + 5 y\right)\left(3 x - 4 y\right)}\]
\end{enumerate}

\vspace{0.2cm}
\subsubsection*{Ejercicio 275 - Aspas}
\textbf{Enunciado:} $ 5 x^{2} + 24 x y - 5 y^{2} $\par
\begin{enumerate}[leftmargin=2em]
\item El metodo de aspas busca dos binomios. Los extremos deben producir el primer y el ultimo termino.
\item Probamos con \((5 x - y)\) y \((x + 5 y)\).
\item Verificamos el termino del medio con los productos cruzados: \(- 5 y^{2}+5 x^{2}=5 x^{2} - 5 y^{2}\).
\item Como los productos cruzados coinciden con el termino central, el resultado es: \[\boxed{\left(5 x - y\right)\left(x + 5 y\right)}\]
\end{enumerate}

\vspace{0.2cm}
\subsubsection*{Ejercicio 276 - Aspas}
\textbf{Enunciado:} $ 4 x^{2} - 8 x y - 12 y^{2} $\par
\begin{enumerate}[leftmargin=2em]
\item El metodo de aspas busca dos binomios. Los extremos deben producir el primer y el ultimo termino.
\item Probamos con \((2 x + 2 y)\) y \((2 x - 6 y)\).
\item Verificamos el termino del medio con los productos cruzados: \(4 x^{2}- 12 y^{2}=4 x^{2} - 12 y^{2}\).
\item Como los productos cruzados coinciden con el termino central, el resultado es: \[\boxed{\left(2 x + 2 y\right)\left(2 x - 6 y\right)}\]
\end{enumerate}

\vspace{0.2cm}
\subsubsection*{Ejercicio 277 - Aspas}
\textbf{Enunciado:} $ 9 x^{2} + 12 x y - 21 y^{2} $\par
\begin{enumerate}[leftmargin=2em]
\item El metodo de aspas busca dos binomios. Los extremos deben producir el primer y el ultimo termino.
\item Probamos con \((3 x - 3 y)\) y \((3 x + 7 y)\).
\item Verificamos el termino del medio con los productos cruzados: \(- 21 y^{2}+9 x^{2}=9 x^{2} - 21 y^{2}\).
\item Como los productos cruzados coinciden con el termino central, el resultado es: \[\boxed{\left(3 x - 3 y\right)\left(3 x + 7 y\right)}\]
\end{enumerate}

\vspace{0.2cm}
\subsubsection*{Ejercicio 278 - Aspas}
\textbf{Enunciado:} $ 4 x^{2} - 4 x y - 8 y^{2} $\par
\begin{enumerate}[leftmargin=2em]
\item El metodo de aspas busca dos binomios. Los extremos deben producir el primer y el ultimo termino.
\item Probamos con \((4 x + 4 y)\) y \((x - 2 y)\).
\item Verificamos el termino del medio con los productos cruzados: \(- 8 x y+4 x y=- 4 x y\).
\item Como los productos cruzados coinciden con el termino central, el resultado es: \[\boxed{\left(4 x + 4 y\right)\left(x - 2 y\right)}\]
\end{enumerate}

\vspace{0.2cm}
\subsubsection*{Ejercicio 279 - Aspas}
\textbf{Enunciado:} $ 10 x^{2} + 5 x y - 15 y^{2} $\par
\begin{enumerate}[leftmargin=2em]
\item El metodo de aspas busca dos binomios. Los extremos deben producir el primer y el ultimo termino.
\item Probamos con \((5 x - 5 y)\) y \((2 x + 3 y)\).
\item Verificamos el termino del medio con los productos cruzados: \(- 15 y^{2}+10 x^{2}=10 x^{2} - 15 y^{2}\).
\item Como los productos cruzados coinciden con el termino central, el resultado es: \[\boxed{\left(5 x - 5 y\right)\left(2 x + 3 y\right)}\]
\end{enumerate}

\vspace{0.2cm}
\subsubsection*{Ejercicio 280 - Aspas}
\textbf{Enunciado:} $ 6 x^{2} - 5 x y - 4 y^{2} $\par
\begin{enumerate}[leftmargin=2em]
\item El metodo de aspas busca dos binomios. Los extremos deben producir el primer y el ultimo termino.
\item Probamos con \((2 x + y)\) y \((3 x - 4 y)\).
\item Verificamos el termino del medio con los productos cruzados: \(3 x y- 8 x y=- 5 x y\).
\item Como los productos cruzados coinciden con el termino central, el resultado es: \[\boxed{\left(2 x + y\right)\left(3 x - 4 y\right)}\]
\end{enumerate}

\vspace{0.2cm}
\nivel{5}{imposible}
\subsubsection*{Ejercicio 281 - Aspas}
\textbf{Enunciado:} $ 9 x^{4} + 6 x^{2} y - 8 y^{2} $\par
\begin{enumerate}[leftmargin=2em]
\item El metodo de aspas busca dos binomios. Los extremos deben producir el primer y el ultimo termino.
\item Probamos con \((3 x^{2} - 2 y)\) y \((3 x^{2} + 4 y)\).
\item Verificamos el termino del medio con los productos cruzados: \(12 x^{2} y- 6 x^{2} y=6 x^{2} y\).
\item Como los productos cruzados coinciden con el termino central, el resultado es: \[\boxed{\left(3 x^{2} - 2 y\right)\left(3 x^{2} + 4 y\right)}\]
\end{enumerate}

\vspace{0.2cm}
\subsubsection*{Ejercicio 282 - Aspas}
\textbf{Enunciado:} $ 16 x^{4} - 8 x^{2} y - 15 y^{2} $\par
\begin{enumerate}[leftmargin=2em]
\item El metodo de aspas busca dos binomios. Los extremos deben producir el primer y el ultimo termino.
\item Probamos con \((4 x^{2} + 3 y)\) y \((4 x^{2} - 5 y)\).
\item Verificamos el termino del medio con los productos cruzados: \(- 20 x^{2} y+12 x^{2} y=- 8 x^{2} y\).
\item Como los productos cruzados coinciden con el termino central, el resultado es: \[\boxed{\left(4 x^{2} + 3 y\right)\left(4 x^{2} - 5 y\right)}\]
\end{enumerate}

\vspace{0.2cm}
\subsubsection*{Ejercicio 283 - Aspas}
\textbf{Enunciado:} $ 25 x^{4} + 10 x^{2} y - 24 y^{2} $\par
\begin{enumerate}[leftmargin=2em]
\item El metodo de aspas busca dos binomios. Los extremos deben producir el primer y el ultimo termino.
\item Probamos con \((5 x^{2} - 4 y)\) y \((5 x^{2} + 6 y)\).
\item Verificamos el termino del medio con los productos cruzados: \(30 x^{2} y- 20 x^{2} y=10 x^{2} y\).
\item Como los productos cruzados coinciden con el termino central, el resultado es: \[\boxed{\left(5 x^{2} - 4 y\right)\left(5 x^{2} + 6 y\right)}\]
\end{enumerate}

\vspace{0.2cm}
\subsubsection*{Ejercicio 284 - Aspas}
\textbf{Enunciado:} $ 12 x^{4} - 32 x^{2} y - 35 y^{2} $\par
\begin{enumerate}[leftmargin=2em]
\item El metodo de aspas busca dos binomios. Los extremos deben producir el primer y el ultimo termino.
\item Probamos con \((6 x^{2} + 5 y)\) y \((2 x^{2} - 7 y)\).
\item Verificamos el termino del medio con los productos cruzados: \(- 42 x^{2} y+10 x^{2} y=- 32 x^{2} y\).
\item Como los productos cruzados coinciden con el termino central, el resultado es: \[\boxed{\left(6 x^{2} + 5 y\right)\left(2 x^{2} - 7 y\right)}\]
\end{enumerate}

\vspace{0.2cm}
\subsubsection*{Ejercicio 285 - Aspas}
\textbf{Enunciado:} $ 6 x^{4} - 2 x^{2} y - 48 y^{2} $\par
\begin{enumerate}[leftmargin=2em]
\item El metodo de aspas busca dos binomios. Los extremos deben producir el primer y el ultimo termino.
\item Probamos con \((2 x^{2} - 6 y)\) y \((3 x^{2} + 8 y)\).
\item Verificamos el termino del medio con los productos cruzados: \(16 x^{2} y- 18 x^{2} y=- 2 x^{2} y\).
\item Como los productos cruzados coinciden con el termino central, el resultado es: \[\boxed{\left(2 x^{2} - 6 y\right)\left(3 x^{2} + 8 y\right)}\]
\end{enumerate}

\vspace{0.2cm}
\subsubsection*{Ejercicio 286 - Aspas}
\textbf{Enunciado:} $ 12 x^{4} - 23 x^{2} y - 9 y^{2} $\par
\begin{enumerate}[leftmargin=2em]
\item El metodo de aspas busca dos binomios. Los extremos deben producir el primer y el ultimo termino.
\item Probamos con \((3 x^{2} + y)\) y \((4 x^{2} - 9 y)\).
\item Verificamos el termino del medio con los productos cruzados: \(- 27 x^{2} y+4 x^{2} y=- 23 x^{2} y\).
\item Como los productos cruzados coinciden con el termino central, el resultado es: \[\boxed{\left(3 x^{2} + y\right)\left(4 x^{2} - 9 y\right)}\]
\end{enumerate}

\vspace{0.2cm}
\subsubsection*{Ejercicio 287 - Aspas}
\textbf{Enunciado:} $ 20 x^{4} + 2 x^{2} y - 6 y^{2} $\par
\begin{enumerate}[leftmargin=2em]
\item El metodo de aspas busca dos binomios. Los extremos deben producir el primer y el ultimo termino.
\item Probamos con \((4 x^{2} - 2 y)\) y \((5 x^{2} + 3 y)\).
\item Verificamos el termino del medio con los productos cruzados: \(12 x^{2} y- 10 x^{2} y=2 x^{2} y\).
\item Como los productos cruzados coinciden con el termino central, el resultado es: \[\boxed{\left(4 x^{2} - 2 y\right)\left(5 x^{2} + 3 y\right)}\]
\end{enumerate}

\vspace{0.2cm}
\subsubsection*{Ejercicio 288 - Aspas}
\textbf{Enunciado:} $ 10 x^{4} - 14 x^{2} y - 12 y^{2} $\par
\begin{enumerate}[leftmargin=2em]
\item El metodo de aspas busca dos binomios. Los extremos deben producir el primer y el ultimo termino.
\item Probamos con \((5 x^{2} + 3 y)\) y \((2 x^{2} - 4 y)\).
\item Verificamos el termino del medio con los productos cruzados: \(- 20 x^{2} y+6 x^{2} y=- 14 x^{2} y\).
\item Como los productos cruzados coinciden con el termino central, el resultado es: \[\boxed{\left(5 x^{2} + 3 y\right)\left(2 x^{2} - 4 y\right)}\]
\end{enumerate}

\vspace{0.2cm}
\subsubsection*{Ejercicio 289 - Aspas}
\textbf{Enunciado:} $ 18 x^{4} + 18 x^{2} y - 20 y^{2} $\par
\begin{enumerate}[leftmargin=2em]
\item El metodo de aspas busca dos binomios. Los extremos deben producir el primer y el ultimo termino.
\item Probamos con \((6 x^{2} - 4 y)\) y \((3 x^{2} + 5 y)\).
\item Verificamos el termino del medio con los productos cruzados: \(30 x^{2} y- 12 x^{2} y=18 x^{2} y\).
\item Como los productos cruzados coinciden con el termino central, el resultado es: \[\boxed{\left(6 x^{2} - 4 y\right)\left(3 x^{2} + 5 y\right)}\]
\end{enumerate}

\vspace{0.2cm}
\subsubsection*{Ejercicio 290 - Aspas}
\textbf{Enunciado:} $ 8 x^{4} + 8 x^{2} y - 30 y^{2} $\par
\begin{enumerate}[leftmargin=2em]
\item El metodo de aspas busca dos binomios. Los extremos deben producir el primer y el ultimo termino.
\item Probamos con \((2 x^{2} + 5 y)\) y \((4 x^{2} - 6 y)\).
\item Verificamos el termino del medio con los productos cruzados: \(- 12 x^{2} y+20 x^{2} y=8 x^{2} y\).
\item Como los productos cruzados coinciden con el termino central, el resultado es: \[\boxed{\left(2 x^{2} + 5 y\right)\left(4 x^{2} - 6 y\right)}\]
\end{enumerate}

\vspace{0.2cm}
\subsubsection*{Ejercicio 291 - Aspas}
\textbf{Enunciado:} $ 15 x^{4} - 9 x^{2} y - 42 y^{2} $\par
\begin{enumerate}[leftmargin=2em]
\item El metodo de aspas busca dos binomios. Los extremos deben producir el primer y el ultimo termino.
\item Probamos con \((3 x^{2} - 6 y)\) y \((5 x^{2} + 7 y)\).
\item Verificamos el termino del medio con los productos cruzados: \(21 x^{2} y- 30 x^{2} y=- 9 x^{2} y\).
\item Como los productos cruzados coinciden con el termino central, el resultado es: \[\boxed{\left(3 x^{2} - 6 y\right)\left(5 x^{2} + 7 y\right)}\]
\end{enumerate}

\vspace{0.2cm}
\subsubsection*{Ejercicio 292 - Aspas}
\textbf{Enunciado:} $ 8 x^{4} - 30 x^{2} y - 8 y^{2} $\par
\begin{enumerate}[leftmargin=2em]
\item El metodo de aspas busca dos binomios. Los extremos deben producir el primer y el ultimo termino.
\item Probamos con \((4 x^{2} + y)\) y \((2 x^{2} - 8 y)\).
\item Verificamos el termino del medio con los productos cruzados: \(- 32 x^{2} y+2 x^{2} y=- 30 x^{2} y\).
\item Como los productos cruzados coinciden con el termino central, el resultado es: \[\boxed{\left(4 x^{2} + y\right)\left(2 x^{2} - 8 y\right)}\]
\end{enumerate}

\vspace{0.2cm}
\subsubsection*{Ejercicio 293 - Aspas}
\textbf{Enunciado:} $ 15 x^{4} + 39 x^{2} y - 18 y^{2} $\par
\begin{enumerate}[leftmargin=2em]
\item El metodo de aspas busca dos binomios. Los extremos deben producir el primer y el ultimo termino.
\item Probamos con \((5 x^{2} - 2 y)\) y \((3 x^{2} + 9 y)\).
\item Verificamos el termino del medio con los productos cruzados: \(45 x^{2} y- 6 x^{2} y=39 x^{2} y\).
\item Como los productos cruzados coinciden con el termino central, el resultado es: \[\boxed{\left(5 x^{2} - 2 y\right)\left(3 x^{2} + 9 y\right)}\]
\end{enumerate}

\vspace{0.2cm}
\subsubsection*{Ejercicio 294 - Aspas}
\textbf{Enunciado:} $ 24 x^{4} - 6 x^{2} y - 9 y^{2} $\par
\begin{enumerate}[leftmargin=2em]
\item El metodo de aspas busca dos binomios. Los extremos deben producir el primer y el ultimo termino.
\item Probamos con \((6 x^{2} + 3 y)\) y \((4 x^{2} - 3 y)\).
\item Verificamos el termino del medio con los productos cruzados: \(- 18 x^{2} y+12 x^{2} y=- 6 x^{2} y\).
\item Como los productos cruzados coinciden con el termino central, el resultado es: \[\boxed{\left(6 x^{2} + 3 y\right)\left(4 x^{2} - 3 y\right)}\]
\end{enumerate}

\vspace{0.2cm}
\subsubsection*{Ejercicio 295 - Aspas}
\textbf{Enunciado:} $ 10 x^{4} - 12 x^{2} y - 16 y^{2} $\par
\begin{enumerate}[leftmargin=2em]
\item El metodo de aspas busca dos binomios. Los extremos deben producir el primer y el ultimo termino.
\item Probamos con \((2 x^{2} - 4 y)\) y \((5 x^{2} + 4 y)\).
\item Verificamos el termino del medio con los productos cruzados: \(8 x^{2} y- 20 x^{2} y=- 12 x^{2} y\).
\item Como los productos cruzados coinciden con el termino central, el resultado es: \[\boxed{\left(2 x^{2} - 4 y\right)\left(5 x^{2} + 4 y\right)}\]
\end{enumerate}

\vspace{0.2cm}
\subsubsection*{Ejercicio 296 - Aspas}
\textbf{Enunciado:} $ 6 x^{4} - 5 x^{2} y - 25 y^{2} $\par
\begin{enumerate}[leftmargin=2em]
\item El metodo de aspas busca dos binomios. Los extremos deben producir el primer y el ultimo termino.
\item Probamos con \((3 x^{2} + 5 y)\) y \((2 x^{2} - 5 y)\).
\item Verificamos el termino del medio con los productos cruzados: \(- 15 x^{2} y+10 x^{2} y=- 5 x^{2} y\).
\item Como los productos cruzados coinciden con el termino central, el resultado es: \[\boxed{\left(3 x^{2} + 5 y\right)\left(2 x^{2} - 5 y\right)}\]
\end{enumerate}

\vspace{0.2cm}
\subsubsection*{Ejercicio 297 - Aspas}
\textbf{Enunciado:} $ 12 x^{4} + 6 x^{2} y - 36 y^{2} $\par
\begin{enumerate}[leftmargin=2em]
\item El metodo de aspas busca dos binomios. Los extremos deben producir el primer y el ultimo termino.
\item Probamos con \((4 x^{2} - 6 y)\) y \((3 x^{2} + 6 y)\).
\item Verificamos el termino del medio con los productos cruzados: \(24 x^{2} y- 18 x^{2} y=6 x^{2} y\).
\item Como los productos cruzados coinciden con el termino central, el resultado es: \[\boxed{\left(4 x^{2} - 6 y\right)\left(3 x^{2} + 6 y\right)}\]
\end{enumerate}

\vspace{0.2cm}
\subsubsection*{Ejercicio 298 - Aspas}
\textbf{Enunciado:} $ 20 x^{4} - 31 x^{2} y - 7 y^{2} $\par
\begin{enumerate}[leftmargin=2em]
\item El metodo de aspas busca dos binomios. Los extremos deben producir el primer y el ultimo termino.
\item Probamos con \((5 x^{2} + y)\) y \((4 x^{2} - 7 y)\).
\item Verificamos el termino del medio con los productos cruzados: \(- 35 x^{2} y+4 x^{2} y=- 31 x^{2} y\).
\item Como los productos cruzados coinciden con el termino central, el resultado es: \[\boxed{\left(5 x^{2} + y\right)\left(4 x^{2} - 7 y\right)}\]
\end{enumerate}

\vspace{0.2cm}
\subsubsection*{Ejercicio 299 - Aspas}
\textbf{Enunciado:} $ 30 x^{4} + 38 x^{2} y - 16 y^{2} $\par
\begin{enumerate}[leftmargin=2em]
\item El metodo de aspas busca dos binomios. Los extremos deben producir el primer y el ultimo termino.
\item Probamos con \((6 x^{2} - 2 y)\) y \((5 x^{2} + 8 y)\).
\item Verificamos el termino del medio con los productos cruzados: \(48 x^{2} y- 10 x^{2} y=38 x^{2} y\).
\item Como los productos cruzados coinciden con el termino central, el resultado es: \[\boxed{\left(6 x^{2} - 2 y\right)\left(5 x^{2} + 8 y\right)}\]
\end{enumerate}

\vspace{0.2cm}
\subsubsection*{Ejercicio 300 - Aspas}
\textbf{Enunciado:} $ 4 x^{4} - 12 x^{2} y - 27 y^{2} $\par
\begin{enumerate}[leftmargin=2em]
\item El metodo de aspas busca dos binomios. Los extremos deben producir el primer y el ultimo termino.
\item Probamos con \((2 x^{2} + 3 y)\) y \((2 x^{2} - 9 y)\).
\item Verificamos el termino del medio con los productos cruzados: \(- 18 x^{2} y+6 x^{2} y=- 12 x^{2} y\).
\item Como los productos cruzados coinciden con el termino central, el resultado es: \[\boxed{\left(2 x^{2} + 3 y\right)\left(2 x^{2} - 9 y\right)}\]
\end{enumerate}

\vspace{0.2cm}
\section{Soluciones: Suma y diferencia de potencias}
\nivel{1}{basico}
\subsubsection*{Ejercicio 301 - Diferencia de potencias}
\textbf{Enunciado:} $ \left(x\right)^{3}-\left(2\right)^{3} $\par
\begin{enumerate}[leftmargin=2em]
\item Identificamos las bases: \(A=x\), \(B=2\) y el exponente \(n=3\).
\item Usamos la formula correspondiente: \[A^n-B^n=(A-B)\sum_{k=0}^{n-1}A^{n-1-k}B^k.\]
\item La letra \(k\) indica que se reemplaza por \(0,1,2,\ldots,n-1\). Eso genera todos los terminos del parentesis largo.
\item Sustituyendo los valores de \(A\), \(B\) y \(n\), queda: \[\boxed{\left(\left(x\right)-\left(2\right)\right)\sum_{k=0}^{2}\left(x\right)^{2-k}\left(2\right)^k}\]
\end{enumerate}

\vspace{0.2cm}
\subsubsection*{Ejercicio 302 - Suma de potencias}
\textbf{Enunciado:} $ \left(x\right)^{3}+\left(3\right)^{3} $\par
\begin{enumerate}[leftmargin=2em]
\item Identificamos las bases: \(A=x\), \(B=3\) y el exponente \(n=3\).
\item Usamos la formula correspondiente: \[A^n+B^n=(A+B)\sum_{k=0}^{n-1}(-1)^kA^{n-1-k}B^k\quad\text{cuando }n\text{ es impar}.\]
\item La letra \(k\) indica que se reemplaza por \(0,1,2,\ldots,n-1\). Eso genera todos los terminos del parentesis largo.
\item Sustituyendo los valores de \(A\), \(B\) y \(n\), queda: \[\boxed{\left(\left(x\right)+\left(3\right)\right)\sum_{k=0}^{2}(-1)^k\left(x\right)^{2-k}\left(3\right)^k}\]
\end{enumerate}

\vspace{0.2cm}
\subsubsection*{Ejercicio 303 - Diferencia de potencias}
\textbf{Enunciado:} $ \left(x\right)^{3}-\left(4\right)^{3} $\par
\begin{enumerate}[leftmargin=2em]
\item Identificamos las bases: \(A=x\), \(B=4\) y el exponente \(n=3\).
\item Usamos la formula correspondiente: \[A^n-B^n=(A-B)\sum_{k=0}^{n-1}A^{n-1-k}B^k.\]
\item La letra \(k\) indica que se reemplaza por \(0,1,2,\ldots,n-1\). Eso genera todos los terminos del parentesis largo.
\item Sustituyendo los valores de \(A\), \(B\) y \(n\), queda: \[\boxed{\left(\left(x\right)-\left(4\right)\right)\sum_{k=0}^{2}\left(x\right)^{2-k}\left(4\right)^k}\]
\end{enumerate}

\vspace{0.2cm}
\subsubsection*{Ejercicio 304 - Suma de potencias}
\textbf{Enunciado:} $ \left(x\right)^{3}+\left(5\right)^{3} $\par
\begin{enumerate}[leftmargin=2em]
\item Identificamos las bases: \(A=x\), \(B=5\) y el exponente \(n=3\).
\item Usamos la formula correspondiente: \[A^n+B^n=(A+B)\sum_{k=0}^{n-1}(-1)^kA^{n-1-k}B^k\quad\text{cuando }n\text{ es impar}.\]
\item La letra \(k\) indica que se reemplaza por \(0,1,2,\ldots,n-1\). Eso genera todos los terminos del parentesis largo.
\item Sustituyendo los valores de \(A\), \(B\) y \(n\), queda: \[\boxed{\left(\left(x\right)+\left(5\right)\right)\sum_{k=0}^{2}(-1)^k\left(x\right)^{2-k}\left(5\right)^k}\]
\end{enumerate}

\vspace{0.2cm}
\subsubsection*{Ejercicio 305 - Diferencia de potencias}
\textbf{Enunciado:} $ \left(x\right)^{3}-\left(1\right)^{3} $\par
\begin{enumerate}[leftmargin=2em]
\item Identificamos las bases: \(A=x\), \(B=1\) y el exponente \(n=3\).
\item Usamos la formula correspondiente: \[A^n-B^n=(A-B)\sum_{k=0}^{n-1}A^{n-1-k}B^k.\]
\item La letra \(k\) indica que se reemplaza por \(0,1,2,\ldots,n-1\). Eso genera todos los terminos del parentesis largo.
\item Sustituyendo los valores de \(A\), \(B\) y \(n\), queda: \[\boxed{\left(\left(x\right)-\left(1\right)\right)\sum_{k=0}^{2}\left(x\right)^{2-k}\left(1\right)^k}\]
\end{enumerate}

\vspace{0.2cm}
\subsubsection*{Ejercicio 306 - Suma de potencias}
\textbf{Enunciado:} $ \left(x\right)^{3}+\left(2\right)^{3} $\par
\begin{enumerate}[leftmargin=2em]
\item Identificamos las bases: \(A=x\), \(B=2\) y el exponente \(n=3\).
\item Usamos la formula correspondiente: \[A^n+B^n=(A+B)\sum_{k=0}^{n-1}(-1)^kA^{n-1-k}B^k\quad\text{cuando }n\text{ es impar}.\]
\item La letra \(k\) indica que se reemplaza por \(0,1,2,\ldots,n-1\). Eso genera todos los terminos del parentesis largo.
\item Sustituyendo los valores de \(A\), \(B\) y \(n\), queda: \[\boxed{\left(\left(x\right)+\left(2\right)\right)\sum_{k=0}^{2}(-1)^k\left(x\right)^{2-k}\left(2\right)^k}\]
\end{enumerate}

\vspace{0.2cm}
\subsubsection*{Ejercicio 307 - Diferencia de potencias}
\textbf{Enunciado:} $ \left(x\right)^{3}-\left(3\right)^{3} $\par
\begin{enumerate}[leftmargin=2em]
\item Identificamos las bases: \(A=x\), \(B=3\) y el exponente \(n=3\).
\item Usamos la formula correspondiente: \[A^n-B^n=(A-B)\sum_{k=0}^{n-1}A^{n-1-k}B^k.\]
\item La letra \(k\) indica que se reemplaza por \(0,1,2,\ldots,n-1\). Eso genera todos los terminos del parentesis largo.
\item Sustituyendo los valores de \(A\), \(B\) y \(n\), queda: \[\boxed{\left(\left(x\right)-\left(3\right)\right)\sum_{k=0}^{2}\left(x\right)^{2-k}\left(3\right)^k}\]
\end{enumerate}

\vspace{0.2cm}
\subsubsection*{Ejercicio 308 - Suma de potencias}
\textbf{Enunciado:} $ \left(x\right)^{3}+\left(4\right)^{3} $\par
\begin{enumerate}[leftmargin=2em]
\item Identificamos las bases: \(A=x\), \(B=4\) y el exponente \(n=3\).
\item Usamos la formula correspondiente: \[A^n+B^n=(A+B)\sum_{k=0}^{n-1}(-1)^kA^{n-1-k}B^k\quad\text{cuando }n\text{ es impar}.\]
\item La letra \(k\) indica que se reemplaza por \(0,1,2,\ldots,n-1\). Eso genera todos los terminos del parentesis largo.
\item Sustituyendo los valores de \(A\), \(B\) y \(n\), queda: \[\boxed{\left(\left(x\right)+\left(4\right)\right)\sum_{k=0}^{2}(-1)^k\left(x\right)^{2-k}\left(4\right)^k}\]
\end{enumerate}

\vspace{0.2cm}
\subsubsection*{Ejercicio 309 - Diferencia de potencias}
\textbf{Enunciado:} $ \left(x\right)^{3}-\left(5\right)^{3} $\par
\begin{enumerate}[leftmargin=2em]
\item Identificamos las bases: \(A=x\), \(B=5\) y el exponente \(n=3\).
\item Usamos la formula correspondiente: \[A^n-B^n=(A-B)\sum_{k=0}^{n-1}A^{n-1-k}B^k.\]
\item La letra \(k\) indica que se reemplaza por \(0,1,2,\ldots,n-1\). Eso genera todos los terminos del parentesis largo.
\item Sustituyendo los valores de \(A\), \(B\) y \(n\), queda: \[\boxed{\left(\left(x\right)-\left(5\right)\right)\sum_{k=0}^{2}\left(x\right)^{2-k}\left(5\right)^k}\]
\end{enumerate}

\vspace{0.2cm}
\subsubsection*{Ejercicio 310 - Suma de potencias}
\textbf{Enunciado:} $ \left(x\right)^{3}+\left(1\right)^{3} $\par
\begin{enumerate}[leftmargin=2em]
\item Identificamos las bases: \(A=x\), \(B=1\) y el exponente \(n=3\).
\item Usamos la formula correspondiente: \[A^n+B^n=(A+B)\sum_{k=0}^{n-1}(-1)^kA^{n-1-k}B^k\quad\text{cuando }n\text{ es impar}.\]
\item La letra \(k\) indica que se reemplaza por \(0,1,2,\ldots,n-1\). Eso genera todos los terminos del parentesis largo.
\item Sustituyendo los valores de \(A\), \(B\) y \(n\), queda: \[\boxed{\left(\left(x\right)+\left(1\right)\right)\sum_{k=0}^{2}(-1)^k\left(x\right)^{2-k}\left(1\right)^k}\]
\end{enumerate}

\vspace{0.2cm}
\subsubsection*{Ejercicio 311 - Diferencia de potencias}
\textbf{Enunciado:} $ \left(x\right)^{3}-\left(2\right)^{3} $\par
\begin{enumerate}[leftmargin=2em]
\item Identificamos las bases: \(A=x\), \(B=2\) y el exponente \(n=3\).
\item Usamos la formula correspondiente: \[A^n-B^n=(A-B)\sum_{k=0}^{n-1}A^{n-1-k}B^k.\]
\item La letra \(k\) indica que se reemplaza por \(0,1,2,\ldots,n-1\). Eso genera todos los terminos del parentesis largo.
\item Sustituyendo los valores de \(A\), \(B\) y \(n\), queda: \[\boxed{\left(\left(x\right)-\left(2\right)\right)\sum_{k=0}^{2}\left(x\right)^{2-k}\left(2\right)^k}\]
\end{enumerate}

\vspace{0.2cm}
\subsubsection*{Ejercicio 312 - Suma de potencias}
\textbf{Enunciado:} $ \left(x\right)^{3}+\left(3\right)^{3} $\par
\begin{enumerate}[leftmargin=2em]
\item Identificamos las bases: \(A=x\), \(B=3\) y el exponente \(n=3\).
\item Usamos la formula correspondiente: \[A^n+B^n=(A+B)\sum_{k=0}^{n-1}(-1)^kA^{n-1-k}B^k\quad\text{cuando }n\text{ es impar}.\]
\item La letra \(k\) indica que se reemplaza por \(0,1,2,\ldots,n-1\). Eso genera todos los terminos del parentesis largo.
\item Sustituyendo los valores de \(A\), \(B\) y \(n\), queda: \[\boxed{\left(\left(x\right)+\left(3\right)\right)\sum_{k=0}^{2}(-1)^k\left(x\right)^{2-k}\left(3\right)^k}\]
\end{enumerate}

\vspace{0.2cm}
\subsubsection*{Ejercicio 313 - Diferencia de potencias}
\textbf{Enunciado:} $ \left(x\right)^{3}-\left(4\right)^{3} $\par
\begin{enumerate}[leftmargin=2em]
\item Identificamos las bases: \(A=x\), \(B=4\) y el exponente \(n=3\).
\item Usamos la formula correspondiente: \[A^n-B^n=(A-B)\sum_{k=0}^{n-1}A^{n-1-k}B^k.\]
\item La letra \(k\) indica que se reemplaza por \(0,1,2,\ldots,n-1\). Eso genera todos los terminos del parentesis largo.
\item Sustituyendo los valores de \(A\), \(B\) y \(n\), queda: \[\boxed{\left(\left(x\right)-\left(4\right)\right)\sum_{k=0}^{2}\left(x\right)^{2-k}\left(4\right)^k}\]
\end{enumerate}

\vspace{0.2cm}
\subsubsection*{Ejercicio 314 - Suma de potencias}
\textbf{Enunciado:} $ \left(x\right)^{3}+\left(5\right)^{3} $\par
\begin{enumerate}[leftmargin=2em]
\item Identificamos las bases: \(A=x\), \(B=5\) y el exponente \(n=3\).
\item Usamos la formula correspondiente: \[A^n+B^n=(A+B)\sum_{k=0}^{n-1}(-1)^kA^{n-1-k}B^k\quad\text{cuando }n\text{ es impar}.\]
\item La letra \(k\) indica que se reemplaza por \(0,1,2,\ldots,n-1\). Eso genera todos los terminos del parentesis largo.
\item Sustituyendo los valores de \(A\), \(B\) y \(n\), queda: \[\boxed{\left(\left(x\right)+\left(5\right)\right)\sum_{k=0}^{2}(-1)^k\left(x\right)^{2-k}\left(5\right)^k}\]
\end{enumerate}

\vspace{0.2cm}
\subsubsection*{Ejercicio 315 - Diferencia de potencias}
\textbf{Enunciado:} $ \left(x\right)^{3}-\left(1\right)^{3} $\par
\begin{enumerate}[leftmargin=2em]
\item Identificamos las bases: \(A=x\), \(B=1\) y el exponente \(n=3\).
\item Usamos la formula correspondiente: \[A^n-B^n=(A-B)\sum_{k=0}^{n-1}A^{n-1-k}B^k.\]
\item La letra \(k\) indica que se reemplaza por \(0,1,2,\ldots,n-1\). Eso genera todos los terminos del parentesis largo.
\item Sustituyendo los valores de \(A\), \(B\) y \(n\), queda: \[\boxed{\left(\left(x\right)-\left(1\right)\right)\sum_{k=0}^{2}\left(x\right)^{2-k}\left(1\right)^k}\]
\end{enumerate}

\vspace{0.2cm}
\subsubsection*{Ejercicio 316 - Suma de potencias}
\textbf{Enunciado:} $ \left(x\right)^{3}+\left(2\right)^{3} $\par
\begin{enumerate}[leftmargin=2em]
\item Identificamos las bases: \(A=x\), \(B=2\) y el exponente \(n=3\).
\item Usamos la formula correspondiente: \[A^n+B^n=(A+B)\sum_{k=0}^{n-1}(-1)^kA^{n-1-k}B^k\quad\text{cuando }n\text{ es impar}.\]
\item La letra \(k\) indica que se reemplaza por \(0,1,2,\ldots,n-1\). Eso genera todos los terminos del parentesis largo.
\item Sustituyendo los valores de \(A\), \(B\) y \(n\), queda: \[\boxed{\left(\left(x\right)+\left(2\right)\right)\sum_{k=0}^{2}(-1)^k\left(x\right)^{2-k}\left(2\right)^k}\]
\end{enumerate}

\vspace{0.2cm}
\subsubsection*{Ejercicio 317 - Diferencia de potencias}
\textbf{Enunciado:} $ \left(x\right)^{3}-\left(3\right)^{3} $\par
\begin{enumerate}[leftmargin=2em]
\item Identificamos las bases: \(A=x\), \(B=3\) y el exponente \(n=3\).
\item Usamos la formula correspondiente: \[A^n-B^n=(A-B)\sum_{k=0}^{n-1}A^{n-1-k}B^k.\]
\item La letra \(k\) indica que se reemplaza por \(0,1,2,\ldots,n-1\). Eso genera todos los terminos del parentesis largo.
\item Sustituyendo los valores de \(A\), \(B\) y \(n\), queda: \[\boxed{\left(\left(x\right)-\left(3\right)\right)\sum_{k=0}^{2}\left(x\right)^{2-k}\left(3\right)^k}\]
\end{enumerate}

\vspace{0.2cm}
\subsubsection*{Ejercicio 318 - Suma de potencias}
\textbf{Enunciado:} $ \left(x\right)^{3}+\left(4\right)^{3} $\par
\begin{enumerate}[leftmargin=2em]
\item Identificamos las bases: \(A=x\), \(B=4\) y el exponente \(n=3\).
\item Usamos la formula correspondiente: \[A^n+B^n=(A+B)\sum_{k=0}^{n-1}(-1)^kA^{n-1-k}B^k\quad\text{cuando }n\text{ es impar}.\]
\item La letra \(k\) indica que se reemplaza por \(0,1,2,\ldots,n-1\). Eso genera todos los terminos del parentesis largo.
\item Sustituyendo los valores de \(A\), \(B\) y \(n\), queda: \[\boxed{\left(\left(x\right)+\left(4\right)\right)\sum_{k=0}^{2}(-1)^k\left(x\right)^{2-k}\left(4\right)^k}\]
\end{enumerate}

\vspace{0.2cm}
\subsubsection*{Ejercicio 319 - Diferencia de potencias}
\textbf{Enunciado:} $ \left(x\right)^{3}-\left(5\right)^{3} $\par
\begin{enumerate}[leftmargin=2em]
\item Identificamos las bases: \(A=x\), \(B=5\) y el exponente \(n=3\).
\item Usamos la formula correspondiente: \[A^n-B^n=(A-B)\sum_{k=0}^{n-1}A^{n-1-k}B^k.\]
\item La letra \(k\) indica que se reemplaza por \(0,1,2,\ldots,n-1\). Eso genera todos los terminos del parentesis largo.
\item Sustituyendo los valores de \(A\), \(B\) y \(n\), queda: \[\boxed{\left(\left(x\right)-\left(5\right)\right)\sum_{k=0}^{2}\left(x\right)^{2-k}\left(5\right)^k}\]
\end{enumerate}

\vspace{0.2cm}
\subsubsection*{Ejercicio 320 - Suma de potencias}
\textbf{Enunciado:} $ \left(x\right)^{3}+\left(1\right)^{3} $\par
\begin{enumerate}[leftmargin=2em]
\item Identificamos las bases: \(A=x\), \(B=1\) y el exponente \(n=3\).
\item Usamos la formula correspondiente: \[A^n+B^n=(A+B)\sum_{k=0}^{n-1}(-1)^kA^{n-1-k}B^k\quad\text{cuando }n\text{ es impar}.\]
\item La letra \(k\) indica que se reemplaza por \(0,1,2,\ldots,n-1\). Eso genera todos los terminos del parentesis largo.
\item Sustituyendo los valores de \(A\), \(B\) y \(n\), queda: \[\boxed{\left(\left(x\right)+\left(1\right)\right)\sum_{k=0}^{2}(-1)^k\left(x\right)^{2-k}\left(1\right)^k}\]
\end{enumerate}

\vspace{0.2cm}
\nivel{2}{intermedio inicial}
\subsubsection*{Ejercicio 321 - Diferencia de potencias}
\textbf{Enunciado:} $ \left(2 x\right)^{3}-\left(2 y\right)^{3} $\par
\begin{enumerate}[leftmargin=2em]
\item Identificamos las bases: \(A=2 x\), \(B=2 y\) y el exponente \(n=3\).
\item Usamos la formula correspondiente: \[A^n-B^n=(A-B)\sum_{k=0}^{n-1}A^{n-1-k}B^k.\]
\item La letra \(k\) indica que se reemplaza por \(0,1,2,\ldots,n-1\). Eso genera todos los terminos del parentesis largo.
\item Sustituyendo los valores de \(A\), \(B\) y \(n\), queda: \[\boxed{\left(\left(2 x\right)-\left(2 y\right)\right)\sum_{k=0}^{2}\left(2 x\right)^{2-k}\left(2 y\right)^k}\]
\end{enumerate}

\vspace{0.2cm}
\subsubsection*{Ejercicio 322 - Suma de potencias}
\textbf{Enunciado:} $ \left(3 x\right)^{5}+\left(3 y\right)^{5} $\par
\begin{enumerate}[leftmargin=2em]
\item Identificamos las bases: \(A=3 x\), \(B=3 y\) y el exponente \(n=5\).
\item Usamos la formula correspondiente: \[A^n+B^n=(A+B)\sum_{k=0}^{n-1}(-1)^kA^{n-1-k}B^k\quad\text{cuando }n\text{ es impar}.\]
\item La letra \(k\) indica que se reemplaza por \(0,1,2,\ldots,n-1\). Eso genera todos los terminos del parentesis largo.
\item Sustituyendo los valores de \(A\), \(B\) y \(n\), queda: \[\boxed{\left(\left(3 x\right)+\left(3 y\right)\right)\sum_{k=0}^{4}(-1)^k\left(3 x\right)^{4-k}\left(3 y\right)^k}\]
\end{enumerate}

\vspace{0.2cm}
\subsubsection*{Ejercicio 323 - Diferencia de potencias}
\textbf{Enunciado:} $ \left(x\right)^{4}-\left(4 y\right)^{4} $\par
\begin{enumerate}[leftmargin=2em]
\item Identificamos las bases: \(A=x\), \(B=4 y\) y el exponente \(n=4\).
\item Usamos la formula correspondiente: \[A^n-B^n=(A-B)\sum_{k=0}^{n-1}A^{n-1-k}B^k.\]
\item La letra \(k\) indica que se reemplaza por \(0,1,2,\ldots,n-1\). Eso genera todos los terminos del parentesis largo.
\item Sustituyendo los valores de \(A\), \(B\) y \(n\), queda: \[\boxed{\left(\left(x\right)-\left(4 y\right)\right)\sum_{k=0}^{3}\left(x\right)^{3-k}\left(4 y\right)^k}\]
\end{enumerate}

\vspace{0.2cm}
\subsubsection*{Ejercicio 324 - Suma de potencias}
\textbf{Enunciado:} $ \left(2 x\right)^{5}+\left(y\right)^{5} $\par
\begin{enumerate}[leftmargin=2em]
\item Identificamos las bases: \(A=2 x\), \(B=y\) y el exponente \(n=5\).
\item Usamos la formula correspondiente: \[A^n+B^n=(A+B)\sum_{k=0}^{n-1}(-1)^kA^{n-1-k}B^k\quad\text{cuando }n\text{ es impar}.\]
\item La letra \(k\) indica que se reemplaza por \(0,1,2,\ldots,n-1\). Eso genera todos los terminos del parentesis largo.
\item Sustituyendo los valores de \(A\), \(B\) y \(n\), queda: \[\boxed{\left(\left(2 x\right)+\left(y\right)\right)\sum_{k=0}^{4}(-1)^k\left(2 x\right)^{4-k}\left(y\right)^k}\]
\end{enumerate}

\vspace{0.2cm}
\subsubsection*{Ejercicio 325 - Diferencia de potencias}
\textbf{Enunciado:} $ \left(3 x\right)^{3}-\left(2 y\right)^{3} $\par
\begin{enumerate}[leftmargin=2em]
\item Identificamos las bases: \(A=3 x\), \(B=2 y\) y el exponente \(n=3\).
\item Usamos la formula correspondiente: \[A^n-B^n=(A-B)\sum_{k=0}^{n-1}A^{n-1-k}B^k.\]
\item La letra \(k\) indica que se reemplaza por \(0,1,2,\ldots,n-1\). Eso genera todos los terminos del parentesis largo.
\item Sustituyendo los valores de \(A\), \(B\) y \(n\), queda: \[\boxed{\left(\left(3 x\right)-\left(2 y\right)\right)\sum_{k=0}^{2}\left(3 x\right)^{2-k}\left(2 y\right)^k}\]
\end{enumerate}

\vspace{0.2cm}
\subsubsection*{Ejercicio 326 - Suma de potencias}
\textbf{Enunciado:} $ \left(x\right)^{5}+\left(3 y\right)^{5} $\par
\begin{enumerate}[leftmargin=2em]
\item Identificamos las bases: \(A=x\), \(B=3 y\) y el exponente \(n=5\).
\item Usamos la formula correspondiente: \[A^n+B^n=(A+B)\sum_{k=0}^{n-1}(-1)^kA^{n-1-k}B^k\quad\text{cuando }n\text{ es impar}.\]
\item La letra \(k\) indica que se reemplaza por \(0,1,2,\ldots,n-1\). Eso genera todos los terminos del parentesis largo.
\item Sustituyendo los valores de \(A\), \(B\) y \(n\), queda: \[\boxed{\left(\left(x\right)+\left(3 y\right)\right)\sum_{k=0}^{4}(-1)^k\left(x\right)^{4-k}\left(3 y\right)^k}\]
\end{enumerate}

\vspace{0.2cm}
\subsubsection*{Ejercicio 327 - Diferencia de potencias}
\textbf{Enunciado:} $ \left(2 x\right)^{3}-\left(4 y\right)^{3} $\par
\begin{enumerate}[leftmargin=2em]
\item Identificamos las bases: \(A=2 x\), \(B=4 y\) y el exponente \(n=3\).
\item Usamos la formula correspondiente: \[A^n-B^n=(A-B)\sum_{k=0}^{n-1}A^{n-1-k}B^k.\]
\item La letra \(k\) indica que se reemplaza por \(0,1,2,\ldots,n-1\). Eso genera todos los terminos del parentesis largo.
\item Sustituyendo los valores de \(A\), \(B\) y \(n\), queda: \[\boxed{\left(\left(2 x\right)-\left(4 y\right)\right)\sum_{k=0}^{2}\left(2 x\right)^{2-k}\left(4 y\right)^k}\]
\end{enumerate}

\vspace{0.2cm}
\subsubsection*{Ejercicio 328 - Suma de potencias}
\textbf{Enunciado:} $ \left(3 x\right)^{5}+\left(y\right)^{5} $\par
\begin{enumerate}[leftmargin=2em]
\item Identificamos las bases: \(A=3 x\), \(B=y\) y el exponente \(n=5\).
\item Usamos la formula correspondiente: \[A^n+B^n=(A+B)\sum_{k=0}^{n-1}(-1)^kA^{n-1-k}B^k\quad\text{cuando }n\text{ es impar}.\]
\item La letra \(k\) indica que se reemplaza por \(0,1,2,\ldots,n-1\). Eso genera todos los terminos del parentesis largo.
\item Sustituyendo los valores de \(A\), \(B\) y \(n\), queda: \[\boxed{\left(\left(3 x\right)+\left(y\right)\right)\sum_{k=0}^{4}(-1)^k\left(3 x\right)^{4-k}\left(y\right)^k}\]
\end{enumerate}

\vspace{0.2cm}
\subsubsection*{Ejercicio 329 - Diferencia de potencias}
\textbf{Enunciado:} $ \left(x\right)^{4}-\left(2 y\right)^{4} $\par
\begin{enumerate}[leftmargin=2em]
\item Identificamos las bases: \(A=x\), \(B=2 y\) y el exponente \(n=4\).
\item Usamos la formula correspondiente: \[A^n-B^n=(A-B)\sum_{k=0}^{n-1}A^{n-1-k}B^k.\]
\item La letra \(k\) indica que se reemplaza por \(0,1,2,\ldots,n-1\). Eso genera todos los terminos del parentesis largo.
\item Sustituyendo los valores de \(A\), \(B\) y \(n\), queda: \[\boxed{\left(\left(x\right)-\left(2 y\right)\right)\sum_{k=0}^{3}\left(x\right)^{3-k}\left(2 y\right)^k}\]
\end{enumerate}

\vspace{0.2cm}
\subsubsection*{Ejercicio 330 - Suma de potencias}
\textbf{Enunciado:} $ \left(2 x\right)^{5}+\left(3 y\right)^{5} $\par
\begin{enumerate}[leftmargin=2em]
\item Identificamos las bases: \(A=2 x\), \(B=3 y\) y el exponente \(n=5\).
\item Usamos la formula correspondiente: \[A^n+B^n=(A+B)\sum_{k=0}^{n-1}(-1)^kA^{n-1-k}B^k\quad\text{cuando }n\text{ es impar}.\]
\item La letra \(k\) indica que se reemplaza por \(0,1,2,\ldots,n-1\). Eso genera todos los terminos del parentesis largo.
\item Sustituyendo los valores de \(A\), \(B\) y \(n\), queda: \[\boxed{\left(\left(2 x\right)+\left(3 y\right)\right)\sum_{k=0}^{4}(-1)^k\left(2 x\right)^{4-k}\left(3 y\right)^k}\]
\end{enumerate}

\vspace{0.2cm}
\subsubsection*{Ejercicio 331 - Diferencia de potencias}
\textbf{Enunciado:} $ \left(3 x\right)^{3}-\left(4 y\right)^{3} $\par
\begin{enumerate}[leftmargin=2em]
\item Identificamos las bases: \(A=3 x\), \(B=4 y\) y el exponente \(n=3\).
\item Usamos la formula correspondiente: \[A^n-B^n=(A-B)\sum_{k=0}^{n-1}A^{n-1-k}B^k.\]
\item La letra \(k\) indica que se reemplaza por \(0,1,2,\ldots,n-1\). Eso genera todos los terminos del parentesis largo.
\item Sustituyendo los valores de \(A\), \(B\) y \(n\), queda: \[\boxed{\left(\left(3 x\right)-\left(4 y\right)\right)\sum_{k=0}^{2}\left(3 x\right)^{2-k}\left(4 y\right)^k}\]
\end{enumerate}

\vspace{0.2cm}
\subsubsection*{Ejercicio 332 - Suma de potencias}
\textbf{Enunciado:} $ \left(x\right)^{5}+\left(y\right)^{5} $\par
\begin{enumerate}[leftmargin=2em]
\item Identificamos las bases: \(A=x\), \(B=y\) y el exponente \(n=5\).
\item Usamos la formula correspondiente: \[A^n+B^n=(A+B)\sum_{k=0}^{n-1}(-1)^kA^{n-1-k}B^k\quad\text{cuando }n\text{ es impar}.\]
\item La letra \(k\) indica que se reemplaza por \(0,1,2,\ldots,n-1\). Eso genera todos los terminos del parentesis largo.
\item Sustituyendo los valores de \(A\), \(B\) y \(n\), queda: \[\boxed{\left(\left(x\right)+\left(y\right)\right)\sum_{k=0}^{4}(-1)^k\left(x\right)^{4-k}\left(y\right)^k}\]
\end{enumerate}

\vspace{0.2cm}
\subsubsection*{Ejercicio 333 - Diferencia de potencias}
\textbf{Enunciado:} $ \left(2 x\right)^{3}-\left(2 y\right)^{3} $\par
\begin{enumerate}[leftmargin=2em]
\item Identificamos las bases: \(A=2 x\), \(B=2 y\) y el exponente \(n=3\).
\item Usamos la formula correspondiente: \[A^n-B^n=(A-B)\sum_{k=0}^{n-1}A^{n-1-k}B^k.\]
\item La letra \(k\) indica que se reemplaza por \(0,1,2,\ldots,n-1\). Eso genera todos los terminos del parentesis largo.
\item Sustituyendo los valores de \(A\), \(B\) y \(n\), queda: \[\boxed{\left(\left(2 x\right)-\left(2 y\right)\right)\sum_{k=0}^{2}\left(2 x\right)^{2-k}\left(2 y\right)^k}\]
\end{enumerate}

\vspace{0.2cm}
\subsubsection*{Ejercicio 334 - Suma de potencias}
\textbf{Enunciado:} $ \left(3 x\right)^{5}+\left(3 y\right)^{5} $\par
\begin{enumerate}[leftmargin=2em]
\item Identificamos las bases: \(A=3 x\), \(B=3 y\) y el exponente \(n=5\).
\item Usamos la formula correspondiente: \[A^n+B^n=(A+B)\sum_{k=0}^{n-1}(-1)^kA^{n-1-k}B^k\quad\text{cuando }n\text{ es impar}.\]
\item La letra \(k\) indica que se reemplaza por \(0,1,2,\ldots,n-1\). Eso genera todos los terminos del parentesis largo.
\item Sustituyendo los valores de \(A\), \(B\) y \(n\), queda: \[\boxed{\left(\left(3 x\right)+\left(3 y\right)\right)\sum_{k=0}^{4}(-1)^k\left(3 x\right)^{4-k}\left(3 y\right)^k}\]
\end{enumerate}

\vspace{0.2cm}
\subsubsection*{Ejercicio 335 - Diferencia de potencias}
\textbf{Enunciado:} $ \left(x\right)^{4}-\left(4 y\right)^{4} $\par
\begin{enumerate}[leftmargin=2em]
\item Identificamos las bases: \(A=x\), \(B=4 y\) y el exponente \(n=4\).
\item Usamos la formula correspondiente: \[A^n-B^n=(A-B)\sum_{k=0}^{n-1}A^{n-1-k}B^k.\]
\item La letra \(k\) indica que se reemplaza por \(0,1,2,\ldots,n-1\). Eso genera todos los terminos del parentesis largo.
\item Sustituyendo los valores de \(A\), \(B\) y \(n\), queda: \[\boxed{\left(\left(x\right)-\left(4 y\right)\right)\sum_{k=0}^{3}\left(x\right)^{3-k}\left(4 y\right)^k}\]
\end{enumerate}

\vspace{0.2cm}
\subsubsection*{Ejercicio 336 - Suma de potencias}
\textbf{Enunciado:} $ \left(2 x\right)^{5}+\left(y\right)^{5} $\par
\begin{enumerate}[leftmargin=2em]
\item Identificamos las bases: \(A=2 x\), \(B=y\) y el exponente \(n=5\).
\item Usamos la formula correspondiente: \[A^n+B^n=(A+B)\sum_{k=0}^{n-1}(-1)^kA^{n-1-k}B^k\quad\text{cuando }n\text{ es impar}.\]
\item La letra \(k\) indica que se reemplaza por \(0,1,2,\ldots,n-1\). Eso genera todos los terminos del parentesis largo.
\item Sustituyendo los valores de \(A\), \(B\) y \(n\), queda: \[\boxed{\left(\left(2 x\right)+\left(y\right)\right)\sum_{k=0}^{4}(-1)^k\left(2 x\right)^{4-k}\left(y\right)^k}\]
\end{enumerate}

\vspace{0.2cm}
\subsubsection*{Ejercicio 337 - Diferencia de potencias}
\textbf{Enunciado:} $ \left(3 x\right)^{3}-\left(2 y\right)^{3} $\par
\begin{enumerate}[leftmargin=2em]
\item Identificamos las bases: \(A=3 x\), \(B=2 y\) y el exponente \(n=3\).
\item Usamos la formula correspondiente: \[A^n-B^n=(A-B)\sum_{k=0}^{n-1}A^{n-1-k}B^k.\]
\item La letra \(k\) indica que se reemplaza por \(0,1,2,\ldots,n-1\). Eso genera todos los terminos del parentesis largo.
\item Sustituyendo los valores de \(A\), \(B\) y \(n\), queda: \[\boxed{\left(\left(3 x\right)-\left(2 y\right)\right)\sum_{k=0}^{2}\left(3 x\right)^{2-k}\left(2 y\right)^k}\]
\end{enumerate}

\vspace{0.2cm}
\subsubsection*{Ejercicio 338 - Suma de potencias}
\textbf{Enunciado:} $ \left(x\right)^{5}+\left(3 y\right)^{5} $\par
\begin{enumerate}[leftmargin=2em]
\item Identificamos las bases: \(A=x\), \(B=3 y\) y el exponente \(n=5\).
\item Usamos la formula correspondiente: \[A^n+B^n=(A+B)\sum_{k=0}^{n-1}(-1)^kA^{n-1-k}B^k\quad\text{cuando }n\text{ es impar}.\]
\item La letra \(k\) indica que se reemplaza por \(0,1,2,\ldots,n-1\). Eso genera todos los terminos del parentesis largo.
\item Sustituyendo los valores de \(A\), \(B\) y \(n\), queda: \[\boxed{\left(\left(x\right)+\left(3 y\right)\right)\sum_{k=0}^{4}(-1)^k\left(x\right)^{4-k}\left(3 y\right)^k}\]
\end{enumerate}

\vspace{0.2cm}
\subsubsection*{Ejercicio 339 - Diferencia de potencias}
\textbf{Enunciado:} $ \left(2 x\right)^{3}-\left(4 y\right)^{3} $\par
\begin{enumerate}[leftmargin=2em]
\item Identificamos las bases: \(A=2 x\), \(B=4 y\) y el exponente \(n=3\).
\item Usamos la formula correspondiente: \[A^n-B^n=(A-B)\sum_{k=0}^{n-1}A^{n-1-k}B^k.\]
\item La letra \(k\) indica que se reemplaza por \(0,1,2,\ldots,n-1\). Eso genera todos los terminos del parentesis largo.
\item Sustituyendo los valores de \(A\), \(B\) y \(n\), queda: \[\boxed{\left(\left(2 x\right)-\left(4 y\right)\right)\sum_{k=0}^{2}\left(2 x\right)^{2-k}\left(4 y\right)^k}\]
\end{enumerate}

\vspace{0.2cm}
\subsubsection*{Ejercicio 340 - Suma de potencias}
\textbf{Enunciado:} $ \left(3 x\right)^{5}+\left(y\right)^{5} $\par
\begin{enumerate}[leftmargin=2em]
\item Identificamos las bases: \(A=3 x\), \(B=y\) y el exponente \(n=5\).
\item Usamos la formula correspondiente: \[A^n+B^n=(A+B)\sum_{k=0}^{n-1}(-1)^kA^{n-1-k}B^k\quad\text{cuando }n\text{ es impar}.\]
\item La letra \(k\) indica que se reemplaza por \(0,1,2,\ldots,n-1\). Eso genera todos los terminos del parentesis largo.
\item Sustituyendo los valores de \(A\), \(B\) y \(n\), queda: \[\boxed{\left(\left(3 x\right)+\left(y\right)\right)\sum_{k=0}^{4}(-1)^k\left(3 x\right)^{4-k}\left(y\right)^k}\]
\end{enumerate}

\vspace{0.2cm}
\nivel{3}{intermedio}
\subsubsection*{Ejercicio 341 - Diferencia de potencias}
\textbf{Enunciado:} $ \left(x + 2\right)^{4}-\left(2 y\right)^{4} $\par
\begin{enumerate}[leftmargin=2em]
\item Identificamos las bases: \(A=x + 2\), \(B=2 y\) y el exponente \(n=4\).
\item Usamos la formula correspondiente: \[A^n-B^n=(A-B)\sum_{k=0}^{n-1}A^{n-1-k}B^k.\]
\item La letra \(k\) indica que se reemplaza por \(0,1,2,\ldots,n-1\). Eso genera todos los terminos del parentesis largo.
\item Sustituyendo los valores de \(A\), \(B\) y \(n\), queda: \[\boxed{\left(\left(x + 2\right)-\left(2 y\right)\right)\sum_{k=0}^{3}\left(x + 2\right)^{3-k}\left(2 y\right)^k}\]
\end{enumerate}

\vspace{0.2cm}
\subsubsection*{Ejercicio 342 - Suma de potencias}
\textbf{Enunciado:} $ \left(x + 3\right)^{5}+\left(3 y\right)^{5} $\par
\begin{enumerate}[leftmargin=2em]
\item Identificamos las bases: \(A=x + 3\), \(B=3 y\) y el exponente \(n=5\).
\item Usamos la formula correspondiente: \[A^n+B^n=(A+B)\sum_{k=0}^{n-1}(-1)^kA^{n-1-k}B^k\quad\text{cuando }n\text{ es impar}.\]
\item La letra \(k\) indica que se reemplaza por \(0,1,2,\ldots,n-1\). Eso genera todos los terminos del parentesis largo.
\item Sustituyendo los valores de \(A\), \(B\) y \(n\), queda: \[\boxed{\left(\left(x + 3\right)+\left(3 y\right)\right)\sum_{k=0}^{4}(-1)^k\left(x + 3\right)^{4-k}\left(3 y\right)^k}\]
\end{enumerate}

\vspace{0.2cm}
\subsubsection*{Ejercicio 343 - Diferencia de potencias}
\textbf{Enunciado:} $ \left(x + 4\right)^{4}-\left(y\right)^{4} $\par
\begin{enumerate}[leftmargin=2em]
\item Identificamos las bases: \(A=x + 4\), \(B=y\) y el exponente \(n=4\).
\item Usamos la formula correspondiente: \[A^n-B^n=(A-B)\sum_{k=0}^{n-1}A^{n-1-k}B^k.\]
\item La letra \(k\) indica que se reemplaza por \(0,1,2,\ldots,n-1\). Eso genera todos los terminos del parentesis largo.
\item Sustituyendo los valores de \(A\), \(B\) y \(n\), queda: \[\boxed{\left(\left(x + 4\right)-\left(y\right)\right)\sum_{k=0}^{3}\left(x + 4\right)^{3-k}\left(y\right)^k}\]
\end{enumerate}

\vspace{0.2cm}
\subsubsection*{Ejercicio 344 - Suma de potencias}
\textbf{Enunciado:} $ \left(x + 1\right)^{5}+\left(2 y\right)^{5} $\par
\begin{enumerate}[leftmargin=2em]
\item Identificamos las bases: \(A=x + 1\), \(B=2 y\) y el exponente \(n=5\).
\item Usamos la formula correspondiente: \[A^n+B^n=(A+B)\sum_{k=0}^{n-1}(-1)^kA^{n-1-k}B^k\quad\text{cuando }n\text{ es impar}.\]
\item La letra \(k\) indica que se reemplaza por \(0,1,2,\ldots,n-1\). Eso genera todos los terminos del parentesis largo.
\item Sustituyendo los valores de \(A\), \(B\) y \(n\), queda: \[\boxed{\left(\left(x + 1\right)+\left(2 y\right)\right)\sum_{k=0}^{4}(-1)^k\left(x + 1\right)^{4-k}\left(2 y\right)^k}\]
\end{enumerate}

\vspace{0.2cm}
\subsubsection*{Ejercicio 345 - Diferencia de potencias}
\textbf{Enunciado:} $ \left(x + 2\right)^{4}-\left(3 y\right)^{4} $\par
\begin{enumerate}[leftmargin=2em]
\item Identificamos las bases: \(A=x + 2\), \(B=3 y\) y el exponente \(n=4\).
\item Usamos la formula correspondiente: \[A^n-B^n=(A-B)\sum_{k=0}^{n-1}A^{n-1-k}B^k.\]
\item La letra \(k\) indica que se reemplaza por \(0,1,2,\ldots,n-1\). Eso genera todos los terminos del parentesis largo.
\item Sustituyendo los valores de \(A\), \(B\) y \(n\), queda: \[\boxed{\left(\left(x + 2\right)-\left(3 y\right)\right)\sum_{k=0}^{3}\left(x + 2\right)^{3-k}\left(3 y\right)^k}\]
\end{enumerate}

\vspace{0.2cm}
\subsubsection*{Ejercicio 346 - Suma de potencias}
\textbf{Enunciado:} $ \left(x + 3\right)^{5}+\left(y\right)^{5} $\par
\begin{enumerate}[leftmargin=2em]
\item Identificamos las bases: \(A=x + 3\), \(B=y\) y el exponente \(n=5\).
\item Usamos la formula correspondiente: \[A^n+B^n=(A+B)\sum_{k=0}^{n-1}(-1)^kA^{n-1-k}B^k\quad\text{cuando }n\text{ es impar}.\]
\item La letra \(k\) indica que se reemplaza por \(0,1,2,\ldots,n-1\). Eso genera todos los terminos del parentesis largo.
\item Sustituyendo los valores de \(A\), \(B\) y \(n\), queda: \[\boxed{\left(\left(x + 3\right)+\left(y\right)\right)\sum_{k=0}^{4}(-1)^k\left(x + 3\right)^{4-k}\left(y\right)^k}\]
\end{enumerate}

\vspace{0.2cm}
\subsubsection*{Ejercicio 347 - Diferencia de potencias}
\textbf{Enunciado:} $ \left(x + 4\right)^{4}-\left(2 y\right)^{4} $\par
\begin{enumerate}[leftmargin=2em]
\item Identificamos las bases: \(A=x + 4\), \(B=2 y\) y el exponente \(n=4\).
\item Usamos la formula correspondiente: \[A^n-B^n=(A-B)\sum_{k=0}^{n-1}A^{n-1-k}B^k.\]
\item La letra \(k\) indica que se reemplaza por \(0,1,2,\ldots,n-1\). Eso genera todos los terminos del parentesis largo.
\item Sustituyendo los valores de \(A\), \(B\) y \(n\), queda: \[\boxed{\left(\left(x + 4\right)-\left(2 y\right)\right)\sum_{k=0}^{3}\left(x + 4\right)^{3-k}\left(2 y\right)^k}\]
\end{enumerate}

\vspace{0.2cm}
\subsubsection*{Ejercicio 348 - Suma de potencias}
\textbf{Enunciado:} $ \left(x + 1\right)^{5}+\left(3 y\right)^{5} $\par
\begin{enumerate}[leftmargin=2em]
\item Identificamos las bases: \(A=x + 1\), \(B=3 y\) y el exponente \(n=5\).
\item Usamos la formula correspondiente: \[A^n+B^n=(A+B)\sum_{k=0}^{n-1}(-1)^kA^{n-1-k}B^k\quad\text{cuando }n\text{ es impar}.\]
\item La letra \(k\) indica que se reemplaza por \(0,1,2,\ldots,n-1\). Eso genera todos los terminos del parentesis largo.
\item Sustituyendo los valores de \(A\), \(B\) y \(n\), queda: \[\boxed{\left(\left(x + 1\right)+\left(3 y\right)\right)\sum_{k=0}^{4}(-1)^k\left(x + 1\right)^{4-k}\left(3 y\right)^k}\]
\end{enumerate}

\vspace{0.2cm}
\subsubsection*{Ejercicio 349 - Diferencia de potencias}
\textbf{Enunciado:} $ \left(x + 2\right)^{4}-\left(y\right)^{4} $\par
\begin{enumerate}[leftmargin=2em]
\item Identificamos las bases: \(A=x + 2\), \(B=y\) y el exponente \(n=4\).
\item Usamos la formula correspondiente: \[A^n-B^n=(A-B)\sum_{k=0}^{n-1}A^{n-1-k}B^k.\]
\item La letra \(k\) indica que se reemplaza por \(0,1,2,\ldots,n-1\). Eso genera todos los terminos del parentesis largo.
\item Sustituyendo los valores de \(A\), \(B\) y \(n\), queda: \[\boxed{\left(\left(x + 2\right)-\left(y\right)\right)\sum_{k=0}^{3}\left(x + 2\right)^{3-k}\left(y\right)^k}\]
\end{enumerate}

\vspace{0.2cm}
\subsubsection*{Ejercicio 350 - Suma de potencias}
\textbf{Enunciado:} $ \left(x + 3\right)^{5}+\left(2 y\right)^{5} $\par
\begin{enumerate}[leftmargin=2em]
\item Identificamos las bases: \(A=x + 3\), \(B=2 y\) y el exponente \(n=5\).
\item Usamos la formula correspondiente: \[A^n+B^n=(A+B)\sum_{k=0}^{n-1}(-1)^kA^{n-1-k}B^k\quad\text{cuando }n\text{ es impar}.\]
\item La letra \(k\) indica que se reemplaza por \(0,1,2,\ldots,n-1\). Eso genera todos los terminos del parentesis largo.
\item Sustituyendo los valores de \(A\), \(B\) y \(n\), queda: \[\boxed{\left(\left(x + 3\right)+\left(2 y\right)\right)\sum_{k=0}^{4}(-1)^k\left(x + 3\right)^{4-k}\left(2 y\right)^k}\]
\end{enumerate}

\vspace{0.2cm}
\subsubsection*{Ejercicio 351 - Diferencia de potencias}
\textbf{Enunciado:} $ \left(x + 4\right)^{4}-\left(3 y\right)^{4} $\par
\begin{enumerate}[leftmargin=2em]
\item Identificamos las bases: \(A=x + 4\), \(B=3 y\) y el exponente \(n=4\).
\item Usamos la formula correspondiente: \[A^n-B^n=(A-B)\sum_{k=0}^{n-1}A^{n-1-k}B^k.\]
\item La letra \(k\) indica que se reemplaza por \(0,1,2,\ldots,n-1\). Eso genera todos los terminos del parentesis largo.
\item Sustituyendo los valores de \(A\), \(B\) y \(n\), queda: \[\boxed{\left(\left(x + 4\right)-\left(3 y\right)\right)\sum_{k=0}^{3}\left(x + 4\right)^{3-k}\left(3 y\right)^k}\]
\end{enumerate}

\vspace{0.2cm}
\subsubsection*{Ejercicio 352 - Suma de potencias}
\textbf{Enunciado:} $ \left(x + 1\right)^{5}+\left(y\right)^{5} $\par
\begin{enumerate}[leftmargin=2em]
\item Identificamos las bases: \(A=x + 1\), \(B=y\) y el exponente \(n=5\).
\item Usamos la formula correspondiente: \[A^n+B^n=(A+B)\sum_{k=0}^{n-1}(-1)^kA^{n-1-k}B^k\quad\text{cuando }n\text{ es impar}.\]
\item La letra \(k\) indica que se reemplaza por \(0,1,2,\ldots,n-1\). Eso genera todos los terminos del parentesis largo.
\item Sustituyendo los valores de \(A\), \(B\) y \(n\), queda: \[\boxed{\left(\left(x + 1\right)+\left(y\right)\right)\sum_{k=0}^{4}(-1)^k\left(x + 1\right)^{4-k}\left(y\right)^k}\]
\end{enumerate}

\vspace{0.2cm}
\subsubsection*{Ejercicio 353 - Diferencia de potencias}
\textbf{Enunciado:} $ \left(x + 2\right)^{4}-\left(2 y\right)^{4} $\par
\begin{enumerate}[leftmargin=2em]
\item Identificamos las bases: \(A=x + 2\), \(B=2 y\) y el exponente \(n=4\).
\item Usamos la formula correspondiente: \[A^n-B^n=(A-B)\sum_{k=0}^{n-1}A^{n-1-k}B^k.\]
\item La letra \(k\) indica que se reemplaza por \(0,1,2,\ldots,n-1\). Eso genera todos los terminos del parentesis largo.
\item Sustituyendo los valores de \(A\), \(B\) y \(n\), queda: \[\boxed{\left(\left(x + 2\right)-\left(2 y\right)\right)\sum_{k=0}^{3}\left(x + 2\right)^{3-k}\left(2 y\right)^k}\]
\end{enumerate}

\vspace{0.2cm}
\subsubsection*{Ejercicio 354 - Suma de potencias}
\textbf{Enunciado:} $ \left(x + 3\right)^{5}+\left(3 y\right)^{5} $\par
\begin{enumerate}[leftmargin=2em]
\item Identificamos las bases: \(A=x + 3\), \(B=3 y\) y el exponente \(n=5\).
\item Usamos la formula correspondiente: \[A^n+B^n=(A+B)\sum_{k=0}^{n-1}(-1)^kA^{n-1-k}B^k\quad\text{cuando }n\text{ es impar}.\]
\item La letra \(k\) indica que se reemplaza por \(0,1,2,\ldots,n-1\). Eso genera todos los terminos del parentesis largo.
\item Sustituyendo los valores de \(A\), \(B\) y \(n\), queda: \[\boxed{\left(\left(x + 3\right)+\left(3 y\right)\right)\sum_{k=0}^{4}(-1)^k\left(x + 3\right)^{4-k}\left(3 y\right)^k}\]
\end{enumerate}

\vspace{0.2cm}
\subsubsection*{Ejercicio 355 - Diferencia de potencias}
\textbf{Enunciado:} $ \left(x + 4\right)^{4}-\left(y\right)^{4} $\par
\begin{enumerate}[leftmargin=2em]
\item Identificamos las bases: \(A=x + 4\), \(B=y\) y el exponente \(n=4\).
\item Usamos la formula correspondiente: \[A^n-B^n=(A-B)\sum_{k=0}^{n-1}A^{n-1-k}B^k.\]
\item La letra \(k\) indica que se reemplaza por \(0,1,2,\ldots,n-1\). Eso genera todos los terminos del parentesis largo.
\item Sustituyendo los valores de \(A\), \(B\) y \(n\), queda: \[\boxed{\left(\left(x + 4\right)-\left(y\right)\right)\sum_{k=0}^{3}\left(x + 4\right)^{3-k}\left(y\right)^k}\]
\end{enumerate}

\vspace{0.2cm}
\subsubsection*{Ejercicio 356 - Suma de potencias}
\textbf{Enunciado:} $ \left(x + 1\right)^{5}+\left(2 y\right)^{5} $\par
\begin{enumerate}[leftmargin=2em]
\item Identificamos las bases: \(A=x + 1\), \(B=2 y\) y el exponente \(n=5\).
\item Usamos la formula correspondiente: \[A^n+B^n=(A+B)\sum_{k=0}^{n-1}(-1)^kA^{n-1-k}B^k\quad\text{cuando }n\text{ es impar}.\]
\item La letra \(k\) indica que se reemplaza por \(0,1,2,\ldots,n-1\). Eso genera todos los terminos del parentesis largo.
\item Sustituyendo los valores de \(A\), \(B\) y \(n\), queda: \[\boxed{\left(\left(x + 1\right)+\left(2 y\right)\right)\sum_{k=0}^{4}(-1)^k\left(x + 1\right)^{4-k}\left(2 y\right)^k}\]
\end{enumerate}

\vspace{0.2cm}
\subsubsection*{Ejercicio 357 - Diferencia de potencias}
\textbf{Enunciado:} $ \left(x + 2\right)^{4}-\left(3 y\right)^{4} $\par
\begin{enumerate}[leftmargin=2em]
\item Identificamos las bases: \(A=x + 2\), \(B=3 y\) y el exponente \(n=4\).
\item Usamos la formula correspondiente: \[A^n-B^n=(A-B)\sum_{k=0}^{n-1}A^{n-1-k}B^k.\]
\item La letra \(k\) indica que se reemplaza por \(0,1,2,\ldots,n-1\). Eso genera todos los terminos del parentesis largo.
\item Sustituyendo los valores de \(A\), \(B\) y \(n\), queda: \[\boxed{\left(\left(x + 2\right)-\left(3 y\right)\right)\sum_{k=0}^{3}\left(x + 2\right)^{3-k}\left(3 y\right)^k}\]
\end{enumerate}

\vspace{0.2cm}
\subsubsection*{Ejercicio 358 - Suma de potencias}
\textbf{Enunciado:} $ \left(x + 3\right)^{5}+\left(y\right)^{5} $\par
\begin{enumerate}[leftmargin=2em]
\item Identificamos las bases: \(A=x + 3\), \(B=y\) y el exponente \(n=5\).
\item Usamos la formula correspondiente: \[A^n+B^n=(A+B)\sum_{k=0}^{n-1}(-1)^kA^{n-1-k}B^k\quad\text{cuando }n\text{ es impar}.\]
\item La letra \(k\) indica que se reemplaza por \(0,1,2,\ldots,n-1\). Eso genera todos los terminos del parentesis largo.
\item Sustituyendo los valores de \(A\), \(B\) y \(n\), queda: \[\boxed{\left(\left(x + 3\right)+\left(y\right)\right)\sum_{k=0}^{4}(-1)^k\left(x + 3\right)^{4-k}\left(y\right)^k}\]
\end{enumerate}

\vspace{0.2cm}
\subsubsection*{Ejercicio 359 - Diferencia de potencias}
\textbf{Enunciado:} $ \left(x + 4\right)^{4}-\left(2 y\right)^{4} $\par
\begin{enumerate}[leftmargin=2em]
\item Identificamos las bases: \(A=x + 4\), \(B=2 y\) y el exponente \(n=4\).
\item Usamos la formula correspondiente: \[A^n-B^n=(A-B)\sum_{k=0}^{n-1}A^{n-1-k}B^k.\]
\item La letra \(k\) indica que se reemplaza por \(0,1,2,\ldots,n-1\). Eso genera todos los terminos del parentesis largo.
\item Sustituyendo los valores de \(A\), \(B\) y \(n\), queda: \[\boxed{\left(\left(x + 4\right)-\left(2 y\right)\right)\sum_{k=0}^{3}\left(x + 4\right)^{3-k}\left(2 y\right)^k}\]
\end{enumerate}

\vspace{0.2cm}
\subsubsection*{Ejercicio 360 - Suma de potencias}
\textbf{Enunciado:} $ \left(x + 1\right)^{5}+\left(3 y\right)^{5} $\par
\begin{enumerate}[leftmargin=2em]
\item Identificamos las bases: \(A=x + 1\), \(B=3 y\) y el exponente \(n=5\).
\item Usamos la formula correspondiente: \[A^n+B^n=(A+B)\sum_{k=0}^{n-1}(-1)^kA^{n-1-k}B^k\quad\text{cuando }n\text{ es impar}.\]
\item La letra \(k\) indica que se reemplaza por \(0,1,2,\ldots,n-1\). Eso genera todos los terminos del parentesis largo.
\item Sustituyendo los valores de \(A\), \(B\) y \(n\), queda: \[\boxed{\left(\left(x + 1\right)+\left(3 y\right)\right)\sum_{k=0}^{4}(-1)^k\left(x + 1\right)^{4-k}\left(3 y\right)^k}\]
\end{enumerate}

\vspace{0.2cm}
\nivel{4}{avanzado}
\subsubsection*{Ejercicio 361 - Diferencia de potencias}
\textbf{Enunciado:} $ \left(2 x - 2 y\right)^{6}-\left(2 x + 2\right)^{6} $\par
\begin{enumerate}[leftmargin=2em]
\item Identificamos las bases: \(A=2 x - 2 y\), \(B=2 x + 2\) y el exponente \(n=6\).
\item Usamos la formula correspondiente: \[A^n-B^n=(A-B)\sum_{k=0}^{n-1}A^{n-1-k}B^k.\]
\item La letra \(k\) indica que se reemplaza por \(0,1,2,\ldots,n-1\). Eso genera todos los terminos del parentesis largo.
\item Sustituyendo los valores de \(A\), \(B\) y \(n\), queda: \[\boxed{\left(\left(2 x - 2 y\right)-\left(2 x + 2\right)\right)\sum_{k=0}^{5}\left(2 x - 2 y\right)^{5-k}\left(2 x + 2\right)^k}\]
\end{enumerate}

\vspace{0.2cm}
\subsubsection*{Ejercicio 362 - Suma de potencias}
\textbf{Enunciado:} $ \left(3 x - 3 y\right)^{5}+\left(3 x + 1\right)^{5} $\par
\begin{enumerate}[leftmargin=2em]
\item Identificamos las bases: \(A=3 x - 3 y\), \(B=3 x + 1\) y el exponente \(n=5\).
\item Usamos la formula correspondiente: \[A^n+B^n=(A+B)\sum_{k=0}^{n-1}(-1)^kA^{n-1-k}B^k\quad\text{cuando }n\text{ es impar}.\]
\item La letra \(k\) indica que se reemplaza por \(0,1,2,\ldots,n-1\). Eso genera todos los terminos del parentesis largo.
\item Sustituyendo los valores de \(A\), \(B\) y \(n\), queda: \[\boxed{\left(\left(3 x - 3 y\right)+\left(3 x + 1\right)\right)\sum_{k=0}^{4}(-1)^k\left(3 x - 3 y\right)^{4-k}\left(3 x + 1\right)^k}\]
\end{enumerate}

\vspace{0.2cm}
\subsubsection*{Ejercicio 363 - Diferencia de potencias}
\textbf{Enunciado:} $ \left(x - 4 y\right)^{6}-\left(4 x + 2\right)^{6} $\par
\begin{enumerate}[leftmargin=2em]
\item Identificamos las bases: \(A=x - 4 y\), \(B=4 x + 2\) y el exponente \(n=6\).
\item Usamos la formula correspondiente: \[A^n-B^n=(A-B)\sum_{k=0}^{n-1}A^{n-1-k}B^k.\]
\item La letra \(k\) indica que se reemplaza por \(0,1,2,\ldots,n-1\). Eso genera todos los terminos del parentesis largo.
\item Sustituyendo los valores de \(A\), \(B\) y \(n\), queda: \[\boxed{\left(\left(x - 4 y\right)-\left(4 x + 2\right)\right)\sum_{k=0}^{5}\left(x - 4 y\right)^{5-k}\left(4 x + 2\right)^k}\]
\end{enumerate}

\vspace{0.2cm}
\subsubsection*{Ejercicio 364 - Suma de potencias}
\textbf{Enunciado:} $ \left(2 x - y\right)^{5}+\left(5 x + 1\right)^{5} $\par
\begin{enumerate}[leftmargin=2em]
\item Identificamos las bases: \(A=2 x - y\), \(B=5 x + 1\) y el exponente \(n=5\).
\item Usamos la formula correspondiente: \[A^n+B^n=(A+B)\sum_{k=0}^{n-1}(-1)^kA^{n-1-k}B^k\quad\text{cuando }n\text{ es impar}.\]
\item La letra \(k\) indica que se reemplaza por \(0,1,2,\ldots,n-1\). Eso genera todos los terminos del parentesis largo.
\item Sustituyendo los valores de \(A\), \(B\) y \(n\), queda: \[\boxed{\left(\left(2 x - y\right)+\left(5 x + 1\right)\right)\sum_{k=0}^{4}(-1)^k\left(2 x - y\right)^{4-k}\left(5 x + 1\right)^k}\]
\end{enumerate}

\vspace{0.2cm}
\subsubsection*{Ejercicio 365 - Diferencia de potencias}
\textbf{Enunciado:} $ \left(3 x - 2 y\right)^{8}-\left(x + 2\right)^{8} $\par
\begin{enumerate}[leftmargin=2em]
\item Identificamos las bases: \(A=3 x - 2 y\), \(B=x + 2\) y el exponente \(n=8\).
\item Usamos la formula correspondiente: \[A^n-B^n=(A-B)\sum_{k=0}^{n-1}A^{n-1-k}B^k.\]
\item La letra \(k\) indica que se reemplaza por \(0,1,2,\ldots,n-1\). Eso genera todos los terminos del parentesis largo.
\item Sustituyendo los valores de \(A\), \(B\) y \(n\), queda: \[\boxed{\left(\left(3 x - 2 y\right)-\left(x + 2\right)\right)\sum_{k=0}^{7}\left(3 x - 2 y\right)^{7-k}\left(x + 2\right)^k}\]
\end{enumerate}

\vspace{0.2cm}
\subsubsection*{Ejercicio 366 - Suma de potencias}
\textbf{Enunciado:} $ \left(x - 3 y\right)^{5}+\left(2 x + 1\right)^{5} $\par
\begin{enumerate}[leftmargin=2em]
\item Identificamos las bases: \(A=x - 3 y\), \(B=2 x + 1\) y el exponente \(n=5\).
\item Usamos la formula correspondiente: \[A^n+B^n=(A+B)\sum_{k=0}^{n-1}(-1)^kA^{n-1-k}B^k\quad\text{cuando }n\text{ es impar}.\]
\item La letra \(k\) indica que se reemplaza por \(0,1,2,\ldots,n-1\). Eso genera todos los terminos del parentesis largo.
\item Sustituyendo los valores de \(A\), \(B\) y \(n\), queda: \[\boxed{\left(\left(x - 3 y\right)+\left(2 x + 1\right)\right)\sum_{k=0}^{4}(-1)^k\left(x - 3 y\right)^{4-k}\left(2 x + 1\right)^k}\]
\end{enumerate}

\vspace{0.2cm}
\subsubsection*{Ejercicio 367 - Diferencia de potencias}
\textbf{Enunciado:} $ \left(2 x - 4 y\right)^{6}-\left(3 x + 2\right)^{6} $\par
\begin{enumerate}[leftmargin=2em]
\item Identificamos las bases: \(A=2 x - 4 y\), \(B=3 x + 2\) y el exponente \(n=6\).
\item Usamos la formula correspondiente: \[A^n-B^n=(A-B)\sum_{k=0}^{n-1}A^{n-1-k}B^k.\]
\item La letra \(k\) indica que se reemplaza por \(0,1,2,\ldots,n-1\). Eso genera todos los terminos del parentesis largo.
\item Sustituyendo los valores de \(A\), \(B\) y \(n\), queda: \[\boxed{\left(\left(2 x - 4 y\right)-\left(3 x + 2\right)\right)\sum_{k=0}^{5}\left(2 x - 4 y\right)^{5-k}\left(3 x + 2\right)^k}\]
\end{enumerate}

\vspace{0.2cm}
\subsubsection*{Ejercicio 368 - Suma de potencias}
\textbf{Enunciado:} $ \left(3 x - y\right)^{5}+\left(4 x + 1\right)^{5} $\par
\begin{enumerate}[leftmargin=2em]
\item Identificamos las bases: \(A=3 x - y\), \(B=4 x + 1\) y el exponente \(n=5\).
\item Usamos la formula correspondiente: \[A^n+B^n=(A+B)\sum_{k=0}^{n-1}(-1)^kA^{n-1-k}B^k\quad\text{cuando }n\text{ es impar}.\]
\item La letra \(k\) indica que se reemplaza por \(0,1,2,\ldots,n-1\). Eso genera todos los terminos del parentesis largo.
\item Sustituyendo los valores de \(A\), \(B\) y \(n\), queda: \[\boxed{\left(\left(3 x - y\right)+\left(4 x + 1\right)\right)\sum_{k=0}^{4}(-1)^k\left(3 x - y\right)^{4-k}\left(4 x + 1\right)^k}\]
\end{enumerate}

\vspace{0.2cm}
\subsubsection*{Ejercicio 369 - Diferencia de potencias}
\textbf{Enunciado:} $ \left(x - 2 y\right)^{6}-\left(5 x + 2\right)^{6} $\par
\begin{enumerate}[leftmargin=2em]
\item Identificamos las bases: \(A=x - 2 y\), \(B=5 x + 2\) y el exponente \(n=6\).
\item Usamos la formula correspondiente: \[A^n-B^n=(A-B)\sum_{k=0}^{n-1}A^{n-1-k}B^k.\]
\item La letra \(k\) indica que se reemplaza por \(0,1,2,\ldots,n-1\). Eso genera todos los terminos del parentesis largo.
\item Sustituyendo los valores de \(A\), \(B\) y \(n\), queda: \[\boxed{\left(\left(x - 2 y\right)-\left(5 x + 2\right)\right)\sum_{k=0}^{5}\left(x - 2 y\right)^{5-k}\left(5 x + 2\right)^k}\]
\end{enumerate}

\vspace{0.2cm}
\subsubsection*{Ejercicio 370 - Suma de potencias}
\textbf{Enunciado:} $ \left(2 x - 3 y\right)^{5}+\left(x + 1\right)^{5} $\par
\begin{enumerate}[leftmargin=2em]
\item Identificamos las bases: \(A=2 x - 3 y\), \(B=x + 1\) y el exponente \(n=5\).
\item Usamos la formula correspondiente: \[A^n+B^n=(A+B)\sum_{k=0}^{n-1}(-1)^kA^{n-1-k}B^k\quad\text{cuando }n\text{ es impar}.\]
\item La letra \(k\) indica que se reemplaza por \(0,1,2,\ldots,n-1\). Eso genera todos los terminos del parentesis largo.
\item Sustituyendo los valores de \(A\), \(B\) y \(n\), queda: \[\boxed{\left(\left(2 x - 3 y\right)+\left(x + 1\right)\right)\sum_{k=0}^{4}(-1)^k\left(2 x - 3 y\right)^{4-k}\left(x + 1\right)^k}\]
\end{enumerate}

\vspace{0.2cm}
\subsubsection*{Ejercicio 371 - Diferencia de potencias}
\textbf{Enunciado:} $ \left(3 x - 4 y\right)^{6}-\left(2 x + 2\right)^{6} $\par
\begin{enumerate}[leftmargin=2em]
\item Identificamos las bases: \(A=3 x - 4 y\), \(B=2 x + 2\) y el exponente \(n=6\).
\item Usamos la formula correspondiente: \[A^n-B^n=(A-B)\sum_{k=0}^{n-1}A^{n-1-k}B^k.\]
\item La letra \(k\) indica que se reemplaza por \(0,1,2,\ldots,n-1\). Eso genera todos los terminos del parentesis largo.
\item Sustituyendo los valores de \(A\), \(B\) y \(n\), queda: \[\boxed{\left(\left(3 x - 4 y\right)-\left(2 x + 2\right)\right)\sum_{k=0}^{5}\left(3 x - 4 y\right)^{5-k}\left(2 x + 2\right)^k}\]
\end{enumerate}

\vspace{0.2cm}
\subsubsection*{Ejercicio 372 - Suma de potencias}
\textbf{Enunciado:} $ \left(x - y\right)^{5}+\left(3 x + 1\right)^{5} $\par
\begin{enumerate}[leftmargin=2em]
\item Identificamos las bases: \(A=x - y\), \(B=3 x + 1\) y el exponente \(n=5\).
\item Usamos la formula correspondiente: \[A^n+B^n=(A+B)\sum_{k=0}^{n-1}(-1)^kA^{n-1-k}B^k\quad\text{cuando }n\text{ es impar}.\]
\item La letra \(k\) indica que se reemplaza por \(0,1,2,\ldots,n-1\). Eso genera todos los terminos del parentesis largo.
\item Sustituyendo los valores de \(A\), \(B\) y \(n\), queda: \[\boxed{\left(\left(x - y\right)+\left(3 x + 1\right)\right)\sum_{k=0}^{4}(-1)^k\left(x - y\right)^{4-k}\left(3 x + 1\right)^k}\]
\end{enumerate}

\vspace{0.2cm}
\subsubsection*{Ejercicio 373 - Diferencia de potencias}
\textbf{Enunciado:} $ \left(2 x - 2 y\right)^{6}-\left(4 x + 2\right)^{6} $\par
\begin{enumerate}[leftmargin=2em]
\item Identificamos las bases: \(A=2 x - 2 y\), \(B=4 x + 2\) y el exponente \(n=6\).
\item Usamos la formula correspondiente: \[A^n-B^n=(A-B)\sum_{k=0}^{n-1}A^{n-1-k}B^k.\]
\item La letra \(k\) indica que se reemplaza por \(0,1,2,\ldots,n-1\). Eso genera todos los terminos del parentesis largo.
\item Sustituyendo los valores de \(A\), \(B\) y \(n\), queda: \[\boxed{\left(\left(2 x - 2 y\right)-\left(4 x + 2\right)\right)\sum_{k=0}^{5}\left(2 x - 2 y\right)^{5-k}\left(4 x + 2\right)^k}\]
\end{enumerate}

\vspace{0.2cm}
\subsubsection*{Ejercicio 374 - Suma de potencias}
\textbf{Enunciado:} $ \left(3 x - 3 y\right)^{5}+\left(5 x + 1\right)^{5} $\par
\begin{enumerate}[leftmargin=2em]
\item Identificamos las bases: \(A=3 x - 3 y\), \(B=5 x + 1\) y el exponente \(n=5\).
\item Usamos la formula correspondiente: \[A^n+B^n=(A+B)\sum_{k=0}^{n-1}(-1)^kA^{n-1-k}B^k\quad\text{cuando }n\text{ es impar}.\]
\item La letra \(k\) indica que se reemplaza por \(0,1,2,\ldots,n-1\). Eso genera todos los terminos del parentesis largo.
\item Sustituyendo los valores de \(A\), \(B\) y \(n\), queda: \[\boxed{\left(\left(3 x - 3 y\right)+\left(5 x + 1\right)\right)\sum_{k=0}^{4}(-1)^k\left(3 x - 3 y\right)^{4-k}\left(5 x + 1\right)^k}\]
\end{enumerate}

\vspace{0.2cm}
\subsubsection*{Ejercicio 375 - Diferencia de potencias}
\textbf{Enunciado:} $ \left(x - 4 y\right)^{8}-\left(x + 2\right)^{8} $\par
\begin{enumerate}[leftmargin=2em]
\item Identificamos las bases: \(A=x - 4 y\), \(B=x + 2\) y el exponente \(n=8\).
\item Usamos la formula correspondiente: \[A^n-B^n=(A-B)\sum_{k=0}^{n-1}A^{n-1-k}B^k.\]
\item La letra \(k\) indica que se reemplaza por \(0,1,2,\ldots,n-1\). Eso genera todos los terminos del parentesis largo.
\item Sustituyendo los valores de \(A\), \(B\) y \(n\), queda: \[\boxed{\left(\left(x - 4 y\right)-\left(x + 2\right)\right)\sum_{k=0}^{7}\left(x - 4 y\right)^{7-k}\left(x + 2\right)^k}\]
\end{enumerate}

\vspace{0.2cm}
\subsubsection*{Ejercicio 376 - Suma de potencias}
\textbf{Enunciado:} $ \left(2 x - y\right)^{5}+\left(2 x + 1\right)^{5} $\par
\begin{enumerate}[leftmargin=2em]
\item Identificamos las bases: \(A=2 x - y\), \(B=2 x + 1\) y el exponente \(n=5\).
\item Usamos la formula correspondiente: \[A^n+B^n=(A+B)\sum_{k=0}^{n-1}(-1)^kA^{n-1-k}B^k\quad\text{cuando }n\text{ es impar}.\]
\item La letra \(k\) indica que se reemplaza por \(0,1,2,\ldots,n-1\). Eso genera todos los terminos del parentesis largo.
\item Sustituyendo los valores de \(A\), \(B\) y \(n\), queda: \[\boxed{\left(\left(2 x - y\right)+\left(2 x + 1\right)\right)\sum_{k=0}^{4}(-1)^k\left(2 x - y\right)^{4-k}\left(2 x + 1\right)^k}\]
\end{enumerate}

\vspace{0.2cm}
\subsubsection*{Ejercicio 377 - Diferencia de potencias}
\textbf{Enunciado:} $ \left(3 x - 2 y\right)^{6}-\left(3 x + 2\right)^{6} $\par
\begin{enumerate}[leftmargin=2em]
\item Identificamos las bases: \(A=3 x - 2 y\), \(B=3 x + 2\) y el exponente \(n=6\).
\item Usamos la formula correspondiente: \[A^n-B^n=(A-B)\sum_{k=0}^{n-1}A^{n-1-k}B^k.\]
\item La letra \(k\) indica que se reemplaza por \(0,1,2,\ldots,n-1\). Eso genera todos los terminos del parentesis largo.
\item Sustituyendo los valores de \(A\), \(B\) y \(n\), queda: \[\boxed{\left(\left(3 x - 2 y\right)-\left(3 x + 2\right)\right)\sum_{k=0}^{5}\left(3 x - 2 y\right)^{5-k}\left(3 x + 2\right)^k}\]
\end{enumerate}

\vspace{0.2cm}
\subsubsection*{Ejercicio 378 - Suma de potencias}
\textbf{Enunciado:} $ \left(x - 3 y\right)^{5}+\left(4 x + 1\right)^{5} $\par
\begin{enumerate}[leftmargin=2em]
\item Identificamos las bases: \(A=x - 3 y\), \(B=4 x + 1\) y el exponente \(n=5\).
\item Usamos la formula correspondiente: \[A^n+B^n=(A+B)\sum_{k=0}^{n-1}(-1)^kA^{n-1-k}B^k\quad\text{cuando }n\text{ es impar}.\]
\item La letra \(k\) indica que se reemplaza por \(0,1,2,\ldots,n-1\). Eso genera todos los terminos del parentesis largo.
\item Sustituyendo los valores de \(A\), \(B\) y \(n\), queda: \[\boxed{\left(\left(x - 3 y\right)+\left(4 x + 1\right)\right)\sum_{k=0}^{4}(-1)^k\left(x - 3 y\right)^{4-k}\left(4 x + 1\right)^k}\]
\end{enumerate}

\vspace{0.2cm}
\subsubsection*{Ejercicio 379 - Diferencia de potencias}
\textbf{Enunciado:} $ \left(2 x - 4 y\right)^{6}-\left(5 x + 2\right)^{6} $\par
\begin{enumerate}[leftmargin=2em]
\item Identificamos las bases: \(A=2 x - 4 y\), \(B=5 x + 2\) y el exponente \(n=6\).
\item Usamos la formula correspondiente: \[A^n-B^n=(A-B)\sum_{k=0}^{n-1}A^{n-1-k}B^k.\]
\item La letra \(k\) indica que se reemplaza por \(0,1,2,\ldots,n-1\). Eso genera todos los terminos del parentesis largo.
\item Sustituyendo los valores de \(A\), \(B\) y \(n\), queda: \[\boxed{\left(\left(2 x - 4 y\right)-\left(5 x + 2\right)\right)\sum_{k=0}^{5}\left(2 x - 4 y\right)^{5-k}\left(5 x + 2\right)^k}\]
\end{enumerate}

\vspace{0.2cm}
\subsubsection*{Ejercicio 380 - Suma de potencias}
\textbf{Enunciado:} $ \left(3 x - y\right)^{5}+\left(x + 1\right)^{5} $\par
\begin{enumerate}[leftmargin=2em]
\item Identificamos las bases: \(A=3 x - y\), \(B=x + 1\) y el exponente \(n=5\).
\item Usamos la formula correspondiente: \[A^n+B^n=(A+B)\sum_{k=0}^{n-1}(-1)^kA^{n-1-k}B^k\quad\text{cuando }n\text{ es impar}.\]
\item La letra \(k\) indica que se reemplaza por \(0,1,2,\ldots,n-1\). Eso genera todos los terminos del parentesis largo.
\item Sustituyendo los valores de \(A\), \(B\) y \(n\), queda: \[\boxed{\left(\left(3 x - y\right)+\left(x + 1\right)\right)\sum_{k=0}^{4}(-1)^k\left(3 x - y\right)^{4-k}\left(x + 1\right)^k}\]
\end{enumerate}

\vspace{0.2cm}
\nivel{5}{imposible}
\subsubsection*{Ejercicio 381 - Diferencia de potencias}
\textbf{Enunciado:} $ \left(2 x^{2} - 2 y\right)^{8}-\left(3 x y + 2\right)^{8} $\par
\begin{enumerate}[leftmargin=2em]
\item Identificamos las bases: \(A=2 x^{2} - 2 y\), \(B=3 x y + 2\) y el exponente \(n=8\).
\item Usamos la formula correspondiente: \[A^n-B^n=(A-B)\sum_{k=0}^{n-1}A^{n-1-k}B^k.\]
\item La letra \(k\) indica que se reemplaza por \(0,1,2,\ldots,n-1\). Eso genera todos los terminos del parentesis largo.
\item Sustituyendo los valores de \(A\), \(B\) y \(n\), queda: \[\boxed{\left(\left(2 x^{2} - 2 y\right)-\left(3 x y + 2\right)\right)\sum_{k=0}^{7}\left(2 x^{2} - 2 y\right)^{7-k}\left(3 x y + 2\right)^k}\]
\end{enumerate}

\vspace{0.2cm}
\subsubsection*{Ejercicio 382 - Suma de potencias}
\textbf{Enunciado:} $ \left(3 x^{2} - 3 y\right)^{7}+\left(4 x y + 3\right)^{7} $\par
\begin{enumerate}[leftmargin=2em]
\item Identificamos las bases: \(A=3 x^{2} - 3 y\), \(B=4 x y + 3\) y el exponente \(n=7\).
\item Usamos la formula correspondiente: \[A^n+B^n=(A+B)\sum_{k=0}^{n-1}(-1)^kA^{n-1-k}B^k\quad\text{cuando }n\text{ es impar}.\]
\item La letra \(k\) indica que se reemplaza por \(0,1,2,\ldots,n-1\). Eso genera todos los terminos del parentesis largo.
\item Sustituyendo los valores de \(A\), \(B\) y \(n\), queda: \[\boxed{\left(\left(3 x^{2} - 3 y\right)+\left(4 x y + 3\right)\right)\sum_{k=0}^{6}(-1)^k\left(3 x^{2} - 3 y\right)^{6-k}\left(4 x y + 3\right)^k}\]
\end{enumerate}

\vspace{0.2cm}
\subsubsection*{Ejercicio 383 - Diferencia de potencias}
\textbf{Enunciado:} $ \left(x^{2} - 4 y\right)^{9}-\left(5 x y + 4\right)^{9} $\par
\begin{enumerate}[leftmargin=2em]
\item Identificamos las bases: \(A=x^{2} - 4 y\), \(B=5 x y + 4\) y el exponente \(n=9\).
\item Usamos la formula correspondiente: \[A^n-B^n=(A-B)\sum_{k=0}^{n-1}A^{n-1-k}B^k.\]
\item La letra \(k\) indica que se reemplaza por \(0,1,2,\ldots,n-1\). Eso genera todos los terminos del parentesis largo.
\item Sustituyendo los valores de \(A\), \(B\) y \(n\), queda: \[\boxed{\left(\left(x^{2} - 4 y\right)-\left(5 x y + 4\right)\right)\sum_{k=0}^{8}\left(x^{2} - 4 y\right)^{8-k}\left(5 x y + 4\right)^k}\]
\end{enumerate}

\vspace{0.2cm}
\subsubsection*{Ejercicio 384 - Suma de potencias}
\textbf{Enunciado:} $ \left(2 x^{2} - y\right)^{7}+\left(2 x y + 5\right)^{7} $\par
\begin{enumerate}[leftmargin=2em]
\item Identificamos las bases: \(A=2 x^{2} - y\), \(B=2 x y + 5\) y el exponente \(n=7\).
\item Usamos la formula correspondiente: \[A^n+B^n=(A+B)\sum_{k=0}^{n-1}(-1)^kA^{n-1-k}B^k\quad\text{cuando }n\text{ es impar}.\]
\item La letra \(k\) indica que se reemplaza por \(0,1,2,\ldots,n-1\). Eso genera todos los terminos del parentesis largo.
\item Sustituyendo los valores de \(A\), \(B\) y \(n\), queda: \[\boxed{\left(\left(2 x^{2} - y\right)+\left(2 x y + 5\right)\right)\sum_{k=0}^{6}(-1)^k\left(2 x^{2} - y\right)^{6-k}\left(2 x y + 5\right)^k}\]
\end{enumerate}

\vspace{0.2cm}
\subsubsection*{Ejercicio 385 - Diferencia de potencias}
\textbf{Enunciado:} $ \left(3 x^{2} - 2 y\right)^{8}-\left(3 x y + 1\right)^{8} $\par
\begin{enumerate}[leftmargin=2em]
\item Identificamos las bases: \(A=3 x^{2} - 2 y\), \(B=3 x y + 1\) y el exponente \(n=8\).
\item Usamos la formula correspondiente: \[A^n-B^n=(A-B)\sum_{k=0}^{n-1}A^{n-1-k}B^k.\]
\item La letra \(k\) indica que se reemplaza por \(0,1,2,\ldots,n-1\). Eso genera todos los terminos del parentesis largo.
\item Sustituyendo los valores de \(A\), \(B\) y \(n\), queda: \[\boxed{\left(\left(3 x^{2} - 2 y\right)-\left(3 x y + 1\right)\right)\sum_{k=0}^{7}\left(3 x^{2} - 2 y\right)^{7-k}\left(3 x y + 1\right)^k}\]
\end{enumerate}

\vspace{0.2cm}
\subsubsection*{Ejercicio 386 - Suma de potencias}
\textbf{Enunciado:} $ \left(x^{2} - 3 y\right)^{7}+\left(4 x y + 2\right)^{7} $\par
\begin{enumerate}[leftmargin=2em]
\item Identificamos las bases: \(A=x^{2} - 3 y\), \(B=4 x y + 2\) y el exponente \(n=7\).
\item Usamos la formula correspondiente: \[A^n+B^n=(A+B)\sum_{k=0}^{n-1}(-1)^kA^{n-1-k}B^k\quad\text{cuando }n\text{ es impar}.\]
\item La letra \(k\) indica que se reemplaza por \(0,1,2,\ldots,n-1\). Eso genera todos los terminos del parentesis largo.
\item Sustituyendo los valores de \(A\), \(B\) y \(n\), queda: \[\boxed{\left(\left(x^{2} - 3 y\right)+\left(4 x y + 2\right)\right)\sum_{k=0}^{6}(-1)^k\left(x^{2} - 3 y\right)^{6-k}\left(4 x y + 2\right)^k}\]
\end{enumerate}

\vspace{0.2cm}
\subsubsection*{Ejercicio 387 - Diferencia de potencias}
\textbf{Enunciado:} $ \left(2 x^{2} - 4 y\right)^{8}-\left(5 x y + 3\right)^{8} $\par
\begin{enumerate}[leftmargin=2em]
\item Identificamos las bases: \(A=2 x^{2} - 4 y\), \(B=5 x y + 3\) y el exponente \(n=8\).
\item Usamos la formula correspondiente: \[A^n-B^n=(A-B)\sum_{k=0}^{n-1}A^{n-1-k}B^k.\]
\item La letra \(k\) indica que se reemplaza por \(0,1,2,\ldots,n-1\). Eso genera todos los terminos del parentesis largo.
\item Sustituyendo los valores de \(A\), \(B\) y \(n\), queda: \[\boxed{\left(\left(2 x^{2} - 4 y\right)-\left(5 x y + 3\right)\right)\sum_{k=0}^{7}\left(2 x^{2} - 4 y\right)^{7-k}\left(5 x y + 3\right)^k}\]
\end{enumerate}

\vspace{0.2cm}
\subsubsection*{Ejercicio 388 - Suma de potencias}
\textbf{Enunciado:} $ \left(3 x^{2} - y\right)^{7}+\left(2 x y + 4\right)^{7} $\par
\begin{enumerate}[leftmargin=2em]
\item Identificamos las bases: \(A=3 x^{2} - y\), \(B=2 x y + 4\) y el exponente \(n=7\).
\item Usamos la formula correspondiente: \[A^n+B^n=(A+B)\sum_{k=0}^{n-1}(-1)^kA^{n-1-k}B^k\quad\text{cuando }n\text{ es impar}.\]
\item La letra \(k\) indica que se reemplaza por \(0,1,2,\ldots,n-1\). Eso genera todos los terminos del parentesis largo.
\item Sustituyendo los valores de \(A\), \(B\) y \(n\), queda: \[\boxed{\left(\left(3 x^{2} - y\right)+\left(2 x y + 4\right)\right)\sum_{k=0}^{6}(-1)^k\left(3 x^{2} - y\right)^{6-k}\left(2 x y + 4\right)^k}\]
\end{enumerate}

\vspace{0.2cm}
\subsubsection*{Ejercicio 389 - Diferencia de potencias}
\textbf{Enunciado:} $ \left(x^{2} - 2 y\right)^{9}-\left(3 x y + 5\right)^{9} $\par
\begin{enumerate}[leftmargin=2em]
\item Identificamos las bases: \(A=x^{2} - 2 y\), \(B=3 x y + 5\) y el exponente \(n=9\).
\item Usamos la formula correspondiente: \[A^n-B^n=(A-B)\sum_{k=0}^{n-1}A^{n-1-k}B^k.\]
\item La letra \(k\) indica que se reemplaza por \(0,1,2,\ldots,n-1\). Eso genera todos los terminos del parentesis largo.
\item Sustituyendo los valores de \(A\), \(B\) y \(n\), queda: \[\boxed{\left(\left(x^{2} - 2 y\right)-\left(3 x y + 5\right)\right)\sum_{k=0}^{8}\left(x^{2} - 2 y\right)^{8-k}\left(3 x y + 5\right)^k}\]
\end{enumerate}

\vspace{0.2cm}
\subsubsection*{Ejercicio 390 - Suma de potencias}
\textbf{Enunciado:} $ \left(2 x^{2} - 3 y\right)^{7}+\left(4 x y + 1\right)^{7} $\par
\begin{enumerate}[leftmargin=2em]
\item Identificamos las bases: \(A=2 x^{2} - 3 y\), \(B=4 x y + 1\) y el exponente \(n=7\).
\item Usamos la formula correspondiente: \[A^n+B^n=(A+B)\sum_{k=0}^{n-1}(-1)^kA^{n-1-k}B^k\quad\text{cuando }n\text{ es impar}.\]
\item La letra \(k\) indica que se reemplaza por \(0,1,2,\ldots,n-1\). Eso genera todos los terminos del parentesis largo.
\item Sustituyendo los valores de \(A\), \(B\) y \(n\), queda: \[\boxed{\left(\left(2 x^{2} - 3 y\right)+\left(4 x y + 1\right)\right)\sum_{k=0}^{6}(-1)^k\left(2 x^{2} - 3 y\right)^{6-k}\left(4 x y + 1\right)^k}\]
\end{enumerate}

\vspace{0.2cm}
\subsubsection*{Ejercicio 391 - Diferencia de potencias}
\textbf{Enunciado:} $ \left(3 x^{2} - 4 y\right)^{8}-\left(5 x y + 2\right)^{8} $\par
\begin{enumerate}[leftmargin=2em]
\item Identificamos las bases: \(A=3 x^{2} - 4 y\), \(B=5 x y + 2\) y el exponente \(n=8\).
\item Usamos la formula correspondiente: \[A^n-B^n=(A-B)\sum_{k=0}^{n-1}A^{n-1-k}B^k.\]
\item La letra \(k\) indica que se reemplaza por \(0,1,2,\ldots,n-1\). Eso genera todos los terminos del parentesis largo.
\item Sustituyendo los valores de \(A\), \(B\) y \(n\), queda: \[\boxed{\left(\left(3 x^{2} - 4 y\right)-\left(5 x y + 2\right)\right)\sum_{k=0}^{7}\left(3 x^{2} - 4 y\right)^{7-k}\left(5 x y + 2\right)^k}\]
\end{enumerate}

\vspace{0.2cm}
\subsubsection*{Ejercicio 392 - Suma de potencias}
\textbf{Enunciado:} $ \left(x^{2} - y\right)^{7}+\left(2 x y + 3\right)^{7} $\par
\begin{enumerate}[leftmargin=2em]
\item Identificamos las bases: \(A=x^{2} - y\), \(B=2 x y + 3\) y el exponente \(n=7\).
\item Usamos la formula correspondiente: \[A^n+B^n=(A+B)\sum_{k=0}^{n-1}(-1)^kA^{n-1-k}B^k\quad\text{cuando }n\text{ es impar}.\]
\item La letra \(k\) indica que se reemplaza por \(0,1,2,\ldots,n-1\). Eso genera todos los terminos del parentesis largo.
\item Sustituyendo los valores de \(A\), \(B\) y \(n\), queda: \[\boxed{\left(\left(x^{2} - y\right)+\left(2 x y + 3\right)\right)\sum_{k=0}^{6}(-1)^k\left(x^{2} - y\right)^{6-k}\left(2 x y + 3\right)^k}\]
\end{enumerate}

\vspace{0.2cm}
\subsubsection*{Ejercicio 393 - Diferencia de potencias}
\textbf{Enunciado:} $ \left(2 x^{2} - 2 y\right)^{8}-\left(3 x y + 4\right)^{8} $\par
\begin{enumerate}[leftmargin=2em]
\item Identificamos las bases: \(A=2 x^{2} - 2 y\), \(B=3 x y + 4\) y el exponente \(n=8\).
\item Usamos la formula correspondiente: \[A^n-B^n=(A-B)\sum_{k=0}^{n-1}A^{n-1-k}B^k.\]
\item La letra \(k\) indica que se reemplaza por \(0,1,2,\ldots,n-1\). Eso genera todos los terminos del parentesis largo.
\item Sustituyendo los valores de \(A\), \(B\) y \(n\), queda: \[\boxed{\left(\left(2 x^{2} - 2 y\right)-\left(3 x y + 4\right)\right)\sum_{k=0}^{7}\left(2 x^{2} - 2 y\right)^{7-k}\left(3 x y + 4\right)^k}\]
\end{enumerate}

\vspace{0.2cm}
\subsubsection*{Ejercicio 394 - Suma de potencias}
\textbf{Enunciado:} $ \left(3 x^{2} - 3 y\right)^{7}+\left(4 x y + 5\right)^{7} $\par
\begin{enumerate}[leftmargin=2em]
\item Identificamos las bases: \(A=3 x^{2} - 3 y\), \(B=4 x y + 5\) y el exponente \(n=7\).
\item Usamos la formula correspondiente: \[A^n+B^n=(A+B)\sum_{k=0}^{n-1}(-1)^kA^{n-1-k}B^k\quad\text{cuando }n\text{ es impar}.\]
\item La letra \(k\) indica que se reemplaza por \(0,1,2,\ldots,n-1\). Eso genera todos los terminos del parentesis largo.
\item Sustituyendo los valores de \(A\), \(B\) y \(n\), queda: \[\boxed{\left(\left(3 x^{2} - 3 y\right)+\left(4 x y + 5\right)\right)\sum_{k=0}^{6}(-1)^k\left(3 x^{2} - 3 y\right)^{6-k}\left(4 x y + 5\right)^k}\]
\end{enumerate}

\vspace{0.2cm}
\subsubsection*{Ejercicio 395 - Diferencia de potencias}
\textbf{Enunciado:} $ \left(x^{2} - 4 y\right)^{9}-\left(5 x y + 1\right)^{9} $\par
\begin{enumerate}[leftmargin=2em]
\item Identificamos las bases: \(A=x^{2} - 4 y\), \(B=5 x y + 1\) y el exponente \(n=9\).
\item Usamos la formula correspondiente: \[A^n-B^n=(A-B)\sum_{k=0}^{n-1}A^{n-1-k}B^k.\]
\item La letra \(k\) indica que se reemplaza por \(0,1,2,\ldots,n-1\). Eso genera todos los terminos del parentesis largo.
\item Sustituyendo los valores de \(A\), \(B\) y \(n\), queda: \[\boxed{\left(\left(x^{2} - 4 y\right)-\left(5 x y + 1\right)\right)\sum_{k=0}^{8}\left(x^{2} - 4 y\right)^{8-k}\left(5 x y + 1\right)^k}\]
\end{enumerate}

\vspace{0.2cm}
\subsubsection*{Ejercicio 396 - Suma de potencias}
\textbf{Enunciado:} $ \left(2 x^{2} - y\right)^{7}+\left(2 x y + 2\right)^{7} $\par
\begin{enumerate}[leftmargin=2em]
\item Identificamos las bases: \(A=2 x^{2} - y\), \(B=2 x y + 2\) y el exponente \(n=7\).
\item Usamos la formula correspondiente: \[A^n+B^n=(A+B)\sum_{k=0}^{n-1}(-1)^kA^{n-1-k}B^k\quad\text{cuando }n\text{ es impar}.\]
\item La letra \(k\) indica que se reemplaza por \(0,1,2,\ldots,n-1\). Eso genera todos los terminos del parentesis largo.
\item Sustituyendo los valores de \(A\), \(B\) y \(n\), queda: \[\boxed{\left(\left(2 x^{2} - y\right)+\left(2 x y + 2\right)\right)\sum_{k=0}^{6}(-1)^k\left(2 x^{2} - y\right)^{6-k}\left(2 x y + 2\right)^k}\]
\end{enumerate}

\vspace{0.2cm}
\subsubsection*{Ejercicio 397 - Diferencia de potencias}
\textbf{Enunciado:} $ \left(3 x^{2} - 2 y\right)^{8}-\left(3 x y + 3\right)^{8} $\par
\begin{enumerate}[leftmargin=2em]
\item Identificamos las bases: \(A=3 x^{2} - 2 y\), \(B=3 x y + 3\) y el exponente \(n=8\).
\item Usamos la formula correspondiente: \[A^n-B^n=(A-B)\sum_{k=0}^{n-1}A^{n-1-k}B^k.\]
\item La letra \(k\) indica que se reemplaza por \(0,1,2,\ldots,n-1\). Eso genera todos los terminos del parentesis largo.
\item Sustituyendo los valores de \(A\), \(B\) y \(n\), queda: \[\boxed{\left(\left(3 x^{2} - 2 y\right)-\left(3 x y + 3\right)\right)\sum_{k=0}^{7}\left(3 x^{2} - 2 y\right)^{7-k}\left(3 x y + 3\right)^k}\]
\end{enumerate}

\vspace{0.2cm}
\subsubsection*{Ejercicio 398 - Suma de potencias}
\textbf{Enunciado:} $ \left(x^{2} - 3 y\right)^{7}+\left(4 x y + 4\right)^{7} $\par
\begin{enumerate}[leftmargin=2em]
\item Identificamos las bases: \(A=x^{2} - 3 y\), \(B=4 x y + 4\) y el exponente \(n=7\).
\item Usamos la formula correspondiente: \[A^n+B^n=(A+B)\sum_{k=0}^{n-1}(-1)^kA^{n-1-k}B^k\quad\text{cuando }n\text{ es impar}.\]
\item La letra \(k\) indica que se reemplaza por \(0,1,2,\ldots,n-1\). Eso genera todos los terminos del parentesis largo.
\item Sustituyendo los valores de \(A\), \(B\) y \(n\), queda: \[\boxed{\left(\left(x^{2} - 3 y\right)+\left(4 x y + 4\right)\right)\sum_{k=0}^{6}(-1)^k\left(x^{2} - 3 y\right)^{6-k}\left(4 x y + 4\right)^k}\]
\end{enumerate}

\vspace{0.2cm}
\subsubsection*{Ejercicio 399 - Diferencia de potencias}
\textbf{Enunciado:} $ \left(2 x^{2} - 4 y\right)^{8}-\left(5 x y + 5\right)^{8} $\par
\begin{enumerate}[leftmargin=2em]
\item Identificamos las bases: \(A=2 x^{2} - 4 y\), \(B=5 x y + 5\) y el exponente \(n=8\).
\item Usamos la formula correspondiente: \[A^n-B^n=(A-B)\sum_{k=0}^{n-1}A^{n-1-k}B^k.\]
\item La letra \(k\) indica que se reemplaza por \(0,1,2,\ldots,n-1\). Eso genera todos los terminos del parentesis largo.
\item Sustituyendo los valores de \(A\), \(B\) y \(n\), queda: \[\boxed{\left(\left(2 x^{2} - 4 y\right)-\left(5 x y + 5\right)\right)\sum_{k=0}^{7}\left(2 x^{2} - 4 y\right)^{7-k}\left(5 x y + 5\right)^k}\]
\end{enumerate}

\vspace{0.2cm}
\subsubsection*{Ejercicio 400 - Suma de potencias}
\textbf{Enunciado:} $ \left(3 x^{2} - y\right)^{7}+\left(2 x y + 1\right)^{7} $\par
\begin{enumerate}[leftmargin=2em]
\item Identificamos las bases: \(A=3 x^{2} - y\), \(B=2 x y + 1\) y el exponente \(n=7\).
\item Usamos la formula correspondiente: \[A^n+B^n=(A+B)\sum_{k=0}^{n-1}(-1)^kA^{n-1-k}B^k\quad\text{cuando }n\text{ es impar}.\]
\item La letra \(k\) indica que se reemplaza por \(0,1,2,\ldots,n-1\). Eso genera todos los terminos del parentesis largo.
\item Sustituyendo los valores de \(A\), \(B\) y \(n\), queda: \[\boxed{\left(\left(3 x^{2} - y\right)+\left(2 x y + 1\right)\right)\sum_{k=0}^{6}(-1)^k\left(3 x^{2} - y\right)^{6-k}\left(2 x y + 1\right)^k}\]
\end{enumerate}

\vspace{0.2cm}
\section{Soluciones: Ruffini}
\nivel{1}{basico}
\subsubsection*{Ejercicio 401 - Ruffini}
\textbf{Enunciado:} $ x^{3} + 3 x^{2} - 4 x - 12 $\par
\begin{enumerate}[leftmargin=2em]
\item Ruffini sirve para dividir un polinomio entre factores de la forma \((x-r)\). Si el resto da cero, entonces \((x-r)\) es factor.
\item Probamos \(r=2\). Coeficientes actuales: \((1, 3, -4, -12)\).
\item La fila final de Ruffini es \((1, 5, 6, 0)\). El resto es \(0\), por eso el factor es \((x-2)\). El cociente queda \(x^{2} + 5 x + 6\).
\item Probamos \(r=-2\). Coeficientes actuales: \((1, 5, 6)\).
\item La fila final de Ruffini es \((1, 3, 0)\). El resto es \(0\), por eso el factor es \((x+2)\). El cociente queda \(x + 3\).
\item Probamos \(r=-3\). Coeficientes actuales: \((1, 3)\).
\item La fila final de Ruffini es \((1, 0)\). El resto es \(0\), por eso el factor es \((x+3)\). El cociente queda \(1\).
\item Juntamos todos los factores encontrados. Resultado: \[\boxed{\left(x - 2\right) \left(x + 2\right) \left(x + 3\right)}\]
\end{enumerate}

\vspace{0.2cm}
\subsubsection*{Ejercicio 402 - Ruffini}
\textbf{Enunciado:} $ x^{3} + 4 x^{2} - 9 x - 36 $\par
\begin{enumerate}[leftmargin=2em]
\item Ruffini sirve para dividir un polinomio entre factores de la forma \((x-r)\). Si el resto da cero, entonces \((x-r)\) es factor.
\item Probamos \(r=3\). Coeficientes actuales: \((1, 4, -9, -36)\).
\item La fila final de Ruffini es \((1, 7, 12, 0)\). El resto es \(0\), por eso el factor es \((x-3)\). El cociente queda \(x^{2} + 7 x + 12\).
\item Probamos \(r=-3\). Coeficientes actuales: \((1, 7, 12)\).
\item La fila final de Ruffini es \((1, 4, 0)\). El resto es \(0\), por eso el factor es \((x+3)\). El cociente queda \(x + 4\).
\item Probamos \(r=-4\). Coeficientes actuales: \((1, 4)\).
\item La fila final de Ruffini es \((1, 0)\). El resto es \(0\), por eso el factor es \((x+4)\). El cociente queda \(1\).
\item Juntamos todos los factores encontrados. Resultado: \[\boxed{\left(x - 3\right) \left(x + 3\right) \left(x + 4\right)}\]
\end{enumerate}

\vspace{0.2cm}
\subsubsection*{Ejercicio 403 - Ruffini}
\textbf{Enunciado:} $ x^{3} + 2 x^{2} - 19 x - 20 $\par
\begin{enumerate}[leftmargin=2em]
\item Ruffini sirve para dividir un polinomio entre factores de la forma \((x-r)\). Si el resto da cero, entonces \((x-r)\) es factor.
\item Probamos \(r=4\). Coeficientes actuales: \((1, 2, -19, -20)\).
\item La fila final de Ruffini es \((1, 6, 5, 0)\). El resto es \(0\), por eso el factor es \((x-4)\). El cociente queda \(x^{2} + 6 x + 5\).
\item Probamos \(r=-1\). Coeficientes actuales: \((1, 6, 5)\).
\item La fila final de Ruffini es \((1, 5, 0)\). El resto es \(0\), por eso el factor es \((x+1)\). El cociente queda \(x + 5\).
\item Probamos \(r=-5\). Coeficientes actuales: \((1, 5)\).
\item La fila final de Ruffini es \((1, 0)\). El resto es \(0\), por eso el factor es \((x+5)\). El cociente queda \(1\).
\item Juntamos todos los factores encontrados. Resultado: \[\boxed{\left(x - 4\right) \left(x + 1\right) \left(x + 5\right)}\]
\end{enumerate}

\vspace{0.2cm}
\subsubsection*{Ejercicio 404 - Ruffini}
\textbf{Enunciado:} $ x^{3} - x^{2} - 16 x - 20 $\par
\begin{enumerate}[leftmargin=2em]
\item Ruffini sirve para dividir un polinomio entre factores de la forma \((x-r)\). Si el resto da cero, entonces \((x-r)\) es factor.
\item Probamos \(r=5\). Coeficientes actuales: \((1, -1, -16, -20)\).
\item La fila final de Ruffini es \((1, 4, 4, 0)\). El resto es \(0\), por eso el factor es \((x-5)\). El cociente queda \(x^{2} + 4 x + 4\).
\item Probamos \(r=-2\). Coeficientes actuales: \((1, 4, 4)\).
\item La fila final de Ruffini es \((1, 2, 0)\). El resto es \(0\), por eso el factor es \((x+2)\). El cociente queda \(x + 2\).
\item Probamos \(r=-2\). Coeficientes actuales: \((1, 2)\).
\item La fila final de Ruffini es \((1, 0)\). El resto es \(0\), por eso el factor es \((x+2)\). El cociente queda \(1\).
\item Juntamos todos los factores encontrados. Resultado: \[\boxed{\left(x - 5\right) \left(x + 2\right)^{2}}\]
\end{enumerate}

\vspace{0.2cm}
\subsubsection*{Ejercicio 405 - Ruffini}
\textbf{Enunciado:} $ x^{3} + 5 x^{2} + 3 x - 9 $\par
\begin{enumerate}[leftmargin=2em]
\item Ruffini sirve para dividir un polinomio entre factores de la forma \((x-r)\). Si el resto da cero, entonces \((x-r)\) es factor.
\item Probamos \(r=1\). Coeficientes actuales: \((1, 5, 3, -9)\).
\item La fila final de Ruffini es \((1, 6, 9, 0)\). El resto es \(0\), por eso el factor es \((x-1)\). El cociente queda \(x^{2} + 6 x + 9\).
\item Probamos \(r=-3\). Coeficientes actuales: \((1, 6, 9)\).
\item La fila final de Ruffini es \((1, 3, 0)\). El resto es \(0\), por eso el factor es \((x+3)\). El cociente queda \(x + 3\).
\item Probamos \(r=-3\). Coeficientes actuales: \((1, 3)\).
\item La fila final de Ruffini es \((1, 0)\). El resto es \(0\), por eso el factor es \((x+3)\). El cociente queda \(1\).
\item Juntamos todos los factores encontrados. Resultado: \[\boxed{\left(x - 1\right) \left(x + 3\right)^{2}}\]
\end{enumerate}

\vspace{0.2cm}
\subsubsection*{Ejercicio 406 - Ruffini}
\textbf{Enunciado:} $ x^{3} + 3 x^{2} - 6 x - 8 $\par
\begin{enumerate}[leftmargin=2em]
\item Ruffini sirve para dividir un polinomio entre factores de la forma \((x-r)\). Si el resto da cero, entonces \((x-r)\) es factor.
\item Probamos \(r=2\). Coeficientes actuales: \((1, 3, -6, -8)\).
\item La fila final de Ruffini es \((1, 5, 4, 0)\). El resto es \(0\), por eso el factor es \((x-2)\). El cociente queda \(x^{2} + 5 x + 4\).
\item Probamos \(r=-1\). Coeficientes actuales: \((1, 5, 4)\).
\item La fila final de Ruffini es \((1, 4, 0)\). El resto es \(0\), por eso el factor es \((x+1)\). El cociente queda \(x + 4\).
\item Probamos \(r=-4\). Coeficientes actuales: \((1, 4)\).
\item La fila final de Ruffini es \((1, 0)\). El resto es \(0\), por eso el factor es \((x+4)\). El cociente queda \(1\).
\item Juntamos todos los factores encontrados. Resultado: \[\boxed{\left(x - 2\right) \left(x + 1\right) \left(x + 4\right)}\]
\end{enumerate}

\vspace{0.2cm}
\subsubsection*{Ejercicio 407 - Ruffini}
\textbf{Enunciado:} $ x^{3} + 4 x^{2} - 11 x - 30 $\par
\begin{enumerate}[leftmargin=2em]
\item Ruffini sirve para dividir un polinomio entre factores de la forma \((x-r)\). Si el resto da cero, entonces \((x-r)\) es factor.
\item Probamos \(r=3\). Coeficientes actuales: \((1, 4, -11, -30)\).
\item La fila final de Ruffini es \((1, 7, 10, 0)\). El resto es \(0\), por eso el factor es \((x-3)\). El cociente queda \(x^{2} + 7 x + 10\).
\item Probamos \(r=-2\). Coeficientes actuales: \((1, 7, 10)\).
\item La fila final de Ruffini es \((1, 5, 0)\). El resto es \(0\), por eso el factor es \((x+2)\). El cociente queda \(x + 5\).
\item Probamos \(r=-5\). Coeficientes actuales: \((1, 5)\).
\item La fila final de Ruffini es \((1, 0)\). El resto es \(0\), por eso el factor es \((x+5)\). El cociente queda \(1\).
\item Juntamos todos los factores encontrados. Resultado: \[\boxed{\left(x - 3\right) \left(x + 2\right) \left(x + 5\right)}\]
\end{enumerate}

\vspace{0.2cm}
\subsubsection*{Ejercicio 408 - Ruffini}
\textbf{Enunciado:} $ x^{3} + x^{2} - 14 x - 24 $\par
\begin{enumerate}[leftmargin=2em]
\item Ruffini sirve para dividir un polinomio entre factores de la forma \((x-r)\). Si el resto da cero, entonces \((x-r)\) es factor.
\item Probamos \(r=4\). Coeficientes actuales: \((1, 1, -14, -24)\).
\item La fila final de Ruffini es \((1, 5, 6, 0)\). El resto es \(0\), por eso el factor es \((x-4)\). El cociente queda \(x^{2} + 5 x + 6\).
\item Probamos \(r=-3\). Coeficientes actuales: \((1, 5, 6)\).
\item La fila final de Ruffini es \((1, 2, 0)\). El resto es \(0\), por eso el factor es \((x+3)\). El cociente queda \(x + 2\).
\item Probamos \(r=-2\). Coeficientes actuales: \((1, 2)\).
\item La fila final de Ruffini es \((1, 0)\). El resto es \(0\), por eso el factor es \((x+2)\). El cociente queda \(1\).
\item Juntamos todos los factores encontrados. Resultado: \[\boxed{\left(x - 4\right) \left(x + 2\right) \left(x + 3\right)}\]
\end{enumerate}

\vspace{0.2cm}
\subsubsection*{Ejercicio 409 - Ruffini}
\textbf{Enunciado:} $ x^{3} - x^{2} - 17 x - 15 $\par
\begin{enumerate}[leftmargin=2em]
\item Ruffini sirve para dividir un polinomio entre factores de la forma \((x-r)\). Si el resto da cero, entonces \((x-r)\) es factor.
\item Probamos \(r=5\). Coeficientes actuales: \((1, -1, -17, -15)\).
\item La fila final de Ruffini es \((1, 4, 3, 0)\). El resto es \(0\), por eso el factor es \((x-5)\). El cociente queda \(x^{2} + 4 x + 3\).
\item Probamos \(r=-1\). Coeficientes actuales: \((1, 4, 3)\).
\item La fila final de Ruffini es \((1, 3, 0)\). El resto es \(0\), por eso el factor es \((x+1)\). El cociente queda \(x + 3\).
\item Probamos \(r=-3\). Coeficientes actuales: \((1, 3)\).
\item La fila final de Ruffini es \((1, 0)\). El resto es \(0\), por eso el factor es \((x+3)\). El cociente queda \(1\).
\item Juntamos todos los factores encontrados. Resultado: \[\boxed{\left(x - 5\right) \left(x + 1\right) \left(x + 3\right)}\]
\end{enumerate}

\vspace{0.2cm}
\subsubsection*{Ejercicio 410 - Ruffini}
\textbf{Enunciado:} $ x^{3} + 5 x^{2} + 2 x - 8 $\par
\begin{enumerate}[leftmargin=2em]
\item Ruffini sirve para dividir un polinomio entre factores de la forma \((x-r)\). Si el resto da cero, entonces \((x-r)\) es factor.
\item Probamos \(r=1\). Coeficientes actuales: \((1, 5, 2, -8)\).
\item La fila final de Ruffini es \((1, 6, 8, 0)\). El resto es \(0\), por eso el factor es \((x-1)\). El cociente queda \(x^{2} + 6 x + 8\).
\item Probamos \(r=-2\). Coeficientes actuales: \((1, 6, 8)\).
\item La fila final de Ruffini es \((1, 4, 0)\). El resto es \(0\), por eso el factor es \((x+2)\). El cociente queda \(x + 4\).
\item Probamos \(r=-4\). Coeficientes actuales: \((1, 4)\).
\item La fila final de Ruffini es \((1, 0)\). El resto es \(0\), por eso el factor es \((x+4)\). El cociente queda \(1\).
\item Juntamos todos los factores encontrados. Resultado: \[\boxed{\left(x - 1\right) \left(x + 2\right) \left(x + 4\right)}\]
\end{enumerate}

\vspace{0.2cm}
\subsubsection*{Ejercicio 411 - Ruffini}
\textbf{Enunciado:} $ x^{3} + 6 x^{2} - x - 30 $\par
\begin{enumerate}[leftmargin=2em]
\item Ruffini sirve para dividir un polinomio entre factores de la forma \((x-r)\). Si el resto da cero, entonces \((x-r)\) es factor.
\item Probamos \(r=2\). Coeficientes actuales: \((1, 6, -1, -30)\).
\item La fila final de Ruffini es \((1, 8, 15, 0)\). El resto es \(0\), por eso el factor es \((x-2)\). El cociente queda \(x^{2} + 8 x + 15\).
\item Probamos \(r=-3\). Coeficientes actuales: \((1, 8, 15)\).
\item La fila final de Ruffini es \((1, 5, 0)\). El resto es \(0\), por eso el factor es \((x+3)\). El cociente queda \(x + 5\).
\item Probamos \(r=-5\). Coeficientes actuales: \((1, 5)\).
\item La fila final de Ruffini es \((1, 0)\). El resto es \(0\), por eso el factor es \((x+5)\). El cociente queda \(1\).
\item Juntamos todos los factores encontrados. Resultado: \[\boxed{\left(x - 2\right) \left(x + 3\right) \left(x + 5\right)}\]
\end{enumerate}

\vspace{0.2cm}
\subsubsection*{Ejercicio 412 - Ruffini}
\textbf{Enunciado:} $ x^{3} - 7 x - 6 $\par
\begin{enumerate}[leftmargin=2em]
\item Ruffini sirve para dividir un polinomio entre factores de la forma \((x-r)\). Si el resto da cero, entonces \((x-r)\) es factor.
\item Probamos \(r=3\). Coeficientes actuales: \((1, 0, -7, -6)\).
\item La fila final de Ruffini es \((1, 3, 2, 0)\). El resto es \(0\), por eso el factor es \((x-3)\). El cociente queda \(x^{2} + 3 x + 2\).
\item Probamos \(r=-1\). Coeficientes actuales: \((1, 3, 2)\).
\item La fila final de Ruffini es \((1, 2, 0)\). El resto es \(0\), por eso el factor es \((x+1)\). El cociente queda \(x + 2\).
\item Probamos \(r=-2\). Coeficientes actuales: \((1, 2)\).
\item La fila final de Ruffini es \((1, 0)\). El resto es \(0\), por eso el factor es \((x+2)\). El cociente queda \(1\).
\item Juntamos todos los factores encontrados. Resultado: \[\boxed{\left(x - 3\right) \left(x + 1\right) \left(x + 2\right)}\]
\end{enumerate}

\vspace{0.2cm}
\subsubsection*{Ejercicio 413 - Ruffini}
\textbf{Enunciado:} $ x^{3} + x^{2} - 14 x - 24 $\par
\begin{enumerate}[leftmargin=2em]
\item Ruffini sirve para dividir un polinomio entre factores de la forma \((x-r)\). Si el resto da cero, entonces \((x-r)\) es factor.
\item Probamos \(r=4\). Coeficientes actuales: \((1, 1, -14, -24)\).
\item La fila final de Ruffini es \((1, 5, 6, 0)\). El resto es \(0\), por eso el factor es \((x-4)\). El cociente queda \(x^{2} + 5 x + 6\).
\item Probamos \(r=-2\). Coeficientes actuales: \((1, 5, 6)\).
\item La fila final de Ruffini es \((1, 3, 0)\). El resto es \(0\), por eso el factor es \((x+2)\). El cociente queda \(x + 3\).
\item Probamos \(r=-3\). Coeficientes actuales: \((1, 3)\).
\item La fila final de Ruffini es \((1, 0)\). El resto es \(0\), por eso el factor es \((x+3)\). El cociente queda \(1\).
\item Juntamos todos los factores encontrados. Resultado: \[\boxed{\left(x - 4\right) \left(x + 2\right) \left(x + 3\right)}\]
\end{enumerate}

\vspace{0.2cm}
\subsubsection*{Ejercicio 414 - Ruffini}
\textbf{Enunciado:} $ x^{3} + 2 x^{2} - 23 x - 60 $\par
\begin{enumerate}[leftmargin=2em]
\item Ruffini sirve para dividir un polinomio entre factores de la forma \((x-r)\). Si el resto da cero, entonces \((x-r)\) es factor.
\item Probamos \(r=5\). Coeficientes actuales: \((1, 2, -23, -60)\).
\item La fila final de Ruffini es \((1, 7, 12, 0)\). El resto es \(0\), por eso el factor es \((x-5)\). El cociente queda \(x^{2} + 7 x + 12\).
\item Probamos \(r=-3\). Coeficientes actuales: \((1, 7, 12)\).
\item La fila final de Ruffini es \((1, 4, 0)\). El resto es \(0\), por eso el factor es \((x+3)\). El cociente queda \(x + 4\).
\item Probamos \(r=-4\). Coeficientes actuales: \((1, 4)\).
\item La fila final de Ruffini es \((1, 0)\). El resto es \(0\), por eso el factor es \((x+4)\). El cociente queda \(1\).
\item Juntamos todos los factores encontrados. Resultado: \[\boxed{\left(x - 5\right) \left(x + 3\right) \left(x + 4\right)}\]
\end{enumerate}

\vspace{0.2cm}
\subsubsection*{Ejercicio 415 - Ruffini}
\textbf{Enunciado:} $ x^{3} + 5 x^{2} - x - 5 $\par
\begin{enumerate}[leftmargin=2em]
\item Ruffini sirve para dividir un polinomio entre factores de la forma \((x-r)\). Si el resto da cero, entonces \((x-r)\) es factor.
\item Probamos \(r=1\). Coeficientes actuales: \((1, 5, -1, -5)\).
\item La fila final de Ruffini es \((1, 6, 5, 0)\). El resto es \(0\), por eso el factor es \((x-1)\). El cociente queda \(x^{2} + 6 x + 5\).
\item Probamos \(r=-1\). Coeficientes actuales: \((1, 6, 5)\).
\item La fila final de Ruffini es \((1, 5, 0)\). El resto es \(0\), por eso el factor es \((x+1)\). El cociente queda \(x + 5\).
\item Probamos \(r=-5\). Coeficientes actuales: \((1, 5)\).
\item La fila final de Ruffini es \((1, 0)\). El resto es \(0\), por eso el factor es \((x+5)\). El cociente queda \(1\).
\item Juntamos todos los factores encontrados. Resultado: \[\boxed{\left(x - 1\right) \left(x + 1\right) \left(x + 5\right)}\]
\end{enumerate}

\vspace{0.2cm}
\subsubsection*{Ejercicio 416 - Ruffini}
\textbf{Enunciado:} $ x^{3} + 2 x^{2} - 4 x - 8 $\par
\begin{enumerate}[leftmargin=2em]
\item Ruffini sirve para dividir un polinomio entre factores de la forma \((x-r)\). Si el resto da cero, entonces \((x-r)\) es factor.
\item Probamos \(r=2\). Coeficientes actuales: \((1, 2, -4, -8)\).
\item La fila final de Ruffini es \((1, 4, 4, 0)\). El resto es \(0\), por eso el factor es \((x-2)\). El cociente queda \(x^{2} + 4 x + 4\).
\item Probamos \(r=-2\). Coeficientes actuales: \((1, 4, 4)\).
\item La fila final de Ruffini es \((1, 2, 0)\). El resto es \(0\), por eso el factor es \((x+2)\). El cociente queda \(x + 2\).
\item Probamos \(r=-2\). Coeficientes actuales: \((1, 2)\).
\item La fila final de Ruffini es \((1, 0)\). El resto es \(0\), por eso el factor es \((x+2)\). El cociente queda \(1\).
\item Juntamos todos los factores encontrados. Resultado: \[\boxed{\left(x - 2\right) \left(x + 2\right)^{2}}\]
\end{enumerate}

\vspace{0.2cm}
\subsubsection*{Ejercicio 417 - Ruffini}
\textbf{Enunciado:} $ x^{3} + 3 x^{2} - 9 x - 27 $\par
\begin{enumerate}[leftmargin=2em]
\item Ruffini sirve para dividir un polinomio entre factores de la forma \((x-r)\). Si el resto da cero, entonces \((x-r)\) es factor.
\item Probamos \(r=3\). Coeficientes actuales: \((1, 3, -9, -27)\).
\item La fila final de Ruffini es \((1, 6, 9, 0)\). El resto es \(0\), por eso el factor es \((x-3)\). El cociente queda \(x^{2} + 6 x + 9\).
\item Probamos \(r=-3\). Coeficientes actuales: \((1, 6, 9)\).
\item La fila final de Ruffini es \((1, 3, 0)\). El resto es \(0\), por eso el factor es \((x+3)\). El cociente queda \(x + 3\).
\item Probamos \(r=-3\). Coeficientes actuales: \((1, 3)\).
\item La fila final de Ruffini es \((1, 0)\). El resto es \(0\), por eso el factor es \((x+3)\). El cociente queda \(1\).
\item Juntamos todos los factores encontrados. Resultado: \[\boxed{\left(x - 3\right) \left(x + 3\right)^{2}}\]
\end{enumerate}

\vspace{0.2cm}
\subsubsection*{Ejercicio 418 - Ruffini}
\textbf{Enunciado:} $ x^{3} + x^{2} - 16 x - 16 $\par
\begin{enumerate}[leftmargin=2em]
\item Ruffini sirve para dividir un polinomio entre factores de la forma \((x-r)\). Si el resto da cero, entonces \((x-r)\) es factor.
\item Probamos \(r=4\). Coeficientes actuales: \((1, 1, -16, -16)\).
\item La fila final de Ruffini es \((1, 5, 4, 0)\). El resto es \(0\), por eso el factor es \((x-4)\). El cociente queda \(x^{2} + 5 x + 4\).
\item Probamos \(r=-1\). Coeficientes actuales: \((1, 5, 4)\).
\item La fila final de Ruffini es \((1, 4, 0)\). El resto es \(0\), por eso el factor es \((x+1)\). El cociente queda \(x + 4\).
\item Probamos \(r=-4\). Coeficientes actuales: \((1, 4)\).
\item La fila final de Ruffini es \((1, 0)\). El resto es \(0\), por eso el factor es \((x+4)\). El cociente queda \(1\).
\item Juntamos todos los factores encontrados. Resultado: \[\boxed{\left(x - 4\right) \left(x + 1\right) \left(x + 4\right)}\]
\end{enumerate}

\vspace{0.2cm}
\subsubsection*{Ejercicio 419 - Ruffini}
\textbf{Enunciado:} $ x^{3} + 2 x^{2} - 25 x - 50 $\par
\begin{enumerate}[leftmargin=2em]
\item Ruffini sirve para dividir un polinomio entre factores de la forma \((x-r)\). Si el resto da cero, entonces \((x-r)\) es factor.
\item Probamos \(r=5\). Coeficientes actuales: \((1, 2, -25, -50)\).
\item La fila final de Ruffini es \((1, 7, 10, 0)\). El resto es \(0\), por eso el factor es \((x-5)\). El cociente queda \(x^{2} + 7 x + 10\).
\item Probamos \(r=-2\). Coeficientes actuales: \((1, 7, 10)\).
\item La fila final de Ruffini es \((1, 5, 0)\). El resto es \(0\), por eso el factor es \((x+2)\). El cociente queda \(x + 5\).
\item Probamos \(r=-5\). Coeficientes actuales: \((1, 5)\).
\item La fila final de Ruffini es \((1, 0)\). El resto es \(0\), por eso el factor es \((x+5)\). El cociente queda \(1\).
\item Juntamos todos los factores encontrados. Resultado: \[\boxed{\left(x - 5\right) \left(x + 2\right) \left(x + 5\right)}\]
\end{enumerate}

\vspace{0.2cm}
\subsubsection*{Ejercicio 420 - Ruffini}
\textbf{Enunciado:} $ x^{3} + 4 x^{2} + x - 6 $\par
\begin{enumerate}[leftmargin=2em]
\item Ruffini sirve para dividir un polinomio entre factores de la forma \((x-r)\). Si el resto da cero, entonces \((x-r)\) es factor.
\item Probamos \(r=1\). Coeficientes actuales: \((1, 4, 1, -6)\).
\item La fila final de Ruffini es \((1, 5, 6, 0)\). El resto es \(0\), por eso el factor es \((x-1)\). El cociente queda \(x^{2} + 5 x + 6\).
\item Probamos \(r=-3\). Coeficientes actuales: \((1, 5, 6)\).
\item La fila final de Ruffini es \((1, 2, 0)\). El resto es \(0\), por eso el factor es \((x+3)\). El cociente queda \(x + 2\).
\item Probamos \(r=-2\). Coeficientes actuales: \((1, 2)\).
\item La fila final de Ruffini es \((1, 0)\). El resto es \(0\), por eso el factor es \((x+2)\). El cociente queda \(1\).
\item Juntamos todos los factores encontrados. Resultado: \[\boxed{\left(x - 1\right) \left(x + 2\right) \left(x + 3\right)}\]
\end{enumerate}

\vspace{0.2cm}
\nivel{2}{intermedio inicial}
\subsubsection*{Ejercicio 421 - Ruffini}
\textbf{Enunciado:} $ x^{3} + 3 x^{2} - 4 x - 12 $\par
\begin{enumerate}[leftmargin=2em]
\item Ruffini sirve para dividir un polinomio entre factores de la forma \((x-r)\). Si el resto da cero, entonces \((x-r)\) es factor.
\item Probamos \(r=-2\). Coeficientes actuales: \((1, 3, -4, -12)\).
\item La fila final de Ruffini es \((1, 1, -6, 0)\). El resto es \(0\), por eso el factor es \((x+2)\). El cociente queda \(x^{2} + x - 6\).
\item Probamos \(r=2\). Coeficientes actuales: \((1, 1, -6)\).
\item La fila final de Ruffini es \((1, 3, 0)\). El resto es \(0\), por eso el factor es \((x-2)\). El cociente queda \(x + 3\).
\item Probamos \(r=-3\). Coeficientes actuales: \((1, 3)\).
\item La fila final de Ruffini es \((1, 0)\). El resto es \(0\), por eso el factor es \((x+3)\). El cociente queda \(1\).
\item Juntamos todos los factores encontrados. Resultado: \[\boxed{\left(x - 2\right) \left(x + 2\right) \left(x + 3\right)}\]
\end{enumerate}

\vspace{0.2cm}
\subsubsection*{Ejercicio 422 - Ruffini}
\textbf{Enunciado:} $ x^{3} + 4 x^{2} - 9 x - 36 $\par
\begin{enumerate}[leftmargin=2em]
\item Ruffini sirve para dividir un polinomio entre factores de la forma \((x-r)\). Si el resto da cero, entonces \((x-r)\) es factor.
\item Probamos \(r=-3\). Coeficientes actuales: \((1, 4, -9, -36)\).
\item La fila final de Ruffini es \((1, 1, -12, 0)\). El resto es \(0\), por eso el factor es \((x+3)\). El cociente queda \(x^{2} + x - 12\).
\item Probamos \(r=3\). Coeficientes actuales: \((1, 1, -12)\).
\item La fila final de Ruffini es \((1, 4, 0)\). El resto es \(0\), por eso el factor es \((x-3)\). El cociente queda \(x + 4\).
\item Probamos \(r=-4\). Coeficientes actuales: \((1, 4)\).
\item La fila final de Ruffini es \((1, 0)\). El resto es \(0\), por eso el factor es \((x+4)\). El cociente queda \(1\).
\item Juntamos todos los factores encontrados. Resultado: \[\boxed{\left(x - 3\right) \left(x + 3\right) \left(x + 4\right)}\]
\end{enumerate}

\vspace{0.2cm}
\subsubsection*{Ejercicio 423 - Ruffini}
\textbf{Enunciado:} $ x^{3} + 2 x^{2} - 16 x - 32 $\par
\begin{enumerate}[leftmargin=2em]
\item Ruffini sirve para dividir un polinomio entre factores de la forma \((x-r)\). Si el resto da cero, entonces \((x-r)\) es factor.
\item Probamos \(r=-4\). Coeficientes actuales: \((1, 2, -16, -32)\).
\item La fila final de Ruffini es \((1, -2, -8, 0)\). El resto es \(0\), por eso el factor es \((x+4)\). El cociente queda \(x^{2} - 2 x - 8\).
\item Probamos \(r=4\). Coeficientes actuales: \((1, -2, -8)\).
\item La fila final de Ruffini es \((1, 2, 0)\). El resto es \(0\), por eso el factor es \((x-4)\). El cociente queda \(x + 2\).
\item Probamos \(r=-2\). Coeficientes actuales: \((1, 2)\).
\item La fila final de Ruffini es \((1, 0)\). El resto es \(0\), por eso el factor es \((x+2)\). El cociente queda \(1\).
\item Juntamos todos los factores encontrados. Resultado: \[\boxed{\left(x - 4\right) \left(x + 2\right) \left(x + 4\right)}\]
\end{enumerate}

\vspace{0.2cm}
\subsubsection*{Ejercicio 424 - Ruffini}
\textbf{Enunciado:} $ x^{3} + 7 x^{2} + 7 x - 15 $\par
\begin{enumerate}[leftmargin=2em]
\item Ruffini sirve para dividir un polinomio entre factores de la forma \((x-r)\). Si el resto da cero, entonces \((x-r)\) es factor.
\item Probamos \(r=-5\). Coeficientes actuales: \((1, 7, 7, -15)\).
\item La fila final de Ruffini es \((1, 2, -3, 0)\). El resto es \(0\), por eso el factor es \((x+5)\). El cociente queda \(x^{2} + 2 x - 3\).
\item Probamos \(r=1\). Coeficientes actuales: \((1, 2, -3)\).
\item La fila final de Ruffini es \((1, 3, 0)\). El resto es \(0\), por eso el factor es \((x-1)\). El cociente queda \(x + 3\).
\item Probamos \(r=-3\). Coeficientes actuales: \((1, 3)\).
\item La fila final de Ruffini es \((1, 0)\). El resto es \(0\), por eso el factor es \((x+3)\). El cociente queda \(1\).
\item Juntamos todos los factores encontrados. Resultado: \[\boxed{\left(x - 1\right) \left(x + 3\right) \left(x + 5\right)}\]
\end{enumerate}

\vspace{0.2cm}
\subsubsection*{Ejercicio 425 - Ruffini}
\textbf{Enunciado:} $ x^{3} + 3 x^{2} - 6 x - 8 $\par
\begin{enumerate}[leftmargin=2em]
\item Ruffini sirve para dividir un polinomio entre factores de la forma \((x-r)\). Si el resto da cero, entonces \((x-r)\) es factor.
\item Probamos \(r=-1\). Coeficientes actuales: \((1, 3, -6, -8)\).
\item La fila final de Ruffini es \((1, 2, -8, 0)\). El resto es \(0\), por eso el factor es \((x+1)\). El cociente queda \(x^{2} + 2 x - 8\).
\item Probamos \(r=2\). Coeficientes actuales: \((1, 2, -8)\).
\item La fila final de Ruffini es \((1, 4, 0)\). El resto es \(0\), por eso el factor es \((x-2)\). El cociente queda \(x + 4\).
\item Probamos \(r=-4\). Coeficientes actuales: \((1, 4)\).
\item La fila final de Ruffini es \((1, 0)\). El resto es \(0\), por eso el factor es \((x+4)\). El cociente queda \(1\).
\item Juntamos todos los factores encontrados. Resultado: \[\boxed{\left(x - 2\right) \left(x + 1\right) \left(x + 4\right)}\]
\end{enumerate}

\vspace{0.2cm}
\subsubsection*{Ejercicio 426 - Ruffini}
\textbf{Enunciado:} $ x^{3} + x^{2} - 8 x - 12 $\par
\begin{enumerate}[leftmargin=2em]
\item Ruffini sirve para dividir un polinomio entre factores de la forma \((x-r)\). Si el resto da cero, entonces \((x-r)\) es factor.
\item Probamos \(r=-2\). Coeficientes actuales: \((1, 1, -8, -12)\).
\item La fila final de Ruffini es \((1, -1, -6, 0)\). El resto es \(0\), por eso el factor es \((x+2)\). El cociente queda \(x^{2} - x - 6\).
\item Probamos \(r=3\). Coeficientes actuales: \((1, -1, -6)\).
\item La fila final de Ruffini es \((1, 2, 0)\). El resto es \(0\), por eso el factor es \((x-3)\). El cociente queda \(x + 2\).
\item Probamos \(r=-2\). Coeficientes actuales: \((1, 2)\).
\item La fila final de Ruffini es \((1, 0)\). El resto es \(0\), por eso el factor es \((x+2)\). El cociente queda \(1\).
\item Juntamos todos los factores encontrados. Resultado: \[\boxed{\left(x - 3\right) \left(x + 2\right)^{2}}\]
\end{enumerate}

\vspace{0.2cm}
\subsubsection*{Ejercicio 427 - Ruffini}
\textbf{Enunciado:} $ x^{3} + 2 x^{2} - 15 x - 36 $\par
\begin{enumerate}[leftmargin=2em]
\item Ruffini sirve para dividir un polinomio entre factores de la forma \((x-r)\). Si el resto da cero, entonces \((x-r)\) es factor.
\item Probamos \(r=-3\). Coeficientes actuales: \((1, 2, -15, -36)\).
\item La fila final de Ruffini es \((1, -1, -12, 0)\). El resto es \(0\), por eso el factor es \((x+3)\). El cociente queda \(x^{2} - x - 12\).
\item Probamos \(r=4\). Coeficientes actuales: \((1, -1, -12)\).
\item La fila final de Ruffini es \((1, 3, 0)\). El resto es \(0\), por eso el factor es \((x-4)\). El cociente queda \(x + 3\).
\item Probamos \(r=-3\). Coeficientes actuales: \((1, 3)\).
\item La fila final de Ruffini es \((1, 0)\). El resto es \(0\), por eso el factor es \((x+3)\). El cociente queda \(1\).
\item Juntamos todos los factores encontrados. Resultado: \[\boxed{\left(x - 4\right) \left(x + 3\right)^{2}}\]
\end{enumerate}

\vspace{0.2cm}
\subsubsection*{Ejercicio 428 - Ruffini}
\textbf{Enunciado:} $ x^{3} + 7 x^{2} + 8 x - 16 $\par
\begin{enumerate}[leftmargin=2em]
\item Ruffini sirve para dividir un polinomio entre factores de la forma \((x-r)\). Si el resto da cero, entonces \((x-r)\) es factor.
\item Probamos \(r=-4\). Coeficientes actuales: \((1, 7, 8, -16)\).
\item La fila final de Ruffini es \((1, 3, -4, 0)\). El resto es \(0\), por eso el factor es \((x+4)\). El cociente queda \(x^{2} + 3 x - 4\).
\item Probamos \(r=1\). Coeficientes actuales: \((1, 3, -4)\).
\item La fila final de Ruffini es \((1, 4, 0)\). El resto es \(0\), por eso el factor es \((x-1)\). El cociente queda \(x + 4\).
\item Probamos \(r=-4\). Coeficientes actuales: \((1, 4)\).
\item La fila final de Ruffini es \((1, 0)\). El resto es \(0\), por eso el factor es \((x+4)\). El cociente queda \(1\).
\item Juntamos todos los factores encontrados. Resultado: \[\boxed{\left(x - 1\right) \left(x + 4\right)^{2}}\]
\end{enumerate}

\vspace{0.2cm}
\subsubsection*{Ejercicio 429 - Ruffini}
\textbf{Enunciado:} $ x^{3} + 5 x^{2} - 4 x - 20 $\par
\begin{enumerate}[leftmargin=2em]
\item Ruffini sirve para dividir un polinomio entre factores de la forma \((x-r)\). Si el resto da cero, entonces \((x-r)\) es factor.
\item Probamos \(r=-5\). Coeficientes actuales: \((1, 5, -4, -20)\).
\item La fila final de Ruffini es \((1, 0, -4, 0)\). El resto es \(0\), por eso el factor es \((x+5)\). El cociente queda \(x^{2} - 4\).
\item Probamos \(r=2\). Coeficientes actuales: \((1, 0, -4)\).
\item La fila final de Ruffini es \((1, 2, 0)\). El resto es \(0\), por eso el factor es \((x-2)\). El cociente queda \(x + 2\).
\item Probamos \(r=-2\). Coeficientes actuales: \((1, 2)\).
\item La fila final de Ruffini es \((1, 0)\). El resto es \(0\), por eso el factor es \((x+2)\). El cociente queda \(1\).
\item Juntamos todos los factores encontrados. Resultado: \[\boxed{\left(x - 2\right) \left(x + 2\right) \left(x + 5\right)}\]
\end{enumerate}

\vspace{0.2cm}
\subsubsection*{Ejercicio 430 - Ruffini}
\textbf{Enunciado:} $ x^{3} + x^{2} - 9 x - 9 $\par
\begin{enumerate}[leftmargin=2em]
\item Ruffini sirve para dividir un polinomio entre factores de la forma \((x-r)\). Si el resto da cero, entonces \((x-r)\) es factor.
\item Probamos \(r=-1\). Coeficientes actuales: \((1, 1, -9, -9)\).
\item La fila final de Ruffini es \((1, 0, -9, 0)\). El resto es \(0\), por eso el factor es \((x+1)\). El cociente queda \(x^{2} - 9\).
\item Probamos \(r=3\). Coeficientes actuales: \((1, 0, -9)\).
\item La fila final de Ruffini es \((1, 3, 0)\). El resto es \(0\), por eso el factor es \((x-3)\). El cociente queda \(x + 3\).
\item Probamos \(r=-3\). Coeficientes actuales: \((1, 3)\).
\item La fila final de Ruffini es \((1, 0)\). El resto es \(0\), por eso el factor es \((x+3)\). El cociente queda \(1\).
\item Juntamos todos los factores encontrados. Resultado: \[\boxed{\left(x - 3\right) \left(x + 1\right) \left(x + 3\right)}\]
\end{enumerate}

\vspace{0.2cm}
\subsubsection*{Ejercicio 431 - Ruffini}
\textbf{Enunciado:} $ x^{3} + 2 x^{2} - 16 x - 32 $\par
\begin{enumerate}[leftmargin=2em]
\item Ruffini sirve para dividir un polinomio entre factores de la forma \((x-r)\). Si el resto da cero, entonces \((x-r)\) es factor.
\item Probamos \(r=-2\). Coeficientes actuales: \((1, 2, -16, -32)\).
\item La fila final de Ruffini es \((1, 0, -16, 0)\). El resto es \(0\), por eso el factor es \((x+2)\). El cociente queda \(x^{2} - 16\).
\item Probamos \(r=4\). Coeficientes actuales: \((1, 0, -16)\).
\item La fila final de Ruffini es \((1, 4, 0)\). El resto es \(0\), por eso el factor es \((x-4)\). El cociente queda \(x + 4\).
\item Probamos \(r=-4\). Coeficientes actuales: \((1, 4)\).
\item La fila final de Ruffini es \((1, 0)\). El resto es \(0\), por eso el factor es \((x+4)\). El cociente queda \(1\).
\item Juntamos todos los factores encontrados. Resultado: \[\boxed{\left(x - 4\right) \left(x + 2\right) \left(x + 4\right)}\]
\end{enumerate}

\vspace{0.2cm}
\subsubsection*{Ejercicio 432 - Ruffini}
\textbf{Enunciado:} $ x^{3} + 4 x^{2} + x - 6 $\par
\begin{enumerate}[leftmargin=2em]
\item Ruffini sirve para dividir un polinomio entre factores de la forma \((x-r)\). Si el resto da cero, entonces \((x-r)\) es factor.
\item Probamos \(r=-3\). Coeficientes actuales: \((1, 4, 1, -6)\).
\item La fila final de Ruffini es \((1, 1, -2, 0)\). El resto es \(0\), por eso el factor es \((x+3)\). El cociente queda \(x^{2} + x - 2\).
\item Probamos \(r=1\). Coeficientes actuales: \((1, 1, -2)\).
\item La fila final de Ruffini es \((1, 2, 0)\). El resto es \(0\), por eso el factor es \((x-1)\). El cociente queda \(x + 2\).
\item Probamos \(r=-2\). Coeficientes actuales: \((1, 2)\).
\item La fila final de Ruffini es \((1, 0)\). El resto es \(0\), por eso el factor es \((x+2)\). El cociente queda \(1\).
\item Juntamos todos los factores encontrados. Resultado: \[\boxed{\left(x - 1\right) \left(x + 2\right) \left(x + 3\right)}\]
\end{enumerate}

\vspace{0.2cm}
\subsubsection*{Ejercicio 433 - Ruffini}
\textbf{Enunciado:} $ x^{3} + 5 x^{2} - 2 x - 24 $\par
\begin{enumerate}[leftmargin=2em]
\item Ruffini sirve para dividir un polinomio entre factores de la forma \((x-r)\). Si el resto da cero, entonces \((x-r)\) es factor.
\item Probamos \(r=-4\). Coeficientes actuales: \((1, 5, -2, -24)\).
\item La fila final de Ruffini es \((1, 1, -6, 0)\). El resto es \(0\), por eso el factor es \((x+4)\). El cociente queda \(x^{2} + x - 6\).
\item Probamos \(r=2\). Coeficientes actuales: \((1, 1, -6)\).
\item La fila final de Ruffini es \((1, 3, 0)\). El resto es \(0\), por eso el factor es \((x-2)\). El cociente queda \(x + 3\).
\item Probamos \(r=-3\). Coeficientes actuales: \((1, 3)\).
\item La fila final de Ruffini es \((1, 0)\). El resto es \(0\), por eso el factor es \((x+3)\). El cociente queda \(1\).
\item Juntamos todos los factores encontrados. Resultado: \[\boxed{\left(x - 2\right) \left(x + 3\right) \left(x + 4\right)}\]
\end{enumerate}

\vspace{0.2cm}
\subsubsection*{Ejercicio 434 - Ruffini}
\textbf{Enunciado:} $ x^{3} + 6 x^{2} - 7 x - 60 $\par
\begin{enumerate}[leftmargin=2em]
\item Ruffini sirve para dividir un polinomio entre factores de la forma \((x-r)\). Si el resto da cero, entonces \((x-r)\) es factor.
\item Probamos \(r=-5\). Coeficientes actuales: \((1, 6, -7, -60)\).
\item La fila final de Ruffini es \((1, 1, -12, 0)\). El resto es \(0\), por eso el factor es \((x+5)\). El cociente queda \(x^{2} + x - 12\).
\item Probamos \(r=3\). Coeficientes actuales: \((1, 1, -12)\).
\item La fila final de Ruffini es \((1, 4, 0)\). El resto es \(0\), por eso el factor es \((x-3)\). El cociente queda \(x + 4\).
\item Probamos \(r=-4\). Coeficientes actuales: \((1, 4)\).
\item La fila final de Ruffini es \((1, 0)\). El resto es \(0\), por eso el factor es \((x+4)\). El cociente queda \(1\).
\item Juntamos todos los factores encontrados. Resultado: \[\boxed{\left(x - 3\right) \left(x + 4\right) \left(x + 5\right)}\]
\end{enumerate}

\vspace{0.2cm}
\subsubsection*{Ejercicio 435 - Ruffini}
\textbf{Enunciado:} $ x^{3} - x^{2} - 10 x - 8 $\par
\begin{enumerate}[leftmargin=2em]
\item Ruffini sirve para dividir un polinomio entre factores de la forma \((x-r)\). Si el resto da cero, entonces \((x-r)\) es factor.
\item Probamos \(r=-1\). Coeficientes actuales: \((1, -1, -10, -8)\).
\item La fila final de Ruffini es \((1, -2, -8, 0)\). El resto es \(0\), por eso el factor es \((x+1)\). El cociente queda \(x^{2} - 2 x - 8\).
\item Probamos \(r=4\). Coeficientes actuales: \((1, -2, -8)\).
\item La fila final de Ruffini es \((1, 2, 0)\). El resto es \(0\), por eso el factor es \((x-4)\). El cociente queda \(x + 2\).
\item Probamos \(r=-2\). Coeficientes actuales: \((1, 2)\).
\item La fila final de Ruffini es \((1, 0)\). El resto es \(0\), por eso el factor es \((x+2)\). El cociente queda \(1\).
\item Juntamos todos los factores encontrados. Resultado: \[\boxed{\left(x - 4\right) \left(x + 1\right) \left(x + 2\right)}\]
\end{enumerate}

\vspace{0.2cm}
\subsubsection*{Ejercicio 436 - Ruffini}
\textbf{Enunciado:} $ x^{3} + 4 x^{2} + x - 6 $\par
\begin{enumerate}[leftmargin=2em]
\item Ruffini sirve para dividir un polinomio entre factores de la forma \((x-r)\). Si el resto da cero, entonces \((x-r)\) es factor.
\item Probamos \(r=-2\). Coeficientes actuales: \((1, 4, 1, -6)\).
\item La fila final de Ruffini es \((1, 2, -3, 0)\). El resto es \(0\), por eso el factor es \((x+2)\). El cociente queda \(x^{2} + 2 x - 3\).
\item Probamos \(r=1\). Coeficientes actuales: \((1, 2, -3)\).
\item La fila final de Ruffini es \((1, 3, 0)\). El resto es \(0\), por eso el factor es \((x-1)\). El cociente queda \(x + 3\).
\item Probamos \(r=-3\). Coeficientes actuales: \((1, 3)\).
\item La fila final de Ruffini es \((1, 0)\). El resto es \(0\), por eso el factor es \((x+3)\). El cociente queda \(1\).
\item Juntamos todos los factores encontrados. Resultado: \[\boxed{\left(x - 1\right) \left(x + 2\right) \left(x + 3\right)}\]
\end{enumerate}

\vspace{0.2cm}
\subsubsection*{Ejercicio 437 - Ruffini}
\textbf{Enunciado:} $ x^{3} + 5 x^{2} - 2 x - 24 $\par
\begin{enumerate}[leftmargin=2em]
\item Ruffini sirve para dividir un polinomio entre factores de la forma \((x-r)\). Si el resto da cero, entonces \((x-r)\) es factor.
\item Probamos \(r=-3\). Coeficientes actuales: \((1, 5, -2, -24)\).
\item La fila final de Ruffini es \((1, 2, -8, 0)\). El resto es \(0\), por eso el factor es \((x+3)\). El cociente queda \(x^{2} + 2 x - 8\).
\item Probamos \(r=2\). Coeficientes actuales: \((1, 2, -8)\).
\item La fila final de Ruffini es \((1, 4, 0)\). El resto es \(0\), por eso el factor es \((x-2)\). El cociente queda \(x + 4\).
\item Probamos \(r=-4\). Coeficientes actuales: \((1, 4)\).
\item La fila final de Ruffini es \((1, 0)\). El resto es \(0\), por eso el factor es \((x+4)\). El cociente queda \(1\).
\item Juntamos todos los factores encontrados. Resultado: \[\boxed{\left(x - 2\right) \left(x + 3\right) \left(x + 4\right)}\]
\end{enumerate}

\vspace{0.2cm}
\subsubsection*{Ejercicio 438 - Ruffini}
\textbf{Enunciado:} $ x^{3} + 3 x^{2} - 10 x - 24 $\par
\begin{enumerate}[leftmargin=2em]
\item Ruffini sirve para dividir un polinomio entre factores de la forma \((x-r)\). Si el resto da cero, entonces \((x-r)\) es factor.
\item Probamos \(r=-4\). Coeficientes actuales: \((1, 3, -10, -24)\).
\item La fila final de Ruffini es \((1, -1, -6, 0)\). El resto es \(0\), por eso el factor es \((x+4)\). El cociente queda \(x^{2} - x - 6\).
\item Probamos \(r=3\). Coeficientes actuales: \((1, -1, -6)\).
\item La fila final de Ruffini es \((1, 2, 0)\). El resto es \(0\), por eso el factor es \((x-3)\). El cociente queda \(x + 2\).
\item Probamos \(r=-2\). Coeficientes actuales: \((1, 2)\).
\item La fila final de Ruffini es \((1, 0)\). El resto es \(0\), por eso el factor es \((x+2)\). El cociente queda \(1\).
\item Juntamos todos los factores encontrados. Resultado: \[\boxed{\left(x - 3\right) \left(x + 2\right) \left(x + 4\right)}\]
\end{enumerate}

\vspace{0.2cm}
\subsubsection*{Ejercicio 439 - Ruffini}
\textbf{Enunciado:} $ x^{3} + 4 x^{2} - 17 x - 60 $\par
\begin{enumerate}[leftmargin=2em]
\item Ruffini sirve para dividir un polinomio entre factores de la forma \((x-r)\). Si el resto da cero, entonces \((x-r)\) es factor.
\item Probamos \(r=-5\). Coeficientes actuales: \((1, 4, -17, -60)\).
\item La fila final de Ruffini es \((1, -1, -12, 0)\). El resto es \(0\), por eso el factor es \((x+5)\). El cociente queda \(x^{2} - x - 12\).
\item Probamos \(r=4\). Coeficientes actuales: \((1, -1, -12)\).
\item La fila final de Ruffini es \((1, 3, 0)\). El resto es \(0\), por eso el factor es \((x-4)\). El cociente queda \(x + 3\).
\item Probamos \(r=-3\). Coeficientes actuales: \((1, 3)\).
\item La fila final de Ruffini es \((1, 0)\). El resto es \(0\), por eso el factor es \((x+3)\). El cociente queda \(1\).
\item Juntamos todos los factores encontrados. Resultado: \[\boxed{\left(x - 4\right) \left(x + 3\right) \left(x + 5\right)}\]
\end{enumerate}

\vspace{0.2cm}
\subsubsection*{Ejercicio 440 - Ruffini}
\textbf{Enunciado:} $ x^{3} + 4 x^{2} - x - 4 $\par
\begin{enumerate}[leftmargin=2em]
\item Ruffini sirve para dividir un polinomio entre factores de la forma \((x-r)\). Si el resto da cero, entonces \((x-r)\) es factor.
\item Probamos \(r=-1\). Coeficientes actuales: \((1, 4, -1, -4)\).
\item La fila final de Ruffini es \((1, 3, -4, 0)\). El resto es \(0\), por eso el factor es \((x+1)\). El cociente queda \(x^{2} + 3 x - 4\).
\item Probamos \(r=1\). Coeficientes actuales: \((1, 3, -4)\).
\item La fila final de Ruffini es \((1, 4, 0)\). El resto es \(0\), por eso el factor es \((x-1)\). El cociente queda \(x + 4\).
\item Probamos \(r=-4\). Coeficientes actuales: \((1, 4)\).
\item La fila final de Ruffini es \((1, 0)\). El resto es \(0\), por eso el factor es \((x+4)\). El cociente queda \(1\).
\item Juntamos todos los factores encontrados. Resultado: \[\boxed{\left(x - 1\right) \left(x + 1\right) \left(x + 4\right)}\]
\end{enumerate}

\vspace{0.2cm}
\nivel{3}{intermedio}
\subsubsection*{Ejercicio 441 - Ruffini}
\textbf{Enunciado:} $ 3 x^{3} - 9 x^{2} - 12 x + 36 $\par
\begin{enumerate}[leftmargin=2em]
\item Ruffini sirve para dividir un polinomio entre factores de la forma \((x-r)\). Si el resto da cero, entonces \((x-r)\) es factor.
\item Probamos \(r=2\). Coeficientes actuales: \((3, -9, -12, 36)\).
\item La fila final de Ruffini es \((3, -3, -18, 0)\). El resto es \(0\), por eso el factor es \((x-2)\). El cociente queda \(3 x^{2} - 3 x - 18\).
\item Probamos \(r=-2\). Coeficientes actuales: \((3, -3, -18)\).
\item La fila final de Ruffini es \((3, -9, 0)\). El resto es \(0\), por eso el factor es \((x+2)\). El cociente queda \(3 x - 9\).
\item Probamos \(r=3\). Coeficientes actuales: \((3, -9)\).
\item La fila final de Ruffini es \((3, 0)\). El resto es \(0\), por eso el factor es \((x-3)\). El cociente queda \(3\).
\item Juntamos todos los factores encontrados. Resultado: \[\boxed{3 \left(x - 3\right) \left(x - 2\right) \left(x + 2\right)}\]
\end{enumerate}

\vspace{0.2cm}
\subsubsection*{Ejercicio 442 - Ruffini}
\textbf{Enunciado:} $ 4 x^{3} - 16 x^{2} - 36 x + 144 $\par
\begin{enumerate}[leftmargin=2em]
\item Ruffini sirve para dividir un polinomio entre factores de la forma \((x-r)\). Si el resto da cero, entonces \((x-r)\) es factor.
\item Probamos \(r=3\). Coeficientes actuales: \((4, -16, -36, 144)\).
\item La fila final de Ruffini es \((4, -4, -48, 0)\). El resto es \(0\), por eso el factor es \((x-3)\). El cociente queda \(4 x^{2} - 4 x - 48\).
\item Probamos \(r=-3\). Coeficientes actuales: \((4, -4, -48)\).
\item La fila final de Ruffini es \((4, -16, 0)\). El resto es \(0\), por eso el factor es \((x+3)\). El cociente queda \(4 x - 16\).
\item Probamos \(r=4\). Coeficientes actuales: \((4, -16)\).
\item La fila final de Ruffini es \((4, 0)\). El resto es \(0\), por eso el factor es \((x-4)\). El cociente queda \(4\).
\item Juntamos todos los factores encontrados. Resultado: \[\boxed{4 \left(x - 4\right) \left(x - 3\right) \left(x + 3\right)}\]
\end{enumerate}

\vspace{0.2cm}
\subsubsection*{Ejercicio 443 - Ruffini}
\textbf{Enunciado:} $ 2 x^{3} - 4 x^{2} - 32 x + 64 $\par
\begin{enumerate}[leftmargin=2em]
\item Ruffini sirve para dividir un polinomio entre factores de la forma \((x-r)\). Si el resto da cero, entonces \((x-r)\) es factor.
\item Probamos \(r=4\). Coeficientes actuales: \((2, -4, -32, 64)\).
\item La fila final de Ruffini es \((2, 4, -16, 0)\). El resto es \(0\), por eso el factor es \((x-4)\). El cociente queda \(2 x^{2} + 4 x - 16\).
\item Probamos \(r=-4\). Coeficientes actuales: \((2, 4, -16)\).
\item La fila final de Ruffini es \((2, -4, 0)\). El resto es \(0\), por eso el factor es \((x+4)\). El cociente queda \(2 x - 4\).
\item Probamos \(r=2\). Coeficientes actuales: \((2, -4)\).
\item La fila final de Ruffini es \((2, 0)\). El resto es \(0\), por eso el factor es \((x-2)\). El cociente queda \(2\).
\item Juntamos todos los factores encontrados. Resultado: \[\boxed{2 \left(x - 4\right) \left(x - 2\right) \left(x + 4\right)}\]
\end{enumerate}

\vspace{0.2cm}
\subsubsection*{Ejercicio 444 - Ruffini}
\textbf{Enunciado:} $ 3 x^{3} + 3 x^{2} - 51 x + 45 $\par
\begin{enumerate}[leftmargin=2em]
\item Ruffini sirve para dividir un polinomio entre factores de la forma \((x-r)\). Si el resto da cero, entonces \((x-r)\) es factor.
\item Probamos \(r=1\). Coeficientes actuales: \((3, 3, -51, 45)\).
\item La fila final de Ruffini es \((3, 6, -45, 0)\). El resto es \(0\), por eso el factor es \((x-1)\). El cociente queda \(3 x^{2} + 6 x - 45\).
\item Probamos \(r=-5\). Coeficientes actuales: \((3, 6, -45)\).
\item La fila final de Ruffini es \((3, -9, 0)\). El resto es \(0\), por eso el factor es \((x+5)\). El cociente queda \(3 x - 9\).
\item Probamos \(r=3\). Coeficientes actuales: \((3, -9)\).
\item La fila final de Ruffini es \((3, 0)\). El resto es \(0\), por eso el factor es \((x-3)\). El cociente queda \(3\).
\item Juntamos todos los factores encontrados. Resultado: \[\boxed{3 \left(x - 3\right) \left(x - 1\right) \left(x + 5\right)}\]
\end{enumerate}

\vspace{0.2cm}
\subsubsection*{Ejercicio 445 - Ruffini}
\textbf{Enunciado:} $ 4 x^{3} - 20 x^{2} + 8 x + 32 $\par
\begin{enumerate}[leftmargin=2em]
\item Ruffini sirve para dividir un polinomio entre factores de la forma \((x-r)\). Si el resto da cero, entonces \((x-r)\) es factor.
\item Probamos \(r=2\). Coeficientes actuales: \((4, -20, 8, 32)\).
\item La fila final de Ruffini es \((4, -12, -16, 0)\). El resto es \(0\), por eso el factor es \((x-2)\). El cociente queda \(4 x^{2} - 12 x - 16\).
\item Probamos \(r=-1\). Coeficientes actuales: \((4, -12, -16)\).
\item La fila final de Ruffini es \((4, -16, 0)\). El resto es \(0\), por eso el factor es \((x+1)\). El cociente queda \(4 x - 16\).
\item Probamos \(r=4\). Coeficientes actuales: \((4, -16)\).
\item La fila final de Ruffini es \((4, 0)\). El resto es \(0\), por eso el factor es \((x-4)\). El cociente queda \(4\).
\item Juntamos todos los factores encontrados. Resultado: \[\boxed{4 \left(x - 4\right) \left(x - 2\right) \left(x + 1\right)}\]
\end{enumerate}

\vspace{0.2cm}
\subsubsection*{Ejercicio 446 - Ruffini}
\textbf{Enunciado:} $ 2 x^{3} - 6 x^{2} - 8 x + 24 $\par
\begin{enumerate}[leftmargin=2em]
\item Ruffini sirve para dividir un polinomio entre factores de la forma \((x-r)\). Si el resto da cero, entonces \((x-r)\) es factor.
\item Probamos \(r=3\). Coeficientes actuales: \((2, -6, -8, 24)\).
\item La fila final de Ruffini es \((2, 0, -8, 0)\). El resto es \(0\), por eso el factor es \((x-3)\). El cociente queda \(2 x^{2} - 8\).
\item Probamos \(r=-2\). Coeficientes actuales: \((2, 0, -8)\).
\item La fila final de Ruffini es \((2, -4, 0)\). El resto es \(0\), por eso el factor es \((x+2)\). El cociente queda \(2 x - 4\).
\item Probamos \(r=2\). Coeficientes actuales: \((2, -4)\).
\item La fila final de Ruffini es \((2, 0)\). El resto es \(0\), por eso el factor es \((x-2)\). El cociente queda \(2\).
\item Juntamos todos los factores encontrados. Resultado: \[\boxed{2 \left(x - 3\right) \left(x - 2\right) \left(x + 2\right)}\]
\end{enumerate}

\vspace{0.2cm}
\subsubsection*{Ejercicio 447 - Ruffini}
\textbf{Enunciado:} $ 3 x^{3} - 12 x^{2} - 27 x + 108 $\par
\begin{enumerate}[leftmargin=2em]
\item Ruffini sirve para dividir un polinomio entre factores de la forma \((x-r)\). Si el resto da cero, entonces \((x-r)\) es factor.
\item Probamos \(r=4\). Coeficientes actuales: \((3, -12, -27, 108)\).
\item La fila final de Ruffini es \((3, 0, -27, 0)\). El resto es \(0\), por eso el factor es \((x-4)\). El cociente queda \(3 x^{2} - 27\).
\item Probamos \(r=-3\). Coeficientes actuales: \((3, 0, -27)\).
\item La fila final de Ruffini es \((3, -9, 0)\). El resto es \(0\), por eso el factor es \((x+3)\). El cociente queda \(3 x - 9\).
\item Probamos \(r=3\). Coeficientes actuales: \((3, -9)\).
\item La fila final de Ruffini es \((3, 0)\). El resto es \(0\), por eso el factor es \((x-3)\). El cociente queda \(3\).
\item Juntamos todos los factores encontrados. Resultado: \[\boxed{3 \left(x - 4\right) \left(x - 3\right) \left(x + 3\right)}\]
\end{enumerate}

\vspace{0.2cm}
\subsubsection*{Ejercicio 448 - Ruffini}
\textbf{Enunciado:} $ 4 x^{3} - 4 x^{2} - 64 x + 64 $\par
\begin{enumerate}[leftmargin=2em]
\item Ruffini sirve para dividir un polinomio entre factores de la forma \((x-r)\). Si el resto da cero, entonces \((x-r)\) es factor.
\item Probamos \(r=1\). Coeficientes actuales: \((4, -4, -64, 64)\).
\item La fila final de Ruffini es \((4, 0, -64, 0)\). El resto es \(0\), por eso el factor es \((x-1)\). El cociente queda \(4 x^{2} - 64\).
\item Probamos \(r=-4\). Coeficientes actuales: \((4, 0, -64)\).
\item La fila final de Ruffini es \((4, -16, 0)\). El resto es \(0\), por eso el factor es \((x+4)\). El cociente queda \(4 x - 16\).
\item Probamos \(r=4\). Coeficientes actuales: \((4, -16)\).
\item La fila final de Ruffini es \((4, 0)\). El resto es \(0\), por eso el factor es \((x-4)\). El cociente queda \(4\).
\item Juntamos todos los factores encontrados. Resultado: \[\boxed{4 \left(x - 4\right) \left(x - 1\right) \left(x + 4\right)}\]
\end{enumerate}

\vspace{0.2cm}
\subsubsection*{Ejercicio 449 - Ruffini}
\textbf{Enunciado:} $ 2 x^{3} + 2 x^{2} - 32 x + 40 $\par
\begin{enumerate}[leftmargin=2em]
\item Ruffini sirve para dividir un polinomio entre factores de la forma \((x-r)\). Si el resto da cero, entonces \((x-r)\) es factor.
\item Probamos \(r=2\). Coeficientes actuales: \((2, 2, -32, 40)\).
\item La fila final de Ruffini es \((2, 6, -20, 0)\). El resto es \(0\), por eso el factor es \((x-2)\). El cociente queda \(2 x^{2} + 6 x - 20\).
\item Probamos \(r=-5\). Coeficientes actuales: \((2, 6, -20)\).
\item La fila final de Ruffini es \((2, -4, 0)\). El resto es \(0\), por eso el factor es \((x+5)\). El cociente queda \(2 x - 4\).
\item Probamos \(r=2\). Coeficientes actuales: \((2, -4)\).
\item La fila final de Ruffini es \((2, 0)\). El resto es \(0\), por eso el factor es \((x-2)\). El cociente queda \(2\).
\item Juntamos todos los factores encontrados. Resultado: \[\boxed{2 \left(x - 2\right)^{2} \left(x + 5\right)}\]
\end{enumerate}

\vspace{0.2cm}
\subsubsection*{Ejercicio 450 - Ruffini}
\textbf{Enunciado:} $ 3 x^{3} - 15 x^{2} + 9 x + 27 $\par
\begin{enumerate}[leftmargin=2em]
\item Ruffini sirve para dividir un polinomio entre factores de la forma \((x-r)\). Si el resto da cero, entonces \((x-r)\) es factor.
\item Probamos \(r=3\). Coeficientes actuales: \((3, -15, 9, 27)\).
\item La fila final de Ruffini es \((3, -6, -9, 0)\). El resto es \(0\), por eso el factor es \((x-3)\). El cociente queda \(3 x^{2} - 6 x - 9\).
\item Probamos \(r=-1\). Coeficientes actuales: \((3, -6, -9)\).
\item La fila final de Ruffini es \((3, -9, 0)\). El resto es \(0\), por eso el factor es \((x+1)\). El cociente queda \(3 x - 9\).
\item Probamos \(r=3\). Coeficientes actuales: \((3, -9)\).
\item La fila final de Ruffini es \((3, 0)\). El resto es \(0\), por eso el factor es \((x-3)\). El cociente queda \(3\).
\item Juntamos todos los factores encontrados. Resultado: \[\boxed{3 \left(x - 3\right)^{2} \left(x + 1\right)}\]
\end{enumerate}

\vspace{0.2cm}
\subsubsection*{Ejercicio 451 - Ruffini}
\textbf{Enunciado:} $ 4 x^{3} - 24 x^{2} + 128 $\par
\begin{enumerate}[leftmargin=2em]
\item Ruffini sirve para dividir un polinomio entre factores de la forma \((x-r)\). Si el resto da cero, entonces \((x-r)\) es factor.
\item Probamos \(r=4\). Coeficientes actuales: \((4, -24, 0, 128)\).
\item La fila final de Ruffini es \((4, -8, -32, 0)\). El resto es \(0\), por eso el factor es \((x-4)\). El cociente queda \(4 x^{2} - 8 x - 32\).
\item Probamos \(r=-2\). Coeficientes actuales: \((4, -8, -32)\).
\item La fila final de Ruffini es \((4, -16, 0)\). El resto es \(0\), por eso el factor es \((x+2)\). El cociente queda \(4 x - 16\).
\item Probamos \(r=4\). Coeficientes actuales: \((4, -16)\).
\item La fila final de Ruffini es \((4, 0)\). El resto es \(0\), por eso el factor es \((x-4)\). El cociente queda \(4\).
\item Juntamos todos los factores encontrados. Resultado: \[\boxed{4 \left(x - 4\right)^{2} \left(x + 2\right)}\]
\end{enumerate}

\vspace{0.2cm}
\subsubsection*{Ejercicio 452 - Ruffini}
\textbf{Enunciado:} $ 2 x^{3} - 14 x + 12 $\par
\begin{enumerate}[leftmargin=2em]
\item Ruffini sirve para dividir un polinomio entre factores de la forma \((x-r)\). Si el resto da cero, entonces \((x-r)\) es factor.
\item Probamos \(r=1\). Coeficientes actuales: \((2, 0, -14, 12)\).
\item La fila final de Ruffini es \((2, 2, -12, 0)\). El resto es \(0\), por eso el factor es \((x-1)\). El cociente queda \(2 x^{2} + 2 x - 12\).
\item Probamos \(r=-3\). Coeficientes actuales: \((2, 2, -12)\).
\item La fila final de Ruffini es \((2, -4, 0)\). El resto es \(0\), por eso el factor es \((x+3)\). El cociente queda \(2 x - 4\).
\item Probamos \(r=2\). Coeficientes actuales: \((2, -4)\).
\item La fila final de Ruffini es \((2, 0)\). El resto es \(0\), por eso el factor es \((x-2)\). El cociente queda \(2\).
\item Juntamos todos los factores encontrados. Resultado: \[\boxed{2 \left(x - 2\right) \left(x - 1\right) \left(x + 3\right)}\]
\end{enumerate}

\vspace{0.2cm}
\subsubsection*{Ejercicio 453 - Ruffini}
\textbf{Enunciado:} $ 3 x^{3} - 3 x^{2} - 42 x + 72 $\par
\begin{enumerate}[leftmargin=2em]
\item Ruffini sirve para dividir un polinomio entre factores de la forma \((x-r)\). Si el resto da cero, entonces \((x-r)\) es factor.
\item Probamos \(r=2\). Coeficientes actuales: \((3, -3, -42, 72)\).
\item La fila final de Ruffini es \((3, 3, -36, 0)\). El resto es \(0\), por eso el factor es \((x-2)\). El cociente queda \(3 x^{2} + 3 x - 36\).
\item Probamos \(r=-4\). Coeficientes actuales: \((3, 3, -36)\).
\item La fila final de Ruffini es \((3, -9, 0)\). El resto es \(0\), por eso el factor es \((x+4)\). El cociente queda \(3 x - 9\).
\item Probamos \(r=3\). Coeficientes actuales: \((3, -9)\).
\item La fila final de Ruffini es \((3, 0)\). El resto es \(0\), por eso el factor es \((x-3)\). El cociente queda \(3\).
\item Juntamos todos los factores encontrados. Resultado: \[\boxed{3 \left(x - 3\right) \left(x - 2\right) \left(x + 4\right)}\]
\end{enumerate}

\vspace{0.2cm}
\subsubsection*{Ejercicio 454 - Ruffini}
\textbf{Enunciado:} $ 4 x^{3} - 8 x^{2} - 92 x + 240 $\par
\begin{enumerate}[leftmargin=2em]
\item Ruffini sirve para dividir un polinomio entre factores de la forma \((x-r)\). Si el resto da cero, entonces \((x-r)\) es factor.
\item Probamos \(r=3\). Coeficientes actuales: \((4, -8, -92, 240)\).
\item La fila final de Ruffini es \((4, 4, -80, 0)\). El resto es \(0\), por eso el factor es \((x-3)\). El cociente queda \(4 x^{2} + 4 x - 80\).
\item Probamos \(r=-5\). Coeficientes actuales: \((4, 4, -80)\).
\item La fila final de Ruffini es \((4, -16, 0)\). El resto es \(0\), por eso el factor es \((x+5)\). El cociente queda \(4 x - 16\).
\item Probamos \(r=4\). Coeficientes actuales: \((4, -16)\).
\item La fila final de Ruffini es \((4, 0)\). El resto es \(0\), por eso el factor es \((x-4)\). El cociente queda \(4\).
\item Juntamos todos los factores encontrados. Resultado: \[\boxed{4 \left(x - 4\right) \left(x - 3\right) \left(x + 5\right)}\]
\end{enumerate}

\vspace{0.2cm}
\subsubsection*{Ejercicio 455 - Ruffini}
\textbf{Enunciado:} $ 2 x^{3} - 10 x^{2} + 4 x + 16 $\par
\begin{enumerate}[leftmargin=2em]
\item Ruffini sirve para dividir un polinomio entre factores de la forma \((x-r)\). Si el resto da cero, entonces \((x-r)\) es factor.
\item Probamos \(r=4\). Coeficientes actuales: \((2, -10, 4, 16)\).
\item La fila final de Ruffini es \((2, -2, -4, 0)\). El resto es \(0\), por eso el factor es \((x-4)\). El cociente queda \(2 x^{2} - 2 x - 4\).
\item Probamos \(r=-1\). Coeficientes actuales: \((2, -2, -4)\).
\item La fila final de Ruffini es \((2, -4, 0)\). El resto es \(0\), por eso el factor es \((x+1)\). El cociente queda \(2 x - 4\).
\item Probamos \(r=2\). Coeficientes actuales: \((2, -4)\).
\item La fila final de Ruffini es \((2, 0)\). El resto es \(0\), por eso el factor es \((x-2)\). El cociente queda \(2\).
\item Juntamos todos los factores encontrados. Resultado: \[\boxed{2 \left(x - 4\right) \left(x - 2\right) \left(x + 1\right)}\]
\end{enumerate}

\vspace{0.2cm}
\subsubsection*{Ejercicio 456 - Ruffini}
\textbf{Enunciado:} $ 3 x^{3} - 6 x^{2} - 15 x + 18 $\par
\begin{enumerate}[leftmargin=2em]
\item Ruffini sirve para dividir un polinomio entre factores de la forma \((x-r)\). Si el resto da cero, entonces \((x-r)\) es factor.
\item Probamos \(r=1\). Coeficientes actuales: \((3, -6, -15, 18)\).
\item La fila final de Ruffini es \((3, -3, -18, 0)\). El resto es \(0\), por eso el factor es \((x-1)\). El cociente queda \(3 x^{2} - 3 x - 18\).
\item Probamos \(r=-2\). Coeficientes actuales: \((3, -3, -18)\).
\item La fila final de Ruffini es \((3, -9, 0)\). El resto es \(0\), por eso el factor es \((x+2)\). El cociente queda \(3 x - 9\).
\item Probamos \(r=3\). Coeficientes actuales: \((3, -9)\).
\item La fila final de Ruffini es \((3, 0)\). El resto es \(0\), por eso el factor es \((x-3)\). El cociente queda \(3\).
\item Juntamos todos los factores encontrados. Resultado: \[\boxed{3 \left(x - 3\right) \left(x - 1\right) \left(x + 2\right)}\]
\end{enumerate}

\vspace{0.2cm}
\subsubsection*{Ejercicio 457 - Ruffini}
\textbf{Enunciado:} $ 4 x^{3} - 12 x^{2} - 40 x + 96 $\par
\begin{enumerate}[leftmargin=2em]
\item Ruffini sirve para dividir un polinomio entre factores de la forma \((x-r)\). Si el resto da cero, entonces \((x-r)\) es factor.
\item Probamos \(r=2\). Coeficientes actuales: \((4, -12, -40, 96)\).
\item La fila final de Ruffini es \((4, -4, -48, 0)\). El resto es \(0\), por eso el factor es \((x-2)\). El cociente queda \(4 x^{2} - 4 x - 48\).
\item Probamos \(r=-3\). Coeficientes actuales: \((4, -4, -48)\).
\item La fila final de Ruffini es \((4, -16, 0)\). El resto es \(0\), por eso el factor es \((x+3)\). El cociente queda \(4 x - 16\).
\item Probamos \(r=4\). Coeficientes actuales: \((4, -16)\).
\item La fila final de Ruffini es \((4, 0)\). El resto es \(0\), por eso el factor es \((x-4)\). El cociente queda \(4\).
\item Juntamos todos los factores encontrados. Resultado: \[\boxed{4 \left(x - 4\right) \left(x - 2\right) \left(x + 3\right)}\]
\end{enumerate}

\vspace{0.2cm}
\subsubsection*{Ejercicio 458 - Ruffini}
\textbf{Enunciado:} $ 2 x^{3} - 2 x^{2} - 28 x + 48 $\par
\begin{enumerate}[leftmargin=2em]
\item Ruffini sirve para dividir un polinomio entre factores de la forma \((x-r)\). Si el resto da cero, entonces \((x-r)\) es factor.
\item Probamos \(r=3\). Coeficientes actuales: \((2, -2, -28, 48)\).
\item La fila final de Ruffini es \((2, 4, -16, 0)\). El resto es \(0\), por eso el factor es \((x-3)\). El cociente queda \(2 x^{2} + 4 x - 16\).
\item Probamos \(r=-4\). Coeficientes actuales: \((2, 4, -16)\).
\item La fila final de Ruffini es \((2, -4, 0)\). El resto es \(0\), por eso el factor es \((x+4)\). El cociente queda \(2 x - 4\).
\item Probamos \(r=2\). Coeficientes actuales: \((2, -4)\).
\item La fila final de Ruffini es \((2, 0)\). El resto es \(0\), por eso el factor es \((x-2)\). El cociente queda \(2\).
\item Juntamos todos los factores encontrados. Resultado: \[\boxed{2 \left(x - 3\right) \left(x - 2\right) \left(x + 4\right)}\]
\end{enumerate}

\vspace{0.2cm}
\subsubsection*{Ejercicio 459 - Ruffini}
\textbf{Enunciado:} $ 3 x^{3} - 6 x^{2} - 69 x + 180 $\par
\begin{enumerate}[leftmargin=2em]
\item Ruffini sirve para dividir un polinomio entre factores de la forma \((x-r)\). Si el resto da cero, entonces \((x-r)\) es factor.
\item Probamos \(r=4\). Coeficientes actuales: \((3, -6, -69, 180)\).
\item La fila final de Ruffini es \((3, 6, -45, 0)\). El resto es \(0\), por eso el factor es \((x-4)\). El cociente queda \(3 x^{2} + 6 x - 45\).
\item Probamos \(r=-5\). Coeficientes actuales: \((3, 6, -45)\).
\item La fila final de Ruffini es \((3, -9, 0)\). El resto es \(0\), por eso el factor es \((x+5)\). El cociente queda \(3 x - 9\).
\item Probamos \(r=3\). Coeficientes actuales: \((3, -9)\).
\item La fila final de Ruffini es \((3, 0)\). El resto es \(0\), por eso el factor es \((x-3)\). El cociente queda \(3\).
\item Juntamos todos los factores encontrados. Resultado: \[\boxed{3 \left(x - 4\right) \left(x - 3\right) \left(x + 5\right)}\]
\end{enumerate}

\vspace{0.2cm}
\subsubsection*{Ejercicio 460 - Ruffini}
\textbf{Enunciado:} $ 4 x^{3} - 16 x^{2} - 4 x + 16 $\par
\begin{enumerate}[leftmargin=2em]
\item Ruffini sirve para dividir un polinomio entre factores de la forma \((x-r)\). Si el resto da cero, entonces \((x-r)\) es factor.
\item Probamos \(r=1\). Coeficientes actuales: \((4, -16, -4, 16)\).
\item La fila final de Ruffini es \((4, -12, -16, 0)\). El resto es \(0\), por eso el factor es \((x-1)\). El cociente queda \(4 x^{2} - 12 x - 16\).
\item Probamos \(r=-1\). Coeficientes actuales: \((4, -12, -16)\).
\item La fila final de Ruffini es \((4, -16, 0)\). El resto es \(0\), por eso el factor es \((x+1)\). El cociente queda \(4 x - 16\).
\item Probamos \(r=4\). Coeficientes actuales: \((4, -16)\).
\item La fila final de Ruffini es \((4, 0)\). El resto es \(0\), por eso el factor es \((x-4)\). El cociente queda \(4\).
\item Juntamos todos los factores encontrados. Resultado: \[\boxed{4 \left(x - 4\right) \left(x - 1\right) \left(x + 1\right)}\]
\end{enumerate}

\vspace{0.2cm}
\nivel{4}{avanzado}
\subsubsection*{Ejercicio 461 - Ruffini}
\textbf{Enunciado:} $ 3 x^{4} - 39 x^{2} + 108 $\par
\begin{enumerate}[leftmargin=2em]
\item Ruffini sirve para dividir un polinomio entre factores de la forma \((x-r)\). Si el resto da cero, entonces \((x-r)\) es factor.
\item Probamos \(r=2\). Coeficientes actuales: \((3, 0, -39, 0, 108)\).
\item La fila final de Ruffini es \((3, 6, -27, -54, 0)\). El resto es \(0\), por eso el factor es \((x-2)\). El cociente queda \(3 x^{3} + 6 x^{2} - 27 x - 54\).
\item Probamos \(r=-2\). Coeficientes actuales: \((3, 6, -27, -54)\).
\item La fila final de Ruffini es \((3, 0, -27, 0)\). El resto es \(0\), por eso el factor es \((x+2)\). El cociente queda \(3 x^{2} - 27\).
\item Probamos \(r=3\). Coeficientes actuales: \((3, 0, -27)\).
\item La fila final de Ruffini es \((3, 9, 0)\). El resto es \(0\), por eso el factor es \((x-3)\). El cociente queda \(3 x + 9\).
\item Probamos \(r=-3\). Coeficientes actuales: \((3, 9)\).
\item La fila final de Ruffini es \((3, 0)\). El resto es \(0\), por eso el factor es \((x+3)\). El cociente queda \(3\).
\item Juntamos todos los factores encontrados. Resultado: \[\boxed{3 \left(x - 3\right) \left(x - 2\right) \left(x + 2\right) \left(x + 3\right)}\]
\end{enumerate}

\vspace{0.2cm}
\subsubsection*{Ejercicio 462 - Ruffini}
\textbf{Enunciado:} $ 4 x^{4} - 8 x^{3} - 68 x^{2} + 72 x + 288 $\par
\begin{enumerate}[leftmargin=2em]
\item Ruffini sirve para dividir un polinomio entre factores de la forma \((x-r)\). Si el resto da cero, entonces \((x-r)\) es factor.
\item Probamos \(r=3\). Coeficientes actuales: \((4, -8, -68, 72, 288)\).
\item La fila final de Ruffini es \((4, 4, -56, -96, 0)\). El resto es \(0\), por eso el factor es \((x-3)\). El cociente queda \(4 x^{3} + 4 x^{2} - 56 x - 96\).
\item Probamos \(r=-3\). Coeficientes actuales: \((4, 4, -56, -96)\).
\item La fila final de Ruffini es \((4, -8, -32, 0)\). El resto es \(0\), por eso el factor es \((x+3)\). El cociente queda \(4 x^{2} - 8 x - 32\).
\item Probamos \(r=4\). Coeficientes actuales: \((4, -8, -32)\).
\item La fila final de Ruffini es \((4, 8, 0)\). El resto es \(0\), por eso el factor es \((x-4)\). El cociente queda \(4 x + 8\).
\item Probamos \(r=-2\). Coeficientes actuales: \((4, 8)\).
\item La fila final de Ruffini es \((4, 0)\). El resto es \(0\), por eso el factor es \((x+2)\). El cociente queda \(4\).
\item Juntamos todos los factores encontrados. Resultado: \[\boxed{4 \left(x - 4\right) \left(x - 3\right) \left(x + 2\right) \left(x + 3\right)}\]
\end{enumerate}

\vspace{0.2cm}
\subsubsection*{Ejercicio 463 - Ruffini}
\textbf{Enunciado:} $ 5 x^{4} + 5 x^{3} - 110 x^{2} - 80 x + 480 $\par
\begin{enumerate}[leftmargin=2em]
\item Ruffini sirve para dividir un polinomio entre factores de la forma \((x-r)\). Si el resto da cero, entonces \((x-r)\) es factor.
\item Probamos \(r=4\). Coeficientes actuales: \((5, 5, -110, -80, 480)\).
\item La fila final de Ruffini es \((5, 25, -10, -120, 0)\). El resto es \(0\), por eso el factor es \((x-4)\). El cociente queda \(5 x^{3} + 25 x^{2} - 10 x - 120\).
\item Probamos \(r=-4\). Coeficientes actuales: \((5, 25, -10, -120)\).
\item La fila final de Ruffini es \((5, 5, -30, 0)\). El resto es \(0\), por eso el factor es \((x+4)\). El cociente queda \(5 x^{2} + 5 x - 30\).
\item Probamos \(r=2\). Coeficientes actuales: \((5, 5, -30)\).
\item La fila final de Ruffini es \((5, 15, 0)\). El resto es \(0\), por eso el factor es \((x-2)\). El cociente queda \(5 x + 15\).
\item Probamos \(r=-3\). Coeficientes actuales: \((5, 15)\).
\item La fila final de Ruffini es \((5, 0)\). El resto es \(0\), por eso el factor es \((x+3)\). El cociente queda \(5\).
\item Juntamos todos los factores encontrados. Resultado: \[\boxed{5 \left(x - 4\right) \left(x - 2\right) \left(x + 3\right) \left(x + 4\right)}\]
\end{enumerate}

\vspace{0.2cm}
\subsubsection*{Ejercicio 464 - Ruffini}
\textbf{Enunciado:} $ 2 x^{4} - 10 x^{3} - 14 x^{2} + 58 x + 60 $\par
\begin{enumerate}[leftmargin=2em]
\item Ruffini sirve para dividir un polinomio entre factores de la forma \((x-r)\). Si el resto da cero, entonces \((x-r)\) es factor.
\item Probamos \(r=5\). Coeficientes actuales: \((2, -10, -14, 58, 60)\).
\item La fila final de Ruffini es \((2, 0, -14, -12, 0)\). El resto es \(0\), por eso el factor es \((x-5)\). El cociente queda \(2 x^{3} - 14 x - 12\).
\item Probamos \(r=-1\). Coeficientes actuales: \((2, 0, -14, -12)\).
\item La fila final de Ruffini es \((2, -2, -12, 0)\). El resto es \(0\), por eso el factor es \((x+1)\). El cociente queda \(2 x^{2} - 2 x - 12\).
\item Probamos \(r=3\). Coeficientes actuales: \((2, -2, -12)\).
\item La fila final de Ruffini es \((2, 4, 0)\). El resto es \(0\), por eso el factor es \((x-3)\). El cociente queda \(2 x + 4\).
\item Probamos \(r=-2\). Coeficientes actuales: \((2, 4)\).
\item La fila final de Ruffini es \((2, 0)\). El resto es \(0\), por eso el factor es \((x+2)\). El cociente queda \(2\).
\item Juntamos todos los factores encontrados. Resultado: \[\boxed{2 \left(x - 5\right) \left(x - 3\right) \left(x + 1\right) \left(x + 2\right)}\]
\end{enumerate}

\vspace{0.2cm}
\subsubsection*{Ejercicio 465 - Ruffini}
\textbf{Enunciado:} $ 3 x^{4} - 45 x^{2} - 30 x + 72 $\par
\begin{enumerate}[leftmargin=2em]
\item Ruffini sirve para dividir un polinomio entre factores de la forma \((x-r)\). Si el resto da cero, entonces \((x-r)\) es factor.
\item Probamos \(r=1\). Coeficientes actuales: \((3, 0, -45, -30, 72)\).
\item La fila final de Ruffini es \((3, 3, -42, -72, 0)\). El resto es \(0\), por eso el factor es \((x-1)\). El cociente queda \(3 x^{3} + 3 x^{2} - 42 x - 72\).
\item Probamos \(r=-2\). Coeficientes actuales: \((3, 3, -42, -72)\).
\item La fila final de Ruffini es \((3, -3, -36, 0)\). El resto es \(0\), por eso el factor es \((x+2)\). El cociente queda \(3 x^{2} - 3 x - 36\).
\item Probamos \(r=4\). Coeficientes actuales: \((3, -3, -36)\).
\item La fila final de Ruffini es \((3, 9, 0)\). El resto es \(0\), por eso el factor es \((x-4)\). El cociente queda \(3 x + 9\).
\item Probamos \(r=-3\). Coeficientes actuales: \((3, 9)\).
\item La fila final de Ruffini es \((3, 0)\). El resto es \(0\), por eso el factor es \((x+3)\). El cociente queda \(3\).
\item Juntamos todos los factores encontrados. Resultado: \[\boxed{3 \left(x - 4\right) \left(x - 1\right) \left(x + 2\right) \left(x + 3\right)}\]
\end{enumerate}

\vspace{0.2cm}
\subsubsection*{Ejercicio 466 - Ruffini}
\textbf{Enunciado:} $ 4 x^{4} + 4 x^{3} - 40 x^{2} - 16 x + 96 $\par
\begin{enumerate}[leftmargin=2em]
\item Ruffini sirve para dividir un polinomio entre factores de la forma \((x-r)\). Si el resto da cero, entonces \((x-r)\) es factor.
\item Probamos \(r=2\). Coeficientes actuales: \((4, 4, -40, -16, 96)\).
\item La fila final de Ruffini es \((4, 12, -16, -48, 0)\). El resto es \(0\), por eso el factor es \((x-2)\). El cociente queda \(4 x^{3} + 12 x^{2} - 16 x - 48\).
\item Probamos \(r=-3\). Coeficientes actuales: \((4, 12, -16, -48)\).
\item La fila final de Ruffini es \((4, 0, -16, 0)\). El resto es \(0\), por eso el factor es \((x+3)\). El cociente queda \(4 x^{2} - 16\).
\item Probamos \(r=2\). Coeficientes actuales: \((4, 0, -16)\).
\item La fila final de Ruffini es \((4, 8, 0)\). El resto es \(0\), por eso el factor es \((x-2)\). El cociente queda \(4 x + 8\).
\item Probamos \(r=-2\). Coeficientes actuales: \((4, 8)\).
\item La fila final de Ruffini es \((4, 0)\). El resto es \(0\), por eso el factor es \((x+2)\). El cociente queda \(4\).
\item Juntamos todos los factores encontrados. Resultado: \[\boxed{4 \left(x - 2\right)^{2} \left(x + 2\right) \left(x + 3\right)}\]
\end{enumerate}

\vspace{0.2cm}
\subsubsection*{Ejercicio 467 - Ruffini}
\textbf{Enunciado:} $ 5 x^{4} + 5 x^{3} - 105 x^{2} - 45 x + 540 $\par
\begin{enumerate}[leftmargin=2em]
\item Ruffini sirve para dividir un polinomio entre factores de la forma \((x-r)\). Si el resto da cero, entonces \((x-r)\) es factor.
\item Probamos \(r=3\). Coeficientes actuales: \((5, 5, -105, -45, 540)\).
\item La fila final de Ruffini es \((5, 20, -45, -180, 0)\). El resto es \(0\), por eso el factor es \((x-3)\). El cociente queda \(5 x^{3} + 20 x^{2} - 45 x - 180\).
\item Probamos \(r=-4\). Coeficientes actuales: \((5, 20, -45, -180)\).
\item La fila final de Ruffini es \((5, 0, -45, 0)\). El resto es \(0\), por eso el factor es \((x+4)\). El cociente queda \(5 x^{2} - 45\).
\item Probamos \(r=3\). Coeficientes actuales: \((5, 0, -45)\).
\item La fila final de Ruffini es \((5, 15, 0)\). El resto es \(0\), por eso el factor es \((x-3)\). El cociente queda \(5 x + 15\).
\item Probamos \(r=-3\). Coeficientes actuales: \((5, 15)\).
\item La fila final de Ruffini es \((5, 0)\). El resto es \(0\), por eso el factor es \((x+3)\). El cociente queda \(5\).
\item Juntamos todos los factores encontrados. Resultado: \[\boxed{5 \left(x - 3\right)^{2} \left(x + 3\right) \left(x + 4\right)}\]
\end{enumerate}

\vspace{0.2cm}
\subsubsection*{Ejercicio 468 - Ruffini}
\textbf{Enunciado:} $ 2 x^{4} - 10 x^{3} - 12 x^{2} + 64 x + 64 $\par
\begin{enumerate}[leftmargin=2em]
\item Ruffini sirve para dividir un polinomio entre factores de la forma \((x-r)\). Si el resto da cero, entonces \((x-r)\) es factor.
\item Probamos \(r=4\). Coeficientes actuales: \((2, -10, -12, 64, 64)\).
\item La fila final de Ruffini es \((2, -2, -20, -16, 0)\). El resto es \(0\), por eso el factor es \((x-4)\). El cociente queda \(2 x^{3} - 2 x^{2} - 20 x - 16\).
\item Probamos \(r=-1\). Coeficientes actuales: \((2, -2, -20, -16)\).
\item La fila final de Ruffini es \((2, -4, -16, 0)\). El resto es \(0\), por eso el factor es \((x+1)\). El cociente queda \(2 x^{2} - 4 x - 16\).
\item Probamos \(r=4\). Coeficientes actuales: \((2, -4, -16)\).
\item La fila final de Ruffini es \((2, 4, 0)\). El resto es \(0\), por eso el factor es \((x-4)\). El cociente queda \(2 x + 4\).
\item Probamos \(r=-2\). Coeficientes actuales: \((2, 4)\).
\item La fila final de Ruffini es \((2, 0)\). El resto es \(0\), por eso el factor es \((x+2)\). El cociente queda \(2\).
\item Juntamos todos los factores encontrados. Resultado: \[\boxed{2 \left(x - 4\right)^{2} \left(x + 1\right) \left(x + 2\right)}\]
\end{enumerate}

\vspace{0.2cm}
\subsubsection*{Ejercicio 469 - Ruffini}
\textbf{Enunciado:} $ 3 x^{4} - 6 x^{3} - 57 x^{2} + 24 x + 180 $\par
\begin{enumerate}[leftmargin=2em]
\item Ruffini sirve para dividir un polinomio entre factores de la forma \((x-r)\). Si el resto da cero, entonces \((x-r)\) es factor.
\item Probamos \(r=5\). Coeficientes actuales: \((3, -6, -57, 24, 180)\).
\item La fila final de Ruffini es \((3, 9, -12, -36, 0)\). El resto es \(0\), por eso el factor es \((x-5)\). El cociente queda \(3 x^{3} + 9 x^{2} - 12 x - 36\).
\item Probamos \(r=-2\). Coeficientes actuales: \((3, 9, -12, -36)\).
\item La fila final de Ruffini es \((3, 3, -18, 0)\). El resto es \(0\), por eso el factor es \((x+2)\). El cociente queda \(3 x^{2} + 3 x - 18\).
\item Probamos \(r=2\). Coeficientes actuales: \((3, 3, -18)\).
\item La fila final de Ruffini es \((3, 9, 0)\). El resto es \(0\), por eso el factor es \((x-2)\). El cociente queda \(3 x + 9\).
\item Probamos \(r=-3\). Coeficientes actuales: \((3, 9)\).
\item La fila final de Ruffini es \((3, 0)\). El resto es \(0\), por eso el factor es \((x+3)\). El cociente queda \(3\).
\item Juntamos todos los factores encontrados. Resultado: \[\boxed{3 \left(x - 5\right) \left(x - 2\right) \left(x + 2\right) \left(x + 3\right)}\]
\end{enumerate}

\vspace{0.2cm}
\subsubsection*{Ejercicio 470 - Ruffini}
\textbf{Enunciado:} $ 4 x^{4} + 4 x^{3} - 44 x^{2} - 36 x + 72 $\par
\begin{enumerate}[leftmargin=2em]
\item Ruffini sirve para dividir un polinomio entre factores de la forma \((x-r)\). Si el resto da cero, entonces \((x-r)\) es factor.
\item Probamos \(r=1\). Coeficientes actuales: \((4, 4, -44, -36, 72)\).
\item La fila final de Ruffini es \((4, 8, -36, -72, 0)\). El resto es \(0\), por eso el factor es \((x-1)\). El cociente queda \(4 x^{3} + 8 x^{2} - 36 x - 72\).
\item Probamos \(r=-3\). Coeficientes actuales: \((4, 8, -36, -72)\).
\item La fila final de Ruffini es \((4, -4, -24, 0)\). El resto es \(0\), por eso el factor es \((x+3)\). El cociente queda \(4 x^{2} - 4 x - 24\).
\item Probamos \(r=3\). Coeficientes actuales: \((4, -4, -24)\).
\item La fila final de Ruffini es \((4, 8, 0)\). El resto es \(0\), por eso el factor es \((x-3)\). El cociente queda \(4 x + 8\).
\item Probamos \(r=-2\). Coeficientes actuales: \((4, 8)\).
\item La fila final de Ruffini es \((4, 0)\). El resto es \(0\), por eso el factor es \((x+2)\). El cociente queda \(4\).
\item Juntamos todos los factores encontrados. Resultado: \[\boxed{4 \left(x - 3\right) \left(x - 1\right) \left(x + 2\right) \left(x + 3\right)}\]
\end{enumerate}

\vspace{0.2cm}
\subsubsection*{Ejercicio 471 - Ruffini}
\textbf{Enunciado:} $ 5 x^{4} + 5 x^{3} - 110 x^{2} - 80 x + 480 $\par
\begin{enumerate}[leftmargin=2em]
\item Ruffini sirve para dividir un polinomio entre factores de la forma \((x-r)\). Si el resto da cero, entonces \((x-r)\) es factor.
\item Probamos \(r=2\). Coeficientes actuales: \((5, 5, -110, -80, 480)\).
\item La fila final de Ruffini es \((5, 15, -80, -240, 0)\). El resto es \(0\), por eso el factor es \((x-2)\). El cociente queda \(5 x^{3} + 15 x^{2} - 80 x - 240\).
\item Probamos \(r=-4\). Coeficientes actuales: \((5, 15, -80, -240)\).
\item La fila final de Ruffini es \((5, -5, -60, 0)\). El resto es \(0\), por eso el factor es \((x+4)\). El cociente queda \(5 x^{2} - 5 x - 60\).
\item Probamos \(r=4\). Coeficientes actuales: \((5, -5, -60)\).
\item La fila final de Ruffini es \((5, 15, 0)\). El resto es \(0\), por eso el factor es \((x-4)\). El cociente queda \(5 x + 15\).
\item Probamos \(r=-3\). Coeficientes actuales: \((5, 15)\).
\item La fila final de Ruffini es \((5, 0)\). El resto es \(0\), por eso el factor es \((x+3)\). El cociente queda \(5\).
\item Juntamos todos los factores encontrados. Resultado: \[\boxed{5 \left(x - 4\right) \left(x - 2\right) \left(x + 3\right) \left(x + 4\right)}\]
\end{enumerate}

\vspace{0.2cm}
\subsubsection*{Ejercicio 472 - Ruffini}
\textbf{Enunciado:} $ 2 x^{4} - 4 x^{3} - 14 x^{2} + 16 x + 24 $\par
\begin{enumerate}[leftmargin=2em]
\item Ruffini sirve para dividir un polinomio entre factores de la forma \((x-r)\). Si el resto da cero, entonces \((x-r)\) es factor.
\item Probamos \(r=3\). Coeficientes actuales: \((2, -4, -14, 16, 24)\).
\item La fila final de Ruffini es \((2, 2, -8, -8, 0)\). El resto es \(0\), por eso el factor es \((x-3)\). El cociente queda \(2 x^{3} + 2 x^{2} - 8 x - 8\).
\item Probamos \(r=-1\). Coeficientes actuales: \((2, 2, -8, -8)\).
\item La fila final de Ruffini es \((2, 0, -8, 0)\). El resto es \(0\), por eso el factor es \((x+1)\). El cociente queda \(2 x^{2} - 8\).
\item Probamos \(r=2\). Coeficientes actuales: \((2, 0, -8)\).
\item La fila final de Ruffini es \((2, 4, 0)\). El resto es \(0\), por eso el factor es \((x-2)\). El cociente queda \(2 x + 4\).
\item Probamos \(r=-2\). Coeficientes actuales: \((2, 4)\).
\item La fila final de Ruffini es \((2, 0)\). El resto es \(0\), por eso el factor es \((x+2)\). El cociente queda \(2\).
\item Juntamos todos los factores encontrados. Resultado: \[\boxed{2 \left(x - 3\right) \left(x - 2\right) \left(x + 1\right) \left(x + 2\right)}\]
\end{enumerate}

\vspace{0.2cm}
\subsubsection*{Ejercicio 473 - Ruffini}
\textbf{Enunciado:} $ 3 x^{4} - 6 x^{3} - 51 x^{2} + 54 x + 216 $\par
\begin{enumerate}[leftmargin=2em]
\item Ruffini sirve para dividir un polinomio entre factores de la forma \((x-r)\). Si el resto da cero, entonces \((x-r)\) es factor.
\item Probamos \(r=4\). Coeficientes actuales: \((3, -6, -51, 54, 216)\).
\item La fila final de Ruffini es \((3, 6, -27, -54, 0)\). El resto es \(0\), por eso el factor es \((x-4)\). El cociente queda \(3 x^{3} + 6 x^{2} - 27 x - 54\).
\item Probamos \(r=-2\). Coeficientes actuales: \((3, 6, -27, -54)\).
\item La fila final de Ruffini es \((3, 0, -27, 0)\). El resto es \(0\), por eso el factor es \((x+2)\). El cociente queda \(3 x^{2} - 27\).
\item Probamos \(r=3\). Coeficientes actuales: \((3, 0, -27)\).
\item La fila final de Ruffini es \((3, 9, 0)\). El resto es \(0\), por eso el factor es \((x-3)\). El cociente queda \(3 x + 9\).
\item Probamos \(r=-3\). Coeficientes actuales: \((3, 9)\).
\item La fila final de Ruffini es \((3, 0)\). El resto es \(0\), por eso el factor es \((x+3)\). El cociente queda \(3\).
\item Juntamos todos los factores encontrados. Resultado: \[\boxed{3 \left(x - 4\right) \left(x - 3\right) \left(x + 2\right) \left(x + 3\right)}\]
\end{enumerate}

\vspace{0.2cm}
\subsubsection*{Ejercicio 474 - Ruffini}
\textbf{Enunciado:} $ 4 x^{4} - 16 x^{3} - 76 x^{2} + 184 x + 480 $\par
\begin{enumerate}[leftmargin=2em]
\item Ruffini sirve para dividir un polinomio entre factores de la forma \((x-r)\). Si el resto da cero, entonces \((x-r)\) es factor.
\item Probamos \(r=5\). Coeficientes actuales: \((4, -16, -76, 184, 480)\).
\item La fila final de Ruffini es \((4, 4, -56, -96, 0)\). El resto es \(0\), por eso el factor es \((x-5)\). El cociente queda \(4 x^{3} + 4 x^{2} - 56 x - 96\).
\item Probamos \(r=-3\). Coeficientes actuales: \((4, 4, -56, -96)\).
\item La fila final de Ruffini es \((4, -8, -32, 0)\). El resto es \(0\), por eso el factor es \((x+3)\). El cociente queda \(4 x^{2} - 8 x - 32\).
\item Probamos \(r=4\). Coeficientes actuales: \((4, -8, -32)\).
\item La fila final de Ruffini es \((4, 8, 0)\). El resto es \(0\), por eso el factor es \((x-4)\). El cociente queda \(4 x + 8\).
\item Probamos \(r=-2\). Coeficientes actuales: \((4, 8)\).
\item La fila final de Ruffini es \((4, 0)\). El resto es \(0\), por eso el factor es \((x+2)\). El cociente queda \(4\).
\item Juntamos todos los factores encontrados. Resultado: \[\boxed{4 \left(x - 5\right) \left(x - 4\right) \left(x + 2\right) \left(x + 3\right)}\]
\end{enumerate}

\vspace{0.2cm}
\subsubsection*{Ejercicio 475 - Ruffini}
\textbf{Enunciado:} $ 5 x^{4} + 20 x^{3} - 35 x^{2} - 110 x + 120 $\par
\begin{enumerate}[leftmargin=2em]
\item Ruffini sirve para dividir un polinomio entre factores de la forma \((x-r)\). Si el resto da cero, entonces \((x-r)\) es factor.
\item Probamos \(r=1\). Coeficientes actuales: \((5, 20, -35, -110, 120)\).
\item La fila final de Ruffini es \((5, 25, -10, -120, 0)\). El resto es \(0\), por eso el factor es \((x-1)\). El cociente queda \(5 x^{3} + 25 x^{2} - 10 x - 120\).
\item Probamos \(r=-4\). Coeficientes actuales: \((5, 25, -10, -120)\).
\item La fila final de Ruffini es \((5, 5, -30, 0)\). El resto es \(0\), por eso el factor es \((x+4)\). El cociente queda \(5 x^{2} + 5 x - 30\).
\item Probamos \(r=2\). Coeficientes actuales: \((5, 5, -30)\).
\item La fila final de Ruffini es \((5, 15, 0)\). El resto es \(0\), por eso el factor es \((x-2)\). El cociente queda \(5 x + 15\).
\item Probamos \(r=-3\). Coeficientes actuales: \((5, 15)\).
\item La fila final de Ruffini es \((5, 0)\). El resto es \(0\), por eso el factor es \((x+3)\). El cociente queda \(5\).
\item Juntamos todos los factores encontrados. Resultado: \[\boxed{5 \left(x - 2\right) \left(x - 1\right) \left(x + 3\right) \left(x + 4\right)}\]
\end{enumerate}

\vspace{0.2cm}
\subsubsection*{Ejercicio 476 - Ruffini}
\textbf{Enunciado:} $ 2 x^{4} - 4 x^{3} - 14 x^{2} + 16 x + 24 $\par
\begin{enumerate}[leftmargin=2em]
\item Ruffini sirve para dividir un polinomio entre factores de la forma \((x-r)\). Si el resto da cero, entonces \((x-r)\) es factor.
\item Probamos \(r=2\). Coeficientes actuales: \((2, -4, -14, 16, 24)\).
\item La fila final de Ruffini es \((2, 0, -14, -12, 0)\). El resto es \(0\), por eso el factor es \((x-2)\). El cociente queda \(2 x^{3} - 14 x - 12\).
\item Probamos \(r=-1\). Coeficientes actuales: \((2, 0, -14, -12)\).
\item La fila final de Ruffini es \((2, -2, -12, 0)\). El resto es \(0\), por eso el factor es \((x+1)\). El cociente queda \(2 x^{2} - 2 x - 12\).
\item Probamos \(r=3\). Coeficientes actuales: \((2, -2, -12)\).
\item La fila final de Ruffini es \((2, 4, 0)\). El resto es \(0\), por eso el factor es \((x-3)\). El cociente queda \(2 x + 4\).
\item Probamos \(r=-2\). Coeficientes actuales: \((2, 4)\).
\item La fila final de Ruffini es \((2, 0)\). El resto es \(0\), por eso el factor es \((x+2)\). El cociente queda \(2\).
\item Juntamos todos los factores encontrados. Resultado: \[\boxed{2 \left(x - 3\right) \left(x - 2\right) \left(x + 1\right) \left(x + 2\right)}\]
\end{enumerate}

\vspace{0.2cm}
\subsubsection*{Ejercicio 477 - Ruffini}
\textbf{Enunciado:} $ 3 x^{4} - 6 x^{3} - 51 x^{2} + 54 x + 216 $\par
\begin{enumerate}[leftmargin=2em]
\item Ruffini sirve para dividir un polinomio entre factores de la forma \((x-r)\). Si el resto da cero, entonces \((x-r)\) es factor.
\item Probamos \(r=3\). Coeficientes actuales: \((3, -6, -51, 54, 216)\).
\item La fila final de Ruffini es \((3, 3, -42, -72, 0)\). El resto es \(0\), por eso el factor es \((x-3)\). El cociente queda \(3 x^{3} + 3 x^{2} - 42 x - 72\).
\item Probamos \(r=-2\). Coeficientes actuales: \((3, 3, -42, -72)\).
\item La fila final de Ruffini es \((3, -3, -36, 0)\). El resto es \(0\), por eso el factor es \((x+2)\). El cociente queda \(3 x^{2} - 3 x - 36\).
\item Probamos \(r=4\). Coeficientes actuales: \((3, -3, -36)\).
\item La fila final de Ruffini es \((3, 9, 0)\). El resto es \(0\), por eso el factor es \((x-4)\). El cociente queda \(3 x + 9\).
\item Probamos \(r=-3\). Coeficientes actuales: \((3, 9)\).
\item La fila final de Ruffini es \((3, 0)\). El resto es \(0\), por eso el factor es \((x+3)\). El cociente queda \(3\).
\item Juntamos todos los factores encontrados. Resultado: \[\boxed{3 \left(x - 4\right) \left(x - 3\right) \left(x + 2\right) \left(x + 3\right)}\]
\end{enumerate}

\vspace{0.2cm}
\subsubsection*{Ejercicio 478 - Ruffini}
\textbf{Enunciado:} $ 4 x^{4} - 4 x^{3} - 64 x^{2} + 16 x + 192 $\par
\begin{enumerate}[leftmargin=2em]
\item Ruffini sirve para dividir un polinomio entre factores de la forma \((x-r)\). Si el resto da cero, entonces \((x-r)\) es factor.
\item Probamos \(r=4\). Coeficientes actuales: \((4, -4, -64, 16, 192)\).
\item La fila final de Ruffini es \((4, 12, -16, -48, 0)\). El resto es \(0\), por eso el factor es \((x-4)\). El cociente queda \(4 x^{3} + 12 x^{2} - 16 x - 48\).
\item Probamos \(r=-3\). Coeficientes actuales: \((4, 12, -16, -48)\).
\item La fila final de Ruffini es \((4, 0, -16, 0)\). El resto es \(0\), por eso el factor es \((x+3)\). El cociente queda \(4 x^{2} - 16\).
\item Probamos \(r=2\). Coeficientes actuales: \((4, 0, -16)\).
\item La fila final de Ruffini es \((4, 8, 0)\). El resto es \(0\), por eso el factor es \((x-2)\). El cociente queda \(4 x + 8\).
\item Probamos \(r=-2\). Coeficientes actuales: \((4, 8)\).
\item La fila final de Ruffini es \((4, 0)\). El resto es \(0\), por eso el factor es \((x+2)\). El cociente queda \(4\).
\item Juntamos todos los factores encontrados. Resultado: \[\boxed{4 \left(x - 4\right) \left(x - 2\right) \left(x + 2\right) \left(x + 3\right)}\]
\end{enumerate}

\vspace{0.2cm}
\subsubsection*{Ejercicio 479 - Ruffini}
\textbf{Enunciado:} $ 5 x^{4} - 5 x^{3} - 145 x^{2} + 45 x + 900 $\par
\begin{enumerate}[leftmargin=2em]
\item Ruffini sirve para dividir un polinomio entre factores de la forma \((x-r)\). Si el resto da cero, entonces \((x-r)\) es factor.
\item Probamos \(r=5\). Coeficientes actuales: \((5, -5, -145, 45, 900)\).
\item La fila final de Ruffini es \((5, 20, -45, -180, 0)\). El resto es \(0\), por eso el factor es \((x-5)\). El cociente queda \(5 x^{3} + 20 x^{2} - 45 x - 180\).
\item Probamos \(r=-4\). Coeficientes actuales: \((5, 20, -45, -180)\).
\item La fila final de Ruffini es \((5, 0, -45, 0)\). El resto es \(0\), por eso el factor es \((x+4)\). El cociente queda \(5 x^{2} - 45\).
\item Probamos \(r=3\). Coeficientes actuales: \((5, 0, -45)\).
\item La fila final de Ruffini es \((5, 15, 0)\). El resto es \(0\), por eso el factor es \((x-3)\). El cociente queda \(5 x + 15\).
\item Probamos \(r=-3\). Coeficientes actuales: \((5, 15)\).
\item La fila final de Ruffini es \((5, 0)\). El resto es \(0\), por eso el factor es \((x+3)\). El cociente queda \(5\).
\item Juntamos todos los factores encontrados. Resultado: \[\boxed{5 \left(x - 5\right) \left(x - 3\right) \left(x + 3\right) \left(x + 4\right)}\]
\end{enumerate}

\vspace{0.2cm}
\subsubsection*{Ejercicio 480 - Ruffini}
\textbf{Enunciado:} $ 2 x^{4} - 4 x^{3} - 18 x^{2} + 4 x + 16 $\par
\begin{enumerate}[leftmargin=2em]
\item Ruffini sirve para dividir un polinomio entre factores de la forma \((x-r)\). Si el resto da cero, entonces \((x-r)\) es factor.
\item Probamos \(r=1\). Coeficientes actuales: \((2, -4, -18, 4, 16)\).
\item La fila final de Ruffini es \((2, -2, -20, -16, 0)\). El resto es \(0\), por eso el factor es \((x-1)\). El cociente queda \(2 x^{3} - 2 x^{2} - 20 x - 16\).
\item Probamos \(r=-1\). Coeficientes actuales: \((2, -2, -20, -16)\).
\item La fila final de Ruffini es \((2, -4, -16, 0)\). El resto es \(0\), por eso el factor es \((x+1)\). El cociente queda \(2 x^{2} - 4 x - 16\).
\item Probamos \(r=4\). Coeficientes actuales: \((2, -4, -16)\).
\item La fila final de Ruffini es \((2, 4, 0)\). El resto es \(0\), por eso el factor es \((x-4)\). El cociente queda \(2 x + 4\).
\item Probamos \(r=-2\). Coeficientes actuales: \((2, 4)\).
\item La fila final de Ruffini es \((2, 0)\). El resto es \(0\), por eso el factor es \((x+2)\). El cociente queda \(2\).
\item Juntamos todos los factores encontrados. Resultado: \[\boxed{2 \left(x - 4\right) \left(x - 1\right) \left(x + 1\right) \left(x + 2\right)}\]
\end{enumerate}

\vspace{0.2cm}
\nivel{5}{imposible}
\subsubsection*{Ejercicio 481 - Ruffini}
\textbf{Enunciado:} $ 3 x^{5} - 9 x^{4} - 39 x^{3} + 117 x^{2} + 108 x - 324 $\par
\begin{enumerate}[leftmargin=2em]
\item Ruffini sirve para dividir un polinomio entre factores de la forma \((x-r)\). Si el resto da cero, entonces \((x-r)\) es factor.
\item Probamos \(r=2\). Coeficientes actuales: \((3, -9, -39, 117, 108, -324)\).
\item La fila final de Ruffini es \((3, -3, -45, 27, 162, 0)\). El resto es \(0\), por eso el factor es \((x-2)\). El cociente queda \(3 x^{4} - 3 x^{3} - 45 x^{2} + 27 x + 162\).
\item Probamos \(r=-2\). Coeficientes actuales: \((3, -3, -45, 27, 162)\).
\item La fila final de Ruffini es \((3, -9, -27, 81, 0)\). El resto es \(0\), por eso el factor es \((x+2)\). El cociente queda \(3 x^{3} - 9 x^{2} - 27 x + 81\).
\item Probamos \(r=3\). Coeficientes actuales: \((3, -9, -27, 81)\).
\item La fila final de Ruffini es \((3, 0, -27, 0)\). El resto es \(0\), por eso el factor es \((x-3)\). El cociente queda \(3 x^{2} - 27\).
\item Probamos \(r=-3\). Coeficientes actuales: \((3, 0, -27)\).
\item La fila final de Ruffini es \((3, -9, 0)\). El resto es \(0\), por eso el factor es \((x+3)\). El cociente queda \(3 x - 9\).
\item Probamos \(r=3\). Coeficientes actuales: \((3, -9)\).
\item La fila final de Ruffini es \((3, 0)\). El resto es \(0\), por eso el factor es \((x-3)\). El cociente queda \(3\).
\item Juntamos todos los factores encontrados. Resultado: \[\boxed{3 \left(x - 3\right)^{2} \left(x - 2\right) \left(x + 2\right) \left(x + 3\right)}\]
\end{enumerate}

\vspace{0.2cm}
\subsubsection*{Ejercicio 482 - Ruffini}
\textbf{Enunciado:} $ 4 x^{5} - 8 x^{4} - 100 x^{3} + 200 x^{2} + 576 x - 1152 $\par
\begin{enumerate}[leftmargin=2em]
\item Ruffini sirve para dividir un polinomio entre factores de la forma \((x-r)\). Si el resto da cero, entonces \((x-r)\) es factor.
\item Probamos \(r=3\). Coeficientes actuales: \((4, -8, -100, 200, 576, -1152)\).
\item La fila final de Ruffini es \((4, 4, -88, -64, 384, 0)\). El resto es \(0\), por eso el factor es \((x-3)\). El cociente queda \(4 x^{4} + 4 x^{3} - 88 x^{2} - 64 x + 384\).
\item Probamos \(r=-3\). Coeficientes actuales: \((4, 4, -88, -64, 384)\).
\item La fila final de Ruffini es \((4, -8, -64, 128, 0)\). El resto es \(0\), por eso el factor es \((x+3)\). El cociente queda \(4 x^{3} - 8 x^{2} - 64 x + 128\).
\item Probamos \(r=4\). Coeficientes actuales: \((4, -8, -64, 128)\).
\item La fila final de Ruffini es \((4, 8, -32, 0)\). El resto es \(0\), por eso el factor es \((x-4)\). El cociente queda \(4 x^{2} + 8 x - 32\).
\item Probamos \(r=-4\). Coeficientes actuales: \((4, 8, -32)\).
\item La fila final de Ruffini es \((4, -8, 0)\). El resto es \(0\), por eso el factor es \((x+4)\). El cociente queda \(4 x - 8\).
\item Probamos \(r=2\). Coeficientes actuales: \((4, -8)\).
\item La fila final de Ruffini es \((4, 0)\). El resto es \(0\), por eso el factor es \((x-2)\). El cociente queda \(4\).
\item Juntamos todos los factores encontrados. Resultado: \[\boxed{4 \left(x - 4\right) \left(x - 3\right) \left(x - 2\right) \left(x + 3\right) \left(x + 4\right)}\]
\end{enumerate}

\vspace{0.2cm}
\subsubsection*{Ejercicio 483 - Ruffini}
\textbf{Enunciado:} $ 5 x^{5} - 30 x^{4} - 85 x^{3} + 630 x^{2} + 80 x - 2400 $\par
\begin{enumerate}[leftmargin=2em]
\item Ruffini sirve para dividir un polinomio entre factores de la forma \((x-r)\). Si el resto da cero, entonces \((x-r)\) es factor.
\item Probamos \(r=4\). Coeficientes actuales: \((5, -30, -85, 630, 80, -2400)\).
\item La fila final de Ruffini es \((5, -10, -125, 130, 600, 0)\). El resto es \(0\), por eso el factor es \((x-4)\). El cociente queda \(5 x^{4} - 10 x^{3} - 125 x^{2} + 130 x + 600\).
\item Probamos \(r=-4\). Coeficientes actuales: \((5, -10, -125, 130, 600)\).
\item La fila final de Ruffini es \((5, -30, -5, 150, 0)\). El resto es \(0\), por eso el factor es \((x+4)\). El cociente queda \(5 x^{3} - 30 x^{2} - 5 x + 150\).
\item Probamos \(r=5\). Coeficientes actuales: \((5, -30, -5, 150)\).
\item La fila final de Ruffini es \((5, -5, -30, 0)\). El resto es \(0\), por eso el factor es \((x-5)\). El cociente queda \(5 x^{2} - 5 x - 30\).
\item Probamos \(r=-2\). Coeficientes actuales: \((5, -5, -30)\).
\item La fila final de Ruffini es \((5, -15, 0)\). El resto es \(0\), por eso el factor es \((x+2)\). El cociente queda \(5 x - 15\).
\item Probamos \(r=3\). Coeficientes actuales: \((5, -15)\).
\item La fila final de Ruffini es \((5, 0)\). El resto es \(0\), por eso el factor es \((x-3)\). El cociente queda \(5\).
\item Juntamos todos los factores encontrados. Resultado: \[\boxed{5 \left(x - 5\right) \left(x - 4\right) \left(x - 3\right) \left(x + 2\right) \left(x + 4\right)}\]
\end{enumerate}

\vspace{0.2cm}
\subsubsection*{Ejercicio 484 - Ruffini}
\textbf{Enunciado:} $ 6 x^{5} - 6 x^{4} - 198 x^{3} + 222 x^{2} + 1200 x - 1800 $\par
\begin{enumerate}[leftmargin=2em]
\item Ruffini sirve para dividir un polinomio entre factores de la forma \((x-r)\). Si el resto da cero, entonces \((x-r)\) es factor.
\item Probamos \(r=5\). Coeficientes actuales: \((6, -6, -198, 222, 1200, -1800)\).
\item La fila final de Ruffini es \((6, 24, -78, -168, 360, 0)\). El resto es \(0\), por eso el factor es \((x-5)\). El cociente queda \(6 x^{4} + 24 x^{3} - 78 x^{2} - 168 x + 360\).
\item Probamos \(r=-5\). Coeficientes actuales: \((6, 24, -78, -168, 360)\).
\item La fila final de Ruffini es \((6, -6, -48, 72, 0)\). El resto es \(0\), por eso el factor es \((x+5)\). El cociente queda \(6 x^{3} - 6 x^{2} - 48 x + 72\).
\item Probamos \(r=2\). Coeficientes actuales: \((6, -6, -48, 72)\).
\item La fila final de Ruffini es \((6, 6, -36, 0)\). El resto es \(0\), por eso el factor es \((x-2)\). El cociente queda \(6 x^{2} + 6 x - 36\).
\item Probamos \(r=-3\). Coeficientes actuales: \((6, 6, -36)\).
\item La fila final de Ruffini es \((6, -12, 0)\). El resto es \(0\), por eso el factor es \((x+3)\). El cociente queda \(6 x - 12\).
\item Probamos \(r=2\). Coeficientes actuales: \((6, -12)\).
\item La fila final de Ruffini es \((6, 0)\). El resto es \(0\), por eso el factor es \((x-2)\). El cociente queda \(6\).
\item Juntamos todos los factores encontrados. Resultado: \[\boxed{6 \left(x - 5\right) \left(x - 2\right)^{2} \left(x + 3\right) \left(x + 5\right)}\]
\end{enumerate}

\vspace{0.2cm}
\subsubsection*{Ejercicio 485 - Ruffini}
\textbf{Enunciado:} $ 2 x^{5} - 14 x^{4} - 22 x^{3} + 246 x^{2} - 180 x - 432 $\par
\begin{enumerate}[leftmargin=2em]
\item Ruffini sirve para dividir un polinomio entre factores de la forma \((x-r)\). Si el resto da cero, entonces \((x-r)\) es factor.
\item Probamos \(r=6\). Coeficientes actuales: \((2, -14, -22, 246, -180, -432)\).
\item La fila final de Ruffini es \((2, -2, -34, 42, 72, 0)\). El resto es \(0\), por eso el factor es \((x-6)\). El cociente queda \(2 x^{4} - 2 x^{3} - 34 x^{2} + 42 x + 72\).
\item Probamos \(r=-1\). Coeficientes actuales: \((2, -2, -34, 42, 72)\).
\item La fila final de Ruffini es \((2, -4, -30, 72, 0)\). El resto es \(0\), por eso el factor es \((x+1)\). El cociente queda \(2 x^{3} - 4 x^{2} - 30 x + 72\).
\item Probamos \(r=3\). Coeficientes actuales: \((2, -4, -30, 72)\).
\item La fila final de Ruffini es \((2, 2, -24, 0)\). El resto es \(0\), por eso el factor es \((x-3)\). El cociente queda \(2 x^{2} + 2 x - 24\).
\item Probamos \(r=-4\). Coeficientes actuales: \((2, 2, -24)\).
\item La fila final de Ruffini es \((2, -6, 0)\). El resto es \(0\), por eso el factor es \((x+4)\). El cociente queda \(2 x - 6\).
\item Probamos \(r=3\). Coeficientes actuales: \((2, -6)\).
\item La fila final de Ruffini es \((2, 0)\). El resto es \(0\), por eso el factor es \((x-3)\). El cociente queda \(2\).
\item Juntamos todos los factores encontrados. Resultado: \[\boxed{2 \left(x - 6\right) \left(x - 3\right)^{2} \left(x + 1\right) \left(x + 4\right)}\]
\end{enumerate}

\vspace{0.2cm}
\subsubsection*{Ejercicio 486 - Ruffini}
\textbf{Enunciado:} $ 3 x^{5} - 9 x^{4} - 30 x^{3} + 60 x^{2} + 72 x - 96 $\par
\begin{enumerate}[leftmargin=2em]
\item Ruffini sirve para dividir un polinomio entre factores de la forma \((x-r)\). Si el resto da cero, entonces \((x-r)\) es factor.
\item Probamos \(r=1\). Coeficientes actuales: \((3, -9, -30, 60, 72, -96)\).
\item La fila final de Ruffini es \((3, -6, -36, 24, 96, 0)\). El resto es \(0\), por eso el factor es \((x-1)\). El cociente queda \(3 x^{4} - 6 x^{3} - 36 x^{2} + 24 x + 96\).
\item Probamos \(r=-2\). Coeficientes actuales: \((3, -6, -36, 24, 96)\).
\item La fila final de Ruffini es \((3, -12, -12, 48, 0)\). El resto es \(0\), por eso el factor es \((x+2)\). El cociente queda \(3 x^{3} - 12 x^{2} - 12 x + 48\).
\item Probamos \(r=4\). Coeficientes actuales: \((3, -12, -12, 48)\).
\item La fila final de Ruffini es \((3, 0, -12, 0)\). El resto es \(0\), por eso el factor es \((x-4)\). El cociente queda \(3 x^{2} - 12\).
\item Probamos \(r=-2\). Coeficientes actuales: \((3, 0, -12)\).
\item La fila final de Ruffini es \((3, -6, 0)\). El resto es \(0\), por eso el factor es \((x+2)\). El cociente queda \(3 x - 6\).
\item Probamos \(r=2\). Coeficientes actuales: \((3, -6)\).
\item La fila final de Ruffini es \((3, 0)\). El resto es \(0\), por eso el factor es \((x-2)\). El cociente queda \(3\).
\item Juntamos todos los factores encontrados. Resultado: \[\boxed{3 \left(x - 4\right) \left(x - 2\right) \left(x - 1\right) \left(x + 2\right)^{2}}\]
\end{enumerate}

\vspace{0.2cm}
\subsubsection*{Ejercicio 487 - Ruffini}
\textbf{Enunciado:} $ 4 x^{5} - 16 x^{4} - 80 x^{3} + 264 x^{2} + 396 x - 1080 $\par
\begin{enumerate}[leftmargin=2em]
\item Ruffini sirve para dividir un polinomio entre factores de la forma \((x-r)\). Si el resto da cero, entonces \((x-r)\) es factor.
\item Probamos \(r=2\). Coeficientes actuales: \((4, -16, -80, 264, 396, -1080)\).
\item La fila final de Ruffini es \((4, -8, -96, 72, 540, 0)\). El resto es \(0\), por eso el factor es \((x-2)\). El cociente queda \(4 x^{4} - 8 x^{3} - 96 x^{2} + 72 x + 540\).
\item Probamos \(r=-3\). Coeficientes actuales: \((4, -8, -96, 72, 540)\).
\item La fila final de Ruffini es \((4, -20, -36, 180, 0)\). El resto es \(0\), por eso el factor es \((x+3)\). El cociente queda \(4 x^{3} - 20 x^{2} - 36 x + 180\).
\item Probamos \(r=5\). Coeficientes actuales: \((4, -20, -36, 180)\).
\item La fila final de Ruffini es \((4, 0, -36, 0)\). El resto es \(0\), por eso el factor es \((x-5)\). El cociente queda \(4 x^{2} - 36\).
\item Probamos \(r=-3\). Coeficientes actuales: \((4, 0, -36)\).
\item La fila final de Ruffini es \((4, -12, 0)\). El resto es \(0\), por eso el factor es \((x+3)\). El cociente queda \(4 x - 12\).
\item Probamos \(r=3\). Coeficientes actuales: \((4, -12)\).
\item La fila final de Ruffini es \((4, 0)\). El resto es \(0\), por eso el factor es \((x-3)\). El cociente queda \(4\).
\item Juntamos todos los factores encontrados. Resultado: \[\boxed{4 \left(x - 5\right) \left(x - 3\right) \left(x - 2\right) \left(x + 3\right)^{2}}\]
\end{enumerate}

\vspace{0.2cm}
\subsubsection*{Ejercicio 488 - Ruffini}
\textbf{Enunciado:} $ 5 x^{5} + 5 x^{4} - 120 x^{3} + 20 x^{2} + 800 x - 960 $\par
\begin{enumerate}[leftmargin=2em]
\item Ruffini sirve para dividir un polinomio entre factores de la forma \((x-r)\). Si el resto da cero, entonces \((x-r)\) es factor.
\item Probamos \(r=3\). Coeficientes actuales: \((5, 5, -120, 20, 800, -960)\).
\item La fila final de Ruffini es \((5, 20, -60, -160, 320, 0)\). El resto es \(0\), por eso el factor es \((x-3)\). El cociente queda \(5 x^{4} + 20 x^{3} - 60 x^{2} - 160 x + 320\).
\item Probamos \(r=-4\). Coeficientes actuales: \((5, 20, -60, -160, 320)\).
\item La fila final de Ruffini es \((5, 0, -60, 80, 0)\). El resto es \(0\), por eso el factor es \((x+4)\). El cociente queda \(5 x^{3} - 60 x + 80\).
\item Probamos \(r=2\). Coeficientes actuales: \((5, 0, -60, 80)\).
\item La fila final de Ruffini es \((5, 10, -40, 0)\). El resto es \(0\), por eso el factor es \((x-2)\). El cociente queda \(5 x^{2} + 10 x - 40\).
\item Probamos \(r=-4\). Coeficientes actuales: \((5, 10, -40)\).
\item La fila final de Ruffini es \((5, -10, 0)\). El resto es \(0\), por eso el factor es \((x+4)\). El cociente queda \(5 x - 10\).
\item Probamos \(r=2\). Coeficientes actuales: \((5, -10)\).
\item La fila final de Ruffini es \((5, 0)\). El resto es \(0\), por eso el factor es \((x-2)\). El cociente queda \(5\).
\item Juntamos todos los factores encontrados. Resultado: \[\boxed{5 \left(x - 3\right) \left(x - 2\right)^{2} \left(x + 4\right)^{2}}\]
\end{enumerate}

\vspace{0.2cm}
\subsubsection*{Ejercicio 489 - Ruffini}
\textbf{Enunciado:} $ 6 x^{5} - 18 x^{4} - 162 x^{3} + 570 x^{2} + 468 x - 2160 $\par
\begin{enumerate}[leftmargin=2em]
\item Ruffini sirve para dividir un polinomio entre factores de la forma \((x-r)\). Si el resto da cero, entonces \((x-r)\) es factor.
\item Probamos \(r=4\). Coeficientes actuales: \((6, -18, -162, 570, 468, -2160)\).
\item La fila final de Ruffini es \((6, 6, -138, 18, 540, 0)\). El resto es \(0\), por eso el factor es \((x-4)\). El cociente queda \(6 x^{4} + 6 x^{3} - 138 x^{2} + 18 x + 540\).
\item Probamos \(r=-5\). Coeficientes actuales: \((6, 6, -138, 18, 540)\).
\item La fila final de Ruffini es \((6, -24, -18, 108, 0)\). El resto es \(0\), por eso el factor es \((x+5)\). El cociente queda \(6 x^{3} - 24 x^{2} - 18 x + 108\).
\item Probamos \(r=3\). Coeficientes actuales: \((6, -24, -18, 108)\).
\item La fila final de Ruffini es \((6, -6, -36, 0)\). El resto es \(0\), por eso el factor es \((x-3)\). El cociente queda \(6 x^{2} - 6 x - 36\).
\item Probamos \(r=-2\). Coeficientes actuales: \((6, -6, -36)\).
\item La fila final de Ruffini es \((6, -18, 0)\). El resto es \(0\), por eso el factor es \((x+2)\). El cociente queda \(6 x - 18\).
\item Probamos \(r=3\). Coeficientes actuales: \((6, -18)\).
\item La fila final de Ruffini es \((6, 0)\). El resto es \(0\), por eso el factor es \((x-3)\). El cociente queda \(6\).
\item Juntamos todos los factores encontrados. Resultado: \[\boxed{6 \left(x - 4\right) \left(x - 3\right)^{2} \left(x + 2\right) \left(x + 5\right)}\]
\end{enumerate}

\vspace{0.2cm}
\subsubsection*{Ejercicio 490 - Ruffini}
\textbf{Enunciado:} $ 2 x^{5} - 14 x^{4} - 6 x^{3} + 158 x^{2} - 92 x - 240 $\par
\begin{enumerate}[leftmargin=2em]
\item Ruffini sirve para dividir un polinomio entre factores de la forma \((x-r)\). Si el resto da cero, entonces \((x-r)\) es factor.
\item Probamos \(r=5\). Coeficientes actuales: \((2, -14, -6, 158, -92, -240)\).
\item La fila final de Ruffini es \((2, -4, -26, 28, 48, 0)\). El resto es \(0\), por eso el factor es \((x-5)\). El cociente queda \(2 x^{4} - 4 x^{3} - 26 x^{2} + 28 x + 48\).
\item Probamos \(r=-1\). Coeficientes actuales: \((2, -4, -26, 28, 48)\).
\item La fila final de Ruffini es \((2, -6, -20, 48, 0)\). El resto es \(0\), por eso el factor es \((x+1)\). El cociente queda \(2 x^{3} - 6 x^{2} - 20 x + 48\).
\item Probamos \(r=4\). Coeficientes actuales: \((2, -6, -20, 48)\).
\item La fila final de Ruffini es \((2, 2, -12, 0)\). El resto es \(0\), por eso el factor es \((x-4)\). El cociente queda \(2 x^{2} + 2 x - 12\).
\item Probamos \(r=-3\). Coeficientes actuales: \((2, 2, -12)\).
\item La fila final de Ruffini es \((2, -4, 0)\). El resto es \(0\), por eso el factor es \((x+3)\). El cociente queda \(2 x - 4\).
\item Probamos \(r=2\). Coeficientes actuales: \((2, -4)\).
\item La fila final de Ruffini es \((2, 0)\). El resto es \(0\), por eso el factor es \((x-2)\). El cociente queda \(2\).
\item Juntamos todos los factores encontrados. Resultado: \[\boxed{2 \left(x - 5\right) \left(x - 4\right) \left(x - 2\right) \left(x + 1\right) \left(x + 3\right)}\]
\end{enumerate}

\vspace{0.2cm}
\subsubsection*{Ejercicio 491 - Ruffini}
\textbf{Enunciado:} $ 3 x^{5} - 24 x^{4} - 39 x^{3} + 528 x^{2} - 108 x - 2160 $\par
\begin{enumerate}[leftmargin=2em]
\item Ruffini sirve para dividir un polinomio entre factores de la forma \((x-r)\). Si el resto da cero, entonces \((x-r)\) es factor.
\item Probamos \(r=6\). Coeficientes actuales: \((3, -24, -39, 528, -108, -2160)\).
\item La fila final de Ruffini es \((3, -6, -75, 78, 360, 0)\). El resto es \(0\), por eso el factor es \((x-6)\). El cociente queda \(3 x^{4} - 6 x^{3} - 75 x^{2} + 78 x + 360\).
\item Probamos \(r=-2\). Coeficientes actuales: \((3, -6, -75, 78, 360)\).
\item La fila final de Ruffini es \((3, -12, -51, 180, 0)\). El resto es \(0\), por eso el factor es \((x+2)\). El cociente queda \(3 x^{3} - 12 x^{2} - 51 x + 180\).
\item Probamos \(r=5\). Coeficientes actuales: \((3, -12, -51, 180)\).
\item La fila final de Ruffini es \((3, 3, -36, 0)\). El resto es \(0\), por eso el factor es \((x-5)\). El cociente queda \(3 x^{2} + 3 x - 36\).
\item Probamos \(r=-4\). Coeficientes actuales: \((3, 3, -36)\).
\item La fila final de Ruffini es \((3, -9, 0)\). El resto es \(0\), por eso el factor es \((x+4)\). El cociente queda \(3 x - 9\).
\item Probamos \(r=3\). Coeficientes actuales: \((3, -9)\).
\item La fila final de Ruffini es \((3, 0)\). El resto es \(0\), por eso el factor es \((x-3)\). El cociente queda \(3\).
\item Juntamos todos los factores encontrados. Resultado: \[\boxed{3 \left(x - 6\right) \left(x - 5\right) \left(x - 3\right) \left(x + 2\right) \left(x + 4\right)}\]
\end{enumerate}

\vspace{0.2cm}
\subsubsection*{Ejercicio 492 - Ruffini}
\textbf{Enunciado:} $ 4 x^{5} - 44 x^{3} + 24 x^{2} + 112 x - 96 $\par
\begin{enumerate}[leftmargin=2em]
\item Ruffini sirve para dividir un polinomio entre factores de la forma \((x-r)\). Si el resto da cero, entonces \((x-r)\) es factor.
\item Probamos \(r=1\). Coeficientes actuales: \((4, 0, -44, 24, 112, -96)\).
\item La fila final de Ruffini es \((4, 4, -40, -16, 96, 0)\). El resto es \(0\), por eso el factor es \((x-1)\). El cociente queda \(4 x^{4} + 4 x^{3} - 40 x^{2} - 16 x + 96\).
\item Probamos \(r=-3\). Coeficientes actuales: \((4, 4, -40, -16, 96)\).
\item La fila final de Ruffini es \((4, -8, -16, 32, 0)\). El resto es \(0\), por eso el factor es \((x+3)\). El cociente queda \(4 x^{3} - 8 x^{2} - 16 x + 32\).
\item Probamos \(r=2\). Coeficientes actuales: \((4, -8, -16, 32)\).
\item La fila final de Ruffini es \((4, 0, -16, 0)\). El resto es \(0\), por eso el factor es \((x-2)\). El cociente queda \(4 x^{2} - 16\).
\item Probamos \(r=-2\). Coeficientes actuales: \((4, 0, -16)\).
\item La fila final de Ruffini es \((4, -8, 0)\). El resto es \(0\), por eso el factor es \((x+2)\). El cociente queda \(4 x - 8\).
\item Probamos \(r=2\). Coeficientes actuales: \((4, -8)\).
\item La fila final de Ruffini es \((4, 0)\). El resto es \(0\), por eso el factor es \((x-2)\). El cociente queda \(4\).
\item Juntamos todos los factores encontrados. Resultado: \[\boxed{4 \left(x - 2\right)^{2} \left(x - 1\right) \left(x + 2\right) \left(x + 3\right)}\]
\end{enumerate}

\vspace{0.2cm}
\subsubsection*{Ejercicio 493 - Ruffini}
\textbf{Enunciado:} $ 5 x^{5} - 5 x^{4} - 115 x^{3} + 165 x^{2} + 630 x - 1080 $\par
\begin{enumerate}[leftmargin=2em]
\item Ruffini sirve para dividir un polinomio entre factores de la forma \((x-r)\). Si el resto da cero, entonces \((x-r)\) es factor.
\item Probamos \(r=2\). Coeficientes actuales: \((5, -5, -115, 165, 630, -1080)\).
\item La fila final de Ruffini es \((5, 5, -105, -45, 540, 0)\). El resto es \(0\), por eso el factor es \((x-2)\). El cociente queda \(5 x^{4} + 5 x^{3} - 105 x^{2} - 45 x + 540\).
\item Probamos \(r=-4\). Coeficientes actuales: \((5, 5, -105, -45, 540)\).
\item La fila final de Ruffini es \((5, -15, -45, 135, 0)\). El resto es \(0\), por eso el factor es \((x+4)\). El cociente queda \(5 x^{3} - 15 x^{2} - 45 x + 135\).
\item Probamos \(r=3\). Coeficientes actuales: \((5, -15, -45, 135)\).
\item La fila final de Ruffini es \((5, 0, -45, 0)\). El resto es \(0\), por eso el factor es \((x-3)\). El cociente queda \(5 x^{2} - 45\).
\item Probamos \(r=-3\). Coeficientes actuales: \((5, 0, -45)\).
\item La fila final de Ruffini es \((5, -15, 0)\). El resto es \(0\), por eso el factor es \((x+3)\). El cociente queda \(5 x - 15\).
\item Probamos \(r=3\). Coeficientes actuales: \((5, -15)\).
\item La fila final de Ruffini es \((5, 0)\). El resto es \(0\), por eso el factor es \((x-3)\). El cociente queda \(5\).
\item Juntamos todos los factores encontrados. Resultado: \[\boxed{5 \left(x - 3\right)^{2} \left(x - 2\right) \left(x + 3\right) \left(x + 4\right)}\]
\end{enumerate}

\vspace{0.2cm}
\subsubsection*{Ejercicio 494 - Ruffini}
\textbf{Enunciado:} $ 6 x^{5} - 210 x^{3} + 180 x^{2} + 1824 x - 2880 $\par
\begin{enumerate}[leftmargin=2em]
\item Ruffini sirve para dividir un polinomio entre factores de la forma \((x-r)\). Si el resto da cero, entonces \((x-r)\) es factor.
\item Probamos \(r=3\). Coeficientes actuales: \((6, 0, -210, 180, 1824, -2880)\).
\item La fila final de Ruffini es \((6, 18, -156, -288, 960, 0)\). El resto es \(0\), por eso el factor es \((x-3)\). El cociente queda \(6 x^{4} + 18 x^{3} - 156 x^{2} - 288 x + 960\).
\item Probamos \(r=-5\). Coeficientes actuales: \((6, 18, -156, -288, 960)\).
\item La fila final de Ruffini es \((6, -12, -96, 192, 0)\). El resto es \(0\), por eso el factor es \((x+5)\). El cociente queda \(6 x^{3} - 12 x^{2} - 96 x + 192\).
\item Probamos \(r=4\). Coeficientes actuales: \((6, -12, -96, 192)\).
\item La fila final de Ruffini es \((6, 12, -48, 0)\). El resto es \(0\), por eso el factor es \((x-4)\). El cociente queda \(6 x^{2} + 12 x - 48\).
\item Probamos \(r=-4\). Coeficientes actuales: \((6, 12, -48)\).
\item La fila final de Ruffini es \((6, -12, 0)\). El resto es \(0\), por eso el factor es \((x+4)\). El cociente queda \(6 x - 12\).
\item Probamos \(r=2\). Coeficientes actuales: \((6, -12)\).
\item La fila final de Ruffini es \((6, 0)\). El resto es \(0\), por eso el factor es \((x-2)\). El cociente queda \(6\).
\item Juntamos todos los factores encontrados. Resultado: \[\boxed{6 \left(x - 4\right) \left(x - 3\right) \left(x - 2\right) \left(x + 4\right) \left(x + 5\right)}\]
\end{enumerate}

\vspace{0.2cm}
\subsubsection*{Ejercicio 495 - Ruffini}
\textbf{Enunciado:} $ 2 x^{5} - 18 x^{4} + 26 x^{3} + 114 x^{2} - 172 x - 240 $\par
\begin{enumerate}[leftmargin=2em]
\item Ruffini sirve para dividir un polinomio entre factores de la forma \((x-r)\). Si el resto da cero, entonces \((x-r)\) es factor.
\item Probamos \(r=4\). Coeficientes actuales: \((2, -18, 26, 114, -172, -240)\).
\item La fila final de Ruffini es \((2, -10, -14, 58, 60, 0)\). El resto es \(0\), por eso el factor es \((x-4)\). El cociente queda \(2 x^{4} - 10 x^{3} - 14 x^{2} + 58 x + 60\).
\item Probamos \(r=-1\). Coeficientes actuales: \((2, -10, -14, 58, 60)\).
\item La fila final de Ruffini es \((2, -12, -2, 60, 0)\). El resto es \(0\), por eso el factor es \((x+1)\). El cociente queda \(2 x^{3} - 12 x^{2} - 2 x + 60\).
\item Probamos \(r=5\). Coeficientes actuales: \((2, -12, -2, 60)\).
\item La fila final de Ruffini es \((2, -2, -12, 0)\). El resto es \(0\), por eso el factor es \((x-5)\). El cociente queda \(2 x^{2} - 2 x - 12\).
\item Probamos \(r=-2\). Coeficientes actuales: \((2, -2, -12)\).
\item La fila final de Ruffini es \((2, -6, 0)\). El resto es \(0\), por eso el factor es \((x+2)\). El cociente queda \(2 x - 6\).
\item Probamos \(r=3\). Coeficientes actuales: \((2, -6)\).
\item La fila final de Ruffini es \((2, 0)\). El resto es \(0\), por eso el factor es \((x-3)\). El cociente queda \(2\).
\item Juntamos todos los factores encontrados. Resultado: \[\boxed{2 \left(x - 5\right) \left(x - 4\right) \left(x - 3\right) \left(x + 1\right) \left(x + 2\right)}\]
\end{enumerate}

\vspace{0.2cm}
\subsubsection*{Ejercicio 496 - Ruffini}
\textbf{Enunciado:} $ 3 x^{5} - 12 x^{4} - 45 x^{3} + 138 x^{2} + 132 x - 360 $\par
\begin{enumerate}[leftmargin=2em]
\item Ruffini sirve para dividir un polinomio entre factores de la forma \((x-r)\). Si el resto da cero, entonces \((x-r)\) es factor.
\item Probamos \(r=5\). Coeficientes actuales: \((3, -12, -45, 138, 132, -360)\).
\item La fila final de Ruffini es \((3, 3, -30, -12, 72, 0)\). El resto es \(0\), por eso el factor es \((x-5)\). El cociente queda \(3 x^{4} + 3 x^{3} - 30 x^{2} - 12 x + 72\).
\item Probamos \(r=-2\). Coeficientes actuales: \((3, 3, -30, -12, 72)\).
\item La fila final de Ruffini es \((3, -3, -24, 36, 0)\). El resto es \(0\), por eso el factor es \((x+2)\). El cociente queda \(3 x^{3} - 3 x^{2} - 24 x + 36\).
\item Probamos \(r=2\). Coeficientes actuales: \((3, -3, -24, 36)\).
\item La fila final de Ruffini es \((3, 3, -18, 0)\). El resto es \(0\), por eso el factor es \((x-2)\). El cociente queda \(3 x^{2} + 3 x - 18\).
\item Probamos \(r=-3\). Coeficientes actuales: \((3, 3, -18)\).
\item La fila final de Ruffini es \((3, -6, 0)\). El resto es \(0\), por eso el factor es \((x+3)\). El cociente queda \(3 x - 6\).
\item Probamos \(r=2\). Coeficientes actuales: \((3, -6)\).
\item La fila final de Ruffini es \((3, 0)\). El resto es \(0\), por eso el factor es \((x-2)\). El cociente queda \(3\).
\item Juntamos todos los factores encontrados. Resultado: \[\boxed{3 \left(x - 5\right) \left(x - 2\right)^{2} \left(x + 2\right) \left(x + 3\right)}\]
\end{enumerate}

\vspace{0.2cm}
\subsubsection*{Ejercicio 497 - Ruffini}
\textbf{Enunciado:} $ 4 x^{5} - 20 x^{4} - 108 x^{3} + 468 x^{2} + 648 x - 2592 $\par
\begin{enumerate}[leftmargin=2em]
\item Ruffini sirve para dividir un polinomio entre factores de la forma \((x-r)\). Si el resto da cero, entonces \((x-r)\) es factor.
\item Probamos \(r=6\). Coeficientes actuales: \((4, -20, -108, 468, 648, -2592)\).
\item La fila final de Ruffini es \((4, 4, -84, -36, 432, 0)\). El resto es \(0\), por eso el factor es \((x-6)\). El cociente queda \(4 x^{4} + 4 x^{3} - 84 x^{2} - 36 x + 432\).
\item Probamos \(r=-3\). Coeficientes actuales: \((4, 4, -84, -36, 432)\).
\item La fila final de Ruffini es \((4, -8, -60, 144, 0)\). El resto es \(0\), por eso el factor es \((x+3)\). El cociente queda \(4 x^{3} - 8 x^{2} - 60 x + 144\).
\item Probamos \(r=3\). Coeficientes actuales: \((4, -8, -60, 144)\).
\item La fila final de Ruffini es \((4, 4, -48, 0)\). El resto es \(0\), por eso el factor es \((x-3)\). El cociente queda \(4 x^{2} + 4 x - 48\).
\item Probamos \(r=-4\). Coeficientes actuales: \((4, 4, -48)\).
\item La fila final de Ruffini es \((4, -12, 0)\). El resto es \(0\), por eso el factor es \((x+4)\). El cociente queda \(4 x - 12\).
\item Probamos \(r=3\). Coeficientes actuales: \((4, -12)\).
\item La fila final de Ruffini es \((4, 0)\). El resto es \(0\), por eso el factor es \((x-3)\). El cociente queda \(4\).
\item Juntamos todos los factores encontrados. Resultado: \[\boxed{4 \left(x - 6\right) \left(x - 3\right)^{2} \left(x + 3\right) \left(x + 4\right)}\]
\end{enumerate}

\vspace{0.2cm}
\subsubsection*{Ejercicio 498 - Ruffini}
\textbf{Enunciado:} $ 5 x^{5} - 5 x^{4} - 100 x^{3} + 100 x^{2} + 320 x - 320 $\par
\begin{enumerate}[leftmargin=2em]
\item Ruffini sirve para dividir un polinomio entre factores de la forma \((x-r)\). Si el resto da cero, entonces \((x-r)\) es factor.
\item Probamos \(r=1\). Coeficientes actuales: \((5, -5, -100, 100, 320, -320)\).
\item La fila final de Ruffini es \((5, 0, -100, 0, 320, 0)\). El resto es \(0\), por eso el factor es \((x-1)\). El cociente queda \(5 x^{4} - 100 x^{2} + 320\).
\item Probamos \(r=-4\). Coeficientes actuales: \((5, 0, -100, 0, 320)\).
\item La fila final de Ruffini es \((5, -20, -20, 80, 0)\). El resto es \(0\), por eso el factor es \((x+4)\). El cociente queda \(5 x^{3} - 20 x^{2} - 20 x + 80\).
\item Probamos \(r=4\). Coeficientes actuales: \((5, -20, -20, 80)\).
\item La fila final de Ruffini es \((5, 0, -20, 0)\). El resto es \(0\), por eso el factor es \((x-4)\). El cociente queda \(5 x^{2} - 20\).
\item Probamos \(r=-2\). Coeficientes actuales: \((5, 0, -20)\).
\item La fila final de Ruffini es \((5, -10, 0)\). El resto es \(0\), por eso el factor es \((x+2)\). El cociente queda \(5 x - 10\).
\item Probamos \(r=2\). Coeficientes actuales: \((5, -10)\).
\item La fila final de Ruffini es \((5, 0)\). El resto es \(0\), por eso el factor es \((x-2)\). El cociente queda \(5\).
\item Juntamos todos los factores encontrados. Resultado: \[\boxed{5 \left(x - 4\right) \left(x - 2\right) \left(x - 1\right) \left(x + 2\right) \left(x + 4\right)}\]
\end{enumerate}

\vspace{0.2cm}
\subsubsection*{Ejercicio 499 - Ruffini}
\textbf{Enunciado:} $ 6 x^{5} - 12 x^{4} - 204 x^{3} + 408 x^{2} + 1350 x - 2700 $\par
\begin{enumerate}[leftmargin=2em]
\item Ruffini sirve para dividir un polinomio entre factores de la forma \((x-r)\). Si el resto da cero, entonces \((x-r)\) es factor.
\item Probamos \(r=2\). Coeficientes actuales: \((6, -12, -204, 408, 1350, -2700)\).
\item La fila final de Ruffini es \((6, 0, -204, 0, 1350, 0)\). El resto es \(0\), por eso el factor es \((x-2)\). El cociente queda \(6 x^{4} - 204 x^{2} + 1350\).
\item Probamos \(r=-5\). Coeficientes actuales: \((6, 0, -204, 0, 1350)\).
\item La fila final de Ruffini es \((6, -30, -54, 270, 0)\). El resto es \(0\), por eso el factor es \((x+5)\). El cociente queda \(6 x^{3} - 30 x^{2} - 54 x + 270\).
\item Probamos \(r=5\). Coeficientes actuales: \((6, -30, -54, 270)\).
\item La fila final de Ruffini es \((6, 0, -54, 0)\). El resto es \(0\), por eso el factor es \((x-5)\). El cociente queda \(6 x^{2} - 54\).
\item Probamos \(r=-3\). Coeficientes actuales: \((6, 0, -54)\).
\item La fila final de Ruffini es \((6, -18, 0)\). El resto es \(0\), por eso el factor es \((x+3)\). El cociente queda \(6 x - 18\).
\item Probamos \(r=3\). Coeficientes actuales: \((6, -18)\).
\item La fila final de Ruffini es \((6, 0)\). El resto es \(0\), por eso el factor es \((x-3)\). El cociente queda \(6\).
\item Juntamos todos los factores encontrados. Resultado: \[\boxed{6 \left(x - 5\right) \left(x - 3\right) \left(x - 2\right) \left(x + 3\right) \left(x + 5\right)}\]
\end{enumerate}

\vspace{0.2cm}
\subsubsection*{Ejercicio 500 - Ruffini}
\textbf{Enunciado:} $ 2 x^{5} - 4 x^{4} - 30 x^{3} + 80 x^{2} + 8 x - 96 $\par
\begin{enumerate}[leftmargin=2em]
\item Ruffini sirve para dividir un polinomio entre factores de la forma \((x-r)\). Si el resto da cero, entonces \((x-r)\) es factor.
\item Probamos \(r=3\). Coeficientes actuales: \((2, -4, -30, 80, 8, -96)\).
\item La fila final de Ruffini es \((2, 2, -24, 8, 32, 0)\). El resto es \(0\), por eso el factor es \((x-3)\). El cociente queda \(2 x^{4} + 2 x^{3} - 24 x^{2} + 8 x + 32\).
\item Probamos \(r=-1\). Coeficientes actuales: \((2, 2, -24, 8, 32)\).
\item La fila final de Ruffini es \((2, 0, -24, 32, 0)\). El resto es \(0\), por eso el factor es \((x+1)\). El cociente queda \(2 x^{3} - 24 x + 32\).
\item Probamos \(r=2\). Coeficientes actuales: \((2, 0, -24, 32)\).
\item La fila final de Ruffini es \((2, 4, -16, 0)\). El resto es \(0\), por eso el factor es \((x-2)\). El cociente queda \(2 x^{2} + 4 x - 16\).
\item Probamos \(r=-4\). Coeficientes actuales: \((2, 4, -16)\).
\item La fila final de Ruffini es \((2, -4, 0)\). El resto es \(0\), por eso el factor es \((x+4)\). El cociente queda \(2 x - 4\).
\item Probamos \(r=2\). Coeficientes actuales: \((2, -4)\).
\item La fila final de Ruffini es \((2, 0)\). El resto es \(0\), por eso el factor es \((x-2)\). El cociente queda \(2\).
\item Juntamos todos los factores encontrados. Resultado: \[\boxed{2 \left(x - 3\right) \left(x - 2\right)^{2} \left(x + 1\right) \left(x + 4\right)}\]
\end{enumerate}

\vspace{0.2cm}
\end{document}
